% Options for packages loaded elsewhere
\PassOptionsToPackage{unicode}{hyperref}
\PassOptionsToPackage{hyphens}{url}
\documentclass[
  spanish,
  11pt,
  a4paper]{article}
\usepackage{xcolor}
\usepackage[top=2.2cm, bottom=2.2cm, left=2.3cm, right=2.3cm]{geometry}
\usepackage{amsmath,amssymb}
\setcounter{secnumdepth}{5}
\usepackage{iftex}
\ifPDFTeX
  \usepackage[T1]{fontenc}
  \usepackage[utf8]{inputenc}
  \usepackage{textcomp} % provide euro and other symbols
\else % if luatex or xetex
  \usepackage{unicode-math} % this also loads fontspec
  \defaultfontfeatures{Scale=MatchLowercase}
  \defaultfontfeatures[\rmfamily]{Ligatures=TeX,Scale=1}
\fi
\usepackage{lmodern}
\ifPDFTeX\else
  % xetex/luatex font selection
\fi
% Use upquote if available, for straight quotes in verbatim environments
\IfFileExists{upquote.sty}{\usepackage{upquote}}{}
\IfFileExists{microtype.sty}{% use microtype if available
  \usepackage[]{microtype}
  \UseMicrotypeSet[protrusion]{basicmath} % disable protrusion for tt fonts
}{}
\usepackage{setspace}
\makeatletter
\@ifundefined{KOMAClassName}{% if non-KOMA class
  \IfFileExists{parskip.sty}{%
    \usepackage{parskip}
  }{% else
    \setlength{\parindent}{0pt}
    \setlength{\parskip}{6pt plus 2pt minus 1pt}}
}{% if KOMA class
  \KOMAoptions{parskip=half}}
\makeatother
\usepackage{color}
\usepackage{fancyvrb}
\newcommand{\VerbBar}{|}
\newcommand{\VERB}{\Verb[commandchars=\\\{\}]}
\DefineVerbatimEnvironment{Highlighting}{Verbatim}{commandchars=\\\{\}}
% Add ',fontsize=\small' for more characters per line
\usepackage{framed}
\definecolor{shadecolor}{RGB}{248,248,248}
\newenvironment{Shaded}{\begin{snugshade}}{\end{snugshade}}
\newcommand{\AlertTok}[1]{\textcolor[rgb]{0.94,0.16,0.16}{#1}}
\newcommand{\AnnotationTok}[1]{\textcolor[rgb]{0.56,0.35,0.01}{\textbf{\textit{#1}}}}
\newcommand{\AttributeTok}[1]{\textcolor[rgb]{0.13,0.29,0.53}{#1}}
\newcommand{\BaseNTok}[1]{\textcolor[rgb]{0.00,0.00,0.81}{#1}}
\newcommand{\BuiltInTok}[1]{#1}
\newcommand{\CharTok}[1]{\textcolor[rgb]{0.31,0.60,0.02}{#1}}
\newcommand{\CommentTok}[1]{\textcolor[rgb]{0.56,0.35,0.01}{\textit{#1}}}
\newcommand{\CommentVarTok}[1]{\textcolor[rgb]{0.56,0.35,0.01}{\textbf{\textit{#1}}}}
\newcommand{\ConstantTok}[1]{\textcolor[rgb]{0.56,0.35,0.01}{#1}}
\newcommand{\ControlFlowTok}[1]{\textcolor[rgb]{0.13,0.29,0.53}{\textbf{#1}}}
\newcommand{\DataTypeTok}[1]{\textcolor[rgb]{0.13,0.29,0.53}{#1}}
\newcommand{\DecValTok}[1]{\textcolor[rgb]{0.00,0.00,0.81}{#1}}
\newcommand{\DocumentationTok}[1]{\textcolor[rgb]{0.56,0.35,0.01}{\textbf{\textit{#1}}}}
\newcommand{\ErrorTok}[1]{\textcolor[rgb]{0.64,0.00,0.00}{\textbf{#1}}}
\newcommand{\ExtensionTok}[1]{#1}
\newcommand{\FloatTok}[1]{\textcolor[rgb]{0.00,0.00,0.81}{#1}}
\newcommand{\FunctionTok}[1]{\textcolor[rgb]{0.13,0.29,0.53}{\textbf{#1}}}
\newcommand{\ImportTok}[1]{#1}
\newcommand{\InformationTok}[1]{\textcolor[rgb]{0.56,0.35,0.01}{\textbf{\textit{#1}}}}
\newcommand{\KeywordTok}[1]{\textcolor[rgb]{0.13,0.29,0.53}{\textbf{#1}}}
\newcommand{\NormalTok}[1]{#1}
\newcommand{\OperatorTok}[1]{\textcolor[rgb]{0.81,0.36,0.00}{\textbf{#1}}}
\newcommand{\OtherTok}[1]{\textcolor[rgb]{0.56,0.35,0.01}{#1}}
\newcommand{\PreprocessorTok}[1]{\textcolor[rgb]{0.56,0.35,0.01}{\textit{#1}}}
\newcommand{\RegionMarkerTok}[1]{#1}
\newcommand{\SpecialCharTok}[1]{\textcolor[rgb]{0.81,0.36,0.00}{\textbf{#1}}}
\newcommand{\SpecialStringTok}[1]{\textcolor[rgb]{0.31,0.60,0.02}{#1}}
\newcommand{\StringTok}[1]{\textcolor[rgb]{0.31,0.60,0.02}{#1}}
\newcommand{\VariableTok}[1]{\textcolor[rgb]{0.00,0.00,0.00}{#1}}
\newcommand{\VerbatimStringTok}[1]{\textcolor[rgb]{0.31,0.60,0.02}{#1}}
\newcommand{\WarningTok}[1]{\textcolor[rgb]{0.56,0.35,0.01}{\textbf{\textit{#1}}}}
\usepackage{longtable,booktabs,array}
\newcounter{none} % for unnumbered tables
\usepackage{calc} % for calculating minipage widths
% Correct order of tables after \paragraph or \subparagraph
\usepackage{etoolbox}
\makeatletter
\patchcmd\longtable{\par}{\if@noskipsec\mbox{}\fi\par}{}{}
\makeatother
% Allow footnotes in longtable head/foot
\IfFileExists{footnotehyper.sty}{\usepackage{footnotehyper}}{\usepackage{footnote}}
\makesavenoteenv{longtable}
\usepackage{graphicx}
\makeatletter
\newsavebox\pandoc@box
\newcommand*\pandocbounded[1]{% scales image to fit in text height/width
  \sbox\pandoc@box{#1}%
  \Gscale@div\@tempa{\textheight}{\dimexpr\ht\pandoc@box+\dp\pandoc@box\relax}%
  \Gscale@div\@tempb{\linewidth}{\wd\pandoc@box}%
  \ifdim\@tempb\p@<\@tempa\p@\let\@tempa\@tempb\fi% select the smaller of both
  \ifdim\@tempa\p@<\p@\scalebox{\@tempa}{\usebox\pandoc@box}%
  \else\usebox{\pandoc@box}%
  \fi%
}
% Set default figure placement to htbp
\def\fps@figure{htbp}
\makeatother
\ifLuaTeX
\usepackage[bidi=basic,shorthands=off]{babel}
\else
\usepackage[bidi=default,shorthands=off]{babel}
\fi
\ifLuaTeX
  \usepackage{selnolig} % disable illegal ligatures
\fi
\setlength{\emergencystretch}{3em} % prevent overfull lines
\providecommand{\tightlist}{%
  \setlength{\itemsep}{0pt}\setlength{\parskip}{0pt}}
\usepackage{microtype}
\usepackage{fancyhdr}
\usepackage{lastpage}
\usepackage{fvextra}
\usepackage{etoolbox}

\setlength{\parindent}{0pt}
\setlength{\parskip}{0.35em}

\fvset{
  breaklines=true,
  breakanywhere=true
}

\AtBeginEnvironment{Highlighting}{\small}

\pagestyle{fancy}
\fancyhf{}
\fancyhead[L]{Manual de referencia de GROMACS}
\fancyhead[R]{\nouppercase{\leftmark}}
\fancyfoot[C]{Página \thepage\ de \pageref{LastPage}}
\setlength{\headheight}{14pt}

\widowpenalty=10000
\clubpenalty=10000
\displaywidowpenalty=10000
\usepackage{bookmark}
\IfFileExists{xurl.sty}{\usepackage{xurl}}{} % add URL line breaks if available
\urlstyle{same}
\hypersetup{
  pdftitle={Manual de referencia de GROMACS},
  pdflang={es},
  hidelinks,
  pdfcreator={LaTeX via pandoc}}

\title{Manual de referencia de GROMACS}
\author{}
\date{\vspace{-2.5em}}

\begin{document}
\maketitle

{
\setcounter{tocdepth}{3}
\tableofcontents
}
\setstretch{1.12}
\section{Protocolos de trabajo con GROMACS
2026.3}\label{protocolos-de-trabajo-con-gromacs-2026.3}

\emph{Dr.~JJ Casal, PhD}

Versión 2023.21.10.9.44

\begin{quote}
The greatest obstacle to discovery is not ignorance -- it is the
illusion of knowledge.

Never tell people how to do things. Tell them what to do and they will
surprise you with their ingenuity.

General George S. Patton

Desordené átomos tuyos para hacerte aparecer.

Puente - Gustavo Cerati

Jerry, just remember, it's not a lie if you believe it.

George Costanza

Sometimes Only sometimes I question everything

Sometimes - Depeche Mode

The laughter penetrates my silence As drunken men find flaws in science

Set The Fire To The Third Bar - Snow Patrol
\end{quote}

\section{Estrategia para la
simulación}\label{estrategia-para-la-simulaciuxf3n}

\section{Instalación de GROMACS}\label{instalaciuxf3n-de-gromacs}

Esta sección corresponde a \textbf{GROMACS 2026.3}. Antes de instalar
conviene decidir qué se necesita:

\begin{itemize}
\tightlist
\item
  \textbf{Paquete de la distribución:} instalación rápida para
  aprendizaje, análisis y pruebas. Puede no ser la versión más reciente
  ni estar optimizado para el hardware.
\item
  \textbf{Compilación desde código fuente:} recomendada para producción,
  GPU, clústeres y control preciso de dependencias.
\item
  \textbf{MPI externo:} necesario principalmente para ejecuciones que
  abarcan varios nodos. Para una sola estación de trabajo, la
  compilación normal con thread-MPI suele ser suficiente.
\item
  \textbf{gmxapi:} interfaz de Python opcional. Requiere primero una
  instalación compatible de GROMACS.
\end{itemize}

La guía oficial vigente requiere CMake 3.28 o posterior, un compilador
C99 y un compilador C++17. Para GROMACS 2026.3, las versiones mínimas
indicadas son GCC 11, Clang 14 o MSVC 2019.

\subsection{Comprobación del sistema}\label{comprobaciuxf3n-del-sistema}

Antes de instalar:

\begin{Shaded}
\begin{Highlighting}[]
\FunctionTok{uname} \AttributeTok{{-}a}
\FunctionTok{cmake} \AttributeTok{{-}{-}version}
\FunctionTok{gcc} \AttributeTok{{-}{-}version}
\ExtensionTok{g++} \AttributeTok{{-}{-}version}
\ExtensionTok{python3} \AttributeTok{{-}{-}version}
\end{Highlighting}
\end{Shaded}

Si se utilizará una GPU NVIDIA:

\begin{Shaded}
\begin{Highlighting}[]
\ExtensionTok{nvidia{-}smi}
\ExtensionTok{nvcc} \AttributeTok{{-}{-}version}
\end{Highlighting}
\end{Shaded}

\textbf{nvidia-smi} informa el controlador instalado y la versión máxima
de CUDA admitida por ese controlador. \textbf{nvcc --version} informa la
versión del toolkit usado para compilar. No son la misma cosa.

Según la compatibilidad menor publicada por NVIDIA:

{\def\LTcaptype{none} % do not increment counter
\begin{longtable}[]{@{}lr@{}}
\toprule\noalign{}
Familia del toolkit & Controlador mínimo \\
\midrule\noalign{}
\endhead
\bottomrule\noalign{}
\endlastfoot
CUDA 13.x & 580 \\
CUDA 12.x & 525 \\
CUDA 11.x & 450 \\
\end{longtable}
}

Un controlador más nuevo puede ejecutar aplicaciones compiladas con una
familia CUDA anterior mediante compatibilidad hacia atrás. La tabla es
un requisito mínimo general; deben revisarse también las notas de la
versión concreta del toolkit y la compatibilidad de la GPU.

\subsection{Instalación mediante el gestor de
paquetes}\label{instalaciuxf3n-mediante-el-gestor-de-paquetes}

\subsubsection{Ubuntu y Debian}\label{ubuntu-y-debian}

\begin{Shaded}
\begin{Highlighting}[]
\FunctionTok{sudo}\NormalTok{ apt update}
\FunctionTok{sudo}\NormalTok{ apt install gromacs}
\end{Highlighting}
\end{Shaded}

Comprobación:

\begin{Shaded}
\begin{Highlighting}[]
\ExtensionTok{gmx} \AttributeTok{{-}{-}version}
\end{Highlighting}
\end{Shaded}

El paquete disponible depende de la versión de Ubuntu o Debian. Esta vía
es adecuada si la versión empaquetada satisface el protocolo. Para
simulaciones de producción conviene revisar en la salida de \textbf{gmx
--version} el soporte SIMD, FFT, MPI y GPU.

\subsubsection{Fedora}\label{fedora}

\begin{Shaded}
\begin{Highlighting}[]
\FunctionTok{sudo}\NormalTok{ dnf install gromacs}
\end{Highlighting}
\end{Shaded}

Algunas variantes empaquetadas pueden instalar ejecutables o módulos
separados para MPI. Compruebe los nombres proporcionados por la versión
de Fedora:

\begin{Shaded}
\begin{Highlighting}[]
\ExtensionTok{rpm} \AttributeTok{{-}ql}\NormalTok{ gromacs }\KeywordTok{|} \FunctionTok{grep}\NormalTok{ /bin/}
\ExtensionTok{gmx} \AttributeTok{{-}{-}version}
\end{Highlighting}
\end{Shaded}

\subsubsection{Arch Linux y Manjaro}\label{arch-linux-y-manjaro}

\begin{Shaded}
\begin{Highlighting}[]
\FunctionTok{sudo}\NormalTok{ pacman }\AttributeTok{{-}S}\NormalTok{ gromacs}
\ExtensionTok{gmx} \AttributeTok{{-}{-}version}
\end{Highlighting}
\end{Shaded}

\subsubsection{macOS con Homebrew}\label{macos-con-homebrew}

\begin{Shaded}
\begin{Highlighting}[]
\ExtensionTok{brew}\NormalTok{ update}
\ExtensionTok{brew}\NormalTok{ install gromacs}
\ExtensionTok{gmx} \AttributeTok{{-}{-}version}
\end{Highlighting}
\end{Shaded}

La aceleración CUDA no está disponible en macOS. En equipos Apple
Silicon, la compilación nativa puede aprovechar SIMD y la GPU solo
mediante backends admitidos por la versión de GROMACS y las herramientas
disponibles; no debe asumirse que una fórmula de Homebrew incluye
aceleración por GPU.

\subsection{Windows y WSL2}\label{windows-y-wsl2}

La ruta más práctica para ejecutar GROMACS en Windows es WSL2 con una
distribución Linux. Después de instalar WSL2 y Ubuntu, los comandos de
instalación y compilación son los mismos que en Linux.

Desde PowerShell con privilegios de administrador:

\begin{Shaded}
\begin{Highlighting}[]
\NormalTok{wsl }\OperatorTok{{-}{-}}\NormalTok{install }\OperatorTok{{-}}\NormalTok{d Ubuntu}
\NormalTok{wsl }\OperatorTok{{-}{-}}\NormalTok{update}
\end{Highlighting}
\end{Shaded}

Dentro de Ubuntu:

\begin{Shaded}
\begin{Highlighting}[]
\FunctionTok{sudo}\NormalTok{ apt update}
\FunctionTok{sudo}\NormalTok{ apt install gromacs}
\ExtensionTok{gmx} \AttributeTok{{-}{-}version}
\end{Highlighting}
\end{Shaded}

Para usar una GPU NVIDIA dentro de WSL2 se necesita un controlador de
Windows compatible con WSL; no debe instalarse un controlador Linux
NVIDIA dentro de la distribución WSL. El toolkit CUDA de usuario puede
instalarse dentro de WSL cuando sea necesario para compilar. Verifique
desde WSL:

\begin{Shaded}
\begin{Highlighting}[]
\ExtensionTok{nvidia{-}smi}
\end{Highlighting}
\end{Shaded}

\subsection{Compilación desde código
fuente}\label{compilaciuxf3n-desde-cuxf3digo-fuente}

\subsubsection{Dependencias en Ubuntu o
Debian}\label{dependencias-en-ubuntu-o-debian}

\begin{Shaded}
\begin{Highlighting}[]
\FunctionTok{sudo}\NormalTok{ apt update}
\FunctionTok{sudo}\NormalTok{ apt install build{-}essential cmake git python3     libfftw3{-}dev libhwloc{-}dev}
\end{Highlighting}
\end{Shaded}

Para una compilación MPI agregue:

\begin{Shaded}
\begin{Highlighting}[]
\FunctionTok{sudo}\NormalTok{ apt install openmpi{-}bin libopenmpi{-}dev}
\end{Highlighting}
\end{Shaded}

\subsubsection{Dependencias en Fedora}\label{dependencias-en-fedora}

\begin{Shaded}
\begin{Highlighting}[]
\FunctionTok{sudo}\NormalTok{ dnf install gcc gcc{-}c++ cmake make git python3     fftw{-}devel hwloc{-}devel}
\end{Highlighting}
\end{Shaded}

Para MPI:

\begin{Shaded}
\begin{Highlighting}[]
\FunctionTok{sudo}\NormalTok{ dnf install openmpi openmpi{-}devel}
\end{Highlighting}
\end{Shaded}

En Fedora puede ser necesario cargar el entorno de OpenMPI según cómo
esté empaquetado:

\begin{Shaded}
\begin{Highlighting}[]
\ExtensionTok{module}\NormalTok{ load mpi/openmpi{-}x86\_64}
\end{Highlighting}
\end{Shaded}

Si el comando \textbf{module} no existe o el módulo tiene otro nombre,
examine los archivos instalados por el paquete OpenMPI.

\subsubsection{Descarga y compilación básica para
CPU}\label{descarga-y-compilaciuxf3n-buxe1sica-para-cpu}

Use una carpeta de compilación separada del código fuente. El prefijo
bajo \texttt{\$HOME/opt} evita requerir permisos de administrador
durante la instalación.

\begin{Shaded}
\begin{Highlighting}[]
\FunctionTok{wget}\NormalTok{ https://ftp.gromacs.org/gromacs/gromacs{-}2026.3.tar.gz}
\FunctionTok{tar}\NormalTok{ xfz gromacs{-}2026.3.tar.gz}
\BuiltInTok{cd}\NormalTok{ gromacs{-}2026.3}

\FunctionTok{mkdir}\NormalTok{ build}
\BuiltInTok{cd}\NormalTok{ build}

\FunctionTok{cmake}\NormalTok{ ..     }\AttributeTok{{-}DGMX\_BUILD\_OWN\_FFTW}\OperatorTok{=}\NormalTok{ON     }\AttributeTok{{-}DREGRESSIONTEST\_DOWNLOAD}\OperatorTok{=}\NormalTok{ON     }\AttributeTok{{-}DCMAKE\_INSTALL\_PREFIX}\OperatorTok{=}\StringTok{"}\VariableTok{$HOME}\StringTok{/opt/gromacs{-}2026.3"}

\FunctionTok{cmake} \AttributeTok{{-}{-}build}\NormalTok{ . }\AttributeTok{{-}{-}parallel}
\ExtensionTok{ctest} \AttributeTok{{-}{-}output{-}on{-}failure}
\FunctionTok{cmake} \AttributeTok{{-}{-}install}\NormalTok{ .}
\end{Highlighting}
\end{Shaded}

Active la instalación:

\begin{Shaded}
\begin{Highlighting}[]
\BuiltInTok{source} \StringTok{"}\VariableTok{$HOME}\StringTok{/opt/gromacs{-}2026.3/bin/GMXRC"}
\ExtensionTok{gmx} \AttributeTok{{-}{-}version}
\end{Highlighting}
\end{Shaded}

Para cargarla automáticamente al iniciar Bash:

\begin{Shaded}
\begin{Highlighting}[]
\BuiltInTok{echo} \StringTok{\textquotesingle{}source "$HOME/opt/gromacs{-}2026.3/bin/GMXRC"\textquotesingle{}} \OperatorTok{\textgreater{}\textgreater{}} \StringTok{"}\VariableTok{$HOME}\StringTok{/.bashrc"}
\end{Highlighting}
\end{Shaded}

\textbf{GMX\_BUILD\_OWN\_FFTW=ON} descarga y compila FFTW. Si existe una
instalación adecuada de FFTW puede omitirse. La FFTW suministrada por
GROMACS es una elección simple y reproducible para una estación de
trabajo.

\subsection{Compilación con GPU
NVIDIA}\label{compilaciuxf3n-con-gpu-nvidia}

Primero instale un controlador NVIDIA compatible y el toolkit CUDA
siguiendo el método correspondiente a la distribución. No mezcle
paquetes CUDA de repositorios incompatibles. Verifique:

\begin{Shaded}
\begin{Highlighting}[]
\ExtensionTok{nvidia{-}smi}
\ExtensionTok{nvcc} \AttributeTok{{-}{-}version}
\end{Highlighting}
\end{Shaded}

Configure GROMACS con el backend CUDA:

\begin{Shaded}
\begin{Highlighting}[]
\BuiltInTok{cd}\NormalTok{ gromacs{-}2026.3}
\FunctionTok{mkdir}\NormalTok{ build{-}cuda}
\BuiltInTok{cd}\NormalTok{ build{-}cuda}

\FunctionTok{cmake}\NormalTok{ ..     }\AttributeTok{{-}DGMX\_BUILD\_OWN\_FFTW}\OperatorTok{=}\NormalTok{ON     }\AttributeTok{{-}DREGRESSIONTEST\_DOWNLOAD}\OperatorTok{=}\NormalTok{ON     }\AttributeTok{{-}DGMX\_GPU}\OperatorTok{=}\NormalTok{CUDA     }\AttributeTok{{-}DCMAKE\_INSTALL\_PREFIX}\OperatorTok{=}\StringTok{"}\VariableTok{$HOME}\StringTok{/opt/gromacs{-}2026.3{-}cuda"}

\FunctionTok{cmake} \AttributeTok{{-}{-}build}\NormalTok{ . }\AttributeTok{{-}{-}parallel}
\ExtensionTok{ctest} \AttributeTok{{-}{-}output{-}on{-}failure}
\FunctionTok{cmake} \AttributeTok{{-}{-}install}\NormalTok{ .}
\end{Highlighting}
\end{Shaded}

Active y compruebe:

\begin{Shaded}
\begin{Highlighting}[]
\BuiltInTok{source} \StringTok{"}\VariableTok{$HOME}\StringTok{/opt/gromacs{-}2026.3{-}cuda/bin/GMXRC"}
\ExtensionTok{gmx} \AttributeTok{{-}{-}version}
\ExtensionTok{gmx}\NormalTok{ mdrun }\AttributeTok{{-}version}
\end{Highlighting}
\end{Shaded}

La salida debe indicar que GROMACS fue compilado con soporte CUDA. Que
CUDA aparezca en \textbf{nvidia-smi} no demuestra que el ejecutable de
GROMACS tenga soporte GPU.

Para comprobar el acceso real al dispositivo:

\begin{Shaded}
\begin{Highlighting}[]
\ExtensionTok{nvidia{-}smi} \AttributeTok{{-}L}
\end{Highlighting}
\end{Shaded}

La aceleración efectiva depende del tamaño del sistema, el modelo de
GPU, CPU, red, configuración PME y opciones de \textbf{gmx mdrun}. No se
debe forzar la descarga de todas las tareas a GPU sin medir el
rendimiento.

\subsection{Compilación con MPI para
clústeres}\label{compilaciuxf3n-con-mpi-para-cluxfasteres}

La compilación MPI externa se usa para distribuir una simulación entre
varios nodos. Puede coexistir con la instalación normal; el ejecutable
suele llamarse \textbf{gmx\_mpi}.

\begin{Shaded}
\begin{Highlighting}[]
\BuiltInTok{cd}\NormalTok{ gromacs{-}2026.3}
\FunctionTok{mkdir}\NormalTok{ build{-}mpi}
\BuiltInTok{cd}\NormalTok{ build{-}mpi}

\FunctionTok{cmake}\NormalTok{ ..     }\AttributeTok{{-}DGMX\_BUILD\_OWN\_FFTW}\OperatorTok{=}\NormalTok{ON     }\AttributeTok{{-}DREGRESSIONTEST\_DOWNLOAD}\OperatorTok{=}\NormalTok{ON     }\AttributeTok{{-}DGMX\_MPI}\OperatorTok{=}\NormalTok{ON     }\AttributeTok{{-}DCMAKE\_INSTALL\_PREFIX}\OperatorTok{=}\StringTok{"}\VariableTok{$HOME}\StringTok{/opt/gromacs{-}2026.3{-}mpi"}

\FunctionTok{cmake} \AttributeTok{{-}{-}build}\NormalTok{ . }\AttributeTok{{-}{-}parallel}
\ExtensionTok{ctest} \AttributeTok{{-}{-}output{-}on{-}failure}
\FunctionTok{cmake} \AttributeTok{{-}{-}install}\NormalTok{ .}
\end{Highlighting}
\end{Shaded}

Comprobación:

\begin{Shaded}
\begin{Highlighting}[]
\BuiltInTok{source} \StringTok{"}\VariableTok{$HOME}\StringTok{/opt/gromacs{-}2026.3{-}mpi/bin/GMXRC"}
\ExtensionTok{gmx\_mpi} \AttributeTok{{-}{-}version}
\ExtensionTok{mpirun} \AttributeTok{{-}np}\NormalTok{ 2 gmx\_mpi }\AttributeTok{{-}{-}version}
\end{Highlighting}
\end{Shaded}

Para combinar MPI y CUDA:

\begin{Shaded}
\begin{Highlighting}[]
\FunctionTok{cmake}\NormalTok{ ..     }\AttributeTok{{-}DGMX\_BUILD\_OWN\_FFTW}\OperatorTok{=}\NormalTok{ON     }\AttributeTok{{-}DREGRESSIONTEST\_DOWNLOAD}\OperatorTok{=}\NormalTok{ON     }\AttributeTok{{-}DGMX\_MPI}\OperatorTok{=}\NormalTok{ON     }\AttributeTok{{-}DGMX\_GPU}\OperatorTok{=}\NormalTok{CUDA     }\AttributeTok{{-}DCMAKE\_INSTALL\_PREFIX}\OperatorTok{=}\StringTok{"}\VariableTok{$HOME}\StringTok{/opt/gromacs{-}2026.3{-}mpi{-}cuda"}
\end{Highlighting}
\end{Shaded}

En un clúster deben usarse el lanzador y las variables indicadas por el
gestor de trabajos. Un comando local con \textbf{mpirun} no sustituye un
script de SLURM, PBS u otro planificador.

\subsection{Opciones CMake relevantes}\label{opciones-cmake-relevantes}

{\def\LTcaptype{none} % do not increment counter
\begin{longtable}[]{@{}
  >{\raggedright\arraybackslash}p{(\linewidth - 2\tabcolsep) * \real{0.2385}}
  >{\raggedright\arraybackslash}p{(\linewidth - 2\tabcolsep) * \real{0.7615}}@{}}
\toprule\noalign{}
\begin{minipage}[b]{\linewidth}\raggedright
Opción
\end{minipage} & \begin{minipage}[b]{\linewidth}\raggedright
Uso
\end{minipage} \\
\midrule\noalign{}
\endhead
\bottomrule\noalign{}
\endlastfoot
\textbf{-DGMX\_MPI=ON} & Compila con MPI externo. \\
\textbf{-DGMX\_GPU=CUDA} & Activa GPU NVIDIA mediante CUDA. \\
\textbf{-DGMX\_GPU=OpenCL} & Activa el backend OpenCL. \\
\textbf{-DGMX\_GPU=SYCL} & Activa el backend SYCL. \\
\textbf{-DGMX\_SIMD=valor} & Selecciona explícitamente SIMD; normalmente
conviene la autodetección. \\
\textbf{-DGMX\_DOUBLE=ON} & Compila en doble precisión; es más lento y
rara vez necesario para dinámica molecular convencional. \\
\textbf{-DGMX\_FFT\_LIBRARY=fftw3} & Selecciona la biblioteca FFT. \\
\textbf{-DBUILD\_SHARED\_LIBS=ON} & Genera bibliotecas compartidas; es
necesario para ciertos clientes, incluido gmxapi. \\
\textbf{-DCMAKE\_INSTALL\_PREFIX=ruta} & Define el directorio de
instalación. \\
\textbf{-DCMAKE\_BUILD\_TYPE=Debug} & Compilación de depuración; no es
adecuada para producción. \\
\end{longtable}
}

Las opciones antiguas \textbf{-DGMX\_GPU=ON} y
\textbf{-DGMX\_USE\_OPENCL=ON} no deben usarse con GROMACS 2026. El
backend se selecciona directamente mediante \textbf{-DGMX\_GPU=CUDA},
\textbf{OpenCL} o \textbf{SYCL}.

\subsection{Compilación reproducible dentro de
Docker}\label{compilaciuxf3n-reproducible-dentro-de-docker}

Un contenedor permite fijar la distribución Linux, compiladores,
bibliotecas, versión de CUDA de usuario y opciones de CMake. Esto mejora
la reproducibilidad del entorno y evita contaminar el sistema anfitrión.
No vuelve al ejecutable universal ni garantiza resultados idénticos bit
a bit: el kernel, el controlador NVIDIA, la arquitectura CPU, la GPU y
la asignación de hilos siguen perteneciendo al anfitrión.

La imagen debe construirse para una versión exacta de GROMACS. No
conviene descargar ``la última versión'' durante cada compilación. Los
ejemplos siguientes fijan \textbf{GROMACS 2026.3}, verifican el MD5
publicado para el archivo fuente y utilizan una construcción multietapa
para que la imagen final no contenga compiladores ni archivos
temporales.

\subsubsection{Requisitos del
anfitrión}\label{requisitos-del-anfitriuxf3n}

Para CPU sólo se necesitan Docker Engine o Docker Desktop y espacio
suficiente para compilar:

\begin{Shaded}
\begin{Highlighting}[]
\ExtensionTok{docker}\NormalTok{ version}
\ExtensionTok{docker}\NormalTok{ info}
\end{Highlighting}
\end{Shaded}

En Linux con una GPU NVIDIA se necesitan además:

\begin{enumerate}
\def\labelenumi{\arabic{enumi}.}
\tightlist
\item
  un controlador NVIDIA instalado en el anfitrión;
\item
  NVIDIA Container Toolkit;
\item
  el runtime de Docker configurado para exponer la GPU.
\end{enumerate}

El controlador se instala únicamente en el anfitrión. La imagen contiene
el toolkit y las bibliotecas CUDA de usuario, pero no reemplaza al
controlador.

Después de instalar NVIDIA Container Toolkit:

\begin{Shaded}
\begin{Highlighting}[]
\FunctionTok{sudo}\NormalTok{ nvidia{-}ctk runtime configure }\AttributeTok{{-}{-}runtime}\OperatorTok{=}\NormalTok{docker}
\FunctionTok{sudo}\NormalTok{ systemctl restart docker}
\end{Highlighting}
\end{Shaded}

Compruebe el acceso antes de compilar GROMACS:

\begin{Shaded}
\begin{Highlighting}[]
\ExtensionTok{docker}\NormalTok{ run }\AttributeTok{{-}{-}rm} \AttributeTok{{-}{-}runtime}\OperatorTok{=}\NormalTok{nvidia }\AttributeTok{{-}{-}gpus}\NormalTok{ all ubuntu nvidia{-}smi}
\end{Highlighting}
\end{Shaded}

Si este comando falla, el problema está en Docker, el runtime NVIDIA o
el controlador del anfitrión; recompilar GROMACS no lo corrige.

\subsubsection{Imagen CPU portable para
x86-64}\label{imagen-cpu-portable-para-x86-64}

Cree un archivo llamado \textbf{Dockerfile.cpu}:

\begin{Shaded}
\begin{Highlighting}[]
\KeywordTok{FROM}\NormalTok{ ubuntu:24.04 }\KeywordTok{AS}\NormalTok{ builder}

\KeywordTok{ARG}\NormalTok{ GROMACS\_VERSION=2026.3}
\KeywordTok{ARG}\NormalTok{ GROMACS\_MD5=7987af0c6ab939ab6e639f32d0dd260f}
\KeywordTok{ARG}\NormalTok{ GMX\_SIMD=SSE2}

\KeywordTok{RUN} \ExtensionTok{apt{-}get}\NormalTok{ update }\KeywordTok{\&\&} \ExtensionTok{apt{-}get}\NormalTok{ install }\AttributeTok{{-}y} \AttributeTok{{-}{-}no{-}install{-}recommends} \DataTypeTok{\textbackslash{}}
\NormalTok{    build{-}essential }\DataTypeTok{\textbackslash{}}
\NormalTok{    ca{-}certificates }\DataTypeTok{\textbackslash{}}
\NormalTok{    cmake }\DataTypeTok{\textbackslash{}}
\NormalTok{    curl }\DataTypeTok{\textbackslash{}}
    \KeywordTok{\&\&} \FunctionTok{rm} \AttributeTok{{-}rf}\NormalTok{ /var/lib/apt/lists/}\PreprocessorTok{*}

\KeywordTok{WORKDIR}\NormalTok{ /tmp/build}

\KeywordTok{RUN} \ExtensionTok{curl} \AttributeTok{{-}fL} \DataTypeTok{\textbackslash{}}
    \StringTok{"https://ftp.gromacs.org/gromacs/gromacs{-}}\VariableTok{$\{GROMACS\_VERSION\}}\StringTok{.tar.gz"} \DataTypeTok{\textbackslash{}}
    \AttributeTok{{-}o}\NormalTok{ gromacs.tar.gz }\DataTypeTok{\textbackslash{}}
    \KeywordTok{\&\&} \BuiltInTok{echo} \StringTok{"}\VariableTok{$\{GROMACS\_MD5\}}\StringTok{  gromacs.tar.gz"} \KeywordTok{|} \FunctionTok{md5sum} \AttributeTok{{-}c} \AttributeTok{{-}} \DataTypeTok{\textbackslash{}}
    \KeywordTok{\&\&} \FunctionTok{tar}\NormalTok{ xzf gromacs.tar.gz }\DataTypeTok{\textbackslash{}}
    \KeywordTok{\&\&} \FunctionTok{cmake} \AttributeTok{{-}S} \StringTok{"gromacs{-}}\VariableTok{$\{GROMACS\_VERSION\}}\StringTok{"} \AttributeTok{{-}B}\NormalTok{ gromacs{-}build }\DataTypeTok{\textbackslash{}}
        \AttributeTok{{-}DGMX\_BUILD\_OWN\_FFTW}\OperatorTok{=}\NormalTok{ON }\DataTypeTok{\textbackslash{}}
        \AttributeTok{{-}DREGRESSIONTEST\_DOWNLOAD}\OperatorTok{=}\NormalTok{ON }\DataTypeTok{\textbackslash{}}
        \AttributeTok{{-}DGMX\_SIMD}\OperatorTok{=}\StringTok{"}\VariableTok{$\{GMX\_SIMD\}}\StringTok{"} \DataTypeTok{\textbackslash{}}
        \AttributeTok{{-}DGMX\_GPU}\OperatorTok{=}\NormalTok{OFF }\DataTypeTok{\textbackslash{}}
        \AttributeTok{{-}DGMX\_MPI}\OperatorTok{=}\NormalTok{OFF }\DataTypeTok{\textbackslash{}}
        \AttributeTok{{-}DCMAKE\_BUILD\_TYPE}\OperatorTok{=}\NormalTok{Release }\DataTypeTok{\textbackslash{}}
        \AttributeTok{{-}DCMAKE\_INSTALL\_PREFIX}\OperatorTok{=}\NormalTok{/opt/gromacs }\DataTypeTok{\textbackslash{}}
    \KeywordTok{\&\&} \FunctionTok{cmake} \AttributeTok{{-}{-}build}\NormalTok{ gromacs{-}build }\AttributeTok{{-}{-}parallel} \DataTypeTok{\textbackslash{}}
    \KeywordTok{\&\&} \ExtensionTok{ctest} \AttributeTok{{-}{-}test{-}dir}\NormalTok{ gromacs{-}build }\AttributeTok{{-}{-}output{-}on{-}failure} \DataTypeTok{\textbackslash{}}
    \KeywordTok{\&\&} \FunctionTok{cmake} \AttributeTok{{-}{-}install}\NormalTok{ gromacs{-}build}

\KeywordTok{FROM}\NormalTok{ ubuntu:24.04 }\KeywordTok{AS}\NormalTok{ runtime}

\KeywordTok{RUN} \ExtensionTok{apt{-}get}\NormalTok{ update }\KeywordTok{\&\&} \ExtensionTok{apt{-}get}\NormalTok{ install }\AttributeTok{{-}y} \AttributeTok{{-}{-}no{-}install{-}recommends} \DataTypeTok{\textbackslash{}}
\NormalTok{    ca{-}certificates }\DataTypeTok{\textbackslash{}}
\NormalTok{    libgomp1 }\DataTypeTok{\textbackslash{}}
    \KeywordTok{\&\&} \FunctionTok{rm} \AttributeTok{{-}rf}\NormalTok{ /var/lib/apt/lists/}\PreprocessorTok{*}

\KeywordTok{COPY} \OperatorTok{{-}{-}from=builder}\NormalTok{ /opt/gromacs /opt/gromacs}

\KeywordTok{ENV}\NormalTok{ PATH=}\StringTok{"/opt/gromacs/bin:$\{PATH\}"}
\KeywordTok{ENV}\NormalTok{ LD\_LIBRARY\_PATH=}\StringTok{"/opt/gromacs/lib:$\{LD\_LIBRARY\_PATH\}"}

\KeywordTok{WORKDIR}\NormalTok{ /work}

\KeywordTok{ENTRYPOINT}\NormalTok{ [}\StringTok{"gmx"}\NormalTok{]}
\KeywordTok{CMD}\NormalTok{ [}\StringTok{"{-}{-}version"}\NormalTok{]}
\end{Highlighting}
\end{Shaded}

Construya la imagen:

\begin{Shaded}
\begin{Highlighting}[]
\ExtensionTok{docker}\NormalTok{ build }\AttributeTok{{-}{-}pull} \AttributeTok{{-}{-}no{-}cache} \DataTypeTok{\textbackslash{}}
  \AttributeTok{{-}f}\NormalTok{ Dockerfile.cpu }\DataTypeTok{\textbackslash{}}
  \AttributeTok{{-}t}\NormalTok{ gromacs:2026.3{-}cpu .}
\end{Highlighting}
\end{Shaded}

El valor \textbf{SSE2} se eligió como mínimo común denominador razonable
para x86-64. Aumenta la posibilidad de ejecutar la misma imagen en
procesadores x86-64 distintos, pero reduce el rendimiento respecto de
AVX2 o AVX-512. Para una imagen destinada a un único nodo o a máquinas
homogéneas puede compilarse otra variante:

\begin{Shaded}
\begin{Highlighting}[]
\ExtensionTok{docker}\NormalTok{ build }\AttributeTok{{-}{-}pull} \DataTypeTok{\textbackslash{}}
  \AttributeTok{{-}{-}build{-}arg}\NormalTok{ GMX\_SIMD=AVX2\_256 }\DataTypeTok{\textbackslash{}}
  \AttributeTok{{-}f}\NormalTok{ Dockerfile.cpu }\DataTypeTok{\textbackslash{}}
  \AttributeTok{{-}t}\NormalTok{ gromacs:2026.3{-}cpu{-}avx2 .}
\end{Highlighting}
\end{Shaded}

Esa imagen fallará o no será apropiada en CPU sin AVX2. Una imagen
construida para \textbf{linux/amd64} tampoco se vuelve compatible
automáticamente con ARM64. En ARM debe realizarse una compilación nativa
con el SIMD correspondiente; la emulación mediante QEMU sirve para
construir o probar, pero no para medir rendimiento de dinámica
molecular.

\subsubsection{Imagen con CUDA para GPU
NVIDIA}\label{imagen-con-cuda-para-gpu-nvidia}

GROMACS 2026.3 requiere CUDA 12.1 o posterior y una GPU con capacidad de
cómputo 5.0 o superior. El ejemplo utiliza CUDA 12.6 sobre Ubuntu 24.04,
una combinación incluida entre las plataformas de prueba declaradas para
esta versión.

Cree \textbf{Dockerfile.cuda}:

\begin{Shaded}
\begin{Highlighting}[]
\KeywordTok{FROM}\NormalTok{ nvidia/cuda:12.6.3{-}devel{-}ubuntu24.04 }\KeywordTok{AS}\NormalTok{ builder}

\KeywordTok{ARG}\NormalTok{ GROMACS\_VERSION=2026.3}
\KeywordTok{ARG}\NormalTok{ GROMACS\_MD5=7987af0c6ab939ab6e639f32d0dd260f}
\KeywordTok{ARG}\NormalTok{ CUDA\_ARCHITECTURES=}\StringTok{"52;60;61;70;75;80;86;89;90"}

\KeywordTok{RUN} \ExtensionTok{apt{-}get}\NormalTok{ update }\KeywordTok{\&\&} \ExtensionTok{apt{-}get}\NormalTok{ install }\AttributeTok{{-}y} \AttributeTok{{-}{-}no{-}install{-}recommends} \DataTypeTok{\textbackslash{}}
\NormalTok{    build{-}essential }\DataTypeTok{\textbackslash{}}
\NormalTok{    ca{-}certificates }\DataTypeTok{\textbackslash{}}
\NormalTok{    cmake }\DataTypeTok{\textbackslash{}}
\NormalTok{    curl }\DataTypeTok{\textbackslash{}}
    \KeywordTok{\&\&} \FunctionTok{rm} \AttributeTok{{-}rf}\NormalTok{ /var/lib/apt/lists/}\PreprocessorTok{*}

\KeywordTok{WORKDIR}\NormalTok{ /tmp/build}

\KeywordTok{RUN} \ExtensionTok{curl} \AttributeTok{{-}fL} \DataTypeTok{\textbackslash{}}
    \StringTok{"https://ftp.gromacs.org/gromacs/gromacs{-}}\VariableTok{$\{GROMACS\_VERSION\}}\StringTok{.tar.gz"} \DataTypeTok{\textbackslash{}}
    \AttributeTok{{-}o}\NormalTok{ gromacs.tar.gz }\DataTypeTok{\textbackslash{}}
    \KeywordTok{\&\&} \BuiltInTok{echo} \StringTok{"}\VariableTok{$\{GROMACS\_MD5\}}\StringTok{  gromacs.tar.gz"} \KeywordTok{|} \FunctionTok{md5sum} \AttributeTok{{-}c} \AttributeTok{{-}} \DataTypeTok{\textbackslash{}}
    \KeywordTok{\&\&} \FunctionTok{tar}\NormalTok{ xzf gromacs.tar.gz }\DataTypeTok{\textbackslash{}}
    \KeywordTok{\&\&} \FunctionTok{cmake} \AttributeTok{{-}S} \StringTok{"gromacs{-}}\VariableTok{$\{GROMACS\_VERSION\}}\StringTok{"} \AttributeTok{{-}B}\NormalTok{ gromacs{-}build }\DataTypeTok{\textbackslash{}}
        \AttributeTok{{-}DGMX\_BUILD\_OWN\_FFTW}\OperatorTok{=}\NormalTok{ON }\DataTypeTok{\textbackslash{}}
        \AttributeTok{{-}DREGRESSIONTEST\_DOWNLOAD}\OperatorTok{=}\NormalTok{ON }\DataTypeTok{\textbackslash{}}
        \AttributeTok{{-}DGMX\_GPU}\OperatorTok{=}\NormalTok{CUDA }\DataTypeTok{\textbackslash{}}
        \AttributeTok{{-}DCUDAToolkit\_ROOT}\OperatorTok{=}\NormalTok{/usr/local/cuda }\DataTypeTok{\textbackslash{}}
        \AttributeTok{{-}DCMAKE\_CUDA\_ARCHITECTURES}\OperatorTok{=}\StringTok{"}\VariableTok{$\{CUDA\_ARCHITECTURES\}}\StringTok{"} \DataTypeTok{\textbackslash{}}
        \AttributeTok{{-}DGMX\_MPI}\OperatorTok{=}\NormalTok{OFF }\DataTypeTok{\textbackslash{}}
        \AttributeTok{{-}DCMAKE\_BUILD\_TYPE}\OperatorTok{=}\NormalTok{Release }\DataTypeTok{\textbackslash{}}
        \AttributeTok{{-}DCMAKE\_INSTALL\_PREFIX}\OperatorTok{=}\NormalTok{/opt/gromacs }\DataTypeTok{\textbackslash{}}
    \KeywordTok{\&\&} \FunctionTok{cmake} \AttributeTok{{-}{-}build}\NormalTok{ gromacs{-}build }\AttributeTok{{-}{-}parallel} \DataTypeTok{\textbackslash{}}
    \KeywordTok{\&\&} \ExtensionTok{ctest} \AttributeTok{{-}{-}test{-}dir}\NormalTok{ gromacs{-}build }\AttributeTok{{-}{-}output{-}on{-}failure} \DataTypeTok{\textbackslash{}}
    \KeywordTok{\&\&} \FunctionTok{cmake} \AttributeTok{{-}{-}install}\NormalTok{ gromacs{-}build}

\KeywordTok{FROM}\NormalTok{ nvidia/cuda:12.6.3{-}runtime{-}ubuntu24.04 }\KeywordTok{AS}\NormalTok{ runtime}

\KeywordTok{RUN} \ExtensionTok{apt{-}get}\NormalTok{ update }\KeywordTok{\&\&} \ExtensionTok{apt{-}get}\NormalTok{ install }\AttributeTok{{-}y} \AttributeTok{{-}{-}no{-}install{-}recommends} \DataTypeTok{\textbackslash{}}
\NormalTok{    ca{-}certificates }\DataTypeTok{\textbackslash{}}
\NormalTok{    libgomp1 }\DataTypeTok{\textbackslash{}}
    \KeywordTok{\&\&} \FunctionTok{rm} \AttributeTok{{-}rf}\NormalTok{ /var/lib/apt/lists/}\PreprocessorTok{*}

\KeywordTok{COPY} \OperatorTok{{-}{-}from=builder}\NormalTok{ /opt/gromacs /opt/gromacs}

\KeywordTok{ENV}\NormalTok{ PATH=}\StringTok{"/opt/gromacs/bin:$\{PATH\}"}
\KeywordTok{ENV}\NormalTok{ LD\_LIBRARY\_PATH=}\StringTok{"/opt/gromacs/lib:$\{LD\_LIBRARY\_PATH\}"}

\KeywordTok{WORKDIR}\NormalTok{ /work}

\KeywordTok{ENTRYPOINT}\NormalTok{ [}\StringTok{"gmx"}\NormalTok{]}
\KeywordTok{CMD}\NormalTok{ [}\StringTok{"{-}{-}version"}\NormalTok{]}
\end{Highlighting}
\end{Shaded}

Construya y verifique:

\begin{Shaded}
\begin{Highlighting}[]
\ExtensionTok{docker}\NormalTok{ build }\AttributeTok{{-}{-}pull} \AttributeTok{{-}{-}no{-}cache} \DataTypeTok{\textbackslash{}}
  \AttributeTok{{-}f}\NormalTok{ Dockerfile.cuda }\DataTypeTok{\textbackslash{}}
  \AttributeTok{{-}t}\NormalTok{ gromacs:2026.3{-}cuda12.6 .}

\ExtensionTok{docker}\NormalTok{ run }\AttributeTok{{-}{-}rm} \AttributeTok{{-}{-}gpus}\NormalTok{ all }\DataTypeTok{\textbackslash{}}
\NormalTok{  gromacs:2026.3{-}cuda12.6 }\AttributeTok{{-}{-}version}

\ExtensionTok{docker}\NormalTok{ run }\AttributeTok{{-}{-}rm} \AttributeTok{{-}{-}gpus}\NormalTok{ all }\DataTypeTok{\textbackslash{}}
\NormalTok{  gromacs:2026.3{-}cuda12.6 mdrun }\AttributeTok{{-}version}
\end{Highlighting}
\end{Shaded}

La lista de \textbf{CMAKE\_CUDA\_ARCHITECTURES} genera código para
varias generaciones y aumenta el tiempo de compilación y el tamaño de la
imagen. Puede reducirse para un parque homogéneo. Por ejemplo, una RTX
3060 utiliza SM 86 y una Tesla P100 utiliza SM 60:

\begin{Shaded}
\begin{Highlighting}[]
\ExtensionTok{docker}\NormalTok{ build }\AttributeTok{{-}{-}pull} \DataTypeTok{\textbackslash{}}
  \AttributeTok{{-}{-}build{-}arg}\NormalTok{ CUDA\_ARCHITECTURES=}\StringTok{"60;86"} \DataTypeTok{\textbackslash{}}
  \AttributeTok{{-}f}\NormalTok{ Dockerfile.cuda }\DataTypeTok{\textbackslash{}}
  \AttributeTok{{-}t}\NormalTok{ gromacs:2026.3{-}cuda12.6{-}sm60{-}sm86 .}
\end{Highlighting}
\end{Shaded}

No agregue una arquitectura que el toolkit seleccionado ya no admita. Si
se busca máxima portabilidad entre GPU NVIDIA, es preferible conservar
la lista predeterminada de arquitecturas generada por GROMACS o definir
explícitamente todas las GPU reales que deberán ejecutar la imagen.

\subsubsection{Compatibilidad entre la imagen CUDA y el
controlador}\label{compatibilidad-entre-la-imagen-cuda-y-el-controlador}

La versión mostrada por \textbf{nvidia-smi} como ``CUDA Version'' es la
versión máxima admitida por el controlador, no el toolkit contenido en
la imagen. Para ejecutar una imagen basada en CUDA 12.x, el controlador
Linux debe cumplir como mínimo el requisito de esa familia; la tabla
general de compatibilidad menor de NVIDIA indica controlador 525 o
posterior para CUDA 12.x. Algunas características que combinan PTX,
bibliotecas nuevas o hardware reciente pueden exigir un controlador más
nuevo.

Verifique ambos lados:

\begin{Shaded}
\begin{Highlighting}[]
\ExtensionTok{nvidia{-}smi}

\ExtensionTok{docker}\NormalTok{ run }\AttributeTok{{-}{-}rm} \AttributeTok{{-}{-}gpus}\NormalTok{ all }\DataTypeTok{\textbackslash{}}
\NormalTok{  gromacs:2026.3{-}cuda12.6 mdrun }\AttributeTok{{-}version}
\end{Highlighting}
\end{Shaded}

La primera orden caracteriza el anfitrión. La segunda confirma cómo fue
compilado GROMACS y si el contenedor ve la GPU.

\subsubsection{Ejecutar una simulación conservando los
archivos}\label{ejecutar-una-simulaciuxf3n-conservando-los-archivos}

Los datos no deben quedar únicamente dentro de la capa efímera del
contenedor. Monte el directorio actual en \textbf{/work} y use el UID y
GID del usuario para evitar archivos propiedad de root:

\begin{Shaded}
\begin{Highlighting}[]
\ExtensionTok{docker}\NormalTok{ run }\AttributeTok{{-}{-}rm} \AttributeTok{{-}it} \DataTypeTok{\textbackslash{}}
  \AttributeTok{{-}{-}user} \StringTok{"}\VariableTok{$(}\FunctionTok{id} \AttributeTok{{-}u}\VariableTok{)}\StringTok{:}\VariableTok{$(}\FunctionTok{id} \AttributeTok{{-}g}\VariableTok{)}\StringTok{"} \DataTypeTok{\textbackslash{}}
  \AttributeTok{{-}{-}volume} \StringTok{"}\VariableTok{$PWD}\StringTok{:/work"} \DataTypeTok{\textbackslash{}}
  \AttributeTok{{-}{-}workdir}\NormalTok{ /work }\DataTypeTok{\textbackslash{}}
\NormalTok{  gromacs:2026.3{-}cpu }\DataTypeTok{\textbackslash{}}
\NormalTok{  grompp }\AttributeTok{{-}f}\NormalTok{ md.mdp }\AttributeTok{{-}c}\NormalTok{ npt.gro }\AttributeTok{{-}t}\NormalTok{ npt.cpt }\DataTypeTok{\textbackslash{}}
  \AttributeTok{{-}p}\NormalTok{ topol.top }\AttributeTok{{-}o}\NormalTok{ md.tpr}
\end{Highlighting}
\end{Shaded}

Para producción con NVIDIA:

\begin{Shaded}
\begin{Highlighting}[]
\ExtensionTok{docker}\NormalTok{ run }\AttributeTok{{-}{-}rm} \AttributeTok{{-}it} \DataTypeTok{\textbackslash{}}
  \AttributeTok{{-}{-}gpus}\NormalTok{ all }\DataTypeTok{\textbackslash{}}
  \AttributeTok{{-}{-}user} \StringTok{"}\VariableTok{$(}\FunctionTok{id} \AttributeTok{{-}u}\VariableTok{)}\StringTok{:}\VariableTok{$(}\FunctionTok{id} \AttributeTok{{-}g}\VariableTok{)}\StringTok{"} \DataTypeTok{\textbackslash{}}
  \AttributeTok{{-}{-}volume} \StringTok{"}\VariableTok{$PWD}\StringTok{:/work"} \DataTypeTok{\textbackslash{}}
  \AttributeTok{{-}{-}workdir}\NormalTok{ /work }\DataTypeTok{\textbackslash{}}
\NormalTok{  gromacs:2026.3{-}cuda12.6 }\DataTypeTok{\textbackslash{}}
\NormalTok{  mdrun }\AttributeTok{{-}deffnm}\NormalTok{ md }\AttributeTok{{-}ntmpi}\NormalTok{ 1 }\AttributeTok{{-}ntomp}\NormalTok{ 8}
\end{Highlighting}
\end{Shaded}

Como el Dockerfile define \textbf{ENTRYPOINT {[}``gmx''{]}}, después del
nombre de la imagen se escribe directamente la suborden, por ejemplo
\textbf{grompp}, \textbf{mdrun} o \textbf{rms}. El directorio montado
conserva TPR, trayectorias, energías, logs y checkpoints cuando se
elimina el contenedor.

Si Docker tiene un límite de CPU o memoria, GROMACS sólo podrá usar los
recursos asignados. Conviene declararlos de forma explícita cuando se
comparan rendimientos:

\begin{Shaded}
\begin{Highlighting}[]
\ExtensionTok{docker}\NormalTok{ run }\AttributeTok{{-}{-}rm} \DataTypeTok{\textbackslash{}}
  \AttributeTok{{-}{-}gpus}\NormalTok{ all }\DataTypeTok{\textbackslash{}}
  \AttributeTok{{-}{-}cpuset{-}cpus}\NormalTok{ 0{-}7 }\DataTypeTok{\textbackslash{}}
  \AttributeTok{{-}{-}memory}\NormalTok{ 24g }\DataTypeTok{\textbackslash{}}
  \AttributeTok{{-}{-}user} \StringTok{"}\VariableTok{$(}\FunctionTok{id} \AttributeTok{{-}u}\VariableTok{)}\StringTok{:}\VariableTok{$(}\FunctionTok{id} \AttributeTok{{-}g}\VariableTok{)}\StringTok{"} \DataTypeTok{\textbackslash{}}
  \AttributeTok{{-}{-}volume} \StringTok{"}\VariableTok{$PWD}\StringTok{:/work"} \DataTypeTok{\textbackslash{}}
  \AttributeTok{{-}{-}workdir}\NormalTok{ /work }\DataTypeTok{\textbackslash{}}
\NormalTok{  gromacs:2026.3{-}cuda12.6 }\DataTypeTok{\textbackslash{}}
\NormalTok{  mdrun }\AttributeTok{{-}deffnm}\NormalTok{ md }\AttributeTok{{-}ntmpi}\NormalTok{ 1 }\AttributeTok{{-}ntomp}\NormalTok{ 8}
\end{Highlighting}
\end{Shaded}

\subsubsection{Validación mínima de la
imagen}\label{validaciuxf3n-muxednima-de-la-imagen}

La compilación ejecuta \textbf{ctest}, pero la imagen GPU se construye
normalmente sin acceso a un dispositivo. Debe realizarse una prueba de
ejecución en cada clase de hardware de destino.

Registre la configuración:

\begin{Shaded}
\begin{Highlighting}[]
\ExtensionTok{docker}\NormalTok{ image inspect gromacs:2026.3{-}cuda12.6 }\OperatorTok{\textgreater{}}\NormalTok{ image{-}inspect.json}

\ExtensionTok{docker}\NormalTok{ run }\AttributeTok{{-}{-}rm} \AttributeTok{{-}{-}gpus}\NormalTok{ all }\DataTypeTok{\textbackslash{}}
\NormalTok{  gromacs:2026.3{-}cuda12.6 mdrun }\AttributeTok{{-}version}
\end{Highlighting}
\end{Shaded}

Después ejecute un sistema pequeño y examine el log:

\begin{Shaded}
\begin{Highlighting}[]
\ExtensionTok{docker}\NormalTok{ run }\AttributeTok{{-}{-}rm} \AttributeTok{{-}{-}gpus}\NormalTok{ all }\DataTypeTok{\textbackslash{}}
  \AttributeTok{{-}{-}user} \StringTok{"}\VariableTok{$(}\FunctionTok{id} \AttributeTok{{-}u}\VariableTok{)}\StringTok{:}\VariableTok{$(}\FunctionTok{id} \AttributeTok{{-}g}\VariableTok{)}\StringTok{"} \DataTypeTok{\textbackslash{}}
  \AttributeTok{{-}{-}volume} \StringTok{"}\VariableTok{$PWD}\StringTok{:/work"} \DataTypeTok{\textbackslash{}}
  \AttributeTok{{-}{-}workdir}\NormalTok{ /work }\DataTypeTok{\textbackslash{}}
\NormalTok{  gromacs:2026.3{-}cuda12.6 }\DataTypeTok{\textbackslash{}}
\NormalTok{  mdrun }\AttributeTok{{-}s}\NormalTok{ test.tpr }\AttributeTok{{-}deffnm}\NormalTok{ test }\AttributeTok{{-}nsteps}\NormalTok{ 1000}
\end{Highlighting}
\end{Shaded}

Compruebe en \textbf{test.log}:

\begin{itemize}
\tightlist
\item
  versión y precisión de GROMACS;
\item
  SIMD detectado;
\item
  backend CUDA;
\item
  GPU seleccionada;
\item
  número de rangos y de hilos;
\item
  ausencia de errores LINCS, NaN o fallos del dispositivo;
\item
  rendimiento coherente con una ejecución nativa equivalente.
\end{itemize}

Una diferencia numérica pequeña entre hardware, número de hilos o
backends no implica por sí sola un error: la dinámica molecular es
caótica y las reducciones en coma flotante no son asociativas. La
validación debe comparar conservación, distribuciones y propiedades
estadísticas, no exigir trayectorias idénticas marco a marco.

\subsubsection{Reproducibilidad de la
imagen}\label{reproducibilidad-de-la-imagen}

Una etiqueta como \textbf{ubuntu:24.04} o
\textbf{nvidia/cuda:12.6.3-runtime-ubuntu24.04} puede apuntar
posteriormente a una imagen base reconstruida. Para congelar una imagen
publicada se debe registrar y usar su digest:

\begin{Shaded}
\begin{Highlighting}[]
\KeywordTok{FROM}\NormalTok{ ubuntu:24.04@sha256:DIGEST\_VERIFICADO }\KeywordTok{AS}\NormalTok{ builder}
\end{Highlighting}
\end{Shaded}

El digest se obtiene del registro utilizado y debe conservarse junto
con:

\begin{itemize}
\tightlist
\item
  Dockerfile;
\item
  versión y suma de comprobación de GROMACS;
\item
  digest de cada imagen base;
\item
  salida de \textbf{docker version};
\item
  salida de \textbf{gmx mdrun -version};
\item
  controlador y modelo de GPU;
\item
  comando exacto de construcción y ejecución.
\end{itemize}

\textbf{--no-cache} fuerza una reconstrucción limpia, pero no garantiza
reproducibilidad si los repositorios APT cambiaron. Para
reconstrucciones archivables se necesitan además repositorios con
instantáneas o una imagen ya construida identificada por digest.

\subsubsection{MPI, clústeres y límites
prácticos}\label{mpi-cluxfasteres-y-luxedmites-pruxe1cticos}

La imagen propuesta usa thread-MPI, apropiado para una estación de
trabajo o un solo nodo. Construir con \textbf{GMX\_MPI=ON} dentro del
contenedor es posible, pero ejecutar eficientemente entre nodos requiere
compatibilidad con el MPI, la red de alta velocidad, UCX/OFED y el
lanzador del clúster. Encapsular una biblioteca MPI arbitraria puede
anular RDMA o generar incompatibilidades con SLURM.

En HPC suele ser más simple construir una imagen OCI validada y
ejecutarla mediante Apptainer/Singularity, o compilar GROMACS contra la
pila MPI suministrada por el centro. Para multinodo, la portabilidad del
contenedor debe validarse con el administrador y con una prueba de
escalamiento; que \textbf{mpirun} funcione dentro de una computadora no
demuestra compatibilidad multinodo.

\subsubsection{Fallos frecuentes}\label{fallos-frecuentes}

\begin{itemize}
\tightlist
\item
  Instalar el controlador NVIDIA dentro de la imagen.
\item
  Confundir el toolkit CUDA de la imagen con el controlador del
  anfitrión.
\item
  Compilar con AVX2 o AVX-512 y asumir que la imagen funcionará en
  cualquier CPU.
\item
  Usar una etiqueta \textbf{latest} para GROMACS, Ubuntu o CUDA.
\item
  No montar el directorio de trabajo y perder los resultados al eliminar
  el contenedor.
\item
  Ejecutar como root y dejar archivos sin permisos para el usuario.
\item
  Omitir \textbf{--gpus all} y concluir que GROMACS fue compilado sin
  CUDA.
\item
  Probar sólo \textbf{gmx --version} sin ejecutar un sistema pequeño.
\item
  Suponer que Docker elimina las diferencias numéricas entre GPU y CPU.
\item
  Copiar una imagen CUDA a un nodo cuyo controlador es demasiado
  antiguo.
\end{itemize}

\subsection{Instalación de gmxapi para
Python}\label{instalaciuxf3n-de-gmxapi-para-python}

gmxapi requiere una instalación previa de GROMACS compilada con
\textbf{GMXAPI=ON} y \textbf{BUILD\_SHARED\_LIBS=ON}. Ambas opciones
suelen estar activadas por defecto, pero conviene declararlas cuando se
prepara una instalación destinada a Python:

\begin{Shaded}
\begin{Highlighting}[]
\FunctionTok{cmake}\NormalTok{ ..     }\AttributeTok{{-}DGMX\_BUILD\_OWN\_FFTW}\OperatorTok{=}\NormalTok{ON     }\AttributeTok{{-}DREGRESSIONTEST\_DOWNLOAD}\OperatorTok{=}\NormalTok{ON     }\AttributeTok{{-}DGMXAPI}\OperatorTok{=}\NormalTok{ON     }\AttributeTok{{-}DBUILD\_SHARED\_LIBS}\OperatorTok{=}\NormalTok{ON     }\AttributeTok{{-}DCMAKE\_INSTALL\_PREFIX}\OperatorTok{=}\StringTok{"}\VariableTok{$HOME}\StringTok{/opt/gromacs{-}2026.3"}
\end{Highlighting}
\end{Shaded}

Después de instalar GROMACS, cree un entorno virtual. gmxapi admite
Python 3.9 o posterior.

\begin{Shaded}
\begin{Highlighting}[]
\ExtensionTok{python3} \AttributeTok{{-}m}\NormalTok{ venv }\StringTok{"}\VariableTok{$HOME}\StringTok{/venvs/gmxapi{-}2026"}
\BuiltInTok{source} \StringTok{"}\VariableTok{$HOME}\StringTok{/venvs/gmxapi{-}2026/bin/activate"}

\ExtensionTok{python} \AttributeTok{{-}m}\NormalTok{ pip install }\AttributeTok{{-}{-}upgrade}\NormalTok{ pip setuptools wheel cmake pybind11}
\BuiltInTok{source} \StringTok{"}\VariableTok{$HOME}\StringTok{/opt/gromacs{-}2026.3/bin/GMXRC"}
\ExtensionTok{python} \AttributeTok{{-}m}\NormalTok{ pip install }\AttributeTok{{-}{-}no{-}cache{-}dir}\NormalTok{ gmxapi}
\end{Highlighting}
\end{Shaded}

Si el instalador no encuentra GROMACS:

\begin{Shaded}
\begin{Highlighting}[]
\VariableTok{gmxapi\_ROOT}\OperatorTok{=}\StringTok{"}\VariableTok{$HOME}\StringTok{/opt/gromacs{-}2026.3"} \ExtensionTok{python} \AttributeTok{{-}m}\NormalTok{ pip install }\AttributeTok{{-}{-}no{-}cache{-}dir}\NormalTok{ gmxapi}
\end{Highlighting}
\end{Shaded}

Comprobación:

\begin{Shaded}
\begin{Highlighting}[]
\ExtensionTok{python} \AttributeTok{{-}c} \StringTok{"import gmxapi; print(gmxapi.\_\_version\_\_)"}
\end{Highlighting}
\end{Shaded}

Después de actualizar o recompilar GROMACS con otro compilador, MPI o
precisión, reinstale gmxapi sin usar la caché. Una rueda compilada
contra otra instalación puede importar incorrectamente o fallar por
símbolos incompatibles.

\subsection{Verificación final}\label{verificaciuxf3n-final}

\begin{Shaded}
\begin{Highlighting}[]
\FunctionTok{which}\NormalTok{ gmx}
\ExtensionTok{gmx} \AttributeTok{{-}{-}version}
\ExtensionTok{gmx}\NormalTok{ mdrun }\AttributeTok{{-}version}
\end{Highlighting}
\end{Shaded}

Compruebe al menos:

\begin{itemize}
\tightlist
\item
  versión exacta de GROMACS;
\item
  compilador y precisión;
\item
  SIMD;
\item
  biblioteca FFT;
\item
  soporte MPI o thread-MPI;
\item
  backend GPU;
\item
  ruta del archivo de datos.
\end{itemize}

Una prueba mínima del ejecutable no reemplaza los tests:

\begin{Shaded}
\begin{Highlighting}[]
\ExtensionTok{gmx}\NormalTok{ help}
\ExtensionTok{gmx}\NormalTok{ check }\AttributeTok{{-}h}
\end{Highlighting}
\end{Shaded}

Para diagnosticar qué ejecutable se está usando cuando existen varias
instalaciones:

\begin{Shaded}
\begin{Highlighting}[]
\BuiltInTok{type} \AttributeTok{{-}a}\NormalTok{ gmx}
\BuiltInTok{echo} \StringTok{"}\VariableTok{$PATH}\StringTok{"}
\end{Highlighting}
\end{Shaded}

Si se activa otro entorno, módulo o instalación, ejecute nuevamente el
\textbf{GMXRC} correspondiente antes de continuar.

\subsection{Fuentes}\label{fuentes}

\begin{enumerate}
\def\labelenumi{\arabic{enumi}.}
\tightlist
\item
  \href{https://manual.gromacs.org/current/install-guide/index.html}{Guía
  de instalación de GROMACS 2026.3}
\item
  \href{https://manual.gromacs.org/current/gmxapi/userguide/install.html}{Instalación
  de gmxapi}
\item
  \href{https://docs.nvidia.com/deploy/cuda-compatibility/minor-version-compatibility.html}{Compatibilidad
  entre CUDA Toolkit y controladores NVIDIA}
\item
  \href{https://docs.nvidia.com/cuda/cuda-installation-guide-linux/}{Guía
  de instalación de CUDA para Linux}
\item
  \href{https://manual.gromacs.org/documentation/2026.3/download.html}{Descarga
  y suma de comprobación de GROMACS 2026.3}
\item
  \href{https://docs.nvidia.com/datacenter/cloud-native/container-toolkit/latest/install-guide.html}{Instalación
  de NVIDIA Container Toolkit}
\item
  \href{https://docs.nvidia.com/datacenter/cloud-native/container-toolkit/latest/sample-workload.html}{Prueba
  de acceso a GPU desde un contenedor}
\item
  \href{https://docs.docker.com/build/building/best-practices/}{Buenas
  prácticas para construir imágenes Docker}
\end{enumerate}

\section{Caso: 1 proteína - 1
ligando}\label{caso-1-proteuxedna---1-ligando}

Este tutorial describe la preparación, simulación y análisis de un
complejo proteína--ligando con \textbf{GROMACS 2026.3}. Se supone que el
ligando ya está colocado en la pose inicial que se desea estudiar.

Los nombres de archivos son ejemplos. Conviene mantener nombres
explícitos y una carpeta por etapa:

\begin{Shaded}
\begin{Highlighting}[]
\FunctionTok{mkdir} \AttributeTok{{-}p}\NormalTok{ 00\_entrada 01\_preparacion 02\_em 03\_nvt 04\_npt 05\_md 06\_analisis}
\end{Highlighting}
\end{Shaded}

\subsection{Objetivo y alcance del
tutorial}\label{objetivo-y-alcance-del-tutorial}

El sistema estudiado contiene una proteína receptora, un ligando,
solvente explícito e iones. El objetivo es construir un modelo
atomístico consistente, eliminar contactos desfavorables, equilibrarlo y
generar un conjunto de conformaciones que permita describir la
estabilidad estructural y las interacciones del complejo.

Este procedimiento puede responder preguntas como:

\begin{itemize}
\tightlist
\item
  si una pose permanece estable durante el intervalo simulado;
\item
  qué residuos mantienen contacto con el ligando;
\item
  qué puentes de hidrógeno, contactos hidrofóbicos o interacciones
  iónicas son persistentes;
\item
  qué regiones de la proteína cambian su movilidad;
\item
  si aparecen reorganizaciones del sitio de unión o entrada de agua.
\end{itemize}

Una trayectoria aislada no demuestra afinidad, mecanismo de acción ni
convergencia termodinámica. Tampoco convierte una pose de docking en una
estructura experimental. La afinidad requiere métodos y diseños
específicos ---por ejemplo FEP, LIE, MM/PBSA o potencial de fuerza
media--- y réplicas suficientes para estimar incertidumbre.

\subsection{Conceptos que deben
distinguirse}\label{conceptos-que-deben-distinguirse}

{\def\LTcaptype{none} % do not increment counter
\begin{longtable}[]{@{}
  >{\raggedright\arraybackslash}p{(\linewidth - 2\tabcolsep) * \real{0.0931}}
  >{\raggedright\arraybackslash}p{(\linewidth - 2\tabcolsep) * \real{0.9069}}@{}}
\toprule\noalign{}
\begin{minipage}[b]{\linewidth}\raggedright
Concepto
\end{minipage} & \begin{minipage}[b]{\linewidth}\raggedright
Definición operativa
\end{minipage} \\
\midrule\noalign{}
\endhead
\bottomrule\noalign{}
\endlastfoot
\textbf{Receptor} & Macromolécula que contiene el sitio de unión. En
este tutorial es una proteína soluble. \\
\textbf{Ligando} & Molécula cuya pose e interacciones con el receptor se
estudian. Puede poseer varios estados de protonación, tautómeros o
estereoisómeros. \\
\textbf{Pose} & Posición, orientación y conformación del ligando dentro
del sitio de unión. \\
\textbf{Microestado} & Combinación concreta de protonación, tautomería
y, cuando corresponda, estereoquímica. Dos microestados del mismo
compuesto son sistemas químicos diferentes para una simulación
clásica. \\
\textbf{Coordenadas} & Posiciones atómicas almacenadas en PDB, GRO u
otro formato. No contienen por sí solas una descripción completa de las
interacciones. \\
\textbf{Topología} & Conectividad, tipos atómicos, cargas, masas y
parámetros que definen cómo se calcula la energía. \\
\textbf{Campo de fuerza} & Conjunto de forma funcional y parámetros
usado para aproximar la energía potencial. \\
\textbf{Protocolo} & Secuencia documentada de preparación, minimización,
equilibración, producción y análisis. \\
\textbf{Réplica} & Simulación independiente, normalmente iniciada con
velocidades aleatorias diferentes y, si corresponde, otra conformación
inicial. \\
\end{longtable}
}

Una estructura químicamente correcta requiere coherencia entre
coordenadas y topología: mismo número y orden de átomos, mismos nombres
de residuos y átomos, misma conectividad, misma carga formal y mismo
microestado.

\subsection{Modelo físico aproximado}\label{modelo-fuxedsico-aproximado}

La dinámica molecular clásica integra las ecuaciones de movimiento de
Newton:

\[
m_i\frac{d^2\mathbf r_i}{dt^2}
=
\mathbf F_i
=
-\nabla_i U(\mathbf r).
\]

La fuerza sobre cada átomo se obtiene del gradiente de una función de
energía potencial. En forma esquemática,

\[
U =
U_{\mathrm{enlaces}}
+U_{\mathrm{ángulos}}
+U_{\mathrm{diedros}}
+U_{\mathrm{impropios}}
+U_{\mathrm{LJ}}
+U_{\mathrm{Coulomb}}
+U_{\mathrm{restricciones}}.
\]

Para un par de átomos, las contribuciones no enlazantes más habituales
son

\[
U_{\mathrm{LJ}}(r)
=
4\varepsilon
\left[
\left(\frac{\sigma}{r}\right)^{12}
-
\left(\frac{\sigma}{r}\right)^6
\right]
\]

y

\[
U_{\mathrm{Coulomb}}(r)
=
\frac{1}{4\pi\varepsilon_0\varepsilon_r}
\frac{q_iq_j}{r}.
\]

El término de Lennard-Jones representa de manera aproximada la repulsión
de corto alcance y la dispersión. El término electrostático depende de
las cargas parciales asignadas por el campo de fuerza. En los campos de
fuerza biomoleculares convencionales estas cargas suelen ser fijas: no
se modela explícitamente la polarización electrónica inducida.

Las ecuaciones anteriores explican por qué una topología no puede
improvisarse a partir de una geometría. Cambiar cargas, tipos atómicos,
reglas de combinación o parámetros torsionales cambia la superficie de
energía potencial y, por tanto, el sistema que se está simulando.

\subsection{Archivos y flujo de
información}\label{archivos-y-flujo-de-informaciuxf3n}

Los archivos cumplen funciones diferentes:

{\def\LTcaptype{none} % do not increment counter
\begin{longtable}[]{@{}
  >{\raggedright\arraybackslash}p{(\linewidth - 2\tabcolsep) * \real{0.0940}}
  >{\raggedright\arraybackslash}p{(\linewidth - 2\tabcolsep) * \real{0.9060}}@{}}
\toprule\noalign{}
\begin{minipage}[b]{\linewidth}\raggedright
Archivo
\end{minipage} & \begin{minipage}[b]{\linewidth}\raggedright
Función
\end{minipage} \\
\midrule\noalign{}
\endhead
\bottomrule\noalign{}
\endlastfoot
\textbf{PDB/GRO} & Coordenadas, nombres de átomos y residuos; GRO
incluye la caja. \\
\textbf{TOP/ITP} & Topología y parámetros del sistema. \\
\textbf{MDP} & Algoritmo, paso temporal, termostato, barostato, cortes y
frecuencia de salida. \\
\textbf{NDX} & Grupos de átomos definidos para acoplamiento,
restricciones o análisis. \\
\textbf{TPR} & Entrada binaria compilada por \textbf{gmx grompp} a
partir de coordenadas, topología, MDP y, si se usa, índice. \\
\textbf{CPT} & Estado completo necesario para continuar una
simulación. \\
\textbf{XTC/TRR} & Trayectoria; XTC contiene coordenadas comprimidas y
TRR puede contener coordenadas, velocidades y fuerzas. \\
\textbf{EDR} & Energías y variables termodinámicas. \\
\textbf{LOG} & Registro de parámetros efectivos, rendimiento,
advertencias y evolución de la corrida. \\
\end{longtable}
}

El flujo mínimo es:

\[
\text{coordenadas}+\text{topología}+\text{MDP}
\xrightarrow{\texttt{gmx grompp}}
\text{TPR}
\xrightarrow{\texttt{gmx mdrun}}
\text{trayectoria}+\text{energías}+\text{checkpoint}.
\]

\textbf{gmx grompp} no es una formalidad: expande directivas del
preprocesador, valida parte de la coherencia del sistema y escribe en el
TPR los parámetros efectivos. Debe conservarse el archivo MDP procesado
generado con \textbf{-po} cuando se desea auditar exactamente qué se
ejecutó.

\subsection{Decisiones previas que no debe tomar
GROMACS}\label{decisiones-previas-que-no-debe-tomar-gromacs}

Antes de ejecutar \textbf{pdb2gmx} deben quedar documentadas:

\begin{enumerate}
\def\labelenumi{\arabic{enumi}.}
\tightlist
\item
  identidad biológica de la estructura, cadenas y ensamblaje
  oligomérico;
\item
  residuos o átomos faltantes y método usado para reconstruirlos;
\item
  conformaciones alternativas retenidas;
\item
  estados de protonación de residuos titulables;
\item
  presencia de enlaces disulfuro, metales, cofactores y aguas
  estructurales;
\item
  microestado, carga formal y estereoquímica del ligando;
\item
  origen de la pose inicial;
\item
  campo de fuerza y modelo de agua;
\item
  condiciones experimentales que se intentan representar: temperatura,
  presión, pH y composición salina.
\end{enumerate}

GROMACS puede detectar inconsistencias sintácticas, pero no puede
decidir si el tautómero, la pose, el ensamblaje biológico o la química
de coordinación elegidos son correctos.

\subsection{Etapas y criterios para
avanzar}\label{etapas-y-criterios-para-avanzar}

{\def\LTcaptype{none} % do not increment counter
\begin{longtable}[]{@{}
  >{\raggedright\arraybackslash}p{(\linewidth - 4\tabcolsep) * \real{0.0952}}
  >{\raggedright\arraybackslash}p{(\linewidth - 4\tabcolsep) * \real{0.3651}}
  >{\raggedright\arraybackslash}p{(\linewidth - 4\tabcolsep) * \real{0.5397}}@{}}
\toprule\noalign{}
\begin{minipage}[b]{\linewidth}\raggedright
Etapa
\end{minipage} & \begin{minipage}[b]{\linewidth}\raggedright
Finalidad
\end{minipage} & \begin{minipage}[b]{\linewidth}\raggedright
Criterio mínimo antes de continuar
\end{minipage} \\
\midrule\noalign{}
\endhead
\bottomrule\noalign{}
\endlastfoot
Preparación & Definir composición, microestados y parámetros &
Coordenadas y topología coherentes; carga y conectividad verificadas \\
Minimización & Eliminar contactos y fuerzas excesivas & Sin NaN; fuerza
máxima y geometría aceptables \\
NVT & Ajustar temperatura a volumen fijo & Temperatura estable y
estructura sin deformaciones \\
NPT & Ajustar densidad y volumen & Densidad y volumen estacionarios;
presión media compatible \\
Producción & Muestrear el conjunto definido & Sin inestabilidad;
duración y réplicas justificadas \\
Análisis & Responder preguntas predefinidas & PBC, ajuste, selección e
intervalo documentados \\
\end{longtable}
}

Una etapa no se valida porque el comando terminó sin error. La salida
estructural, el log y las variables relevantes deben inspeccionarse
antes de usarla como entrada de la etapa siguiente.

\subsection{Requisitos básicos}\label{requisitos-buxe1sicos}

Archivos iniciales:

\begin{itemize}
\tightlist
\item
  \textbf{protein.pdb:} estructura de la proteína.
\item
  \textbf{ligand.pdb} o \textbf{ligand.mol2:} ligando con geometría,
  estereoquímica, protonación y carga formal revisadas.
\item
  \textbf{ligand.itp:} topología del ligando compatible con el campo de
  fuerza de la proteína.
\item
  \textbf{ions.mdp}, \textbf{em.mdp}, \textbf{nvt.mdp}, \textbf{npt.mdp}
  y \textbf{md.mdp:} parámetros de cada etapa.
\end{itemize}

Antes de procesar la proteína deben revisarse residuos faltantes,
conformaciones alternativas, enlaces disulfuro, metales, cofactores,
aguas estructurales y estados de protonación. Eliminar automáticamente
todos los heteroátomos puede destruir información necesaria.

El ligando debe conservar la pose obtenida por cristalografía, docking u
otro procedimiento. Un archivo MOL2 puede contener tipos atómicos y
cargas, pero eso no demuestra que sean compatibles con el campo de
fuerza seleccionado.

\subsection{Unidades utilizadas por
GROMACS}\label{unidades-utilizadas-por-gromacs}

{\def\LTcaptype{none} % do not increment counter
\begin{longtable}[]{@{}ll@{}}
\toprule\noalign{}
Magnitud & Unidad habitual \\
\midrule\noalign{}
\endhead
\bottomrule\noalign{}
\endlastfoot
Distancia & nm \\
Tiempo & ps \\
Temperatura & K \\
Presión & bar \\
Energía & kJ mol⁻¹ \\
Fuerza & kJ mol⁻¹ nm⁻¹ \\
Velocidad & nm ps⁻¹ \\
Constante de fuerza de restricción posicional & kJ mol⁻¹ nm⁻² \\
\end{longtable}
}

Equivalencias útiles:

\[
1\ \mathrm{nm}=10\ \text{Å}
\]

\[
1\ \mathrm{ps}=10^{-12}\ \mathrm{s},\qquad
1\ \mathrm{ns}=1000\ \mathrm{ps}
\]

\[
1\ \mathrm{kcal\ mol^{-1}}=4.184\ \mathrm{kJ\ mol^{-1}}
\]

No convierta distancias de Å a nm de forma implícita. Una distancia de
corte de 1.0 en un archivo MDP significa 1.0 nm, es decir, 10 Å.

\subsection{Preparación del ligando}\label{preparaciuxf3n-del-ligando}

La protonación debe corresponder al pH y al microentorno que se pretende
simular. Añadir hidrógenos con una herramienta de conversión no
sustituye esta decisión química.

Conversión básica con Open Babel:

\begin{Shaded}
\begin{Highlighting}[]
\ExtensionTok{obabel} \AttributeTok{{-}ipdb}\NormalTok{ ligand.pdb }\AttributeTok{{-}omol2} \AttributeTok{{-}O}\NormalTok{ ligand.mol2 }\AttributeTok{{-}h}
\end{Highlighting}
\end{Shaded}

Protonación aproximada a pH 7.4:

\begin{Shaded}
\begin{Highlighting}[]
\ExtensionTok{obabel} \AttributeTok{{-}ipdb}\NormalTok{ ligand.pdb }\AttributeTok{{-}omol2} \AttributeTok{{-}O}\NormalTok{ ligand\_pH7.4.mol2 }\AttributeTok{{-}p}\NormalTok{ 7.4}
\end{Highlighting}
\end{Shaded}

La opción \textbf{-p 7.4} predice una forma de protonación. Para
moléculas con tautomería, centros ionizables acoplados o metales, el
resultado debe verificarse. También deben revisarse carga formal, orden
de enlaces y estereoquímica.

GROMACS no parametriza automáticamente moléculas orgánicas arbitrarias.
La topología debe obtenerse con una herramienta apropiada para la
familia del campo de fuerza:

\begin{itemize}
\tightlist
\item
  CHARMM: parámetros compatibles con CGenFF/CHARMM.
\item
  AMBER: parámetros compatibles con GAFF u otra metodología AMBER.
\item
  OPLS-AA: tipos y cargas consistentes con OPLS.
\item
  GROMOS: parámetros desarrollados bajo las convenciones GROMOS.
\end{itemize}

No deben mezclarse tipos atómicos, reglas de combinación, cargas o
términos de CHARMM y AMBER sin una transformación validada.

\subsubsection{Parametrización del ligando con
ACPYPE}\label{parametrizaciuxf3n-del-ligando-con-acpype}

\textbf{ACPYPE} (\emph{AnteChamber PYthon Parser interfacE}) automatiza
la parametrización mediante herramientas de AmberTools y convierte el
resultado a formatos utilizables por GROMACS. Para una molécula pequeña
ejecuta, en esencia:

\begin{enumerate}
\def\labelenumi{\arabic{enumi}.}
\tightlist
\item
  \textbf{Antechamber}, que asigna tipos atómicos GAFF o GAFF2 y cargas
  parciales;
\item
  \textbf{parmchk2}, que busca parámetros ausentes y genera las
  adiciones necesarias;
\item
  \textbf{tleap}, que construye la topología y las coordenadas AMBER;
\item
  la conversión de la topología AMBER al formato de GROMACS.
\end{enumerate}

\paragraph{Correspondencia con el campo de
fuerza}\label{correspondencia-con-el-campo-de-fuerza}

Los parámetros producidos de esta manera pertenecen a la familia
\textbf{AMBER}:

{\def\LTcaptype{none} % do not increment counter
\begin{longtable}[]{@{}
  >{\raggedright\arraybackslash}p{(\linewidth - 4\tabcolsep) * \real{0.0842}}
  >{\raggedright\arraybackslash}p{(\linewidth - 4\tabcolsep) * \real{0.3053}}
  >{\raggedright\arraybackslash}p{(\linewidth - 4\tabcolsep) * \real{0.6105}}@{}}
\toprule\noalign{}
\begin{minipage}[b]{\linewidth}\raggedright
Opción de ACPYPE
\end{minipage} & \begin{minipage}[b]{\linewidth}\raggedright
Parametrización
\end{minipage} & \begin{minipage}[b]{\linewidth}\raggedright
Uso previsto
\end{minipage} \\
\midrule\noalign{}
\endhead
\bottomrule\noalign{}
\endlastfoot
\texttt{-a\ gaff} & General Amber Force Field, GAFF & Moléculas
orgánicas pequeñas dentro de un sistema AMBER compatible. \\
\texttt{-a\ gaff2} & General Amber Force Field 2, GAFF2 & Opción
preferente para moléculas pequeñas cuando se utiliza una parametrización
AMBER moderna. \\
\texttt{-a\ amber} & Tipos AMBER para residuos reconocidos y GAFF como
respaldo & Péptidos o sistemas que pueden compararse con plantillas
AMBER; no es la opción habitual para un ligando arbitrario. \\
\texttt{-a\ amber2} & Tipos AMBER y GAFF2 como respaldo & Variante
equivalente basada en GAFF2. \\
\end{longtable}
}

La opción \textbf{\texttt{-o\ gmx}} solicita archivos con sintaxis de
GROMACS; \textbf{no convierte GAFF2 en CHARMM, CGenFF, OPLS-AA o
GROMOS}. Del mismo modo, que ACPYPE pueda escribir un archivo denominado
CHARMM significa que cambia el formato, no que los parámetros hayan sido
desarrollados bajo CGenFF.

Una combinación coherente es, por ejemplo:

\begin{Shaded}
\begin{Highlighting}[]
\NormalTok{proteína: AMBER14SB o un port AMBER validado para GROMACS}
\NormalTok{ligando:  GAFF2 generado mediante ACPYPE/Antechamber}
\NormalTok{agua:     modelo recomendado para el campo AMBER elegido}
\end{Highlighting}
\end{Shaded}

No mezcle directamente una proteína CHARMM36 con un ligando GAFF2. Las
familias pueden diferir en cargas, torsiones, reglas de combinación y
factores de escala para interacciones 1--4. Un TPR que compila sin
errores no demuestra compatibilidad física.

En los ports AMBER para GROMACS deben conservarse las convenciones AMBER
declaradas en \texttt{{[}\ defaults\ {]}}, incluidas las reglas de
combinación y los factores \textbf{fudgeLJ/fudgeQQ}. No copie el bloque
\texttt{{[}\ defaults\ {]}} de la topología completa de ACPYPE dentro de
una topología ya creada por \textbf{pdb2gmx}.

\paragraph{Instalación reproducible}\label{instalaciuxf3n-reproducible}

Una instalación aislada mediante conda puede prepararse con:

\begin{Shaded}
\begin{Highlighting}[]
\ExtensionTok{conda}\NormalTok{ create }\AttributeTok{{-}n}\NormalTok{ acpype }\AttributeTok{{-}c}\NormalTok{ conda{-}forge }\DataTypeTok{\textbackslash{}}
\NormalTok{    acpype ambertools openbabel}

\ExtensionTok{conda}\NormalTok{ activate acpype}
\ExtensionTok{acpype} \AttributeTok{{-}h}
\ExtensionTok{antechamber} \AttributeTok{{-}h}
\ExtensionTok{obabel} \AttributeTok{{-}V}
\end{Highlighting}
\end{Shaded}

Registre las versiones exactas:

\begin{Shaded}
\begin{Highlighting}[]
\ExtensionTok{conda}\NormalTok{ list }\KeywordTok{|} \FunctionTok{grep} \AttributeTok{{-}E} \StringTok{\textquotesingle{}acpype|ambertools|openbabel\textquotesingle{}}
\end{Highlighting}
\end{Shaded}

La distribución actual de ACPYPE también puede instalarse mediante
\texttt{pip} en plataformas para las que existan ruedas con las
herramientas de AmberTools incorporadas. La disponibilidad depende del
sistema operativo y de la arquitectura; para un protocolo reproducible
conviene registrar tanto la versión de ACPYPE como la de AmberTools
realmente utilizada.

\paragraph{Preparación del archivo de
entrada}\label{preparaciuxf3n-del-archivo-de-entrada}

Para la parametrización es preferible un archivo MOL2 o SDF con:

\begin{itemize}
\tightlist
\item
  orden de enlaces correcto;
\item
  aromaticidad revisada;
\item
  protonación y tautomería definidas;
\item
  estereoquímica correcta;
\item
  coordenadas tridimensionales razonables;
\item
  carga formal conocida.
\end{itemize}

Un PDB no representa de manera inequívoca órdenes de enlace ni carga
formal. Si solo se dispone de PDB, la conversión mediante Open Babel
debe revisarse antes de ejecutar ACPYPE:

\begin{Shaded}
\begin{Highlighting}[]
\ExtensionTok{obabel}\NormalTok{ ligand.pdb }\DataTypeTok{\textbackslash{}}
    \AttributeTok{{-}O}\NormalTok{ LIG.mol2 }\DataTypeTok{\textbackslash{}}
    \AttributeTok{{-}h}
\end{Highlighting}
\end{Shaded}

Si se parte de SDF:

\begin{Shaded}
\begin{Highlighting}[]
\ExtensionTok{obabel}\NormalTok{ ligand.sdf }\DataTypeTok{\textbackslash{}}
    \AttributeTok{{-}O}\NormalTok{ LIG.mol2}
\end{Highlighting}
\end{Shaded}

No use automáticamente \textbf{\texttt{-\/-gen3d}} sobre una geometría
experimental o una pose de docking que se desea conservar. La generación
tridimensional crea otra conformación.

\paragraph{Ejecución con GAFF2}\label{ejecuciuxf3n-con-gaff2}

Para un ligando neutro llamado \textbf{LIG}:

\begin{Shaded}
\begin{Highlighting}[]
\ExtensionTok{acpype} \DataTypeTok{\textbackslash{}}
    \AttributeTok{{-}i}\NormalTok{ LIG.mol2 }\DataTypeTok{\textbackslash{}}
    \AttributeTok{{-}b}\NormalTok{ LIG }\DataTypeTok{\textbackslash{}}
    \AttributeTok{{-}c}\NormalTok{ abcg2 }\DataTypeTok{\textbackslash{}}
    \AttributeTok{{-}n}\NormalTok{ 0 }\DataTypeTok{\textbackslash{}}
    \AttributeTok{{-}a}\NormalTok{ gaff2 }\DataTypeTok{\textbackslash{}}
    \AttributeTok{{-}o}\NormalTok{ gmx}
\end{Highlighting}
\end{Shaded}

La documentación actual de ACPYPE recomienda \textbf{\texttt{abcg2}}
como método de cargas para GAFF2. El método AM1-BCC clásico puede
solicitarse con:

\begin{Shaded}
\begin{Highlighting}[]
\ExtensionTok{acpype} \DataTypeTok{\textbackslash{}}
    \AttributeTok{{-}i}\NormalTok{ LIG.mol2 }\DataTypeTok{\textbackslash{}}
    \AttributeTok{{-}b}\NormalTok{ LIG }\DataTypeTok{\textbackslash{}}
    \AttributeTok{{-}c}\NormalTok{ bcc }\DataTypeTok{\textbackslash{}}
    \AttributeTok{{-}n}\NormalTok{ 0 }\DataTypeTok{\textbackslash{}}
    \AttributeTok{{-}a}\NormalTok{ gaff2 }\DataTypeTok{\textbackslash{}}
    \AttributeTok{{-}o}\NormalTok{ gmx}
\end{Highlighting}
\end{Shaded}

No combine cargas generadas mediante métodos diferentes dentro de una
misma serie de ligandos si se pretende comparar resultados. El método,
la versión de AmberTools y la geometría usada para calcular las cargas
forman parte del protocolo.

La opción \textbf{\texttt{-n}} es la carga neta entera del microestado
simulado. Ejemplos:

\begin{Shaded}
\begin{Highlighting}[]
\CommentTok{\# Catión monovalente}
\ExtensionTok{acpype} \AttributeTok{{-}i}\NormalTok{ LIG\_plus.mol2 }\AttributeTok{{-}b}\NormalTok{ LIP }\AttributeTok{{-}c}\NormalTok{ abcg2 }\AttributeTok{{-}n}\NormalTok{ 1 }\AttributeTok{{-}a}\NormalTok{ gaff2 }\AttributeTok{{-}o}\NormalTok{ gmx}

\CommentTok{\# Anión monovalente}
\ExtensionTok{acpype} \AttributeTok{{-}i}\NormalTok{ LIG\_minus.mol2 }\AttributeTok{{-}b}\NormalTok{ LIM }\AttributeTok{{-}c}\NormalTok{ abcg2 }\AttributeTok{{-}n} \AttributeTok{{-}1} \AttributeTok{{-}a}\NormalTok{ gaff2 }\AttributeTok{{-}o}\NormalTok{ gmx}
\end{Highlighting}
\end{Shaded}

No utilice \texttt{-n\ 0} porque sea el valor del ejemplo. Una carga
neta incorrecta cambia todas las cargas parciales y la electrostática
del ligando.

Para cargas obtenidas por otro procedimiento cuántico o ajustadas
externamente, prepare un MOL2 que ya contenga esas cargas y consulte la
opción:

\begin{Shaded}
\begin{Highlighting}[]
\ExtensionTok{acpype} \AttributeTok{{-}i}\NormalTok{ LIG\_charged.mol2 }\AttributeTok{{-}b}\NormalTok{ LIG }\AttributeTok{{-}c}\NormalTok{ user }\AttributeTok{{-}n}\NormalTok{ 0 }\AttributeTok{{-}a}\NormalTok{ gaff2 }\AttributeTok{{-}o}\NormalTok{ gmx}
\end{Highlighting}
\end{Shaded}

La carga indicada mediante \textbf{\texttt{-n}} debe coincidir con la
suma de las cargas suministradas.

\paragraph{Archivos generados}\label{archivos-generados}

La ejecución crea normalmente el directorio
\textbf{\texttt{LIG.acpype/}}. Entre los archivos relevantes pueden
encontrarse:

\begin{Shaded}
\begin{Highlighting}[]
\NormalTok{LIG\_GMX.gro       coordenadas para GROMACS}
\NormalTok{LIG\_GMX.itp       tipos, cargas y parámetros del ligando}
\NormalTok{LIG\_GMX.top       topología completa para probar el ligando aislado}
\NormalTok{LIG\_AC.mol2       molécula procesada por Antechamber}
\NormalTok{LIG\_AC.frcmod     parámetros complementarios producidos por parmchk2}
\end{Highlighting}
\end{Shaded}

Los nombres exactos pueden variar entre versiones. Conserve también la
salida estándar y los archivos de registro: allí aparecen advertencias
de Antechamber, convergencia de cargas y parámetros inferidos por
analogía.

\paragraph{Validación del resultado}\label{validaciuxf3n-del-resultado}

Antes de incorporarlo al complejo, revise:

\begin{itemize}
\tightlist
\item
  suma de cargas parciales y carga neta;
\item
  tipos atómicos asignados;
\item
  orden de enlaces y aromaticidad;
\item
  geometría de anillos, amidas y centros quirales;
\item
  parámetros informados como ausentes o inferidos;
\item
  nombres y orden de átomos entre GRO, MOL2 e ITP;
\item
  energía y geometría después de una minimización corta;
\item
  presencia de fuerzas excesivas, NaN o advertencias de \textbf{grompp}.
\end{itemize}

Prueba independiente del ligando:

\begin{Shaded}
\begin{Highlighting}[]
\BuiltInTok{cd}\NormalTok{ LIG.acpype}

\ExtensionTok{gmx}\NormalTok{ grompp }\DataTypeTok{\textbackslash{}}
    \AttributeTok{{-}f}\NormalTok{ ../em\_ligand.mdp }\DataTypeTok{\textbackslash{}}
    \AttributeTok{{-}c}\NormalTok{ LIG\_GMX.gro }\DataTypeTok{\textbackslash{}}
    \AttributeTok{{-}p}\NormalTok{ LIG\_GMX.top }\DataTypeTok{\textbackslash{}}
    \AttributeTok{{-}pp}\NormalTok{ LIG\_processed.top }\DataTypeTok{\textbackslash{}}
    \AttributeTok{{-}o}\NormalTok{ LIG\_em.tpr}

\ExtensionTok{gmx}\NormalTok{ mdrun }\DataTypeTok{\textbackslash{}}
    \AttributeTok{{-}deffnm}\NormalTok{ LIG\_em }\DataTypeTok{\textbackslash{}}
    \AttributeTok{{-}v}
\end{Highlighting}
\end{Shaded}

Una minimización exitosa solo verifica coherencia mecánica básica. No
valida cargas, torsiones, estados de protonación ni la capacidad del
modelo para reproducir conformaciones o propiedades experimentales.

\paragraph{Incorporación en una topología AMBER de
proteína--ligando}\label{incorporaciuxf3n-en-una-topologuxeda-amber-de-proteuxednaligando}

No incluya \textbf{\texttt{LIG\_GMX.top}} completo dentro de
\textbf{topol.top}, porque contiene secciones globales que ya existen.
Utilice el ITP del ligando y asegúrese de que sus tipos atómicos se
definan antes de ser usados.

Esquema conceptual:

\begin{Shaded}
\begin{Highlighting}[]
\CommentTok{; Campo de fuerza AMBER seleccionado mediante pdb2gmx}
\CommentTok{\#include "amber14sb.ff/forcefield.itp"}

\CommentTok{; El ITP de ACPYPE puede contener [ atomtypes ]; debe aparecer}
\CommentTok{; antes de cualquier [ moleculetype ] que use esos tipos.}
\CommentTok{\#include "LIG\_GMX.itp"}

\CommentTok{; Topología de la proteína}
\CommentTok{\#include "topol\_Protein.itp"}

\CommentTok{; Agua e iones compatibles con el campo de fuerza}
\CommentTok{\#include "amber14sb.ff/tip3p.itp"}
\CommentTok{\#include "amber14sb.ff/ions.itp"}

\KeywordTok{[ system ]}
\DataTypeTok{Complejo proteína–ligando}

\KeywordTok{[ molecules ]}
\DataTypeTok{Protein    1}
\DataTypeTok{LIG        1}
\end{Highlighting}
\end{Shaded}

El identificador \texttt{amber14sb.ff} es ilustrativo: debe coincidir
con el directorio realmente instalado. Si el ITP de ACPYPE contiene
\texttt{{[}\ atomtypes\ {]}} y se incluye después de la topología de la
proteína, \textbf{grompp} puede detenerse con un error de orden de
directivas. Como alternativa, separe los tipos del ligando en
\textbf{ligand\_atomtypes.itp}, inclúyalos inmediatamente después de
\textbf{forcefield.itp} y deje las secciones
\texttt{{[}\ moleculetype\ {]}}, \texttt{{[}\ atoms\ {]}}, enlaces,
ángulos y diedros en \textbf{LIG.itp}.

El orden de \texttt{{[}\ molecules\ {]}} debe reproducir el orden de las
moléculas en el archivo de coordenadas. El nombre \textbf{LIG} debe
coincidir exactamente con el valor definido en
\texttt{{[}\ moleculetype\ {]}}, no solo con el nombre de residuo
mostrado en el GRO.

Genere siempre una topología preprocesada para auditar la integración:

\begin{Shaded}
\begin{Highlighting}[]
\ExtensionTok{gmx}\NormalTok{ grompp }\DataTypeTok{\textbackslash{}}
    \AttributeTok{{-}f}\NormalTok{ em.mdp }\DataTypeTok{\textbackslash{}}
    \AttributeTok{{-}c}\NormalTok{ complex.gro }\DataTypeTok{\textbackslash{}}
    \AttributeTok{{-}p}\NormalTok{ topol.top }\DataTypeTok{\textbackslash{}}
    \AttributeTok{{-}pp}\NormalTok{ complex\_processed.top }\DataTypeTok{\textbackslash{}}
    \AttributeTok{{-}o}\NormalTok{ complex\_em.tpr}
\end{Highlighting}
\end{Shaded}

No use \textbf{\texttt{-maxwarn}} para ignorar cargas inesperadas,
átomos ausentes, directivas fuera de orden o parámetros duplicados.

\paragraph{Cuándo ACPYPE no es
suficiente}\label{cuuxe1ndo-acpype-no-es-suficiente}

ACPYPE y GAFF2 requieren una revisión adicional o una estrategia
diferente para:

\begin{itemize}
\tightlist
\item
  metales coordinados y compuestos organometálicos;
\item
  enlaces covalentes entre ligando y proteína;
\item
  estados electrónicos abiertos;
\item
  moléculas con química poco representada por GAFF2;
\item
  sistemas para los que se necesita compatibilidad estricta con
  CHARMM/CGenFF, OPLS o GROMOS;
\item
  torsiones críticas que deben validarse contra cálculos cuánticos;
\item
  transformaciones alquímicas donde la correspondencia atómica y la
  topología híbrida requieren construcción específica.
\end{itemize}

Para una proteína parametrizada con CHARMM36, use una ruta CGenFF
compatible. Para OPLS-AA o GROMOS, emplee parámetros desarrollados y
validados bajo las convenciones respectivas. Convertir el archivo final
a sintaxis de GROMACS no resuelve la incompatibilidad entre familias.

Ejemplo mínimo de un archivo ITP:

\begin{Shaded}
\begin{Highlighting}[]
\KeywordTok{[ moleculetype ]}
\CommentTok{; nombre   nrexcl}
\DataTypeTok{LIG        3}

\KeywordTok{[ atoms ]}
\CommentTok{; nr  type  resnr  residuo  átomo  cgnr  carga  masa}
\DataTypeTok{1     CG2R61  1     LIG      C1     1     0.12   12.011}
\DataTypeTok{2     NG2R50  1     LIG      N1     2    {-}0.32   14.007}
\CommentTok{; ...}
\end{Highlighting}
\end{Shaded}

Las líneas iniciadas por punto y coma son comentarios. El nombre
definido en \texttt{{[}\ moleculetype\ {]}} debe coincidir exactamente
con el utilizado en \texttt{{[}\ molecules\ {]}}. La suma de las cargas
parciales debe reproducir la carga formal esperada dentro de la
precisión numérica del modelo.

El valor de \textbf{nrexcl} depende de la convención del campo de
fuerza. No debe cambiarse solo para eliminar errores de
preprocesamiento.

\subsection{Preparación de la
proteína}\label{preparaciuxf3n-de-la-proteuxedna}

Ejecute \textbf{pdb2gmx} sobre una copia de la estructura original:

\begin{Shaded}
\begin{Highlighting}[]
\BuiltInTok{cd}\NormalTok{ 01\_preparacion}

\ExtensionTok{gmx}\NormalTok{ pdb2gmx     }\AttributeTok{{-}f}\NormalTok{ ../00\_entrada/protein.pdb     }\AttributeTok{{-}o}\NormalTok{ protein\_processed.gro     }\AttributeTok{{-}p}\NormalTok{ topol.top     }\AttributeTok{{-}i}\NormalTok{ posre\_protein.itp     }\AttributeTok{{-}water}\NormalTok{ tip3p}
\end{Highlighting}
\end{Shaded}

Seleccione el campo de fuerza de manera interactiva o especifíquelo
mediante \textbf{-ff} si el identificador instalado está documentado:

\begin{Shaded}
\begin{Highlighting}[]
\ExtensionTok{gmx}\NormalTok{ pdb2gmx     }\AttributeTok{{-}f}\NormalTok{ ../00\_entrada/protein.pdb     }\AttributeTok{{-}o}\NormalTok{ protein\_processed.gro     }\AttributeTok{{-}p}\NormalTok{ topol.top     }\AttributeTok{{-}i}\NormalTok{ posre\_protein.itp     }\AttributeTok{{-}ff}\NormalTok{ charmm36{-}jul2022     }\AttributeTok{{-}water}\NormalTok{ tip3p}
\end{Highlighting}
\end{Shaded}

El nombre disponible puede diferir según los campos de fuerza
instalados. Consulte:

\begin{Shaded}
\begin{Highlighting}[]
\ExtensionTok{gmx}\NormalTok{ pdb2gmx }\AttributeTok{{-}h}
\end{Highlighting}
\end{Shaded}

Evite \textbf{-ignh} como opción automática: descarta los hidrógenos
presentes y los regenera según las plantillas del campo de fuerza. Puede
ser útil, pero también elimina estados de protonación definidos
previamente.

Compruebe la estructura:

\begin{Shaded}
\begin{Highlighting}[]
\ExtensionTok{gmx}\NormalTok{ check }\AttributeTok{{-}f}\NormalTok{ protein\_processed.gro}
\end{Highlighting}
\end{Shaded}

Revise en la salida de \textbf{pdb2gmx}:

\begin{itemize}
\tightlist
\item
  residuos y átomos no reconocidos;
\item
  carga total de la proteína;
\item
  enlaces disulfuro;
\item
  terminales asignados;
\item
  histidinas y otros residuos titulables;
\item
  moléculas separadas por cadena.
\end{itemize}

\subsection{Construcción del
complejo}\label{construcciuxf3n-del-complejo}

Las coordenadas finales deben contener primero la proteína y después el
ligando si ese es el orden declarado en la topología. No concatene
directamente dos archivos GRO completos: cada uno contiene título,
número de átomos y caja.

Use un editor molecular o un script validado para combinar las
coordenadas sin alterar la pose. Después inspeccione el complejo
visualmente y descarte solapamientos.

En \textbf{topol.top}, incluya la topología del ligando después de los
parámetros generales del campo de fuerza y antes de las topologías de
agua e iones:

\begin{Shaded}
\begin{Highlighting}[]
\CommentTok{; Parámetros generales}
\CommentTok{\#include "charmm36{-}jul2022.ff/forcefield.itp"}

\CommentTok{; Tipos o parámetros adicionales del ligando, si existen}
\CommentTok{\#include "ligand\_atomtypes.itp"}

\CommentTok{; Topología molecular del ligando}
\CommentTok{\#include "ligand.itp"}

\CommentTok{; Topología de la proteína generada por pdb2gmx}
\CommentTok{\#include "topol\_Protein\_chain\_A.itp"}
\end{Highlighting}
\end{Shaded}

Una sección \texttt{{[}\ atomtypes\ {]}} debe aparecer antes de
cualquier \texttt{{[}\ moleculetype\ {]}} que utilice esos tipos. No
incluya el mismo bloque de tipos atómicos más de una vez.

Al final de \textbf{topol.top}:

\begin{Shaded}
\begin{Highlighting}[]
\KeywordTok{[ system ]}
\DataTypeTok{Complejo proteína–ligando}

\KeywordTok{[ molecules ]}
\CommentTok{; molécula          cantidad}
\DataTypeTok{Protein\_chain\_A     1}
\DataTypeTok{LIG                 1}
\end{Highlighting}
\end{Shaded}

El orden en \texttt{{[}\ molecules\ {]}} debe coincidir con el orden de
las moléculas en el archivo de coordenadas. La cantidad representa
moléculas, no átomos ni residuos.

\subsection{Caja periódica}\label{caja-periuxf3dica}

Para una proteína soluble aproximadamente globular, una caja
dodecaédrica reduce el número de moléculas de agua respecto de una caja
cúbica:

\begin{Shaded}
\begin{Highlighting}[]
\ExtensionTok{gmx}\NormalTok{ editconf     }\AttributeTok{{-}f}\NormalTok{ complex.gro     }\AttributeTok{{-}o}\NormalTok{ complex\_box.gro     }\AttributeTok{{-}c}     \AttributeTok{{-}d}\NormalTok{ 1.0     }\AttributeTok{{-}bt}\NormalTok{ dodecahedron}
\end{Highlighting}
\end{Shaded}

\textbf{-d 1.0} establece una distancia mínima de 1.0 nm entre el soluto
y el límite de la caja. No es una constante universal. Debe ser
compatible con los radios de corte y con los movimientos esperados del
sistema.

Compruebe el volumen y los vectores de caja:

\begin{Shaded}
\begin{Highlighting}[]
\ExtensionTok{gmx}\NormalTok{ editconf }\AttributeTok{{-}f}\NormalTok{ complex\_box.gro}
\end{Highlighting}
\end{Shaded}

\subsection{Solvatación}\label{solvataciuxf3n}

\begin{Shaded}
\begin{Highlighting}[]
\ExtensionTok{gmx}\NormalTok{ solvate     }\AttributeTok{{-}cp}\NormalTok{ complex\_box.gro     }\AttributeTok{{-}cs}\NormalTok{ spc216.gro     }\AttributeTok{{-}o}\NormalTok{ complex\_solv.gro     }\AttributeTok{{-}p}\NormalTok{ topol.top}
\end{Highlighting}
\end{Shaded}

El nombre \textbf{spc216.gro} identifica una configuración
preequilibrada distribuida con GROMACS. La topología final del agua está
determinada por el modelo elegido en \textbf{pdb2gmx}, por lo que debe
mantenerse la compatibilidad entre campo de fuerza, archivo de agua y
topología.

\textbf{gmx solvate} actualiza automáticamente el número de moléculas de
solvente en \textbf{topol.top}. Revise la sección
\texttt{{[}\ molecules\ {]}} después del comando.

\subsection{Neutralización y concentración
salina}\label{neutralizaciuxf3n-y-concentraciuxf3n-salina}

Archivo \textbf{ions.mdp} mínimo:

\begin{Shaded}
\begin{Highlighting}[]
\DataTypeTok{integrator      }\OtherTok{=}\StringTok{ steep}
\DataTypeTok{nsteps          }\OtherTok{=}\StringTok{ }\DecValTok{0}
\DataTypeTok{emtol           }\OtherTok{=}\StringTok{ }\FloatTok{1000.0}

\DataTypeTok{cutoff{-}scheme   }\OtherTok{=}\StringTok{ Verlet}
\DataTypeTok{coulombtype     }\OtherTok{=}\StringTok{ PME}
\DataTypeTok{rcoulomb        }\OtherTok{=}\StringTok{ }\FloatTok{1.0}
\DataTypeTok{rvdw            }\OtherTok{=}\StringTok{ }\FloatTok{1.0}
\DataTypeTok{pbc             }\OtherTok{=}\StringTok{ xyz}
\end{Highlighting}
\end{Shaded}

Genere un TPR temporal:

\begin{Shaded}
\begin{Highlighting}[]
\ExtensionTok{gmx}\NormalTok{ grompp     }\AttributeTok{{-}f}\NormalTok{ ions.mdp     }\AttributeTok{{-}c}\NormalTok{ complex\_solv.gro     }\AttributeTok{{-}p}\NormalTok{ topol.top     }\AttributeTok{{-}o}\NormalTok{ ions.tpr}
\end{Highlighting}
\end{Shaded}

Añada NaCl, neutralice la carga neta y solicite una concentración
nominal de 0.15 mol L⁻¹:

\begin{Shaded}
\begin{Highlighting}[]
\ExtensionTok{gmx}\NormalTok{ genion     }\AttributeTok{{-}s}\NormalTok{ ions.tpr     }\AttributeTok{{-}o}\NormalTok{ complex\_solv\_ions.gro     }\AttributeTok{{-}p}\NormalTok{ topol.top     }\AttributeTok{{-}pname}\NormalTok{ NA     }\AttributeTok{{-}nname}\NormalTok{ CL     }\AttributeTok{{-}neutral}     \AttributeTok{{-}conc}\NormalTok{ 0.15}
\end{Highlighting}
\end{Shaded}

Seleccione el grupo de solvente, normalmente \textbf{SOL}, cuando el
programa pregunte qué moléculas reemplazar. No use un número de grupo
copiado de otro sistema.

La concentración se relaciona con el número de pares iónicos mediante:

\[
N_{\mathrm{pares}} \approx c\,N_A\,V
\]

donde \textbf{c} es la concentración en mol L⁻¹, \(N_A\) es la constante
de Avogadro y \textbf{V} es el volumen en litros. Debido a que el número
de iones debe ser entero y la caja es pequeña, la concentración efectiva
puede diferir del valor solicitado.

La opción \textbf{-neutral} agrega los contraiones necesarios para
llevar la carga neta a cero. \textbf{-conc 0.15} agrega además la sal
correspondiente a la concentración solicitada. Revise la cantidad final
de NA y CL en \textbf{topol.top}.

Concentración de sal y fuerza iónica no son siempre equivalentes. La
fuerza iónica se define como

\[
I=\frac{1}{2}\sum_i c_i z_i^2,
\]

donde \(c_i\) es la concentración molar de la especie iónica \(i\) y
\(z_i\) su número de carga. Para NaCl ideal, 0.15 mol L⁻¹ corresponde
aproximadamente a \(I=0.15\) mol L⁻¹. La presencia de contraiones
añadidos para neutralizar una proteína cargada modifica la composición y
la fuerza iónica efectiva; el efecto es más visible en cajas pequeñas o
con solutos muy cargados.

\subsection{Grupos de índice}\label{grupos-de-uxedndice}

Los números de grupo cambian con la composición del sistema. Cree los
grupos necesarios y documente las selecciones:

\begin{Shaded}
\begin{Highlighting}[]
\ExtensionTok{gmx}\NormalTok{ make\_ndx }\AttributeTok{{-}f}\NormalTok{ complex\_solv\_ions.gro }\AttributeTok{{-}o}\NormalTok{ index.ndx}
\end{Highlighting}
\end{Shaded}

Ejemplo interactivo:

\begin{Shaded}
\begin{Highlighting}[]
\NormalTok{r LIG}
\NormalTok{"Protein" | "LIG"}
\NormalTok{name 18 Protein\_LIG}
\NormalTok{q}
\end{Highlighting}
\end{Shaded}

El número 18 es ilustrativo: debe reemplazarse por el número asignado
durante esa sesión. El grupo \textbf{Protein\_LIG} resulta útil para
centrar y visualizar el complejo, pero no crea una única molécula
física.

Puede examinar selecciones modernas con:

\begin{Shaded}
\begin{Highlighting}[]
\ExtensionTok{gmx}\NormalTok{ select     }\AttributeTok{{-}s}\NormalTok{ complex\_solv\_ions.gro     }\AttributeTok{{-}select} \StringTok{\textquotesingle{}group "Protein" or resname LIG\textquotesingle{}}
\end{Highlighting}
\end{Shaded}

\subsection{Restricciones de posición del
ligando}\label{restricciones-de-posiciuxf3n-del-ligando}

Genere las restricciones usando una estructura que contenga solamente el
ligando y cuyo orden atómico coincida con \textbf{ligand.itp}:

\begin{Shaded}
\begin{Highlighting}[]
\ExtensionTok{gmx}\NormalTok{ genrestr     }\AttributeTok{{-}f}\NormalTok{ ligand.gro     }\AttributeTok{{-}o}\NormalTok{ posre\_ligand.itp     }\AttributeTok{{-}fc}\NormalTok{ 1000 1000 1000}
\end{Highlighting}
\end{Shaded}

Seleccione \textbf{LIG}. Los índices del archivo de restricciones son
locales al \texttt{{[}\ moleculetype\ {]}}. Si se usa la estructura
completa y se generan índices globales, las restricciones pueden apuntar
a átomos incorrectos.

Incluya el archivo inmediatamente después de la topología del ligando:

\begin{Shaded}
\begin{Highlighting}[]
\CommentTok{\#include "ligand.itp"}

\CommentTok{\#ifdef POSRES\_LIG}
\CommentTok{\#include "posre\_ligand.itp"}
\CommentTok{\#endif}
\end{Highlighting}
\end{Shaded}

El valor 1000 corresponde a 1000 kJ mol⁻¹ nm⁻² en cada eje. La energía
armónica de una restricción unidimensional es:

\[
V(x)=\frac{1}{2}k(x-x_0)^2
\]

donde \textbf{k} es la constante de fuerza y \(x_0\) la posición de
referencia.

Active las restricciones desde el MDP:

\begin{Shaded}
\begin{Highlighting}[]
\DataTypeTok{define }\OtherTok{=}\StringTok{ {-}DPOSRES {-}DPOSRES\_LIG}
\end{Highlighting}
\end{Shaded}

El archivo indicado mediante \textbf{gmx grompp -r} proporciona las
coordenadas de referencia. Las restricciones se usan normalmente durante
la equilibración y se eliminan en producción.

\subsection{Minimización de
energía}\label{minimizaciuxf3n-de-energuxeda}

Archivo \textbf{em.mdp}:

\begin{Shaded}
\begin{Highlighting}[]
\DataTypeTok{title            }\OtherTok{=}\StringTok{ Minimización de energía}
\DataTypeTok{integrator       }\OtherTok{=}\StringTok{ steep}
\DataTypeTok{nsteps           }\OtherTok{=}\StringTok{ }\DecValTok{50000}
\DataTypeTok{emtol            }\OtherTok{=}\StringTok{ }\FloatTok{1000.0}
\DataTypeTok{emstep           }\OtherTok{=}\StringTok{ }\FloatTok{0.01}

\DataTypeTok{cutoff{-}scheme    }\OtherTok{=}\StringTok{ Verlet}
\DataTypeTok{nstlist          }\OtherTok{=}\StringTok{ }\DecValTok{20}
\DataTypeTok{rlist            }\OtherTok{=}\StringTok{ }\FloatTok{1.0}
\DataTypeTok{coulombtype      }\OtherTok{=}\StringTok{ PME}
\DataTypeTok{rcoulomb         }\OtherTok{=}\StringTok{ }\FloatTok{1.0}
\DataTypeTok{vdwtype          }\OtherTok{=}\StringTok{ Cut{-}off}
\DataTypeTok{rvdw             }\OtherTok{=}\StringTok{ }\FloatTok{1.0}
\DataTypeTok{pbc              }\OtherTok{=}\StringTok{ xyz}
\end{Highlighting}
\end{Shaded}

Preprocese y ejecute:

\begin{Shaded}
\begin{Highlighting}[]
\FunctionTok{mkdir} \AttributeTok{{-}p}\NormalTok{ ../02\_em}
\BuiltInTok{cd}\NormalTok{ ../02\_em}

\ExtensionTok{gmx}\NormalTok{ grompp     }\AttributeTok{{-}f}\NormalTok{ ../01\_preparacion/em.mdp     }\AttributeTok{{-}c}\NormalTok{ ../01\_preparacion/complex\_solv\_ions.gro     }\AttributeTok{{-}p}\NormalTok{ ../01\_preparacion/topol.top     }\AttributeTok{{-}o}\NormalTok{ em.tpr}

\ExtensionTok{gmx}\NormalTok{ mdrun }\AttributeTok{{-}deffnm}\NormalTok{ em }\AttributeTok{{-}v}
\end{Highlighting}
\end{Shaded}

El criterio \textbf{emtol = 1000} significa que la minimización puede
finalizar cuando la fuerza máxima sea menor que 1000 kJ mol⁻¹ nm⁻¹:

\[
F_{\max}=\max_i\left|-\nabla_i U\right|.
\]

Este valor es habitual antes de equilibrar, pero no garantiza que la
estructura represente un mínimo profundo ni que el sistema esté
equilibrado. La minimización desplaza coordenadas cuesta abajo sobre la
superficie de energía potencial; no genera un conjunto termodinámico y
no sustituye NVT o NPT.

Extraiga la energía potencial:

\begin{Shaded}
\begin{Highlighting}[]
\KeywordTok{(}\BuiltInTok{echo}\NormalTok{ Potential}\KeywordTok{;} \BuiltInTok{echo}\NormalTok{ 0}\KeywordTok{)} \KeywordTok{|}
\ExtensionTok{gmx}\NormalTok{ energy }\AttributeTok{{-}f}\NormalTok{ em.edr }\AttributeTok{{-}o}\NormalTok{ ../06\_analisis/em\_potential.xvg}
\end{Highlighting}
\end{Shaded}

Compruebe:

\begin{itemize}
\tightlist
\item
  descenso de la energía potencial;
\item
  ausencia de valores NaN;
\item
  fuerza máxima final;
\item
  átomo sobre el cual actúa la fuerza máxima;
\item
  advertencias de LINCS o contactos anómalos.
\end{itemize}

No use \textbf{-maxwarn} para ocultar advertencias de \textbf{grompp}.
Debe entenderse la causa antes de continuar.

\subsection{Equilibración NVT}\label{equilibraciuxf3n-nvt}

En el conjunto canónico NVT permanecen fijos el número de partículas
\(N\), el volumen \(V\) y la temperatura objetivo \(T\). El termostato
modifica la dinámica para muestrear la distribución correspondiente; no
debe interpretarse como una simple corrección periódica de velocidades.

La probabilidad de un microestado de energía \(E\) es proporcional a

\[
P(E)\propto \exp\left(-\frac{E}{k_{\mathrm B}T}\right).
\]

En esta etapa se estabiliza la temperatura manteniendo fijo el volumen.
Ejemplo a 300 K durante 100 ps:

\begin{Shaded}
\begin{Highlighting}[]
\DataTypeTok{title                    }\OtherTok{=}\StringTok{ Equilibración NVT}
\DataTypeTok{define                   }\OtherTok{=}\StringTok{ {-}DPOSRES {-}DPOSRES\_LIG}

\DataTypeTok{integrator               }\OtherTok{=}\StringTok{ md}
\DataTypeTok{dt                       }\OtherTok{=}\StringTok{ }\FloatTok{0.002}
\DataTypeTok{nsteps                   }\OtherTok{=}\StringTok{ }\DecValTok{50000}
\DataTypeTok{continuation             }\OtherTok{=}\StringTok{ }\KeywordTok{no}

\DataTypeTok{gen{-}vel                  }\OtherTok{=}\StringTok{ }\KeywordTok{yes}
\DataTypeTok{gen{-}temp                 }\OtherTok{=}\StringTok{ }\DecValTok{300}
\DataTypeTok{gen{-}seed                 }\OtherTok{=}\StringTok{ }\DecValTok{{-}1}

\DataTypeTok{constraints              }\OtherTok{=}\StringTok{ h{-}bonds}
\DataTypeTok{constraint{-}algorithm     }\OtherTok{=}\StringTok{ lincs}

\DataTypeTok{cutoff{-}scheme            }\OtherTok{=}\StringTok{ Verlet}
\DataTypeTok{nstlist                  }\OtherTok{=}\StringTok{ }\DecValTok{20}
\DataTypeTok{rlist                    }\OtherTok{=}\StringTok{ }\FloatTok{1.0}
\DataTypeTok{coulombtype              }\OtherTok{=}\StringTok{ PME}
\DataTypeTok{rcoulomb                 }\OtherTok{=}\StringTok{ }\FloatTok{1.0}
\DataTypeTok{vdwtype                  }\OtherTok{=}\StringTok{ Cut{-}off}
\DataTypeTok{rvdw                     }\OtherTok{=}\StringTok{ }\FloatTok{1.0}

\DataTypeTok{tcoupl                   }\OtherTok{=}\StringTok{ V{-}rescale}
\DataTypeTok{tc{-}grps                  }\OtherTok{=}\StringTok{ Protein\_LIG Water\_and\_ions}
\DataTypeTok{tau{-}t                    }\OtherTok{=}\StringTok{ 1.0 1.0}
\DataTypeTok{ref{-}t                    }\OtherTok{=}\StringTok{ 300 300}

\DataTypeTok{pcoupl                   }\OtherTok{=}\StringTok{ }\KeywordTok{no}
\DataTypeTok{pbc                      }\OtherTok{=}\StringTok{ xyz}

\DataTypeTok{nstxout{-}compressed       }\OtherTok{=}\StringTok{ }\DecValTok{500}
\DataTypeTok{compressed{-}x{-}precision   }\OtherTok{=}\StringTok{ }\DecValTok{1000}
\DataTypeTok{nstenergy                }\OtherTok{=}\StringTok{ }\DecValTok{500}
\DataTypeTok{nstlog                   }\OtherTok{=}\StringTok{ }\DecValTok{500}
\end{Highlighting}
\end{Shaded}

El tiempo simulado es:

\[
t_{\mathrm{total}}=\mathrm{dt}\times\mathrm{nsteps}
\]

En este ejemplo:

\[
0.002\ \mathrm{ps}\times 50000=100\ \mathrm{ps}
\]

Con restricciones sobre enlaces con hidrógeno, un paso de 0.002 ps
equivale a 2 fs.

Ejecución:

\begin{Shaded}
\begin{Highlighting}[]
\FunctionTok{mkdir} \AttributeTok{{-}p}\NormalTok{ ../03\_nvt}
\BuiltInTok{cd}\NormalTok{ ../03\_nvt}

\ExtensionTok{gmx}\NormalTok{ grompp     }\AttributeTok{{-}f}\NormalTok{ ../01\_preparacion/nvt.mdp     }\AttributeTok{{-}c}\NormalTok{ ../02\_em/em.gro     }\AttributeTok{{-}r}\NormalTok{ ../02\_em/em.gro     }\AttributeTok{{-}p}\NormalTok{ ../01\_preparacion/topol.top     }\AttributeTok{{-}n}\NormalTok{ ../01\_preparacion/index.ndx     }\AttributeTok{{-}o}\NormalTok{ nvt.tpr}

\ExtensionTok{gmx}\NormalTok{ mdrun }\AttributeTok{{-}deffnm}\NormalTok{ nvt }\AttributeTok{{-}v}
\end{Highlighting}
\end{Shaded}

Las velocidades se generan una sola vez. Para las etapas posteriores se
conserva el estado mediante el checkpoint.

Análisis de temperatura:

\begin{Shaded}
\begin{Highlighting}[]
\KeywordTok{(}\BuiltInTok{echo}\NormalTok{ Temperature}\KeywordTok{;} \BuiltInTok{echo}\NormalTok{ 0}\KeywordTok{)} \KeywordTok{|}
\ExtensionTok{gmx}\NormalTok{ energy }\AttributeTok{{-}f}\NormalTok{ nvt.edr }\AttributeTok{{-}o}\NormalTok{ ../06\_analisis/nvt\_temperature.xvg}
\end{Highlighting}
\end{Shaded}

La temperatura debe evaluarse como serie temporal y promedio, no por un
único valor final.

\subsection{Equilibración NPT}\label{equilibraciuxf3n-npt}

En el conjunto isotérmico-isobárico NPT permanecen fijos \(N\), la
temperatura objetivo \(T\) y la presión objetivo \(P\), mientras el
volumen fluctúa. El barostato modifica los vectores de caja y las
coordenadas para permitir que la densidad se relaje.

La presión instantánea es una magnitud ruidosa, especialmente en cajas
pequeñas. El criterio práctico no es obtener una línea plana en 1 bar,
sino comprobar que volumen y densidad alcanzaron un régimen estacionario
y que la presión promedio es compatible con el objetivo dentro de su
incertidumbre.

Ejemplo de 500 ps con barostato C-rescale:

\begin{Shaded}
\begin{Highlighting}[]
\DataTypeTok{title                    }\OtherTok{=}\StringTok{ Equilibración NPT}
\DataTypeTok{define                   }\OtherTok{=}\StringTok{ {-}DPOSRES {-}DPOSRES\_LIG}

\DataTypeTok{integrator               }\OtherTok{=}\StringTok{ md}
\DataTypeTok{dt                       }\OtherTok{=}\StringTok{ }\FloatTok{0.002}
\DataTypeTok{nsteps                   }\OtherTok{=}\StringTok{ }\DecValTok{250000}
\DataTypeTok{continuation             }\OtherTok{=}\StringTok{ }\KeywordTok{yes}
\DataTypeTok{gen{-}vel                  }\OtherTok{=}\StringTok{ }\KeywordTok{no}

\DataTypeTok{constraints              }\OtherTok{=}\StringTok{ h{-}bonds}
\DataTypeTok{constraint{-}algorithm     }\OtherTok{=}\StringTok{ lincs}

\DataTypeTok{cutoff{-}scheme            }\OtherTok{=}\StringTok{ Verlet}
\DataTypeTok{nstlist                  }\OtherTok{=}\StringTok{ }\DecValTok{20}
\DataTypeTok{rlist                    }\OtherTok{=}\StringTok{ }\FloatTok{1.0}
\DataTypeTok{coulombtype              }\OtherTok{=}\StringTok{ PME}
\DataTypeTok{rcoulomb                 }\OtherTok{=}\StringTok{ }\FloatTok{1.0}
\DataTypeTok{vdwtype                  }\OtherTok{=}\StringTok{ Cut{-}off}
\DataTypeTok{rvdw                     }\OtherTok{=}\StringTok{ }\FloatTok{1.0}

\DataTypeTok{tcoupl                   }\OtherTok{=}\StringTok{ V{-}rescale}
\DataTypeTok{tc{-}grps                  }\OtherTok{=}\StringTok{ Protein\_LIG Water\_and\_ions}
\DataTypeTok{tau{-}t                    }\OtherTok{=}\StringTok{ 1.0 1.0}
\DataTypeTok{ref{-}t                    }\OtherTok{=}\StringTok{ 300 300}

\DataTypeTok{pcoupl                   }\OtherTok{=}\StringTok{ C{-}rescale}
\DataTypeTok{pcoupltype               }\OtherTok{=}\StringTok{ isotropic}
\DataTypeTok{tau{-}p                    }\OtherTok{=}\StringTok{ }\FloatTok{5.0}
\DataTypeTok{ref{-}p                    }\OtherTok{=}\StringTok{ }\FloatTok{1.0}
\DataTypeTok{compressibility          }\OtherTok{=}\StringTok{ 4.5e{-}5}

\DataTypeTok{pbc                      }\OtherTok{=}\StringTok{ xyz}

\DataTypeTok{nstxout{-}compressed       }\OtherTok{=}\StringTok{ }\DecValTok{500}
\DataTypeTok{compressed{-}x{-}precision   }\OtherTok{=}\StringTok{ }\DecValTok{1000}
\DataTypeTok{nstenergy                }\OtherTok{=}\StringTok{ }\DecValTok{500}
\DataTypeTok{nstlog                   }\OtherTok{=}\StringTok{ }\DecValTok{500}
\end{Highlighting}
\end{Shaded}

\textbf{ref-p = 1.0} significa 1 bar. \textbf{compressibility = 4.5e-5
bar⁻¹} es un valor habitual para agua líquida cerca de condiciones
ambientales; debe adaptarse si el medio no es agua.

Ejecución:

\begin{Shaded}
\begin{Highlighting}[]
\FunctionTok{mkdir} \AttributeTok{{-}p}\NormalTok{ ../04\_npt}
\BuiltInTok{cd}\NormalTok{ ../04\_npt}

\ExtensionTok{gmx}\NormalTok{ grompp     }\AttributeTok{{-}f}\NormalTok{ ../01\_preparacion/npt.mdp     }\AttributeTok{{-}c}\NormalTok{ ../03\_nvt/nvt.gro     }\AttributeTok{{-}r}\NormalTok{ ../03\_nvt/nvt.gro     }\AttributeTok{{-}t}\NormalTok{ ../03\_nvt/nvt.cpt     }\AttributeTok{{-}p}\NormalTok{ ../01\_preparacion/topol.top     }\AttributeTok{{-}n}\NormalTok{ ../01\_preparacion/index.ndx     }\AttributeTok{{-}o}\NormalTok{ npt.tpr}

\ExtensionTok{gmx}\NormalTok{ mdrun }\AttributeTok{{-}deffnm}\NormalTok{ npt }\AttributeTok{{-}v}
\end{Highlighting}
\end{Shaded}

Extraiga presión y densidad:

\begin{Shaded}
\begin{Highlighting}[]
\KeywordTok{(}\BuiltInTok{echo}\NormalTok{ Pressure}\KeywordTok{;} \BuiltInTok{echo}\NormalTok{ Density}\KeywordTok{;} \BuiltInTok{echo}\NormalTok{ 0}\KeywordTok{)} \KeywordTok{|}
\ExtensionTok{gmx}\NormalTok{ energy }\AttributeTok{{-}f}\NormalTok{ npt.edr }\AttributeTok{{-}o}\NormalTok{ ../06\_analisis/npt\_pressure\_density.xvg}
\end{Highlighting}
\end{Shaded}

La presión instantánea fluctúa mucho en sistemas pequeños. Evalúe
promedios por bloques, densidad y volumen. Si existe una deriva
sistemática, prolongue la equilibración desde el checkpoint.

\subsection{Dinámica molecular de
producción}\label{dinuxe1mica-molecular-de-producciuxf3n}

Durante producción se retiran las restricciones posicionales, salvo que
formen parte explícita del protocolo. El siguiente MDP describe 100 ns:

\begin{Shaded}
\begin{Highlighting}[]
\DataTypeTok{title                    }\OtherTok{=}\StringTok{ Producción NPT}

\DataTypeTok{integrator               }\OtherTok{=}\StringTok{ md}
\DataTypeTok{dt                       }\OtherTok{=}\StringTok{ }\FloatTok{0.002}
\DataTypeTok{nsteps                   }\OtherTok{=}\StringTok{ }\DecValTok{50000000}
\DataTypeTok{continuation             }\OtherTok{=}\StringTok{ }\KeywordTok{yes}
\DataTypeTok{gen{-}vel                  }\OtherTok{=}\StringTok{ }\KeywordTok{no}

\DataTypeTok{constraints              }\OtherTok{=}\StringTok{ h{-}bonds}
\DataTypeTok{constraint{-}algorithm     }\OtherTok{=}\StringTok{ lincs}

\DataTypeTok{cutoff{-}scheme            }\OtherTok{=}\StringTok{ Verlet}
\DataTypeTok{nstlist                  }\OtherTok{=}\StringTok{ }\DecValTok{20}
\DataTypeTok{rlist                    }\OtherTok{=}\StringTok{ }\FloatTok{1.0}
\DataTypeTok{coulombtype              }\OtherTok{=}\StringTok{ PME}
\DataTypeTok{rcoulomb                 }\OtherTok{=}\StringTok{ }\FloatTok{1.0}
\DataTypeTok{vdwtype                  }\OtherTok{=}\StringTok{ Cut{-}off}
\DataTypeTok{rvdw                     }\OtherTok{=}\StringTok{ }\FloatTok{1.0}

\DataTypeTok{tcoupl                   }\OtherTok{=}\StringTok{ V{-}rescale}
\DataTypeTok{tc{-}grps                  }\OtherTok{=}\StringTok{ Protein\_LIG Water\_and\_ions}
\DataTypeTok{tau{-}t                    }\OtherTok{=}\StringTok{ 1.0 1.0}
\DataTypeTok{ref{-}t                    }\OtherTok{=}\StringTok{ 300 300}

\DataTypeTok{pcoupl                   }\OtherTok{=}\StringTok{ Parrinello{-}Rahman}
\DataTypeTok{pcoupltype               }\OtherTok{=}\StringTok{ isotropic}
\DataTypeTok{tau{-}p                    }\OtherTok{=}\StringTok{ }\FloatTok{5.0}
\DataTypeTok{ref{-}p                    }\OtherTok{=}\StringTok{ }\FloatTok{1.0}
\DataTypeTok{compressibility          }\OtherTok{=}\StringTok{ 4.5e{-}5}

\DataTypeTok{pbc                      }\OtherTok{=}\StringTok{ xyz}

\DataTypeTok{nstxout{-}compressed       }\OtherTok{=}\StringTok{ }\DecValTok{5000}
\DataTypeTok{compressed{-}x{-}precision   }\OtherTok{=}\StringTok{ }\DecValTok{1000}
\DataTypeTok{nstenergy                }\OtherTok{=}\StringTok{ }\DecValTok{5000}
\DataTypeTok{nstlog                   }\OtherTok{=}\StringTok{ }\DecValTok{5000}
\end{Highlighting}
\end{Shaded}

Con 2 fs por paso:

\[
50000000\times 0.002\ \mathrm{ps}
=100000\ \mathrm{ps}
=100\ \mathrm{ns}
\]

Ejecución:

\begin{Shaded}
\begin{Highlighting}[]
\FunctionTok{mkdir} \AttributeTok{{-}p}\NormalTok{ ../05\_md}
\BuiltInTok{cd}\NormalTok{ ../05\_md}

\ExtensionTok{gmx}\NormalTok{ grompp     }\AttributeTok{{-}f}\NormalTok{ ../01\_preparacion/md.mdp     }\AttributeTok{{-}c}\NormalTok{ ../04\_npt/npt.gro     }\AttributeTok{{-}t}\NormalTok{ ../04\_npt/npt.cpt     }\AttributeTok{{-}p}\NormalTok{ ../01\_preparacion/topol.top     }\AttributeTok{{-}n}\NormalTok{ ../01\_preparacion/index.ndx     }\AttributeTok{{-}o}\NormalTok{ md.tpr}

\ExtensionTok{gmx}\NormalTok{ mdrun }\AttributeTok{{-}deffnm}\NormalTok{ md }\AttributeTok{{-}v}
\end{Highlighting}
\end{Shaded}

Archivos principales:

{\def\LTcaptype{none} % do not increment counter
\begin{longtable}[]{@{}
  >{\raggedright\arraybackslash}p{(\linewidth - 2\tabcolsep) * \real{0.1163}}
  >{\raggedright\arraybackslash}p{(\linewidth - 2\tabcolsep) * \real{0.8837}}@{}}
\toprule\noalign{}
\begin{minipage}[b]{\linewidth}\raggedright
Archivo
\end{minipage} & \begin{minipage}[b]{\linewidth}\raggedright
Contenido
\end{minipage} \\
\midrule\noalign{}
\endhead
\bottomrule\noalign{}
\endlastfoot
\textbf{md.tpr} & Topología, parámetros, coordenadas y estado de
entrada. \\
\textbf{md.xtc} & Coordenadas comprimidas. \\
\textbf{md.edr} & Energías y variables termodinámicas. \\
\textbf{md.log} & Registro detallado de la ejecución. \\
\textbf{md.cpt} & Checkpoint para continuar. \\
\textbf{md.gro} & Coordenadas finales. \\
\textbf{md.trr} & Coordenadas, velocidades o fuerzas en precisión
completa, si se solicitaron. \\
\end{longtable}
}

No es necesario generar TRR si el análisis no requiere velocidades,
fuerzas o coordenadas de precisión completa.

\subsection{Continuación y
extensión}\label{continuaciuxf3n-y-extensiuxf3n}

Para reanudar una ejecución interrumpida:

\begin{Shaded}
\begin{Highlighting}[]
\ExtensionTok{gmx}\NormalTok{ mdrun }\AttributeTok{{-}deffnm}\NormalTok{ md }\AttributeTok{{-}cpi}\NormalTok{ md.cpt }\AttributeTok{{-}append}
\end{Highlighting}
\end{Shaded}

\textbf{-append} verifica y continúa los archivos existentes. Use
\textbf{-noappend} solo cuando necesite segmentos separados.

Si el TPR agotó el número de pasos, extiéndalo. \textbf{-extend} se
expresa en picosegundos:

\begin{Shaded}
\begin{Highlighting}[]
\ExtensionTok{gmx}\NormalTok{ convert{-}tpr     }\AttributeTok{{-}s}\NormalTok{ md.tpr     }\AttributeTok{{-}extend}\NormalTok{ 50000     }\AttributeTok{{-}o}\NormalTok{ md\_extended.tpr}

\ExtensionTok{gmx}\NormalTok{ mdrun     }\AttributeTok{{-}s}\NormalTok{ md\_extended.tpr     }\AttributeTok{{-}deffnm}\NormalTok{ md     }\AttributeTok{{-}cpi}\NormalTok{ md.cpt     }\AttributeTok{{-}append}
\end{Highlighting}
\end{Shaded}

Aquí se agregan 50000 ps, equivalentes a 50 ns.

Concatenación de segmentos XTC:

\begin{Shaded}
\begin{Highlighting}[]
\ExtensionTok{gmx}\NormalTok{ trjcat     }\AttributeTok{{-}f}\NormalTok{ md.part0001.xtc md.part0002.xtc     }\AttributeTok{{-}o}\NormalTok{ md\_complete.xtc}
\end{Highlighting}
\end{Shaded}

Revise el orden y los tiempos. La concatenación no corrige
superposiciones temporales ni discontinuidades físicas.

\subsection{Corrección de condiciones
periódicas}\label{correcciuxf3n-de-condiciones-periuxf3dicas}

Las condiciones periódicas pueden separar visualmente moléculas que
continúan próximas. No existe una única secuencia válida para todos los
sistemas. Para un complejo soluble:

\begin{Shaded}
\begin{Highlighting}[]
\BuiltInTok{echo}\NormalTok{ System }\KeywordTok{|}
\ExtensionTok{gmx}\NormalTok{ trjconv     }\AttributeTok{{-}s}\NormalTok{ md.tpr     }\AttributeTok{{-}f}\NormalTok{ md.xtc     }\AttributeTok{{-}o}\NormalTok{ md\_whole.xtc     }\AttributeTok{{-}pbc}\NormalTok{ whole}
\end{Highlighting}
\end{Shaded}

Trayectoria sin saltos, útil para difusión:

\begin{Shaded}
\begin{Highlighting}[]
\BuiltInTok{echo}\NormalTok{ System }\KeywordTok{|}
\ExtensionTok{gmx}\NormalTok{ trjconv     }\AttributeTok{{-}s}\NormalTok{ md.tpr     }\AttributeTok{{-}f}\NormalTok{ md\_whole.xtc     }\AttributeTok{{-}o}\NormalTok{ md\_nojump.xtc     }\AttributeTok{{-}pbc}\NormalTok{ nojump}
\end{Highlighting}
\end{Shaded}

Centre el complejo y lleve las moléculas a una caja compacta:

\begin{Shaded}
\begin{Highlighting}[]
\KeywordTok{(}\BuiltInTok{echo}\NormalTok{ Protein\_LIG}\KeywordTok{;} \BuiltInTok{echo}\NormalTok{ System}\KeywordTok{)} \KeywordTok{|}
\ExtensionTok{gmx}\NormalTok{ trjconv     }\AttributeTok{{-}s}\NormalTok{ md.tpr     }\AttributeTok{{-}f}\NormalTok{ md\_whole.xtc     }\AttributeTok{{-}o}\NormalTok{ md\_center.xtc     }\AttributeTok{{-}center}     \AttributeTok{{-}pbc}\NormalTok{ mol     }\AttributeTok{{-}ur}\NormalTok{ compact     }\AttributeTok{{-}n}\NormalTok{ index.ndx}
\end{Highlighting}
\end{Shaded}

Elimine rotación y traslación ajustando el backbone:

\begin{Shaded}
\begin{Highlighting}[]
\KeywordTok{(}\BuiltInTok{echo}\NormalTok{ Backbone}\KeywordTok{;} \BuiltInTok{echo}\NormalTok{ System}\KeywordTok{)} \KeywordTok{|}
\ExtensionTok{gmx}\NormalTok{ trjconv     }\AttributeTok{{-}s}\NormalTok{ md.tpr     }\AttributeTok{{-}f}\NormalTok{ md\_center.xtc     }\AttributeTok{{-}o}\NormalTok{ md\_fit.xtc     }\AttributeTok{{-}fit}\NormalTok{ rot+trans     }\AttributeTok{{-}n}\NormalTok{ index.ndx}
\end{Highlighting}
\end{Shaded}

El orden importa: no aplique \textbf{-pbc nojump} después de centrar.
Inspeccione visualmente la trayectoria final.

Extracción de una estructura a 50 ns:

\begin{Shaded}
\begin{Highlighting}[]
\BuiltInTok{echo}\NormalTok{ Protein\_LIG }\KeywordTok{|}
\ExtensionTok{gmx}\NormalTok{ trjconv     }\AttributeTok{{-}s}\NormalTok{ md.tpr     }\AttributeTok{{-}f}\NormalTok{ md\_fit.xtc     }\AttributeTok{{-}o}\NormalTok{ frame\_50ns.pdb     }\AttributeTok{{-}dump}\NormalTok{ 50     }\AttributeTok{{-}tu}\NormalTok{ ns     }\AttributeTok{{-}n}\NormalTok{ index.ndx}
\end{Highlighting}
\end{Shaded}

\subsection{Análisis de resultados}\label{anuxe1lisis-de-resultados}

Use el mismo intervalo de producción, tratamiento de PBC y selección en
todas las réplicas. Los archivos XVG son texto y pueden analizarse con
Grace, Python, R u otro programa.

\subsubsection{RMSD}\label{rmsd}

El RMSD mide la desviación respecto de una referencia después del
ajuste:

\[
\mathrm{RMSD}(t)=
\sqrt{\frac{1}{M}\sum_i m_i
\left\|\mathbf r_i(t)-\mathbf r_i^{\mathrm{ref}}\right\|^2}
\]

donde \textbf{M} es la masa total de los átomos seleccionados. Una
meseta indica estabilidad relativa frente a esa referencia, no
convergencia termodinámica.

\begin{Shaded}
\begin{Highlighting}[]
\FunctionTok{mkdir} \AttributeTok{{-}p}\NormalTok{ ../06\_analisis/rmsd}

\KeywordTok{(}\BuiltInTok{echo}\NormalTok{ Backbone}\KeywordTok{;} \BuiltInTok{echo}\NormalTok{ Backbone}\KeywordTok{)} \KeywordTok{|}
\ExtensionTok{gmx}\NormalTok{ rms     }\AttributeTok{{-}s}\NormalTok{ md.tpr     }\AttributeTok{{-}f}\NormalTok{ md\_fit.xtc     }\AttributeTok{{-}n}\NormalTok{ index.ndx     }\AttributeTok{{-}o}\NormalTok{ ../06\_analisis/rmsd/rmsd\_backbone.xvg     }\AttributeTok{{-}tu}\NormalTok{ ns}
\end{Highlighting}
\end{Shaded}

RMSD del ligando después de ajustar la proteína:

\begin{Shaded}
\begin{Highlighting}[]
\KeywordTok{(}\BuiltInTok{echo}\NormalTok{ Backbone}\KeywordTok{;} \BuiltInTok{echo}\NormalTok{ LIG}\KeywordTok{)} \KeywordTok{|}
\ExtensionTok{gmx}\NormalTok{ rms     }\AttributeTok{{-}s}\NormalTok{ md.tpr     }\AttributeTok{{-}f}\NormalTok{ md\_fit.xtc     }\AttributeTok{{-}n}\NormalTok{ index.ndx     }\AttributeTok{{-}o}\NormalTok{ ../06\_analisis/rmsd/rmsd\_ligand\_fit\_protein.xvg     }\AttributeTok{{-}tu}\NormalTok{ ns}
\end{Highlighting}
\end{Shaded}

\subsubsection{Radio de giro}\label{radio-de-giro}

\begin{Shaded}
\begin{Highlighting}[]
\ExtensionTok{gmx}\NormalTok{ gyrate     }\AttributeTok{{-}s}\NormalTok{ md.tpr     }\AttributeTok{{-}f}\NormalTok{ md\_fit.xtc     }\AttributeTok{{-}n}\NormalTok{ index.ndx     }\AttributeTok{{-}sel} \StringTok{\textquotesingle{}group "Protein"\textquotesingle{}}     \AttributeTok{{-}o}\NormalTok{ ../06\_analisis/gyrate\_protein.xvg}
\end{Highlighting}
\end{Shaded}

El radio de giro informa sobre la compacidad global. Debe interpretarse
junto con RMSD, estructura secundaria y contactos internos.

\subsubsection{Distancias y contactos
proteína--ligando}\label{distancias-y-contactos-proteuxednaligando}

Distancia entre el centro geométrico del ligando y un átomo de
referencia:

\begin{Shaded}
\begin{Highlighting}[]
\ExtensionTok{gmx}\NormalTok{ distance     }\AttributeTok{{-}s}\NormalTok{ md.tpr     }\AttributeTok{{-}f}\NormalTok{ md\_fit.xtc     }\AttributeTok{{-}n}\NormalTok{ index.ndx     }\AttributeTok{{-}select} \StringTok{\textquotesingle{}com of group "LIG" plus com of resid 123 and name CA\textquotesingle{}}     \AttributeTok{{-}oall}\NormalTok{ ../06\_analisis/dist\_lig\_res123.xvg}
\end{Highlighting}
\end{Shaded}

El residuo 123 es un ejemplo y debe reemplazarse. Para distancia mínima
y número de contactos:

\begin{Shaded}
\begin{Highlighting}[]
\KeywordTok{(}\BuiltInTok{echo}\NormalTok{ Protein}\KeywordTok{;} \BuiltInTok{echo}\NormalTok{ LIG}\KeywordTok{)} \KeywordTok{|}
\ExtensionTok{gmx}\NormalTok{ mindist     }\AttributeTok{{-}s}\NormalTok{ md.tpr     }\AttributeTok{{-}f}\NormalTok{ md\_fit.xtc     }\AttributeTok{{-}n}\NormalTok{ index.ndx     }\AttributeTok{{-}od}\NormalTok{ ../06\_analisis/mindist\_protein\_lig.xvg     }\AttributeTok{{-}on}\NormalTok{ ../06\_analisis/contacts\_protein\_lig.xvg     }\AttributeTok{{-}d}\NormalTok{ 0.4}
\end{Highlighting}
\end{Shaded}

\textbf{-d 0.4} establece un umbral de contacto de 0.4 nm, equivalente a
4 Å. Debe informarse el umbral utilizado.

\subsubsection{Puentes de hidrógeno}\label{puentes-de-hidruxf3geno}

GROMACS 2026 incluye la implementación moderna de \textbf{gmx hbond},
incorporada inicialmente en GROMACS 2024:

\begin{Shaded}
\begin{Highlighting}[]
\ExtensionTok{gmx}\NormalTok{ hbond     }\AttributeTok{{-}s}\NormalTok{ md.tpr     }\AttributeTok{{-}f}\NormalTok{ md\_fit.xtc     }\AttributeTok{{-}n}\NormalTok{ index.ndx     }\AttributeTok{{-}r} \StringTok{\textquotesingle{}group "Protein"\textquotesingle{}}     \AttributeTok{{-}t} \StringTok{\textquotesingle{}group "LIG"\textquotesingle{}}     \AttributeTok{{-}num}\NormalTok{ ../06\_analisis/hbonds\_protein\_lig.xvg}
\end{Highlighting}
\end{Shaded}

Las selecciones de referencia y objetivo deben ser idénticas o no
solaparse. Los valores recomendados por la herramienta son 0.35 nm para
distancia y 30 grados para el criterio angular. Si se modifican, deben
reportarse.

\subsubsection{Contactos iónicos}\label{contactos-iuxf3nicos}

Un contacto corto entre grupos cargados puede estudiarse mediante
selecciones explícitas:

\begin{Shaded}
\begin{Highlighting}[]
\ExtensionTok{gmx}\NormalTok{ pairdist     }\AttributeTok{{-}s}\NormalTok{ md.tpr     }\AttributeTok{{-}f}\NormalTok{ md\_fit.xtc     }\AttributeTok{{-}n}\NormalTok{ index.ndx     }\AttributeTok{{-}ref} \StringTok{\textquotesingle{}group "Protein\_charged"\textquotesingle{}}     \AttributeTok{{-}sel} \StringTok{\textquotesingle{}group "LIG\_charged"\textquotesingle{}}     \AttributeTok{{-}type}\NormalTok{ min     }\AttributeTok{{-}o}\NormalTok{ ../06\_analisis/ion\_pairs.xvg}
\end{Highlighting}
\end{Shaded}

Los grupos \textbf{Protein\_charged} y \textbf{LIG\_charged} deben
construirse previamente según los átomos efectivamente cargados. Una
distancia corta no demuestra por sí sola una interacción energéticamente
favorable.

\subsubsection{Área accesible al
solvente}\label{uxe1rea-accesible-al-solvente}

\begin{Shaded}
\begin{Highlighting}[]
\ExtensionTok{gmx}\NormalTok{ sasa     }\AttributeTok{{-}s}\NormalTok{ md.tpr     }\AttributeTok{{-}f}\NormalTok{ md\_fit.xtc     }\AttributeTok{{-}n}\NormalTok{ index.ndx     }\AttributeTok{{-}surface} \StringTok{\textquotesingle{}group "Protein\_LIG"\textquotesingle{}}     \AttributeTok{{-}output} \StringTok{\textquotesingle{}group "Protein\_LIG"\textquotesingle{}}     \AttributeTok{{-}o}\NormalTok{ ../06\_analisis/sasa\_complex.xvg     }\AttributeTok{{-}or}\NormalTok{ ../06\_analisis/sasa\_per\_residue.xvg}
\end{Highlighting}
\end{Shaded}

SASA depende de la selección, los radios atómicos y la sonda. Una
disminución del área expuesta del ligando puede acompañar su
enterramiento, pero no equivale a energía de unión.

\subsubsection{Variables
termodinámicas}\label{variables-termodinuxe1micas}

\begin{Shaded}
\begin{Highlighting}[]
\KeywordTok{(}\BuiltInTok{echo}\NormalTok{ Temperature}\KeywordTok{;} \BuiltInTok{echo}\NormalTok{ Pressure}\KeywordTok{;} \BuiltInTok{echo}\NormalTok{ Density}\KeywordTok{;} \BuiltInTok{echo}\NormalTok{ Volume}\KeywordTok{;} \BuiltInTok{echo}\NormalTok{ Potential}\KeywordTok{;} \BuiltInTok{echo}\NormalTok{ 0}\KeywordTok{)} \KeywordTok{|}
\ExtensionTok{gmx}\NormalTok{ energy     }\AttributeTok{{-}f}\NormalTok{ md.edr     }\AttributeTok{{-}o}\NormalTok{ ../06\_analisis/thermodynamics.xvg}
\end{Highlighting}
\end{Shaded}

Los nombres disponibles dependen del contenido de EDR. Evalúe promedios
por bloques y deriva temporal. La presión instantánea suele presentar
fluctuaciones grandes.

\subsubsection{Desplazamiento cuadrático
medio}\label{desplazamiento-cuadruxe1tico-medio}

\textbf{gmx msd} usa selecciones modernas:

\begin{Shaded}
\begin{Highlighting}[]
\ExtensionTok{gmx}\NormalTok{ msd     }\AttributeTok{{-}s}\NormalTok{ md.tpr     }\AttributeTok{{-}f}\NormalTok{ md\_nojump.xtc     }\AttributeTok{{-}n}\NormalTok{ index.ndx     }\AttributeTok{{-}sel} \StringTok{\textquotesingle{}group "LIG"\textquotesingle{}}     \AttributeTok{{-}o}\NormalTok{ ../06\_analisis/msd\_ligand.xvg}
\end{Highlighting}
\end{Shaded}

Difusión lateral en el plano XY:

\begin{Shaded}
\begin{Highlighting}[]
\ExtensionTok{gmx}\NormalTok{ msd     }\AttributeTok{{-}s}\NormalTok{ md.tpr     }\AttributeTok{{-}f}\NormalTok{ md\_nojump.xtc     }\AttributeTok{{-}n}\NormalTok{ index.ndx     }\AttributeTok{{-}sel} \StringTok{\textquotesingle{}group "LIG"\textquotesingle{}}     \AttributeTok{{-}lateral}\NormalTok{ z     }\AttributeTok{{-}o}\NormalTok{ ../06\_analisis/msd\_ligand\_xy.xvg}
\end{Highlighting}
\end{Shaded}

La estimación del coeficiente de difusión se basa en la región lineal de
la relación de Einstein:

\[
\left\langle |\mathbf r(t)-\mathbf r(0)|^2\right\rangle=2dDt
\]

donde \textbf{d} es la dimensionalidad: 3 para difusión tridimensional y
2 para difusión lateral. Debe elegirse un intervalo de ajuste en el
régimen difusivo; el tramo inicial balístico y las regiones con muestreo
insuficiente sesgan el resultado.

\subsubsection{Mapas de densidad}\label{mapas-de-densidad}

\begin{Shaded}
\begin{Highlighting}[]
\BuiltInTok{echo}\NormalTok{ LIG }\KeywordTok{|}
\ExtensionTok{gmx}\NormalTok{ densmap     }\AttributeTok{{-}s}\NormalTok{ md.tpr     }\AttributeTok{{-}f}\NormalTok{ md\_fit.xtc     }\AttributeTok{{-}n}\NormalTok{ index.ndx     }\AttributeTok{{-}od}\NormalTok{ ../06\_analisis/density\_ligand.xpm     }\AttributeTok{{-}aver}\NormalTok{ z}
\end{Highlighting}
\end{Shaded}

El sistema debe estar alineado previamente. Informe el eje promediado,
la resolución de la grilla y la selección.

\subsubsection{Fluctuación de los
residuos}\label{fluctuaciuxf3n-de-los-residuos}

La RMSF mide la fluctuación de cada átomo alrededor de su posición
promedio:

\[
\mathrm{RMSF}_i=
\sqrt{
\left\langle
\left\|\mathbf r_i(t)-\langle\mathbf r_i\rangle\right\|^2
\right\rangle}
\]

Primero ajuste la trayectoria sobre una región estructuralmente estable:

\begin{Shaded}
\begin{Highlighting}[]
\FunctionTok{mkdir} \AttributeTok{{-}p}\NormalTok{ ../06\_analisis/rmsf}

\BuiltInTok{echo}\NormalTok{ Backbone }\KeywordTok{|}
\ExtensionTok{gmx}\NormalTok{ trjconv     }\AttributeTok{{-}s}\NormalTok{ md.tpr     }\AttributeTok{{-}f}\NormalTok{ md\_center.xtc     }\AttributeTok{{-}o}\NormalTok{ ../06\_analisis/rmsf/md\_fit\_backbone.xtc     }\AttributeTok{{-}fit}\NormalTok{ rot+trans     }\AttributeTok{{-}n}\NormalTok{ index.ndx}
\end{Highlighting}
\end{Shaded}

Calcule RMSF por residuo usando C-alpha:

\begin{Shaded}
\begin{Highlighting}[]
\BuiltInTok{echo}\NormalTok{ C{-}alpha }\KeywordTok{|}
\ExtensionTok{gmx}\NormalTok{ rmsf     }\AttributeTok{{-}s}\NormalTok{ md.tpr     }\AttributeTok{{-}f}\NormalTok{ ../06\_analisis/rmsf/md\_fit\_backbone.xtc     }\AttributeTok{{-}n}\NormalTok{ index.ndx     }\AttributeTok{{-}o}\NormalTok{ ../06\_analisis/rmsf/rmsf\_calpha.xvg     }\AttributeTok{{-}res}     \AttributeTok{{-}oq}\NormalTok{ ../06\_analisis/rmsf/rmsf\_calpha\_bfactor.pdb}
\end{Highlighting}
\end{Shaded}

\textbf{-oq} escribe los valores convertidos al campo B del PDB. La
relación isotrópica es:

\[
B_i=\frac{8\pi^2}{3}\mathrm{RMSF}_i^2
\]

Para evaluar estabilidad temporal, compare bloques de igual duración:

\begin{Shaded}
\begin{Highlighting}[]
\BuiltInTok{echo}\NormalTok{ C{-}alpha }\KeywordTok{|}
\ExtensionTok{gmx}\NormalTok{ rmsf     }\AttributeTok{{-}s}\NormalTok{ md.tpr     }\AttributeTok{{-}f}\NormalTok{ ../06\_analisis/rmsf/md\_fit\_backbone.xtc     }\AttributeTok{{-}n}\NormalTok{ index.ndx     }\AttributeTok{{-}b}\NormalTok{ 20     }\AttributeTok{{-}e}\NormalTok{ 40     }\AttributeTok{{-}tu}\NormalTok{ ns     }\AttributeTok{{-}res}     \AttributeTok{{-}o}\NormalTok{ ../06\_analisis/rmsf/rmsf\_20\_40ns.xvg}

\BuiltInTok{echo}\NormalTok{ C{-}alpha }\KeywordTok{|}
\ExtensionTok{gmx}\NormalTok{ rmsf     }\AttributeTok{{-}s}\NormalTok{ md.tpr     }\AttributeTok{{-}f}\NormalTok{ ../06\_analisis/rmsf/md\_fit\_backbone.xtc     }\AttributeTok{{-}n}\NormalTok{ index.ndx     }\AttributeTok{{-}b}\NormalTok{ 40     }\AttributeTok{{-}e}\NormalTok{ 60     }\AttributeTok{{-}tu}\NormalTok{ ns     }\AttributeTok{{-}res}     \AttributeTok{{-}o}\NormalTok{ ../06\_analisis/rmsf/rmsf\_40\_60ns.xvg}
\end{Highlighting}
\end{Shaded}

Picos persistentes suelen corresponder a terminales, bucles o regiones
expuestas. Cambios localizados cerca del sitio de unión pueden sugerir
estabilización o reorganización, pero deben contrastarse con contactos,
estructura secundaria, RMSD y réplicas independientes.

No compare RMSF obtenidas con distinta alineación, selección, duración o
intervalo temporal.

\subsection{Control de calidad y
reproducibilidad}\label{control-de-calidad-y-reproducibilidad}

Conserve:

\begin{itemize}
\tightlist
\item
  versión exacta de GROMACS y salida de \textbf{gmx --version};
\item
  estructuras iniciales y decisiones de protonación;
\item
  campo de fuerza, modelo de agua y método de parametrización del
  ligando;
\item
  archivos MDP, TOP, ITP y NDX;
\item
  TPR, checkpoints y registros;
\item
  comandos ejecutados;
\item
  semillas y configuración de hardware;
\item
  intervalo y selección utilizados en cada análisis.
\end{itemize}

Use réplicas independientes cuando la conclusión dependa del muestreo.
RMSD, RMSF, contactos, puentes de hidrógeno y distancias no son
estimaciones directas de afinidad ni energía libre de unión.

\subsection{Fuentes}\label{fuentes-1}

\begin{enumerate}
\def\labelenumi{\arabic{enumi}.}
\tightlist
\item
  \href{https://manual.gromacs.org/current/reference-manual/file-formats.html}{Formatos
  de archivo de GROMACS 2026.3}
\item
  \href{https://manual.gromacs.org/current/reference-manual/topologies/topologies.html}{Topologías
  de GROMACS}
\item
  \href{https://manual.gromacs.org/current/reference-manual/analysis/analysis.html}{Manual
  de análisis}
\item
  \href{https://manual.gromacs.org/current/user-guide/mdp-options.html}{Opciones
  de archivos MDP}
\item
  \href{https://manual.gromacs.org/current/onlinehelp/gmx.html}{Referencia
  de comandos de GROMACS}
\item
  \href{https://manual.gromacs.org/current/user-guide/flow.html}{Flujo
  general de una simulación en GROMACS}
\item
  \href{https://manual.gromacs.org/current/reference-manual/definitions.html}{Definiciones
  y unidades}
\item
  \href{https://manual.gromacs.org/current/reference-manual/algorithms/molecular-dynamics.html}{Dinámica
  molecular}
\item
  \href{https://manual.gromacs.org/current/reference-manual/algorithms/periodic-boundary-conditions.html}{Condiciones
  periódicas de contorno}
\item
  \href{https://manual.gromacs.org/current/reference-manual/algorithms/constraint-algorithms.html}{Algoritmos
  de restricciones}
\item
  \href{https://alanwilter.github.io/acpype/}{Documentación oficial de
  ACPYPE}
\item
  \href{https://github.com/alanwilter/acpype}{Repositorio oficial de
  ACPYPE}
\item
  \href{https://doi.org/10.1186/1756-0500-5-367}{Sousa da Silva y
  Vranken, ACPYPE --- AnteChamber PYthon Parser interfacE}
\item
  \href{https://doi.org/10.1002/jcc.20035}{Wang et al., Development and
  testing of a General Amber Force Field}
\item
  \href{https://ambermd.org/AmberTools.php}{AmberTools}
\end{enumerate}

\section{Simulación de una caja de
agua}\label{simulaciuxf3n-de-una-caja-de-agua}

Vamos a simular un sistema compuesto por agua. Puede ser que nos sirva
para simular sistemas sencillos y así testear la instalación de GROMACS
o estudiar variables.

\subsection{Múltiples moléculas de ligandos y de
proteínas}\label{muxfaltiples-moluxe9culas-de-ligandos-y-de-proteuxednas}

Los sistemas con varias cadenas proteicas y varios ligandos requieren
distinguir tres conceptos:

\begin{enumerate}
\def\labelenumi{\arabic{enumi}.}
\tightlist
\item
  \textbf{Copia molecular:} instancia concreta con coordenadas propias.
\item
  \textbf{Tipo molecular:} definición de la molécula en una sección
  \texttt{{[}\ moleculetype\ {]}}.
\item
  \textbf{Nombre de residuo:} etiqueta empleada en las coordenadas y en
  las selecciones.
\end{enumerate}

Dos ligandos químicamente idénticos, con los mismos átomos, orden
atómico, enlaces, cargas y parámetros, pueden ser dos copias del mismo
tipo molecular. Sus posiciones y conformaciones iniciales pueden ser
diferentes. En cambio, ligandos distintos o copias que deban recibir
topologías diferentes necesitan tipos moleculares separados.

\subsubsection{Caso 1: varias copias de un ligando
idéntico}\label{caso-1-varias-copias-de-un-ligando-iduxe9ntico}

La estructura PDB 1OKE contiene dos cadenas proteicas y varias moléculas
no proteicas, entre ellas dos moléculas BOG. Si ambas copias de BOG
tienen la misma composición y parametrización, se necesita una sola
topología molecular para BOG.

Antes de parametrizar:

\begin{itemize}
\tightlist
\item
  extraiga una copia completa del ligando;
\item
  revise protonación, carga formal, estereoquímica y orden de enlaces;
\item
  confirme que ambas copias tienen los mismos átomos;
\item
  conserve las coordenadas originales de cada copia;
\item
  use el mismo orden atómico en las coordenadas y en la topología.
\end{itemize}

Puede asignarse \textbf{BOG} como nombre de residuo y como nombre del
tipo molecular. Las dos copias se distinguen mediante su número de
residuo y, si el formato lo permite, su cadena.

Ejemplo simplificado de coordenadas:

\begin{Shaded}
\begin{Highlighting}[]
\NormalTok{HETATM    1  C1  BOG C   1     {-}14.258  79.953  45.302  1.00  0.00           C}
\NormalTok{HETATM    2  O1  BOG C   1     {-}13.074  79.344  45.814  1.00  0.00           O}
\NormalTok{...}
\NormalTok{TER}
\NormalTok{HETATM   49  C1  BOG D   2     {-}17.418  58.852   3.720  1.00  0.00           C}
\NormalTok{HETATM   50  O1  BOG D   2     {-}16.140  59.295   3.262  1.00  0.00           O}
\NormalTok{...}
\end{Highlighting}
\end{Shaded}

Los números de átomo deben ser únicos dentro del PDB. Los números de
residuo o identificadores de cadena diferentes facilitan el análisis,
aunque GROMACS asigna la topología principalmente según el orden de las
moléculas.

Topología molecular:

\begin{Shaded}
\begin{Highlighting}[]
\KeywordTok{[ moleculetype ]}
\CommentTok{; nombre   nrexcl}
\DataTypeTok{BOG        3}

\KeywordTok{[ atoms ]}
\CommentTok{; nr  tipo  resnr  residuo  átomo  cgnr  carga  masa}
\DataTypeTok{1     ...   1      BOG      C1     1     ...    12.011}
\DataTypeTok{2     ...   1      BOG      O1     2     ...    15.999}
\CommentTok{; ...}
\end{Highlighting}
\end{Shaded}

En \textbf{topol.top} se incluye una sola vez:

\begin{Shaded}
\begin{Highlighting}[]
\CommentTok{; Parámetros generales}
\CommentTok{\#include "charmm36{-}jul2022.ff/forcefield.itp"}

\CommentTok{; Parámetros adicionales del ligando, si corresponden}
\CommentTok{\#include "bog\_atomtypes.itp"}

\CommentTok{; Definición del tipo molecular BOG}
\CommentTok{\#include "bog.itp"}

\CommentTok{; Cadenas proteicas generadas por pdb2gmx}
\CommentTok{\#include "topol\_Protein\_chain\_A.itp"}
\CommentTok{\#include "topol\_Protein\_chain\_B.itp"}
\end{Highlighting}
\end{Shaded}

Al final del archivo:

\begin{Shaded}
\begin{Highlighting}[]
\KeywordTok{[ system ]}
\DataTypeTok{Proteína 1OKE con dos moléculas BOG}

\KeywordTok{[ molecules ]}
\CommentTok{; tipo molecular       cantidad}
\DataTypeTok{Protein\_chain\_A         1}
\DataTypeTok{Protein\_chain\_B         1}
\DataTypeTok{BOG                     2}
\end{Highlighting}
\end{Shaded}

La sección \texttt{{[}\ molecules\ {]}} indica que existen dos
instancias de BOG. No deben incluirse dos copias idénticas de
\textbf{bog.itp} porque se redefinirían los mismos tipos o el mismo
\texttt{{[}\ moleculetype\ {]}}.

El orden de las coordenadas debe ser:

\begin{Shaded}
\begin{Highlighting}[]
\NormalTok{Protein\_chain\_A}
\NormalTok{Protein\_chain\_B}
\NormalTok{BOG, copia 1}
\NormalTok{BOG, copia 2}
\end{Highlighting}
\end{Shaded}

Este orden debe coincidir exactamente con \texttt{{[}\ molecules\ {]}}.

\subsubsection{Cuándo son necesarias topologías
separadas}\label{cuuxe1ndo-son-necesarias-topologuxedas-separadas}

Use tipos moleculares diferentes, por ejemplo \textbf{LIGA} y
\textbf{LIGB}, cuando:

\begin{itemize}
\tightlist
\item
  los ligandos son químicamente distintos;
\item
  tienen diferente protonación o carga;
\item
  una copia contiene átomos o enlaces diferentes;
\item
  se parametrizaron con términos diferentes;
\item
  se aplicarán restricciones posicionales distintas;
\item
  se realizarán transformaciones alquí\-micas o cálculos de energía
  libre diferentes por copia.
\end{itemize}

En ese caso:

\begin{Shaded}
\begin{Highlighting}[]
\CommentTok{\#include "ligand\_A.itp"}
\CommentTok{\#include "ligand\_B.itp"}

\KeywordTok{[ molecules ]}
\DataTypeTok{Protein\_chain\_A    1}
\DataTypeTok{Protein\_chain\_B    1}
\DataTypeTok{LIGA               1}
\DataTypeTok{LIGB               1}
\end{Highlighting}
\end{Shaded}

Duplicar una topología solo para cambiar el nombre del residuo suele ser
innecesario. Si se crean \textbf{LIGA} y \textbf{LIGB}, ambos archivos
deben tener nombres de \texttt{{[}\ moleculetype\ {]}} diferentes y
todos sus índices deben seguir siendo locales a cada molécula.

\subsubsection{Sistemas con muchas cadenas: ejemplo
6BKK}\label{sistemas-con-muchas-cadenas-ejemplo-6bkk}

La estructura 6BKK corresponde al dominio transmembrana M2 de influenza
A unido a amantadina. El sistema biológico y la unidad cristalográfica
deben revisarse antes de simular: no debe suponerse que todas las
cadenas del PDB forman la unidad funcional.

Prepare la proteína sin eliminar ligandos o cofactores hasta haber
registrado su identidad y posición. Luego genere la topología de las
cadenas:

\begin{Shaded}
\begin{Highlighting}[]
\ExtensionTok{gmx}\NormalTok{ pdb2gmx     }\AttributeTok{{-}f}\NormalTok{ protein\_only.pdb     }\AttributeTok{{-}o}\NormalTok{ protein\_processed.gro     }\AttributeTok{{-}p}\NormalTok{ topol.top     }\AttributeTok{{-}i}\NormalTok{ posre\_protein.itp     }\AttributeTok{{-}water}\NormalTok{ tip3p}
\end{Highlighting}
\end{Shaded}

Seleccione un campo de fuerza compatible con la topología de amantadina.
No use automáticamente CHARMM27 por aparecer en protocolos antiguos;
registre el campo de fuerza concreto disponible en la instalación.

Si \textbf{pdb2gmx} separa las cadenas en archivos ITP,
\textbf{topol.top} puede contener:

\begin{Shaded}
\begin{Highlighting}[]
\CommentTok{\#include "topol\_Protein\_chain\_A.itp"}
\CommentTok{\#include "topol\_Protein\_chain\_B.itp"}
\CommentTok{\#include "topol\_Protein\_chain\_C.itp"}
\CommentTok{\#include "topol\_Protein\_chain\_D.itp"}
\CommentTok{\#include "topol\_Protein\_chain\_E.itp"}
\CommentTok{\#include "topol\_Protein\_chain\_F.itp"}
\CommentTok{\#include "topol\_Protein\_chain\_G.itp"}
\CommentTok{\#include "topol\_Protein\_chain\_H.itp"}
\end{Highlighting}
\end{Shaded}

La cantidad y el orden deben verificarse en la salida real. No copie
esta lista a otro sistema.

Si existen dos amantadinas idénticas:

\begin{Shaded}
\begin{Highlighting}[]
\KeywordTok{[ molecules ]}
\DataTypeTok{Protein\_chain\_A    1}
\DataTypeTok{Protein\_chain\_B    1}
\DataTypeTok{Protein\_chain\_C    1}
\DataTypeTok{Protein\_chain\_D    1}
\DataTypeTok{Protein\_chain\_E    1}
\DataTypeTok{Protein\_chain\_F    1}
\DataTypeTok{Protein\_chain\_G    1}
\DataTypeTok{Protein\_chain\_H    1}
\DataTypeTok{AMA                2}
\end{Highlighting}
\end{Shaded}

Cada copia debe tener coordenadas propias, números de átomo válidos y el
mismo orden interno que \textbf{ama.itp}.

\subsubsection{Construcción y validación de las
coordenadas}\label{construcciuxf3n-y-validaciuxf3n-de-las-coordenadas}

Después de combinar proteína y ligandos:

\begin{Shaded}
\begin{Highlighting}[]
\ExtensionTok{gmx}\NormalTok{ editconf     }\AttributeTok{{-}f}\NormalTok{ protein\_ligands.pdb     }\AttributeTok{{-}o}\NormalTok{ complex.gro}
\end{Highlighting}
\end{Shaded}

Revise el archivo:

\begin{Shaded}
\begin{Highlighting}[]
\ExtensionTok{gmx}\NormalTok{ check }\AttributeTok{{-}f}\NormalTok{ complex.gro}
\end{Highlighting}
\end{Shaded}

Genere un TPR de validación antes de solvatar. Puede usarse un MDP
mínimo:

\begin{Shaded}
\begin{Highlighting}[]
\DataTypeTok{integrator      }\OtherTok{=}\StringTok{ steep}
\DataTypeTok{nsteps          }\OtherTok{=}\StringTok{ }\DecValTok{0}
\DataTypeTok{cutoff{-}scheme   }\OtherTok{=}\StringTok{ Verlet}
\DataTypeTok{coulombtype     }\OtherTok{=}\StringTok{ PME}
\DataTypeTok{rcoulomb        }\OtherTok{=}\StringTok{ }\FloatTok{1.0}
\DataTypeTok{rvdw            }\OtherTok{=}\StringTok{ }\FloatTok{1.0}
\DataTypeTok{pbc             }\OtherTok{=}\StringTok{ xyz}
\end{Highlighting}
\end{Shaded}

\begin{Shaded}
\begin{Highlighting}[]
\ExtensionTok{gmx}\NormalTok{ grompp     }\AttributeTok{{-}f}\NormalTok{ validate.mdp     }\AttributeTok{{-}c}\NormalTok{ complex.gro     }\AttributeTok{{-}p}\NormalTok{ topol.top     }\AttributeTok{{-}o}\NormalTok{ validate.tpr     }\AttributeTok{{-}pp}\NormalTok{ processed.top}
\end{Highlighting}
\end{Shaded}

La opción \textbf{-pp processed.top} guarda la topología ya expandida
por el preprocesador. Resulta útil para comprobar el orden de las
inclusiones, macros y restricciones activadas.

No continúe si aparece alguno de estos problemas:

\begin{itemize}
\tightlist
\item
  número de coordenadas distinto del número de átomos de la topología;
\item
  tipos atómicos desconocidos;
\item
  parámetros enlazados ausentes;
\item
  carga total inesperada;
\item
  redefiniciones de \texttt{{[}\ atomtypes\ {]}};
\item
  nombres distintos entre \texttt{{[}\ molecules\ {]}} y
  \texttt{{[}\ moleculetype\ {]}}.
\end{itemize}

No use \textbf{-maxwarn} para forzar el TPR sin comprender cada
advertencia.

\subsubsection{Caja, solvatación e
iones}\label{caja-solvataciuxf3n-e-iones}

Para una proteína soluble:

\begin{Shaded}
\begin{Highlighting}[]
\ExtensionTok{gmx}\NormalTok{ editconf     }\AttributeTok{{-}f}\NormalTok{ complex.gro     }\AttributeTok{{-}o}\NormalTok{ complex\_box.gro     }\AttributeTok{{-}c}     \AttributeTok{{-}d}\NormalTok{ 1.0     }\AttributeTok{{-}bt}\NormalTok{ dodecahedron}
\end{Highlighting}
\end{Shaded}

Para un sistema de membrana no debe construirse una caja acuosa genérica
alrededor de la proteína. Primero debe definirse la membrana, su
orientación, composición y espesor mediante un protocolo específico.

Solvatación:

\begin{Shaded}
\begin{Highlighting}[]
\ExtensionTok{gmx}\NormalTok{ solvate     }\AttributeTok{{-}cp}\NormalTok{ complex\_box.gro     }\AttributeTok{{-}cs}\NormalTok{ spc216.gro     }\AttributeTok{{-}o}\NormalTok{ complex\_solv.gro     }\AttributeTok{{-}p}\NormalTok{ topol.top}
\end{Highlighting}
\end{Shaded}

Preprocesamiento para iones:

\begin{Shaded}
\begin{Highlighting}[]
\ExtensionTok{gmx}\NormalTok{ grompp     }\AttributeTok{{-}f}\NormalTok{ ions.mdp     }\AttributeTok{{-}c}\NormalTok{ complex\_solv.gro     }\AttributeTok{{-}p}\NormalTok{ topol.top     }\AttributeTok{{-}o}\NormalTok{ ions.tpr}
\end{Highlighting}
\end{Shaded}

Neutralización y NaCl 0.15 mol L⁻¹:

\begin{Shaded}
\begin{Highlighting}[]
\ExtensionTok{gmx}\NormalTok{ genion     }\AttributeTok{{-}s}\NormalTok{ ions.tpr     }\AttributeTok{{-}o}\NormalTok{ complex\_solv\_ions.gro     }\AttributeTok{{-}p}\NormalTok{ topol.top     }\AttributeTok{{-}pname}\NormalTok{ NA     }\AttributeTok{{-}nname}\NormalTok{ CL     }\AttributeTok{{-}neutral}     \AttributeTok{{-}conc}\NormalTok{ 0.15}
\end{Highlighting}
\end{Shaded}

Seleccione el grupo de solvente, normalmente \textbf{SOL}. El número del
grupo depende del sistema.

\subsubsection{Restricciones de
posición}\label{restricciones-de-posiciuxf3n}

\paragraph{Múltiples copias del mismo tipo
molecular}\label{muxfaltiples-copias-del-mismo-tipo-molecular}

Si dos ligandos son instancias del mismo
\texttt{{[}\ moleculetype\ {]}}, un único archivo de restricciones
incluido dentro de esa topología se aplica a todas las copias:

\begin{Shaded}
\begin{Highlighting}[]
\ExtensionTok{gmx}\NormalTok{ genrestr     }\AttributeTok{{-}f}\NormalTok{ bog\_only.gro     }\AttributeTok{{-}o}\NormalTok{ posre\_bog.itp     }\AttributeTok{{-}fc}\NormalTok{ 1000 1000 1000}
\end{Highlighting}
\end{Shaded}

Incluya las restricciones inmediatamente después de la definición de
BOG:

\begin{Shaded}
\begin{Highlighting}[]
\CommentTok{\#include "bog.itp"}

\CommentTok{\#ifdef POSRES\_BOG}
\CommentTok{\#include "posre\_bog.itp"}
\CommentTok{\#endif}
\end{Highlighting}
\end{Shaded}

Active la macro durante la equilibración:

\begin{Shaded}
\begin{Highlighting}[]
\DataTypeTok{define }\OtherTok{=}\StringTok{ {-}DPOSRES {-}DPOSRES\_BOG}
\end{Highlighting}
\end{Shaded}

No genere \textbf{posre\_BOG1.itp} y \textbf{posre\_BOG2.itp} para
incluirlos sobre un único tipo molecular. Los índices de las
restricciones son locales al tipo molecular, por lo que ambas copias
usarán la misma selección local.

\paragraph{Restricciones diferentes para cada
copia}\label{restricciones-diferentes-para-cada-copia}

Si BOG1 y BOG2 deben tener restricciones diferentes, defina dos tipos
moleculares:

\begin{Shaded}
\begin{Highlighting}[]
\CommentTok{\#include "bog1.itp"}
\CommentTok{\#ifdef POSRES\_BOG1}
\CommentTok{\#include "posre\_bog1.itp"}
\CommentTok{\#endif}

\CommentTok{\#include "bog2.itp"}
\CommentTok{\#ifdef POSRES\_BOG2}
\CommentTok{\#include "posre\_bog2.itp"}
\CommentTok{\#endif}
\end{Highlighting}
\end{Shaded}

\begin{Shaded}
\begin{Highlighting}[]
\KeywordTok{[ molecules ]}
\DataTypeTok{Protein\_chain\_A    1}
\DataTypeTok{Protein\_chain\_B    1}
\DataTypeTok{BOG1               1}
\DataTypeTok{BOG2               1}
\end{Highlighting}
\end{Shaded}

Esta duplicación aumenta el mantenimiento y solo se justifica si las
copias requieren un tratamiento realmente diferente.

\paragraph{Cadenas proteicas}\label{cadenas-proteicas}

Las restricciones de cada cadena deben estar dentro del ámbito del
\texttt{{[}\ moleculetype\ {]}} correspondiente. Los ITP generados por
\textbf{pdb2gmx} suelen incluir su propio archivo de restricciones
mediante la macro \textbf{POSRES}. Revise cada archivo antes de añadir
nuevas inclusiones.

Ejemplo dentro de \textbf{topol\_Protein\_chain\_A.itp}:

\begin{Shaded}
\begin{Highlighting}[]
\KeywordTok{[ moleculetype ]}
\DataTypeTok{Protein\_chain\_A    3}

\CommentTok{; átomos, enlaces y otros términos}
\CommentTok{; ...}

\CommentTok{\#ifdef POSRES}
\CommentTok{\#include "posre\_Protein\_chain\_A.itp"}
\CommentTok{\#endif}
\end{Highlighting}
\end{Shaded}

No reúna al final de \textbf{topol.top} restricciones de varias
moléculas como si utilizaran índices globales.

\subsubsection{Grupos de índice para varias cadenas y
ligandos}\label{grupos-de-uxedndice-para-varias-cadenas-y-ligandos}

Cree grupos explícitos:

\begin{Shaded}
\begin{Highlighting}[]
\ExtensionTok{gmx}\NormalTok{ make\_ndx     }\AttributeTok{{-}f}\NormalTok{ complex\_solv\_ions.gro     }\AttributeTok{{-}o}\NormalTok{ index.ndx}
\end{Highlighting}
\end{Shaded}

Ejemplo interactivo para dos ligandos con residuo BOG:

\begin{Shaded}
\begin{Highlighting}[]
\NormalTok{r BOG}
\NormalTok{name 18 BOG\_all}
\NormalTok{"Protein" | 18}
\NormalTok{name 19 Protein\_BOG}
\NormalTok{q}
\end{Highlighting}
\end{Shaded}

Los números 18 y 19 son ejemplos. Use los asignados en su sesión.

Las dos copias pueden seleccionarse por número de residuo, cadena o
índice atómico. La sintaxis exacta depende de la información conservada
en el archivo:

\begin{Shaded}
\begin{Highlighting}[]
\ExtensionTok{gmx}\NormalTok{ select     }\AttributeTok{{-}s}\NormalTok{ complex\_solv\_ions.gro     }\AttributeTok{{-}n}\NormalTok{ index.ndx     }\AttributeTok{{-}select} \StringTok{\textquotesingle{}resname BOG\textquotesingle{}}
\end{Highlighting}
\end{Shaded}

Para separar las copias por número de residuo:

\begin{Shaded}
\begin{Highlighting}[]
\ExtensionTok{gmx}\NormalTok{ select     }\AttributeTok{{-}s}\NormalTok{ complex\_solv\_ions.gro     }\AttributeTok{{-}select} \StringTok{\textquotesingle{}resname BOG and resid 1\textquotesingle{}}     \AttributeTok{{-}on}\NormalTok{ bog\_1.ndx}

\ExtensionTok{gmx}\NormalTok{ select     }\AttributeTok{{-}s}\NormalTok{ complex\_solv\_ions.gro     }\AttributeTok{{-}select} \StringTok{\textquotesingle{}resname BOG and resid 2\textquotesingle{}}     \AttributeTok{{-}on}\NormalTok{ bog\_2.ndx}
\end{Highlighting}
\end{Shaded}

Si los números de residuo no son únicos entre cadenas, seleccione además
por identificador de cadena cuando esté disponible o utilice rangos de
índices verificados.

\subsubsection{Minimización y
equilibración}\label{minimizaciuxf3n-y-equilibraciuxf3n}

La minimización sigue el flujo general:

\begin{Shaded}
\begin{Highlighting}[]
\ExtensionTok{gmx}\NormalTok{ grompp     }\AttributeTok{{-}f}\NormalTok{ em.mdp     }\AttributeTok{{-}c}\NormalTok{ complex\_solv\_ions.gro     }\AttributeTok{{-}p}\NormalTok{ topol.top     }\AttributeTok{{-}n}\NormalTok{ index.ndx     }\AttributeTok{{-}o}\NormalTok{ em.tpr}

\ExtensionTok{gmx}\NormalTok{ mdrun }\AttributeTok{{-}deffnm}\NormalTok{ em }\AttributeTok{{-}v}
\end{Highlighting}
\end{Shaded}

NVT con restricciones:

\begin{Shaded}
\begin{Highlighting}[]
\ExtensionTok{gmx}\NormalTok{ grompp     }\AttributeTok{{-}f}\NormalTok{ nvt.mdp     }\AttributeTok{{-}c}\NormalTok{ em.gro     }\AttributeTok{{-}r}\NormalTok{ em.gro     }\AttributeTok{{-}p}\NormalTok{ topol.top     }\AttributeTok{{-}n}\NormalTok{ index.ndx     }\AttributeTok{{-}o}\NormalTok{ nvt.tpr}

\ExtensionTok{gmx}\NormalTok{ mdrun }\AttributeTok{{-}deffnm}\NormalTok{ nvt }\AttributeTok{{-}v}
\end{Highlighting}
\end{Shaded}

NPT continuando coordenadas y velocidades:

\begin{Shaded}
\begin{Highlighting}[]
\ExtensionTok{gmx}\NormalTok{ grompp     }\AttributeTok{{-}f}\NormalTok{ npt.mdp     }\AttributeTok{{-}c}\NormalTok{ nvt.gro     }\AttributeTok{{-}r}\NormalTok{ nvt.gro     }\AttributeTok{{-}t}\NormalTok{ nvt.cpt     }\AttributeTok{{-}p}\NormalTok{ topol.top     }\AttributeTok{{-}n}\NormalTok{ index.ndx     }\AttributeTok{{-}o}\NormalTok{ npt.tpr}

\ExtensionTok{gmx}\NormalTok{ mdrun }\AttributeTok{{-}deffnm}\NormalTok{ npt }\AttributeTok{{-}v}
\end{Highlighting}
\end{Shaded}

No use el barostato Berendsen para obtener un ensamble de producción.
Para equilibración, \textbf{C-rescale} permite controlar la presión y
genera el ensamble correcto. Para producción suele utilizarse
\textbf{Parrinello-Rahman}, según el sistema y el protocolo.

Ejemplo de acoplamiento durante NPT:

\begin{Shaded}
\begin{Highlighting}[]
\DataTypeTok{tcoupl           }\OtherTok{=}\StringTok{ V{-}rescale}
\DataTypeTok{tc{-}grps          }\OtherTok{=}\StringTok{ Protein\_LIG Water\_and\_ions}
\DataTypeTok{tau{-}t            }\OtherTok{=}\StringTok{ 1.0 1.0}
\DataTypeTok{ref{-}t            }\OtherTok{=}\StringTok{ 300 300}

\DataTypeTok{pcoupl           }\OtherTok{=}\StringTok{ C{-}rescale}
\DataTypeTok{pcoupltype       }\OtherTok{=}\StringTok{ isotropic}
\DataTypeTok{tau{-}p            }\OtherTok{=}\StringTok{ }\FloatTok{5.0}
\DataTypeTok{ref{-}p            }\OtherTok{=}\StringTok{ }\FloatTok{1.0}
\DataTypeTok{compressibility  }\OtherTok{=}\StringTok{ 4.5e{-}5}
\end{Highlighting}
\end{Shaded}

\textbf{tau-t} y \textbf{tau-p} se expresan en ps; \textbf{ref-t}, en K;
\textbf{ref-p}, en bar; la compresibilidad, en bar⁻¹.

El tiempo de simulación se calcula como:

\begin{Shaded}
\begin{Highlighting}[]
\NormalTok{tiempo = dt × nsteps}
\end{Highlighting}
\end{Shaded}

Con \textbf{dt = 0.002 ps} y \textbf{nsteps = 500000}, el tiempo es 1000
ps, equivalente a 1 ns. En el texto anterior, \textbf{nsteps = 20000} se
interpretaba erróneamente como 60 ns: en realidad corresponde a 40 ps
con un paso de 2 fs.

\subsubsection{Producción}\label{producciuxf3n}

Genere el TPR sin las macros de restricciones, salvo que formen parte
deliberada del experimento:

\begin{Shaded}
\begin{Highlighting}[]
\ExtensionTok{gmx}\NormalTok{ grompp     }\AttributeTok{{-}f}\NormalTok{ md.mdp     }\AttributeTok{{-}c}\NormalTok{ npt.gro     }\AttributeTok{{-}t}\NormalTok{ npt.cpt     }\AttributeTok{{-}p}\NormalTok{ topol.top     }\AttributeTok{{-}n}\NormalTok{ index.ndx     }\AttributeTok{{-}o}\NormalTok{ md.tpr}

\ExtensionTok{gmx}\NormalTok{ mdrun }\AttributeTok{{-}deffnm}\NormalTok{ md }\AttributeTok{{-}v}
\end{Highlighting}
\end{Shaded}

Compruebe que \textbf{md.mdp} contiene \textbf{continuation = yes} y
\textbf{gen-vel = no}. La producción debe continuar las velocidades de
NPT.

Para revisar métodos y parámetros contenidos en el TPR:

\begin{Shaded}
\begin{Highlighting}[]
\ExtensionTok{gmx}\NormalTok{ dump }\AttributeTok{{-}s}\NormalTok{ md.tpr }\OperatorTok{\textgreater{}}\NormalTok{ md\_tpr\_dump.txt}
\ExtensionTok{gmx}\NormalTok{ report{-}methods }\AttributeTok{{-}s}\NormalTok{ md.tpr }\AttributeTok{{-}o}\NormalTok{ methods.tex}
\end{Highlighting}
\end{Shaded}

\textbf{gmx report-methods} genera una descripción de métodos a partir
del TPR. Debe revisarse antes de incorporarla a un manuscrito.

\subsubsection{Preparación de la
trayectoria}\label{preparaciuxf3n-de-la-trayectoria}

Haga las moléculas completas:

\begin{Shaded}
\begin{Highlighting}[]
\BuiltInTok{echo}\NormalTok{ System }\KeywordTok{|}
\ExtensionTok{gmx}\NormalTok{ trjconv     }\AttributeTok{{-}s}\NormalTok{ md.tpr     }\AttributeTok{{-}f}\NormalTok{ md.xtc     }\AttributeTok{{-}o}\NormalTok{ md\_whole.xtc     }\AttributeTok{{-}pbc}\NormalTok{ whole}
\end{Highlighting}
\end{Shaded}

Centre el complejo completo:

\begin{Shaded}
\begin{Highlighting}[]
\KeywordTok{(}\BuiltInTok{echo}\NormalTok{ Protein\_LIG}\KeywordTok{;} \BuiltInTok{echo}\NormalTok{ System}\KeywordTok{)} \KeywordTok{|}
\ExtensionTok{gmx}\NormalTok{ trjconv     }\AttributeTok{{-}s}\NormalTok{ md.tpr     }\AttributeTok{{-}f}\NormalTok{ md\_whole.xtc     }\AttributeTok{{-}o}\NormalTok{ md\_center.xtc     }\AttributeTok{{-}center}     \AttributeTok{{-}pbc}\NormalTok{ mol     }\AttributeTok{{-}ur}\NormalTok{ compact     }\AttributeTok{{-}n}\NormalTok{ index.ndx}
\end{Highlighting}
\end{Shaded}

Ajuste rotación y traslación respecto del backbone de todas las cadenas:

\begin{Shaded}
\begin{Highlighting}[]
\KeywordTok{(}\BuiltInTok{echo}\NormalTok{ Backbone}\KeywordTok{;} \BuiltInTok{echo}\NormalTok{ System}\KeywordTok{)} \KeywordTok{|}
\ExtensionTok{gmx}\NormalTok{ trjconv     }\AttributeTok{{-}s}\NormalTok{ md.tpr     }\AttributeTok{{-}f}\NormalTok{ md\_center.xtc     }\AttributeTok{{-}o}\NormalTok{ md\_fit.xtc     }\AttributeTok{{-}fit}\NormalTok{ rot+trans     }\AttributeTok{{-}n}\NormalTok{ index.ndx}
\end{Highlighting}
\end{Shaded}

Para oligómeros, inspeccione que el grupo usado para centrar contiene
todas las cadenas funcionales y todos los ligandos relevantes.

\subsubsection{RMSD de cada ligando}\label{rmsd-de-cada-ligando}

Si ambas copias están en un mismo grupo, \textbf{gmx rms} calcula un
único RMSD sobre el conjunto completo; no produce automáticamente una
curva independiente por molécula.

Cree grupos separados para BOG1 y BOG2 y calcule cada curva después de
ajustar sobre la proteína:

\begin{Shaded}
\begin{Highlighting}[]
\KeywordTok{(}\BuiltInTok{echo}\NormalTok{ Backbone}\KeywordTok{;} \BuiltInTok{echo}\NormalTok{ BOG1}\KeywordTok{)} \KeywordTok{|}
\ExtensionTok{gmx}\NormalTok{ rms     }\AttributeTok{{-}s}\NormalTok{ md.tpr     }\AttributeTok{{-}f}\NormalTok{ md\_fit.xtc     }\AttributeTok{{-}n}\NormalTok{ index.ndx     }\AttributeTok{{-}o}\NormalTok{ rmsd\_bog1.xvg     }\AttributeTok{{-}tu}\NormalTok{ ns}

\KeywordTok{(}\BuiltInTok{echo}\NormalTok{ Backbone}\KeywordTok{;} \BuiltInTok{echo}\NormalTok{ BOG2}\KeywordTok{)} \KeywordTok{|}
\ExtensionTok{gmx}\NormalTok{ rms     }\AttributeTok{{-}s}\NormalTok{ md.tpr     }\AttributeTok{{-}f}\NormalTok{ md\_fit.xtc     }\AttributeTok{{-}n}\NormalTok{ index.ndx     }\AttributeTok{{-}o}\NormalTok{ rmsd\_bog2.xvg     }\AttributeTok{{-}tu}\NormalTok{ ns}
\end{Highlighting}
\end{Shaded}

El primer grupo selecciona los átomos usados para el ajuste; el segundo,
los átomos cuyo RMSD se calcula. Para ligandos simétricos, el RMSD
convencional puede mostrar saltos por permutaciones de átomos
equivalentes y debe interpretarse con cuidado.

Comando extra para medir la distancia mínima de cada ligando a la
proteína:

\begin{Shaded}
\begin{Highlighting}[]
\ExtensionTok{gmx}\NormalTok{ pairdist     }\AttributeTok{{-}s}\NormalTok{ md.tpr     }\AttributeTok{{-}f}\NormalTok{ md\_fit.xtc     }\AttributeTok{{-}n}\NormalTok{ index.ndx     }\AttributeTok{{-}ref} \StringTok{\textquotesingle{}group "Protein"\textquotesingle{}}     \AttributeTok{{-}sel} \StringTok{\textquotesingle{}group "BOG1"\textquotesingle{}} \StringTok{\textquotesingle{}group "BOG2"\textquotesingle{}}     \AttributeTok{{-}type}\NormalTok{ min     }\AttributeTok{{-}o}\NormalTok{ ligand\_protein\_mindist.xvg}
\end{Highlighting}
\end{Shaded}

\subsubsection{RMSD de cadenas
individuales}\label{rmsd-de-cadenas-individuales}

Cree grupos para cada cadena y use el oligómero completo como grupo de
ajuste si desea comparar movimientos internos bajo una referencia común:

\begin{Shaded}
\begin{Highlighting}[]
\KeywordTok{(}\BuiltInTok{echo}\NormalTok{ Backbone}\KeywordTok{;} \BuiltInTok{echo}\NormalTok{ Chain\_A}\KeywordTok{)} \KeywordTok{|}
\ExtensionTok{gmx}\NormalTok{ rms     }\AttributeTok{{-}s}\NormalTok{ md.tpr     }\AttributeTok{{-}f}\NormalTok{ md\_fit.xtc     }\AttributeTok{{-}n}\NormalTok{ index.ndx     }\AttributeTok{{-}o}\NormalTok{ rmsd\_chain\_A.xvg     }\AttributeTok{{-}tu}\NormalTok{ ns}

\KeywordTok{(}\BuiltInTok{echo}\NormalTok{ Backbone}\KeywordTok{;} \BuiltInTok{echo}\NormalTok{ Chain\_B}\KeywordTok{)} \KeywordTok{|}
\ExtensionTok{gmx}\NormalTok{ rms     }\AttributeTok{{-}s}\NormalTok{ md.tpr     }\AttributeTok{{-}f}\NormalTok{ md\_fit.xtc     }\AttributeTok{{-}n}\NormalTok{ index.ndx     }\AttributeTok{{-}o}\NormalTok{ rmsd\_chain\_B.xvg     }\AttributeTok{{-}tu}\NormalTok{ ns}
\end{Highlighting}
\end{Shaded}

Si cada cadena se ajusta sobre sí misma, se elimina su movimiento
relativo respecto del oligómero. Ambas estrategias responden preguntas
diferentes y no deben mezclarse en una misma comparación.

\subsubsection{RMSD de la proteína, todas las cadenas
juntas}\label{rmsd-de-la-proteuxedna-todas-las-cadenas-juntas}

Después de corregir PBC y ajustar la trayectoria:

\begin{Shaded}
\begin{Highlighting}[]
\KeywordTok{(}\BuiltInTok{echo}\NormalTok{ Backbone}\KeywordTok{;} \BuiltInTok{echo}\NormalTok{ Backbone}\KeywordTok{)} \KeywordTok{|}
\ExtensionTok{gmx}\NormalTok{ rms     }\AttributeTok{{-}s}\NormalTok{ md.tpr     }\AttributeTok{{-}f}\NormalTok{ md\_fit.xtc     }\AttributeTok{{-}n}\NormalTok{ index.ndx     }\AttributeTok{{-}o}\NormalTok{ md\_rmsd\_protein\_all\_chains.xvg     }\AttributeTok{{-}tu}\NormalTok{ ns}
\end{Highlighting}
\end{Shaded}

Este RMSD describe el cambio del backbone del conjunto de cadenas
respecto de la referencia seleccionada. Puede aumentar por
reorganización cuaternaria aunque cada cadena conserve su estructura
interna. Por eso conviene compararlo con los RMSD por cadena y con
distancias entre centros de masa:

\begin{Shaded}
\begin{Highlighting}[]
\ExtensionTok{gmx}\NormalTok{ distance     }\AttributeTok{{-}s}\NormalTok{ md.tpr     }\AttributeTok{{-}f}\NormalTok{ md\_fit.xtc     }\AttributeTok{{-}n}\NormalTok{ index.ndx     }\AttributeTok{{-}select} \StringTok{\textquotesingle{}com of group "Chain\_A" plus com of group "Chain\_B"\textquotesingle{}}     \AttributeTok{{-}oall}\NormalTok{ chain\_A\_chain\_B\_distance.xvg}
\end{Highlighting}
\end{Shaded}

\subsection{Dinámica de una proteína en
agua}\label{dinuxe1mica-de-una-proteuxedna-en-agua}

\subsubsection{Caso: proteína
monomérica}\label{caso-proteuxedna-monomuxe9rica}

Este sistema contiene una proteína formada por aminoácidos estándar,
agua e iones. Es útil para estudiar estabilidad conformacional,
flexibilidad, compactación, estructura secundaria, exposición al
solvente y movimientos colectivos. También sirve para relajar una
estructura modelada, pero una dinámica molecular breve no corrige
automáticamente errores de secuencia, plegamiento, protonación o
ensamblaje.

El flujo general es:

\begin{Shaded}
\begin{Highlighting}[]
\NormalTok{estructura inicial}
\NormalTok{    ↓}
\NormalTok{revisión y preparación}
\NormalTok{    ↓}
\NormalTok{topología con pdb2gmx}
\NormalTok{    ↓}
\NormalTok{caja periódica}
\NormalTok{    ↓}
\NormalTok{solvatación e iones}
\NormalTok{    ↓}
\NormalTok{minimización de energía}
\NormalTok{    ↓}
\NormalTok{equilibración NVT}
\NormalTok{    ↓}
\NormalTok{equilibración NPT}
\NormalTok{    ↓}
\NormalTok{dinámica de producción}
\NormalTok{    ↓}
\NormalTok{corrección de PBC y análisis}
\end{Highlighting}
\end{Shaded}

\subsubsection{1. Definir qué estructura se
simulará}\label{definir-quuxe9-estructura-se-simularuxe1}

Antes de ejecutar GROMACS, determine si la estructura representa
realmente un monómero biológico. Una cadena aislada de un PDB puede ser:

\begin{itemize}
\tightlist
\item
  una proteína monomérica funcional;
\item
  una subunidad extraída de un oligómero;
\item
  una construcción truncada;
\item
  una estructura con mutaciones o etiquetas;
\item
  un modelo incompleto;
\item
  una cadena estabilizada por contactos cristalográficos.
\end{itemize}

Simular una única cadena de una proteína oligomérica puede producir
exposición artificial de superficies hidrofóbicas, pérdida de estructura
o movimientos que no ocurren en la unidad biológica.

Revise como mínimo:

\begin{itemize}
\tightlist
\item
  residuos y segmentos faltantes;
\item
  átomos con ocupación alternativa;
\item
  residuos no estándar;
\item
  terminales reales de la construcción;
\item
  enlaces disulfuro;
\item
  histidinas y otros grupos titulables;
\item
  metales, cofactores, ligandos y aguas estructurales;
\item
  mutaciones y modificaciones postraduccionales;
\item
  orientación y entorno experimental de la proteína.
\end{itemize}

Conserve el archivo original y trabaje sobre una copia:

\begin{Shaded}
\begin{Highlighting}[]
\FunctionTok{mkdir} \AttributeTok{{-}p}\NormalTok{ 00\_entrada 01\_preparacion 02\_em 03\_nvt 04\_npt 05\_md 06\_analisis}
\FunctionTok{cp}\NormalTok{ protein.pdb 00\_entrada/protein\_original.pdb}
\end{Highlighting}
\end{Shaded}

\subsubsection{2. Campo de fuerza y modelo de
agua}\label{campo-de-fuerza-y-modelo-de-agua}

No existe un campo de fuerza universalmente superior para todas las
proteínas y propiedades. La elección debe justificarse por:

\begin{itemize}
\tightlist
\item
  tipo de proteína y entorno;
\item
  propiedad que se medirá;
\item
  calidad de la parametrización disponible;
\item
  compatibilidad con cofactores o ligandos;
\item
  antecedentes metodológicos comparables;
\item
  modelo de agua utilizado durante el desarrollo o la validación del
  campo de fuerza.
\end{itemize}

No seleccione OPLS-AA, CHARMM, AMBER o GROMOS solo por el número que
ocupa en el menú de \textbf{pdb2gmx}. Los números y campos instalados
pueden cambiar entre equipos.

Liste las opciones disponibles:

\begin{Shaded}
\begin{Highlighting}[]
\ExtensionTok{gmx}\NormalTok{ pdb2gmx }\AttributeTok{{-}h}
\end{Highlighting}
\end{Shaded}

También puede iniciar \textbf{pdb2gmx} sin \textbf{-ff} para seleccionar
el campo de fuerza interactivamente.

\subsubsection{3. Limpieza de la
estructura}\label{limpieza-de-la-estructura}

Elimine únicamente componentes cuya ausencia esté justificada. Las aguas
cristalográficas alejadas suelen descartarse, pero una molécula de agua
conservada en un sitio catalítico o en una red de puentes de hidrógeno
puede ser relevante.

Las conformaciones alternativas deben resolverse antes de
\textbf{pdb2gmx}. No deben quedar dos posiciones incompatibles para el
mismo átomo.

Verificación inicial:

\begin{Shaded}
\begin{Highlighting}[]
\FunctionTok{grep} \StringTok{\textquotesingle{}\^{}ATOM\textbackslash{}|\^{}HETATM\textbackslash{}|\^{}TER\textbackslash{}|\^{}SSBOND\textquotesingle{}}\NormalTok{ 00\_entrada/protein\_original.pdb     }\OperatorTok{\textgreater{}}\NormalTok{ 01\_preparacion/structure\_records.txt}
\end{Highlighting}
\end{Shaded}

Este comando solo extrae registros para inspección; no genera por sí
mismo un PDB listo para simular.

\subsubsection{4. Generación de la
topología}\label{generaciuxf3n-de-la-topologuxeda}

Ejemplo interactivo:

\begin{Shaded}
\begin{Highlighting}[]
\BuiltInTok{cd}\NormalTok{ 01\_preparacion}

\ExtensionTok{gmx}\NormalTok{ pdb2gmx     }\AttributeTok{{-}f}\NormalTok{ ../00\_entrada/protein\_original.pdb     }\AttributeTok{{-}o}\NormalTok{ protein\_processed.gro     }\AttributeTok{{-}p}\NormalTok{ topol.top     }\AttributeTok{{-}i}\NormalTok{ posre\_protein.itp}
\end{Highlighting}
\end{Shaded}

Ejemplo con campo de fuerza y agua indicados explícitamente:

\begin{Shaded}
\begin{Highlighting}[]
\ExtensionTok{gmx}\NormalTok{ pdb2gmx     }\AttributeTok{{-}f}\NormalTok{ ../00\_entrada/protein\_original.pdb     }\AttributeTok{{-}o}\NormalTok{ protein\_processed.gro     }\AttributeTok{{-}p}\NormalTok{ topol.top     }\AttributeTok{{-}i}\NormalTok{ posre\_protein.itp     }\AttributeTok{{-}ff}\NormalTok{ charmm36{-}jul2022     }\AttributeTok{{-}water}\NormalTok{ tip3p}
\end{Highlighting}
\end{Shaded}

El identificador \textbf{charmm36-jul2022} es un ejemplo. Debe coincidir
con un campo instalado y con el protocolo elegido.

Opciones útiles:

{\def\LTcaptype{none} % do not increment counter
\begin{longtable}[]{@{}
  >{\raggedright\arraybackslash}p{(\linewidth - 2\tabcolsep) * \real{0.1461}}
  >{\raggedright\arraybackslash}p{(\linewidth - 2\tabcolsep) * \real{0.8539}}@{}}
\toprule\noalign{}
\begin{minipage}[b]{\linewidth}\raggedright
Opción
\end{minipage} & \begin{minipage}[b]{\linewidth}\raggedright
Uso
\end{minipage} \\
\midrule\noalign{}
\endhead
\bottomrule\noalign{}
\endlastfoot
\textbf{-ff} & Selecciona el campo de fuerza. \\
\textbf{-water} & Selecciona el modelo de agua. \\
\textbf{-ter} & Permite elegir estados de los terminales. \\
\textbf{-his} & Permite seleccionar estados de protonación de
histidinas. \\
\textbf{-ss} & Permite elegir interactivamente enlaces disulfuro. \\
\textbf{-ignh} & Ignora hidrógenos de entrada y los reconstruye. Debe
usarse deliberadamente. \\
\textbf{-chainsep} & Controla cuándo separar cadenas en tipos
moleculares. \\
\textbf{-merge} & Controla la fusión de cadenas en un único tipo
molecular. \\
\end{longtable}
}

Consulte las opciones exactas de la instalación:

\begin{Shaded}
\begin{Highlighting}[]
\ExtensionTok{gmx}\NormalTok{ pdb2gmx }\AttributeTok{{-}h}
\end{Highlighting}
\end{Shaded}

\textbf{pdb2gmx} no asigna estados de protonación mediante una
simulación de pH constante. Aplica plantillas y decisiones del usuario.
Para histidina deben evaluarse las formas protonadas en Nδ, Nε o en
ambos nitrógenos según su entorno.

Revise cuidadosamente la salida. Debe comprobar:

\begin{itemize}
\tightlist
\item
  residuos reconocidos;
\item
  terminales asignados;
\item
  carga total;
\item
  enlaces disulfuro;
\item
  átomos añadidos o eliminados;
\item
  advertencias;
\item
  nombre del tipo molecular generado.
\end{itemize}

Compruebe el archivo:

\begin{Shaded}
\begin{Highlighting}[]
\ExtensionTok{gmx}\NormalTok{ check }\AttributeTok{{-}f}\NormalTok{ protein\_processed.gro}
\end{Highlighting}
\end{Shaded}

\subsubsection{5. Contenido esperado de la
topología}\label{contenido-esperado-de-la-topologuxeda}

Un \textbf{topol.top} típico incluye:

\begin{Shaded}
\begin{Highlighting}[]
\CommentTok{; Campo de fuerza}
\CommentTok{\#include "charmm36{-}jul2022.ff/forcefield.itp"}

\CommentTok{; Topología de la proteína}
\CommentTok{\#include "topol\_Protein.itp"}

\CommentTok{; Agua}
\CommentTok{\#include "charmm36{-}jul2022.ff/tip3p.itp"}

\CommentTok{; Iones}
\CommentTok{\#include "charmm36{-}jul2022.ff/ions.itp"}

\KeywordTok{[ system ]}
\DataTypeTok{Proteína monomérica en agua}

\KeywordTok{[ molecules ]}
\DataTypeTok{Protein    1}
\end{Highlighting}
\end{Shaded}

La estructura exacta depende de la salida de \textbf{pdb2gmx}. No
reescriba manualmente nombres sin cambiar también la definición
\texttt{{[}\ moleculetype\ {]}} correspondiente.

El archivo de restricciones suele estar incluido dentro del ITP de la
proteína:

\begin{Shaded}
\begin{Highlighting}[]
\CommentTok{\#ifdef POSRES}
\CommentTok{\#include "posre\_protein.itp"}
\CommentTok{\#endif}
\end{Highlighting}
\end{Shaded}

La inclusión debe permanecer dentro del ámbito del tipo molecular al que
pertenecen sus índices.

\subsubsection{6. Caja periódica}\label{caja-periuxf3dica-1}

Para una proteína soluble aproximadamente globular:

\begin{Shaded}
\begin{Highlighting}[]
\ExtensionTok{gmx}\NormalTok{ editconf     }\AttributeTok{{-}f}\NormalTok{ protein\_processed.gro     }\AttributeTok{{-}o}\NormalTok{ protein\_box.gro     }\AttributeTok{{-}c}     \AttributeTok{{-}d}\NormalTok{ 1.0     }\AttributeTok{{-}bt}\NormalTok{ dodecahedron}
\end{Highlighting}
\end{Shaded}

\textbf{-d 1.0} establece una distancia mínima de 1.0 nm, equivalente a
10 Å, entre el soluto y la caja. Debe ser compatible con los radios de
corte y con el tamaño y movimiento esperados de la proteína.

Una caja dodecaédrica suele contener menos agua que una cúbica. Una caja
cúbica puede ser más simple de visualizar, pero normalmente aumenta el
número de átomos:

\begin{Shaded}
\begin{Highlighting}[]
\ExtensionTok{gmx}\NormalTok{ editconf     }\AttributeTok{{-}f}\NormalTok{ protein\_processed.gro     }\AttributeTok{{-}o}\NormalTok{ protein\_box\_cubic.gro     }\AttributeTok{{-}c}     \AttributeTok{{-}d}\NormalTok{ 1.0     }\AttributeTok{{-}bt}\NormalTok{ cubic}
\end{Highlighting}
\end{Shaded}

No use este protocolo de caja acuosa para una proteína transmembrana:
necesita una bicapa, orientación y composición lipídica apropiadas.

\subsubsection{7. Solvatación}\label{solvataciuxf3n-1}

\begin{Shaded}
\begin{Highlighting}[]
\ExtensionTok{gmx}\NormalTok{ solvate     }\AttributeTok{{-}cp}\NormalTok{ protein\_box.gro     }\AttributeTok{{-}cs}\NormalTok{ spc216.gro     }\AttributeTok{{-}o}\NormalTok{ protein\_solv.gro     }\AttributeTok{{-}p}\NormalTok{ topol.top}
\end{Highlighting}
\end{Shaded}

\textbf{gmx solvate} actualiza la cantidad de solvente en
\texttt{{[}\ molecules\ {]}}. Revise el final de \textbf{topol.top} y
confirme que el agua utilizada es compatible con la topología elegida.

Compruebe que la proteína no cruza de forma problemática la caja y que
no existen cavidades o solapamientos extraños mediante una inspección
visual.

\subsubsection{8. Neutralización y fuerza
iónica}\label{neutralizaciuxf3n-y-fuerza-iuxf3nica}

Archivo mínimo \textbf{ions.mdp}:

\begin{Shaded}
\begin{Highlighting}[]
\DataTypeTok{integrator      }\OtherTok{=}\StringTok{ steep}
\DataTypeTok{nsteps          }\OtherTok{=}\StringTok{ }\DecValTok{0}
\DataTypeTok{emtol           }\OtherTok{=}\StringTok{ }\FloatTok{1000.0}

\DataTypeTok{cutoff{-}scheme   }\OtherTok{=}\StringTok{ Verlet}
\DataTypeTok{coulombtype     }\OtherTok{=}\StringTok{ PME}
\DataTypeTok{rcoulomb        }\OtherTok{=}\StringTok{ }\FloatTok{1.0}
\DataTypeTok{rvdw            }\OtherTok{=}\StringTok{ }\FloatTok{1.0}
\DataTypeTok{pbc             }\OtherTok{=}\StringTok{ xyz}
\end{Highlighting}
\end{Shaded}

Genere el TPR:

\begin{Shaded}
\begin{Highlighting}[]
\ExtensionTok{gmx}\NormalTok{ grompp     }\AttributeTok{{-}f}\NormalTok{ ions.mdp     }\AttributeTok{{-}c}\NormalTok{ protein\_solv.gro     }\AttributeTok{{-}p}\NormalTok{ topol.top     }\AttributeTok{{-}o}\NormalTok{ ions.tpr}
\end{Highlighting}
\end{Shaded}

Neutralización con NaCl y concentración nominal de 0.15 mol L⁻¹:

\begin{Shaded}
\begin{Highlighting}[]
\ExtensionTok{gmx}\NormalTok{ genion     }\AttributeTok{{-}s}\NormalTok{ ions.tpr     }\AttributeTok{{-}o}\NormalTok{ protein\_solv\_ions.gro     }\AttributeTok{{-}p}\NormalTok{ topol.top     }\AttributeTok{{-}pname}\NormalTok{ NA     }\AttributeTok{{-}nname}\NormalTok{ CL     }\AttributeTok{{-}neutral}     \AttributeTok{{-}conc}\NormalTok{ 0.15}
\end{Highlighting}
\end{Shaded}

Seleccione el grupo de agua, normalmente \textbf{SOL}. No memorice su
número: depende del sistema.

La opción \textbf{-neutral} incorpora los contraiones necesarios para
que la carga total sea cero. \textbf{-conc 0.15} agrega sal hasta
aproximar la concentración solicitada. Debido al volumen finito de la
caja, la concentración efectiva puede diferir ligeramente.

Si el experimento requiere otra sal, deben existir parámetros
compatibles para cada especie. Un ion metálico coordinado en un sitio
activo no debe tratarse como un contraion difusible común.

\subsubsection{9. Validación antes de
minimizar}\label{validaciuxf3n-antes-de-minimizar}

Genere una topología expandida:

\begin{Shaded}
\begin{Highlighting}[]
\ExtensionTok{gmx}\NormalTok{ grompp     }\AttributeTok{{-}f}\NormalTok{ em.mdp     }\AttributeTok{{-}c}\NormalTok{ protein\_solv\_ions.gro     }\AttributeTok{{-}p}\NormalTok{ topol.top     }\AttributeTok{{-}o}\NormalTok{ em\_test.tpr     }\AttributeTok{{-}pp}\NormalTok{ processed.top}
\end{Highlighting}
\end{Shaded}

\textbf{processed.top} permite revisar las inclusiones y macros después
del preprocesamiento.

No continúe si existen:

\begin{itemize}
\tightlist
\item
  diferencias entre el número de átomos y coordenadas;
\item
  carga total inesperada;
\item
  residuos sin parametrizar;
\item
  parámetros faltantes;
\item
  nombres moleculares inconsistentes;
\item
  advertencias no comprendidas.
\end{itemize}

No use \textbf{-maxwarn} para forzar la creación del TPR.

\subsubsection{10. Minimización de
energía}\label{minimizaciuxf3n-de-energuxeda-1}

Archivo \textbf{em.mdp}:

\begin{Shaded}
\begin{Highlighting}[]
\DataTypeTok{title            }\OtherTok{=}\StringTok{ Minimización de proteína en agua}
\DataTypeTok{integrator       }\OtherTok{=}\StringTok{ steep}
\DataTypeTok{nsteps           }\OtherTok{=}\StringTok{ }\DecValTok{50000}
\DataTypeTok{emtol            }\OtherTok{=}\StringTok{ }\FloatTok{1000.0}
\DataTypeTok{emstep           }\OtherTok{=}\StringTok{ }\FloatTok{0.01}

\DataTypeTok{cutoff{-}scheme    }\OtherTok{=}\StringTok{ Verlet}
\DataTypeTok{nstlist          }\OtherTok{=}\StringTok{ }\DecValTok{20}
\DataTypeTok{rlist            }\OtherTok{=}\StringTok{ }\FloatTok{1.0}
\DataTypeTok{coulombtype      }\OtherTok{=}\StringTok{ PME}
\DataTypeTok{rcoulomb         }\OtherTok{=}\StringTok{ }\FloatTok{1.0}
\DataTypeTok{vdwtype          }\OtherTok{=}\StringTok{ Cut{-}off}
\DataTypeTok{rvdw             }\OtherTok{=}\StringTok{ }\FloatTok{1.0}
\DataTypeTok{pbc              }\OtherTok{=}\StringTok{ xyz}
\end{Highlighting}
\end{Shaded}

Ejecución:

\begin{Shaded}
\begin{Highlighting}[]
\FunctionTok{mkdir} \AttributeTok{{-}p}\NormalTok{ ../02\_em}
\BuiltInTok{cd}\NormalTok{ ../02\_em}

\ExtensionTok{gmx}\NormalTok{ grompp     }\AttributeTok{{-}f}\NormalTok{ ../01\_preparacion/em.mdp     }\AttributeTok{{-}c}\NormalTok{ ../01\_preparacion/protein\_solv\_ions.gro     }\AttributeTok{{-}p}\NormalTok{ ../01\_preparacion/topol.top     }\AttributeTok{{-}o}\NormalTok{ em.tpr}

\ExtensionTok{gmx}\NormalTok{ mdrun }\AttributeTok{{-}deffnm}\NormalTok{ em }\AttributeTok{{-}v}
\end{Highlighting}
\end{Shaded}

Extraiga la energía potencial:

\begin{Shaded}
\begin{Highlighting}[]
\KeywordTok{(}\BuiltInTok{echo}\NormalTok{ Potential}\KeywordTok{;} \BuiltInTok{echo}\NormalTok{ 0}\KeywordTok{)} \KeywordTok{|}
\ExtensionTok{gmx}\NormalTok{ energy     }\AttributeTok{{-}f}\NormalTok{ em.edr     }\AttributeTok{{-}o}\NormalTok{ ../06\_analisis/em\_potential.xvg}
\end{Highlighting}
\end{Shaded}

La minimización elimina contactos desfavorables; no equilibra
temperatura, presión ni distribución conformacional. Revise energía
potencial, fuerza máxima, átomo asociado a esa fuerza y posibles valores
NaN.

\subsubsection{11. Equilibración NVT}\label{equilibraciuxf3n-nvt-1}

NVT estabiliza la temperatura con volumen fijo. Use restricciones
posicionales inicialmente si necesita evitar una relajación brusca del
soluto mientras se reorganizan agua e iones.

Parámetros principales:

\begin{Shaded}
\begin{Highlighting}[]
\DataTypeTok{title         }\OtherTok{=}\StringTok{ Equilibración NVT}
\DataTypeTok{define        }\OtherTok{=}\StringTok{ {-}DPOSRES}

\DataTypeTok{integrator    }\OtherTok{=}\StringTok{ md}
\DataTypeTok{dt            }\OtherTok{=}\StringTok{ }\FloatTok{0.002}
\DataTypeTok{nsteps        }\OtherTok{=}\StringTok{ }\DecValTok{50000}
\DataTypeTok{continuation  }\OtherTok{=}\StringTok{ }\KeywordTok{no}

\DataTypeTok{gen{-}vel       }\OtherTok{=}\StringTok{ }\KeywordTok{yes}
\DataTypeTok{gen{-}temp      }\OtherTok{=}\StringTok{ }\DecValTok{300}
\DataTypeTok{gen{-}seed      }\OtherTok{=}\StringTok{ }\DecValTok{{-}1}

\DataTypeTok{constraints   }\OtherTok{=}\StringTok{ h{-}bonds}

\DataTypeTok{tcoupl        }\OtherTok{=}\StringTok{ V{-}rescale}
\DataTypeTok{tc{-}grps       }\OtherTok{=}\StringTok{ Protein Water\_and\_ions}
\DataTypeTok{tau{-}t         }\OtherTok{=}\StringTok{ 1.0 1.0}
\DataTypeTok{ref{-}t         }\OtherTok{=}\StringTok{ 300 300}

\DataTypeTok{pcoupl        }\OtherTok{=}\StringTok{ }\KeywordTok{no}
\DataTypeTok{pbc           }\OtherTok{=}\StringTok{ xyz}
\end{Highlighting}
\end{Shaded}

Con \textbf{dt = 0.002 ps} y \textbf{nsteps = 50000}, la duración es 100
ps. El paso de 0.002 ps equivale a 2 fs.

\begin{Shaded}
\begin{Highlighting}[]
\FunctionTok{mkdir} \AttributeTok{{-}p}\NormalTok{ ../03\_nvt}
\BuiltInTok{cd}\NormalTok{ ../03\_nvt}

\ExtensionTok{gmx}\NormalTok{ grompp     }\AttributeTok{{-}f}\NormalTok{ ../01\_preparacion/nvt.mdp     }\AttributeTok{{-}c}\NormalTok{ ../02\_em/em.gro     }\AttributeTok{{-}r}\NormalTok{ ../02\_em/em.gro     }\AttributeTok{{-}p}\NormalTok{ ../01\_preparacion/topol.top     }\AttributeTok{{-}o}\NormalTok{ nvt.tpr}

\ExtensionTok{gmx}\NormalTok{ mdrun }\AttributeTok{{-}deffnm}\NormalTok{ nvt }\AttributeTok{{-}v}
\end{Highlighting}
\end{Shaded}

Analice la temperatura:

\begin{Shaded}
\begin{Highlighting}[]
\KeywordTok{(}\BuiltInTok{echo}\NormalTok{ Temperature}\KeywordTok{;} \BuiltInTok{echo}\NormalTok{ 0}\KeywordTok{)} \KeywordTok{|}
\ExtensionTok{gmx}\NormalTok{ energy     }\AttributeTok{{-}f}\NormalTok{ nvt.edr     }\AttributeTok{{-}o}\NormalTok{ ../06\_analisis/nvt\_temperature.xvg}
\end{Highlighting}
\end{Shaded}

\subsubsection{12. Equilibración NPT}\label{equilibraciuxf3n-npt-1}

NPT ajusta presión, densidad y volumen. Continúe desde el checkpoint de
NVT y no regenere velocidades:

\begin{Shaded}
\begin{Highlighting}[]
\DataTypeTok{title            }\OtherTok{=}\StringTok{ Equilibración NPT}
\DataTypeTok{define           }\OtherTok{=}\StringTok{ {-}DPOSRES}

\DataTypeTok{integrator       }\OtherTok{=}\StringTok{ md}
\DataTypeTok{dt               }\OtherTok{=}\StringTok{ }\FloatTok{0.002}
\DataTypeTok{nsteps           }\OtherTok{=}\StringTok{ }\DecValTok{250000}
\DataTypeTok{continuation     }\OtherTok{=}\StringTok{ }\KeywordTok{yes}
\DataTypeTok{gen{-}vel          }\OtherTok{=}\StringTok{ }\KeywordTok{no}

\DataTypeTok{constraints      }\OtherTok{=}\StringTok{ h{-}bonds}

\DataTypeTok{tcoupl           }\OtherTok{=}\StringTok{ V{-}rescale}
\DataTypeTok{tc{-}grps          }\OtherTok{=}\StringTok{ Protein Water\_and\_ions}
\DataTypeTok{tau{-}t            }\OtherTok{=}\StringTok{ 1.0 1.0}
\DataTypeTok{ref{-}t            }\OtherTok{=}\StringTok{ 300 300}

\DataTypeTok{pcoupl           }\OtherTok{=}\StringTok{ C{-}rescale}
\DataTypeTok{pcoupltype       }\OtherTok{=}\StringTok{ isotropic}
\DataTypeTok{tau{-}p            }\OtherTok{=}\StringTok{ }\FloatTok{5.0}
\DataTypeTok{ref{-}p            }\OtherTok{=}\StringTok{ }\FloatTok{1.0}
\DataTypeTok{compressibility  }\OtherTok{=}\StringTok{ 4.5e{-}5}

\DataTypeTok{pbc              }\OtherTok{=}\StringTok{ xyz}
\end{Highlighting}
\end{Shaded}

Aquí se simulan 500 ps. \textbf{ref-p} se expresa en bar y la
compresibilidad en bar⁻¹. El valor 4.5 × 10⁻⁵ bar⁻¹ es habitual para
agua líquida cerca de condiciones ambientales.

\begin{Shaded}
\begin{Highlighting}[]
\FunctionTok{mkdir} \AttributeTok{{-}p}\NormalTok{ ../04\_npt}
\BuiltInTok{cd}\NormalTok{ ../04\_npt}

\ExtensionTok{gmx}\NormalTok{ grompp     }\AttributeTok{{-}f}\NormalTok{ ../01\_preparacion/npt.mdp     }\AttributeTok{{-}c}\NormalTok{ ../03\_nvt/nvt.gro     }\AttributeTok{{-}r}\NormalTok{ ../03\_nvt/nvt.gro     }\AttributeTok{{-}t}\NormalTok{ ../03\_nvt/nvt.cpt     }\AttributeTok{{-}p}\NormalTok{ ../01\_preparacion/topol.top     }\AttributeTok{{-}o}\NormalTok{ npt.tpr}

\ExtensionTok{gmx}\NormalTok{ mdrun }\AttributeTok{{-}deffnm}\NormalTok{ npt }\AttributeTok{{-}v}
\end{Highlighting}
\end{Shaded}

Analice temperatura, presión, densidad y volumen:

\begin{Shaded}
\begin{Highlighting}[]
\KeywordTok{(}\BuiltInTok{echo}\NormalTok{ Temperature}\KeywordTok{;} \BuiltInTok{echo}\NormalTok{ Pressure}\KeywordTok{;} \BuiltInTok{echo}\NormalTok{ Density}\KeywordTok{;} \BuiltInTok{echo}\NormalTok{ Volume}\KeywordTok{;} \BuiltInTok{echo}\NormalTok{ 0}\KeywordTok{)} \KeywordTok{|}
\ExtensionTok{gmx}\NormalTok{ energy     }\AttributeTok{{-}f}\NormalTok{ npt.edr     }\AttributeTok{{-}o}\NormalTok{ ../06\_analisis/npt\_thermodynamics.xvg}
\end{Highlighting}
\end{Shaded}

La presión instantánea fluctúa intensamente. Evalúe promedios por
bloques, densidad y deriva temporal.

\subsubsection{13. Dinámica de
producción}\label{dinuxe1mica-de-producciuxf3n}

Retire \textbf{define = -DPOSRES} salvo que mantener restricciones forme
parte explícita del experimento.

Parámetros principales:

\begin{Shaded}
\begin{Highlighting}[]
\DataTypeTok{title            }\OtherTok{=}\StringTok{ Producción de proteína en agua}

\DataTypeTok{integrator       }\OtherTok{=}\StringTok{ md}
\DataTypeTok{dt               }\OtherTok{=}\StringTok{ }\FloatTok{0.002}
\DataTypeTok{nsteps           }\OtherTok{=}\StringTok{ }\DecValTok{50000000}
\DataTypeTok{continuation     }\OtherTok{=}\StringTok{ }\KeywordTok{yes}
\DataTypeTok{gen{-}vel          }\OtherTok{=}\StringTok{ }\KeywordTok{no}

\DataTypeTok{constraints      }\OtherTok{=}\StringTok{ h{-}bonds}

\DataTypeTok{tcoupl           }\OtherTok{=}\StringTok{ V{-}rescale}
\DataTypeTok{tc{-}grps          }\OtherTok{=}\StringTok{ Protein Water\_and\_ions}
\DataTypeTok{tau{-}t            }\OtherTok{=}\StringTok{ 1.0 1.0}
\DataTypeTok{ref{-}t            }\OtherTok{=}\StringTok{ 300 300}

\DataTypeTok{pcoupl           }\OtherTok{=}\StringTok{ Parrinello{-}Rahman}
\DataTypeTok{pcoupltype       }\OtherTok{=}\StringTok{ isotropic}
\DataTypeTok{tau{-}p            }\OtherTok{=}\StringTok{ }\FloatTok{5.0}
\DataTypeTok{ref{-}p            }\OtherTok{=}\StringTok{ }\FloatTok{1.0}
\DataTypeTok{compressibility  }\OtherTok{=}\StringTok{ 4.5e{-}5}

\DataTypeTok{pbc              }\OtherTok{=}\StringTok{ xyz}
\end{Highlighting}
\end{Shaded}

Con 50000000 pasos de 0.002 ps se simulan 100 ns. Agregue los parámetros
de PME, Verlet y control de salida validados en el tutorial principal;
no mezcle MDP de campos de fuerza diferentes sin revisar cortes y
modificadores.

\begin{Shaded}
\begin{Highlighting}[]
\FunctionTok{mkdir} \AttributeTok{{-}p}\NormalTok{ ../05\_md}
\BuiltInTok{cd}\NormalTok{ ../05\_md}

\ExtensionTok{gmx}\NormalTok{ grompp     }\AttributeTok{{-}f}\NormalTok{ ../01\_preparacion/md.mdp     }\AttributeTok{{-}c}\NormalTok{ ../04\_npt/npt.gro     }\AttributeTok{{-}t}\NormalTok{ ../04\_npt/npt.cpt     }\AttributeTok{{-}p}\NormalTok{ ../01\_preparacion/topol.top     }\AttributeTok{{-}o}\NormalTok{ md.tpr}

\ExtensionTok{gmx}\NormalTok{ mdrun }\AttributeTok{{-}deffnm}\NormalTok{ md }\AttributeTok{{-}v}
\end{Highlighting}
\end{Shaded}

Compruebe en el registro:

\begin{itemize}
\tightlist
\item
  ausencia de errores LINCS;
\item
  temperatura, presión y densidad estables;
\item
  rendimiento y uso esperado de CPU/GPU;
\item
  frecuencia y tamaño de los archivos de salida;
\item
  escritura periódica del checkpoint.
\end{itemize}

\subsubsection{14. Controles estructurales
mínimos}\label{controles-estructurales-muxednimos}

Corrija PBC antes de analizar:

\begin{Shaded}
\begin{Highlighting}[]
\BuiltInTok{echo}\NormalTok{ Protein }\KeywordTok{|}
\ExtensionTok{gmx}\NormalTok{ trjconv     }\AttributeTok{{-}s}\NormalTok{ md.tpr     }\AttributeTok{{-}f}\NormalTok{ md.xtc     }\AttributeTok{{-}o}\NormalTok{ md\_protein\_center.gro     }\AttributeTok{{-}center}     \AttributeTok{{-}pbc}\NormalTok{ mol     }\AttributeTok{{-}ur}\NormalTok{ compact}
\end{Highlighting}
\end{Shaded}

Para conservar una trayectoria XTC:

\begin{Shaded}
\begin{Highlighting}[]
\KeywordTok{(}\BuiltInTok{echo}\NormalTok{ Protein}\KeywordTok{;} \BuiltInTok{echo}\NormalTok{ Protein}\KeywordTok{)} \KeywordTok{|}
\ExtensionTok{gmx}\NormalTok{ trjconv     }\AttributeTok{{-}s}\NormalTok{ md.tpr     }\AttributeTok{{-}f}\NormalTok{ md.xtc     }\AttributeTok{{-}o}\NormalTok{ md\_protein\_center.xtc     }\AttributeTok{{-}center}     \AttributeTok{{-}pbc}\NormalTok{ mol     }\AttributeTok{{-}ur}\NormalTok{ compact}
\end{Highlighting}
\end{Shaded}

Ajuste rotación y traslación:

\begin{Shaded}
\begin{Highlighting}[]
\KeywordTok{(}\BuiltInTok{echo}\NormalTok{ Backbone}\KeywordTok{;} \BuiltInTok{echo}\NormalTok{ Protein}\KeywordTok{)} \KeywordTok{|}
\ExtensionTok{gmx}\NormalTok{ trjconv     }\AttributeTok{{-}s}\NormalTok{ md.tpr     }\AttributeTok{{-}f}\NormalTok{ md\_protein\_center.xtc     }\AttributeTok{{-}o}\NormalTok{ md\_protein\_fit.xtc     }\AttributeTok{{-}fit}\NormalTok{ rot+trans}
\end{Highlighting}
\end{Shaded}

RMSD del backbone:

\begin{Shaded}
\begin{Highlighting}[]
\KeywordTok{(}\BuiltInTok{echo}\NormalTok{ Backbone}\KeywordTok{;} \BuiltInTok{echo}\NormalTok{ Backbone}\KeywordTok{)} \KeywordTok{|}
\ExtensionTok{gmx}\NormalTok{ rms     }\AttributeTok{{-}s}\NormalTok{ md.tpr     }\AttributeTok{{-}f}\NormalTok{ md\_protein\_fit.xtc     }\AttributeTok{{-}o}\NormalTok{ ../06\_analisis/rmsd\_backbone.xvg     }\AttributeTok{{-}tu}\NormalTok{ ns}
\end{Highlighting}
\end{Shaded}

RMSF por residuo:

\begin{Shaded}
\begin{Highlighting}[]
\BuiltInTok{echo}\NormalTok{ C{-}alpha }\KeywordTok{|}
\ExtensionTok{gmx}\NormalTok{ rmsf     }\AttributeTok{{-}s}\NormalTok{ md.tpr     }\AttributeTok{{-}f}\NormalTok{ md\_protein\_fit.xtc     }\AttributeTok{{-}o}\NormalTok{ ../06\_analisis/rmsf\_calpha.xvg     }\AttributeTok{{-}res}
\end{Highlighting}
\end{Shaded}

Radio de giro:

\begin{Shaded}
\begin{Highlighting}[]
\ExtensionTok{gmx}\NormalTok{ gyrate     }\AttributeTok{{-}s}\NormalTok{ md.tpr     }\AttributeTok{{-}f}\NormalTok{ md\_protein\_fit.xtc     }\AttributeTok{{-}sel} \StringTok{\textquotesingle{}group "Protein"\textquotesingle{}}     \AttributeTok{{-}o}\NormalTok{ ../06\_analisis/gyrate\_protein.xvg}
\end{Highlighting}
\end{Shaded}

Estructura secundaria:

\begin{Shaded}
\begin{Highlighting}[]
\ExtensionTok{gmx}\NormalTok{ dssp     }\AttributeTok{{-}s}\NormalTok{ md.tpr     }\AttributeTok{{-}f}\NormalTok{ md\_protein\_fit.xtc     }\AttributeTok{{-}sel} \StringTok{\textquotesingle{}group "Protein"\textquotesingle{}}     \AttributeTok{{-}o}\NormalTok{ ../06\_analisis/secondary\_structure.dat}
\end{Highlighting}
\end{Shaded}

La disponibilidad y las opciones de \textbf{gmx dssp} deben comprobarse
con:

\begin{Shaded}
\begin{Highlighting}[]
\ExtensionTok{gmx}\NormalTok{ dssp }\AttributeTok{{-}h}
\end{Highlighting}
\end{Shaded}

Una meseta de RMSD no demuestra convergencia. Compare bloques
temporales, réplicas independientes y observables complementarios.

\subsubsection{15. Diseñar el análisis antes de calcular
observables}\label{diseuxf1ar-el-anuxe1lisis-antes-de-calcular-observables}

El análisis debe responder una pregunta física. Acumular RMSD, RMSF,
SASA, puentes de hidrógeno y PCA sin definir qué hipótesis evalúa cada
uno produce figuras, pero no necesariamente evidencia.

Antes de analizar, defina:

\begin{itemize}
\tightlist
\item
  qué intervalo se considera producción equilibrada;
\item
  qué átomos representan mejor el fenómeno: Cα, backbone, todos los
  átomos pesados, sitio activo o dominios;
\item
  qué estructura se usará como referencia;
\item
  qué observable se espera que cambie y en qué dirección;
\item
  cuántas réplicas independientes se compararán;
\item
  cómo se estimará la incertidumbre.
\end{itemize}

El descarte inicial no debe elegirse porque ``mejora'' una curva.
Examine temperatura, densidad, energía, RMSD y los observables
relevantes por bloques. Si una transición lenta ocurre una sola vez,
eliminar la parte anterior puede borrar un estado real en lugar de
retirar equilibración.

Divida la producción en bloques temporales comparables. Por ejemplo,
para una trayectoria de 100 ns puede contrastar 0--25, 25--50, 50--75 y
75--100 ns, además de repetir los análisis en réplicas independientes.
La estabilidad del promedio entre bloques es más informativa que una
línea aparentemente plana.

\subsubsection{16. Exposición al solvente, enlaces y geometría
funcional}\label{exposiciuxf3n-al-solvente-enlaces-y-geometruxeda-funcional}

El área accesible al solvente, SASA, depende de la superficie molecular
y de una sonda que representa aproximadamente una molécula de agua. No
equivale al área geométrica total ni mide por sí sola desplegamiento.

\begin{Shaded}
\begin{Highlighting}[]
\ExtensionTok{gmx}\NormalTok{ sasa }\DataTypeTok{\textbackslash{}}
    \AttributeTok{{-}s}\NormalTok{ md.tpr }\DataTypeTok{\textbackslash{}}
    \AttributeTok{{-}f}\NormalTok{ md\_protein\_fit.xtc }\DataTypeTok{\textbackslash{}}
    \AttributeTok{{-}surface} \StringTok{\textquotesingle{}group "Protein"\textquotesingle{}} \DataTypeTok{\textbackslash{}}
    \AttributeTok{{-}output} \StringTok{\textquotesingle{}group "Protein"\textquotesingle{}} \DataTypeTok{\textbackslash{}}
    \AttributeTok{{-}o}\NormalTok{ ../06\_analisis/sasa\_total.xvg }\DataTypeTok{\textbackslash{}}
    \AttributeTok{{-}or}\NormalTok{ ../06\_analisis/sasa\_residue.xvg}
\end{Highlighting}
\end{Shaded}

Compruebe la sintaxis de la versión instalada con \textbf{gmx sasa -h}.
Para interpretar cambios conviene separar, cuando sea pertinente,
superficie hidrofóbica e hidrofílica mediante grupos definidos y
documentados.

Los puentes de hidrógeno intraproteicos pueden calcularse con:

\begin{Shaded}
\begin{Highlighting}[]
\ExtensionTok{gmx}\NormalTok{ hbond }\DataTypeTok{\textbackslash{}}
    \AttributeTok{{-}s}\NormalTok{ md.tpr }\DataTypeTok{\textbackslash{}}
    \AttributeTok{{-}f}\NormalTok{ md\_protein\_fit.xtc }\DataTypeTok{\textbackslash{}}
    \AttributeTok{{-}r} \StringTok{\textquotesingle{}group "Protein"\textquotesingle{}} \DataTypeTok{\textbackslash{}}
    \AttributeTok{{-}t} \StringTok{\textquotesingle{}group "Protein"\textquotesingle{}} \DataTypeTok{\textbackslash{}}
    \AttributeTok{{-}num}\NormalTok{ ../06\_analisis/hbonds\_intraprotein.xvg}
\end{Highlighting}
\end{Shaded}

La interfaz de \textbf{gmx hbond} cambió entre versiones. Revise
\textbf{gmx hbond -h} y registre los cortes geométricos utilizados. Un
mayor número de puentes no significa automáticamente una estructura más
estable: importa qué interacciones aparecen, su ocupación y si conectan
regiones funcionales.

Para una distancia entre dominios, residuos catalíticos o extremos de
una compuerta, cree grupos específicos en \textbf{index.ndx} y use:

\begin{Shaded}
\begin{Highlighting}[]
\ExtensionTok{gmx}\NormalTok{ distance }\DataTypeTok{\textbackslash{}}
    \AttributeTok{{-}s}\NormalTok{ md.tpr }\DataTypeTok{\textbackslash{}}
    \AttributeTok{{-}f}\NormalTok{ md\_protein\_fit.xtc }\DataTypeTok{\textbackslash{}}
    \AttributeTok{{-}n}\NormalTok{ index.ndx }\DataTypeTok{\textbackslash{}}
    \AttributeTok{{-}select} \StringTok{\textquotesingle{}com of group "Domain\_A" plus com of group "Domain\_B"\textquotesingle{}} \DataTypeTok{\textbackslash{}}
    \AttributeTok{{-}oall}\NormalTok{ ../06\_analisis/domain\_distance.xvg}
\end{Highlighting}
\end{Shaded}

Una distancia entre centros de masa puede ocultar rotaciones. Cuando la
geometría sea funcional, combine distancias, ángulos y contactos
definidos a partir de residuos concretos.

\subsubsection{17. Ángulos diedros, contactos y estructura
secundaria}\label{uxe1ngulos-diedros-contactos-y-estructura-secundaria}

La estabilidad del plegamiento no queda descrita por un RMSD global. Los
cambios locales pueden examinarse mediante ángulos φ/ψ, matrices de
distancia y evolución de la estructura secundaria.

Diagrama de Ramachandran muestreado durante la trayectoria:

\begin{Shaded}
\begin{Highlighting}[]
\ExtensionTok{gmx}\NormalTok{ rama }\DataTypeTok{\textbackslash{}}
    \AttributeTok{{-}s}\NormalTok{ md.tpr }\DataTypeTok{\textbackslash{}}
    \AttributeTok{{-}f}\NormalTok{ md\_protein\_fit.xtc }\DataTypeTok{\textbackslash{}}
    \AttributeTok{{-}o}\NormalTok{ ../06\_analisis/ramachandran.xvg}
\end{Highlighting}
\end{Shaded}

Matriz media de distancias entre residuos:

\begin{Shaded}
\begin{Highlighting}[]
\BuiltInTok{echo}\NormalTok{ Protein }\KeywordTok{|} \DataTypeTok{\textbackslash{}}
\ExtensionTok{gmx}\NormalTok{ mdmat }\DataTypeTok{\textbackslash{}}
    \AttributeTok{{-}s}\NormalTok{ md.tpr }\DataTypeTok{\textbackslash{}}
    \AttributeTok{{-}f}\NormalTok{ md\_protein\_fit.xtc }\DataTypeTok{\textbackslash{}}
    \AttributeTok{{-}mean}\NormalTok{ ../06\_analisis/mean\_distance\_map.xpm }\DataTypeTok{\textbackslash{}}
    \AttributeTok{{-}frames}\NormalTok{ ../06\_analisis/distance\_map\_frames.xpm}
\end{Highlighting}
\end{Shaded}

La opción \textbf{-frames} puede generar archivos grandes; úsela solo
cuando se necesite evolución temporal. Para contactos funcionales es
preferible definir pares de residuos y calcular ocupaciones con un corte
explícito, en lugar de interpretar visualmente toda la matriz.

La estructura secundaria ya puede obtenerse con \textbf{gmx dssp}.
Resuma la ocupación por residuo y por bloque, no solo una representación
coloreada. Una hélice presente el 60 \% del tiempo no es equivalente a
una transición irreversible de hélice a coil.

\subsubsection{18. Fundamento de PCA o dinámica
esencial}\label{fundamento-de-pca-o-dinuxe1mica-esencial}

El análisis de componentes principales, PCA, intenta separar movimientos
colectivos de gran amplitud de fluctuaciones locales. Después de
eliminar traslación y rotación, se construye la matriz de covarianza de
las coordenadas:

\[
C_{ij}=\left\langle
\left(x_i-\langle x_i\rangle\right)
\left(x_j-\langle x_j\rangle\right)
\right\rangle
\]

Su diagonalización produce autovectores \(\mathbf{v}_k\) y autovalores
\(\lambda_k\):

\[
\mathbf{C}\mathbf{v}_k=\lambda_k\mathbf{v}_k
\]

\begin{itemize}
\tightlist
\item
  \(\mathbf{v}_k\) define una dirección colectiva en el espacio de
  coordenadas;
\item
  \(\lambda_k\) es la varianza a lo largo de esa dirección;
\item
  los componentes se ordenan de mayor a menor varianza;
\item
  la fracción de fluctuación explicada por el componente \(k\) es:
\end{itemize}

\[
f_k=\frac{\lambda_k}{\sum_j\lambda_j}
\]

La proyección del fotograma \(t\) sobre el componente \(k\) es:

\[
p_k(t)=\mathbf{v}_k^{\mathrm{T}}
\left[\mathbf{x}(t)-\langle\mathbf{x}\rangle\right]
\]

``Principal'' significa mayor varianza, no mayor importancia biológica.
PC1 puede representar una relajación inicial, difusión no convergida o
un movimiento inducido por una mala preparación. La interpretación
funcional requiere localizar qué dominios se mueven, comprobar
recurrencia y contrastar réplicas o datos experimentales.

PCA cartesiano depende de la selección atómica y del ajuste. Para
dinámica global suele ser razonable usar backbone o Cα. Incluir cadenas
laterales e hidrógenos aumenta dimensionalidad y puede hacer que
movimientos locales oculten los modos colectivos.

\subsubsection{19. PCA con gmx covar y gmx
anaeig}\label{pca-con-gmx-covar-y-gmx-anaeig}

Use una trayectoria con la proteína entera y sin saltos periódicos. No
es necesario ajustar previamente la trayectoria: \textbf{gmx covar}
realiza un ajuste de mínimos cuadrados usando el primer grupo
solicitado. Así se evita ajustar dos veces con criterios diferentes.

Construcción y diagonalización de la matriz de covarianza:

\begin{Shaded}
\begin{Highlighting}[]
\ExtensionTok{gmx}\NormalTok{ covar }\DataTypeTok{\textbackslash{}}
    \AttributeTok{{-}s}\NormalTok{ md.tpr }\DataTypeTok{\textbackslash{}}
    \AttributeTok{{-}f}\NormalTok{ md\_protein\_center.xtc }\DataTypeTok{\textbackslash{}}
    \AttributeTok{{-}o}\NormalTok{ ../06\_analisis/pca\_eigenval.xvg }\DataTypeTok{\textbackslash{}}
    \AttributeTok{{-}v}\NormalTok{ ../06\_analisis/pca\_eigenvec.trr }\DataTypeTok{\textbackslash{}}
    \AttributeTok{{-}av}\NormalTok{ ../06\_analisis/pca\_average.pdb }\DataTypeTok{\textbackslash{}}
    \AttributeTok{{-}l}\NormalTok{ ../06\_analisis/pca\_covar.log}
\end{Highlighting}
\end{Shaded}

Seleccione \textbf{Backbone} para el ajuste y \textbf{Backbone} para la
matriz. Si se elige otro grupo, debe conservarse exactamente la misma
selección en las proyecciones posteriores.

La covarianza cartesiana estándar no está ponderada por masa salvo que
se solicite la opción correspondiente. Para Cα la diferencia es
irrelevante porque todos los átomos seleccionados son carbonos; para
selecciones heterogéneas debe declararse si se utilizó ponderación por
masa.

Proyección sobre PC1 y PC2:

\begin{Shaded}
\begin{Highlighting}[]
\ExtensionTok{gmx}\NormalTok{ anaeig }\DataTypeTok{\textbackslash{}}
    \AttributeTok{{-}v}\NormalTok{ ../06\_analisis/pca\_eigenvec.trr }\DataTypeTok{\textbackslash{}}
    \AttributeTok{{-}s}\NormalTok{ md.tpr }\DataTypeTok{\textbackslash{}}
    \AttributeTok{{-}f}\NormalTok{ md\_protein\_center.xtc }\DataTypeTok{\textbackslash{}}
    \AttributeTok{{-}first}\NormalTok{ 1 }\DataTypeTok{\textbackslash{}}
    \AttributeTok{{-}last}\NormalTok{ 2 }\DataTypeTok{\textbackslash{}}
    \AttributeTok{{-}2d}\NormalTok{ ../06\_analisis/pca\_pc1\_pc2.xvg}
\end{Highlighting}
\end{Shaded}

Evolución temporal de PC1:

\begin{Shaded}
\begin{Highlighting}[]
\ExtensionTok{gmx}\NormalTok{ anaeig }\DataTypeTok{\textbackslash{}}
    \AttributeTok{{-}v}\NormalTok{ ../06\_analisis/pca\_eigenvec.trr }\DataTypeTok{\textbackslash{}}
    \AttributeTok{{-}s}\NormalTok{ md.tpr }\DataTypeTok{\textbackslash{}}
    \AttributeTok{{-}f}\NormalTok{ md\_protein\_center.xtc }\DataTypeTok{\textbackslash{}}
    \AttributeTok{{-}first}\NormalTok{ 1 }\DataTypeTok{\textbackslash{}}
    \AttributeTok{{-}last}\NormalTok{ 1 }\DataTypeTok{\textbackslash{}}
    \AttributeTok{{-}proj}\NormalTok{ ../06\_analisis/pca\_pc1\_time.xvg}
\end{Highlighting}
\end{Shaded}

Estructuras extremas interpoladas a lo largo de PC1:

\begin{Shaded}
\begin{Highlighting}[]
\ExtensionTok{gmx}\NormalTok{ anaeig }\DataTypeTok{\textbackslash{}}
    \AttributeTok{{-}v}\NormalTok{ ../06\_analisis/pca\_eigenvec.trr }\DataTypeTok{\textbackslash{}}
    \AttributeTok{{-}s}\NormalTok{ md.tpr }\DataTypeTok{\textbackslash{}}
    \AttributeTok{{-}f}\NormalTok{ md\_protein\_center.xtc }\DataTypeTok{\textbackslash{}}
    \AttributeTok{{-}first}\NormalTok{ 1 }\DataTypeTok{\textbackslash{}}
    \AttributeTok{{-}last}\NormalTok{ 1 }\DataTypeTok{\textbackslash{}}
    \AttributeTok{{-}extr}\NormalTok{ ../06\_analisis/pca\_pc1\_extremes.pdb }\DataTypeTok{\textbackslash{}}
    \AttributeTok{{-}nframes}\NormalTok{ 30}
\end{Highlighting}
\end{Shaded}

Las estructuras intermedias generadas con \textbf{-extr} son una
interpolación para visualizar el vector. No constituyen una trayectoria
física ni prueban que la proteína recorra suavemente ese camino.

\subsubsection{20. Interpretación y convergencia de
PCA}\label{interpretaciuxf3n-y-convergencia-de-pca}

Examine en conjunto:

\begin{enumerate}
\def\labelenumi{\arabic{enumi}.}
\tightlist
\item
  espectro de autovalores;
\item
  fracción acumulada de varianza;
\item
  proyecciones de cada componente en función del tiempo;
\item
  mapa PC1--PC2 coloreado por tiempo o réplica;
\item
  movimiento estructural de los primeros componentes;
\item
  estabilidad de los autovectores entre bloques y réplicas.
\end{enumerate}

Un mapa PC1--PC2 con varias nubes puede indicar estados
conformacionales, pero también muestreo insuficiente. La conexión
temporal entre puntos y la revisita de cada región son esenciales. Una
trayectoria que avanza una sola vez de izquierda a derecha no demuestra
equilibrio entre dos estados.

El contenido coseno ayuda a detectar modos parecidos a una difusión
aleatoria finita:

\begin{Shaded}
\begin{Highlighting}[]
\ExtensionTok{gmx}\NormalTok{ analyze }\DataTypeTok{\textbackslash{}}
    \AttributeTok{{-}f}\NormalTok{ ../06\_analisis/pca\_pc1\_time.xvg }\DataTypeTok{\textbackslash{}}
    \AttributeTok{{-}cc}\NormalTok{ ../06\_analisis/pca\_pc1\_cosine.xvg}
\end{Highlighting}
\end{Shaded}

Un contenido coseno alto sugiere muestreo insuficiente del modo, pero no
demuestra que la dirección carezca de significado físico. Compare además
PCA calculados por separado sobre mitades de la trayectoria y sobre
réplicas. El solapamiento entre subespacios es más exigente que comparar
solamente autovalores.

Para comparar dos sistemas ---por ejemplo proteína libre y complejo---
no calcule dos PCA independientes y llame ``PC1'' al mismo movimiento.
Los ejes son específicos de cada conjunto. Construya una base común con
trayectorias concatenadas o proyecte ambas trayectorias sobre los
autovectores obtenidos de un único conjunto de referencia. Los sistemas
deben contener los mismos átomos, en el mismo orden, y usar el mismo
ajuste.

\subsubsection{21. Agrupamiento y paisaje
conformacional}\label{agrupamiento-y-paisaje-conformacional}

El agrupamiento busca estructuras representativas; PCA reduce
dimensionalidad. No son equivalentes. Un análisis puede agrupar por
RMSD, contactos o distancias funcionales, según la pregunta.

Ejemplo de agrupamiento por RMSD del backbone:

\begin{Shaded}
\begin{Highlighting}[]
\BuiltInTok{echo}\NormalTok{ Backbone }\KeywordTok{|} \DataTypeTok{\textbackslash{}}
\ExtensionTok{gmx}\NormalTok{ cluster }\DataTypeTok{\textbackslash{}}
    \AttributeTok{{-}s}\NormalTok{ md.tpr }\DataTypeTok{\textbackslash{}}
    \AttributeTok{{-}f}\NormalTok{ md\_protein\_fit.xtc }\DataTypeTok{\textbackslash{}}
    \AttributeTok{{-}method}\NormalTok{ gromos }\DataTypeTok{\textbackslash{}}
    \AttributeTok{{-}cutoff}\NormalTok{ 0.20 }\DataTypeTok{\textbackslash{}}
    \AttributeTok{{-}o}\NormalTok{ ../06\_analisis/cluster\_rmsd.xpm }\DataTypeTok{\textbackslash{}}
    \AttributeTok{{-}g}\NormalTok{ ../06\_analisis/cluster.log }\DataTypeTok{\textbackslash{}}
    \AttributeTok{{-}sz}\NormalTok{ ../06\_analisis/cluster\_sizes.xvg }\DataTypeTok{\textbackslash{}}
    \AttributeTok{{-}cl}\NormalTok{ ../06\_analisis/cluster\_representatives.pdb}
\end{Highlighting}
\end{Shaded}

El corte de 0.20 nm es solo un punto de partida. Debe evaluarse su
sensibilidad: un corte pequeño fragmenta el conjunto y uno grande mezcla
estados distintos. Informe selección, métrica, método y corte.

Un paisaje de energía libre aparente en dos coordenadas \(q_1,q_2\) se
obtiene de la probabilidad muestreada:

\[
F(q_1,q_2)=-k_{\mathrm B}T\ln P(q_1,q_2)+C
\]

Con PC1 y PC2 puede estimarse mediante:

\begin{Shaded}
\begin{Highlighting}[]
\ExtensionTok{gmx}\NormalTok{ sham }\DataTypeTok{\textbackslash{}}
    \AttributeTok{{-}f}\NormalTok{ ../06\_analisis/pca\_pc1\_pc2.xvg }\DataTypeTok{\textbackslash{}}
    \AttributeTok{{-}ls}\NormalTok{ ../06\_analisis/pca\_free\_energy.xpm}
\end{Highlighting}
\end{Shaded}

El resultado depende del binning, el muestreo y las coordenadas
elegidas. Un mínimo oscuro no es una energía absoluta ni demuestra un
estado termodinámico convergido. Regiones no visitadas no tienen
probabilidad estimable; no deben interpretarse como barreras
cuantificadas.

\subsubsection{22. Comparación con información
experimental}\label{comparaciuxf3n-con-informaciuxf3n-experimental}

Una simulación gana valor cuando sus observables se conectan con
mediciones independientes. Según el sistema pueden compararse:

\begin{itemize}
\tightlist
\item
  estructuras de rayos X o cryo-EM, considerando contactos cristalinos y
  resolución;
\item
  factores B, con cautela por desorden estático y efectos del cristal;
\item
  ensambles de RMN, NOE, parámetros de orden o acoplamientos residuales;
\item
  datos SAXS, FRET, HDX-MS o mutagénesis;
\item
  distancias conocidas del sitio activo;
\item
  estados abiertos y cerrados depositados en PDB.
\end{itemize}

La relación idealizada entre desplazamiento cuadrático medio isotrópico
y factor B es:

\[
B=\frac{8\pi^2}{3}\langle u^2\rangle
\]

No compare directamente RMSF de solución con factores B cristalográficos
como si fueran la misma magnitud. El factor B incluye contribuciones del
cristal, refinamiento y desorden, mientras que RMSF depende del
alineamiento, la ventana temporal y el campo de fuerza. Son comparables
principalmente como perfiles cualitativos o mediante un modelo
explícito.

Para proyectar un conjunto de estructuras experimentales sobre una base
PCA deben alinearse, contener una selección atómica común y respetar el
mismo orden. La coincidencia parcial indica que la simulación visita
movimientos compatibles; la ausencia de coincidencia puede reflejar
tiempo insuficiente, ligandos ausentes, condiciones diferentes o sesgo
del campo de fuerza.

\subsubsection{23. Aplicación del esquema a otros tipos de
simulación}\label{aplicaciuxf3n-del-esquema-a-otros-tipos-de-simulaciuxf3n}

La lógica general ---corregir PBC, seleccionar un marco de referencia,
comprobar estabilidad, analizar coordenadas funcionales, estimar
incertidumbre y comparar réplicas--- se transfiere a otros sistemas. Los
comandos no pueden copiarse sin modificar selecciones y significado
físico.

{\def\LTcaptype{none} % do not increment counter
\begin{longtable}[]{@{}
  >{\raggedright\arraybackslash}p{(\linewidth - 4\tabcolsep) * \real{0.1148}}
  >{\raggedright\arraybackslash}p{(\linewidth - 4\tabcolsep) * \real{0.2131}}
  >{\raggedright\arraybackslash}p{(\linewidth - 4\tabcolsep) * \real{0.6721}}@{}}
\toprule\noalign{}
\begin{minipage}[b]{\linewidth}\raggedright
Sistema
\end{minipage} & \begin{minipage}[b]{\linewidth}\raggedright
Parte reutilizable
\end{minipage} & \begin{minipage}[b]{\linewidth}\raggedright
Adaptación necesaria
\end{minipage} \\
\midrule\noalign{}
\endhead
\bottomrule\noalign{}
\endlastfoot
Proteína--ligando & RMSD/RMSF de proteína, PCA, clustering & Ajustar por
proteína; analizar pose, contactos y distancias del ligando por
separado. \\
Oligómero & PCA, contactos, distancias y clustering & Decidir si el
ajuste usa una subunidad o el ensamblado; separar movimientos internos
de reorganización cuaternaria. \\
Membrana--proteína & Análisis de proteína y modos colectivos & Hacer
moléculas enteras; no usar trayectoria ajustada para difusión lipídica,
área o fluctuaciones de caja. \\
Proteína con cofactor & Controles estructurales y PCA & Incluir el
cofactor solo si su parametrización y la selección responden a la
pregunta; controlar geometría de coordinación. \\
Coarse-grained & PCA, clustering y contactos & Usar partículas
equivalentes, interpretar otra resolución y no comparar amplitudes
directamente con all-atom. \\
Réplicas o variantes & Base PCA común y análisis por bloques & Mantener
idénticos átomos, orden, referencia y preprocesamiento; cuantificar
variabilidad entre réplicas. \\
\end{longtable}
}

PCA no debe aplicarse indiscriminadamente a todas las coordenadas del
sistema. Incluir agua, iones o lípidos junto con la proteína genera una
matriz dominada por difusión y permutación de moléculas equivalentes.
Para solvente y membrana son preferibles densidades, difusión,
orientación, contactos o variables colectivas específicas.

\subsubsection{24. Proteínas modeladas o inicialmente
inestables}\label{proteuxednas-modeladas-o-inicialmente-inestables}

Si la estructura proviene de homología, predicción o modelado de bucles:

\begin{itemize}
\tightlist
\item
  revise regiones de baja confianza;
\item
  minimice contactos locales antes de interpretar movimientos;
\item
  use calentamiento gradual si el sistema presenta inestabilidad
  inicial;
\item
  considere un paso de integración menor, por ejemplo 1 fs, durante las
  primeras etapas;
\item
  reduzca las restricciones por etapas en vez de retirarlas
  abruptamente;
\item
  no interprete la relajación del modelo como un cambio biológico.
\end{itemize}

Ejemplo de liberación gradual:

\begin{Shaded}
\begin{Highlighting}[]
\CommentTok{; Etapa inicial}
\DataTypeTok{define }\OtherTok{=}\StringTok{ {-}DPOSRES}
\end{Highlighting}
\end{Shaded}

Puede generar archivos de restricciones con constantes decrecientes, por
ejemplo 1000, 500, 100 y 0 kJ mol⁻¹ nm⁻², manteniendo la duración y los
criterios documentados. Cada etapa debe continuar desde las coordenadas
y velocidades de la anterior.

\subsubsection{25. Cofactores, metales y residuos no
estándar}\label{cofactores-metales-y-residuos-no-estuxe1ndar}

Una proteína que contiene solo residuos reconocidos por el campo de
fuerza puede procesarse directamente mediante \textbf{pdb2gmx}. Un
cofactor orgánico, grupo prostético o residuo modificado requiere
parámetros compatibles.

No todos los cofactores deben tratarse como ligandos independientes.
Existen varios casos:

\begin{itemize}
\tightlist
\item
  \textbf{cofactor no covalente:} puede definirse como otro
  \texttt{{[}\ moleculetype\ {]}};
\item
  \textbf{grupo covalente:} requiere enlaces y parámetros entre proteína
  y cofactor;
\item
  \textbf{residuo modificado:} puede necesitar una entrada RTP y reglas
  de enlace;
\item
  \textbf{metal estructural o catalítico:} exige un modelo específico de
  coordinación;
\item
  \textbf{ion difusible:} puede usar la topología iónica del campo de
  fuerza.
\end{itemize}

Un metal coordinado no debe reemplazarse por un ion genérico sin evaluar
geometría, estado de oxidación, coordinación y transferencia de carga.

Valide siempre que el orden de las coordenadas coincida con
\texttt{{[}\ molecules\ {]}} y que la carga total sea la esperada.

\subsubsection{26. Monómeros, oligómeros y estabilidad
artificial}\label{monuxf3meros-oliguxf3meros-y-estabilidad-artificial}

Cambiar el identificador de cadena o agrupar átomos en
\textbf{index.ndx} no crea enlaces ni estabiliza físicamente un
oligómero. Los grupos sirven para selecciones, acoplamiento, salida o
análisis.

Si una estructura polimérica se desarma durante la simulación,
investigue primero:

\begin{itemize}
\tightlist
\item
  si se simuló la unidad biológica correcta;
\item
  si faltan cadenas, lípidos, ligandos o cofactores;
\item
  si las interfaces dependen del pH o fuerza iónica;
\item
  si se perdieron enlaces covalentes o disulfuro;
\item
  si la orientación inicial es correcta;
\item
  si existen errores de PBC o visualización;
\item
  si la escala temporal observada es compatible con el fenómeno.
\end{itemize}

Las restricciones de posición pueden mantener una geometría, pero
también impedir la dinámica que se pretende estudiar. Para conservar
contactos específicos pueden emplearse restricciones de distancia,
siempre que exista una justificación física o experimental. Congelar
grupos completos altera más drásticamente la dinámica y debe evitarse
salvo casos técnicos muy controlados.

Las restricciones aplicadas durante producción deben informarse porque
modifican el ensamble y limitan la interpretación de fluctuaciones,
RMSD, RMSF y transiciones conformacionales.

\subsubsection{27. Criterio para
continuar}\label{criterio-para-continuar}

Proceda a EM, NVT, NPT y producción únicamente cuando:

\begin{itemize}
\tightlist
\item
  la estructura inicial represente la especie biológica correcta;
\item
  todos los componentes estén parametrizados;
\item
  la topología y las coordenadas tengan el mismo número de átomos;
\item
  la carga total sea razonable;
\item
  las advertencias de \textbf{grompp} estén resueltas;
\item
  caja, solvente e iones sean apropiados;
\item
  se haya definido de antemano qué propiedades se analizarán.
\end{itemize}

Para una proteína monomérica formada exclusivamente por aminoácidos
estándar, el flujo EM → NVT → NPT → MD suele ser suficiente. La
presencia de cofactores, modificaciones, metales, membranas o interfaces
oligoméricas requiere un protocolo específico; no debe resolverse
mediante restricciones genéricas sin modelar antes la química faltante.

\subsubsection{Fuentes}\label{fuentes-2}

\begin{enumerate}
\def\labelenumi{\arabic{enumi}.}
\tightlist
\item
  \href{https://manual.gromacs.org/current/user-guide/system-preparation.html}{Preparación
  de sistemas en GROMACS 2026.3}
\item
  \href{https://manual.gromacs.org/current/onlinehelp/gmx-pdb2gmx.html}{Referencia
  de gmx pdb2gmx}
\item
  \href{https://manual.gromacs.org/current/user-guide/force-fields.html}{Campos
  de fuerza en GROMACS}
\item
  \href{https://manual.gromacs.org/current/onlinehelp/gmx-solvate.html}{Referencia
  de gmx solvate}
\item
  \href{https://manual.gromacs.org/current/user-guide/mdp-options.html}{Opciones
  de archivos MDP}
\item
  \href{https://biophys.uni-saarland.de/courses/comp-mol-bio/pract06/}{Tutorial
  práctico de PCA de trayectorias moleculares, Universidad del Sarre}
\item
  \href{https://manual.gromacs.org/current/onlinehelp/gmx-covar.html\#gmx-covar}{Referencia
  oficial de gmx covar}
\item
  \href{https://manual.gromacs.org/current/onlinehelp/gmx-anaeig.html\#gmx-anaeig}{Referencia
  oficial de gmx anaeig}
\item
  \href{https://pubs.acs.org/doi/10.1021/acs.jpcb.4c04050}{Tutorial
  introductorio de simulación de proteínas con GROMACS, J. Phys. Chem.
  B}
\item
  \href{https://www.tandfonline.com/doi/full/10.1080/23746149.2021.2006080}{Revisión
  sobre aprendizaje automático en el análisis de simulaciones
  biomoleculares}
\item
  \href{https://stemskillslab.com/how-to-do-pca-gromacs-trajectory-essential-dynamics/}{Guía
  práctica complementaria de dinámica esencial con GROMACS}
\end{enumerate}

\subsection{Sistema bifásico}\label{sistema-bifuxe1sico}

\subsubsection{Modelo de dos solventes:
1-octanol--agua}\label{modelo-de-dos-solventes-1-octanolagua}

Un sistema bifásico permite estudiar la distribución de moléculas entre
dos medios líquidos, la estructura de la interfaz, la penetración mutua
de los solventes y, con un protocolo específico, propiedades como el
coeficiente de partición. El sistema no debe tratarse como una caja
acuosa convencional: contiene dos fases condensadas, dos interfaces por
las condiciones periódicas de contorno y una composición que cambia
durante la equilibración.

Este ejemplo describe la construcción de una capa de 1-octanol en
contacto con agua. Los nombres de archivos, residuos y tipos atómicos
son ejemplos; deben coincidir exactamente con la parametrización
utilizada.

El flujo recomendado es:

\begin{Shaded}
\begin{Highlighting}[]
\NormalTok{parametrización del 1{-}octanol}
\NormalTok{    ↓}
\NormalTok{validación de una molécula aislada}
\NormalTok{    ↓}
\NormalTok{construcción del líquido de 1{-}octanol}
\NormalTok{    ↓}
\NormalTok{EM y equilibración del líquido puro}
\NormalTok{    ↓}
\NormalTok{ampliación de la caja en z}
\NormalTok{    ↓}
\NormalTok{incorporación del agua}
\NormalTok{    ↓}
\NormalTok{EM → NVT → NPT o NVT de interfaz}
\NormalTok{    ↓}
\NormalTok{producción y perfiles a lo largo de z}
\end{Highlighting}
\end{Shaded}

\subsubsection{1. Decidir qué sistema físico se quiere
representar}\label{decidir-quuxe9-sistema-fuxedsico-se-quiere-representar}

Antes de construir la caja, defina:

\begin{itemize}
\tightlist
\item
  temperatura y presión;
\item
  campo de fuerza y modelo de agua;
\item
  dimensiones laterales de la interfaz;
\item
  espesor mínimo de cada fase;
\item
  composición inicial;
\item
  presencia de solutos, iones o contraiones;
\item
  duración de la equilibración;
\item
  observable que se calculará.
\end{itemize}

El sistema inicial puede contener 1-octanol puro y agua pura, pero las
fases de equilibrio no serán químicamente puras: parte del agua
ingresará en la región rica en octanol y parte del octanol ingresará en
la región acuosa. Para estudios de partición conviene considerar fases
previamente saturadas o equilibrar el sistema durante un tiempo
suficiente. Una configuración visualmente separada no demuestra que se
haya alcanzado el equilibrio composicional.

Debido a la periodicidad, una lámina de octanol rodeada por agua genera
normalmente dos interfaces aproximadamente paralelas al plano xy. El
tamaño de la caja en z debe evitar que ambas interfaces interactúen de
manera artificial.

\subsubsection{2. Parametrización del
1-octanol}\label{parametrizaciuxf3n-del-1-octanol}

El 1-octanol no es un residuo proteico estándar y no debe procesarse con
\textbf{gmx pdb2gmx} como si fuera un aminoácido. Se necesitan:

\begin{itemize}
\tightlist
\item
  coordenadas tridimensionales;
\item
  topología molecular;
\item
  tipos atómicos;
\item
  cargas parciales;
\item
  parámetros enlazados;
\item
  parámetros de Lennard-Jones compatibles con el campo de fuerza;
\item
  reglas de combinación coherentes con el resto del sistema.
\end{itemize}

SMILES del 1-octanol:

\begin{Shaded}
\begin{Highlighting}[]
\NormalTok{CCCCCCCCO}
\end{Highlighting}
\end{Shaded}

SwissParam genera parámetros compatibles con la familia CHARMM y puede
ser útil para una prueba inicial. No convierte automáticamente esos
parámetros en una parametrización validada para propiedades de partición
o de interfaz. Para un trabajo cuantitativo deben comprobarse, al menos:

\begin{itemize}
\tightlist
\item
  carga total igual a cero;
\item
  geometría y conformaciones;
\item
  densidad del líquido;
\item
  entalpía de vaporización, si corresponde;
\item
  solubilidad mutua con agua;
\item
  distribución de cargas;
\item
  compatibilidad exacta con la versión del campo de fuerza.
\end{itemize}

No mezcle una topología generada para CHARMM con un campo AMBER, GROMOS
u OPLS. Tampoco copie bloques \texttt{{[}\ atomtypes\ {]}} sin comprobar
si los nombres ya existen: dos tipos con el mismo nombre y parámetros
diferentes invalidan la topología.

Una organización sencilla es:

\begin{Shaded}
\begin{Highlighting}[]
\NormalTok{00\_parametros/}
\NormalTok{    octanol.gro}
\NormalTok{    octanol.itp}
\NormalTok{    octanol\_atomtypes.itp}
\NormalTok{01\_octanol\_liquido/}
\NormalTok{02\_interfaz/}
\NormalTok{03\_em/}
\NormalTok{04\_nvt/}
\NormalTok{05\_npt/}
\NormalTok{06\_md/}
\NormalTok{07\_analisis/}
\NormalTok{topol.top}
\end{Highlighting}
\end{Shaded}

\subsubsection{3. Revisar la topología
molecular}\label{revisar-la-topologuxeda-molecular}

El archivo \textbf{octanol.itp} debe contener un único
\texttt{{[}\ moleculetype\ {]}} y las secciones moleculares
correspondientes:

\begin{Shaded}
\begin{Highlighting}[]
\KeywordTok{[ moleculetype ]}
\CommentTok{; nombre    nrexcl}
\DataTypeTok{OCT         3}

\KeywordTok{[ atoms ]}
\CommentTok{; nr  tipo  resnr  residuo  átomo  cgnr  carga  masa}
\CommentTok{; ...}

\KeywordTok{[ bonds ]}
\CommentTok{; ...}

\KeywordTok{[ pairs ]}
\CommentTok{; ...}

\KeywordTok{[ angles ]}
\CommentTok{; ...}

\KeywordTok{[ dihedrals ]}
\CommentTok{; ...}
\end{Highlighting}
\end{Shaded}

Si el generador entrega tipos atómicos nuevos, colóquelos en un archivo
separado que se incluya inmediatamente después del campo de fuerza y
antes de \textbf{octanol.itp}:

\begin{Shaded}
\begin{Highlighting}[]
\CommentTok{\#include "campo\_de\_fuerza.ff/forcefield.itp"}
\CommentTok{\#include "00\_parametros/octanol\_atomtypes.itp"}
\CommentTok{\#include "00\_parametros/octanol.itp"}
\end{Highlighting}
\end{Shaded}

No coloque \texttt{{[}\ atomtypes\ {]}} después de haber comenzado una
definición \texttt{{[}\ moleculetype\ {]}}. El preprocesador de
topologías exige un orden específico de directivas.

Compruebe que la coordenada de una molécula aislada tenga exactamente
los mismos átomos, nombres y orden que la sección
\texttt{{[}\ atoms\ {]}}:

\begin{Shaded}
\begin{Highlighting}[]
\ExtensionTok{gmx}\NormalTok{ check }\AttributeTok{{-}f}\NormalTok{ 00\_parametros/octanol.gro}
\end{Highlighting}
\end{Shaded}

Para inspeccionar la topología expandida:

\begin{Shaded}
\begin{Highlighting}[]
\ExtensionTok{gmx}\NormalTok{ grompp }\DataTypeTok{\textbackslash{}}
    \AttributeTok{{-}f}\NormalTok{ em\_single.mdp }\DataTypeTok{\textbackslash{}}
    \AttributeTok{{-}c}\NormalTok{ 00\_parametros/octanol.gro }\DataTypeTok{\textbackslash{}}
    \AttributeTok{{-}p}\NormalTok{ topol\_single.top }\DataTypeTok{\textbackslash{}}
    \AttributeTok{{-}o}\NormalTok{ octanol\_single.tpr }\DataTypeTok{\textbackslash{}}
    \AttributeTok{{-}pp}\NormalTok{ octanol\_single\_processed.top}
\end{Highlighting}
\end{Shaded}

No use \textbf{-maxwarn} para ocultar incompatibilidades.

\subsubsection{4. Calcular el número inicial de
moléculas}\label{calcular-el-nuxfamero-inicial-de-moluxe9culas}

El número de moléculas no debe elegirse de manera arbitraria. A partir
de una densidad objetivo:

\[
N = \frac{\rho V N_\mathrm{A}}{M}
\]

donde:

\begin{itemize}
\tightlist
\item
  \(N\) es el número de moléculas;
\item
  \(\rho\) es la densidad;
\item
  \(V\) es el volumen;
\item
  \(N_\mathrm{A}\) es la constante de Avogadro;
\item
  \(M\) es la masa molar.
\end{itemize}

Para usar \(\rho\) en g·cm⁻³, \(V\) en nm³ y \(M\) en g·mol⁻¹:

\[
N = \frac{\rho\,V\,10^{-21}\,N_\mathrm{A}}{M}
\]

porque:

\[
1\ \mathrm{nm^3}=10^{-21}\ \mathrm{cm^3}
\]

Para 1-octanol, \(M = 130.23\ \mathrm{g\,mol^{-1}}\). Usando como
ejemplo \(\rho \approx 0.827\ \mathrm{g\,cm^{-3}}\) y una caja de 5 × 5
× 5 nm:

\[
V=125\ \mathrm{nm^3}
\]

\[
N \approx 478\ \text{moléculas}
\]

Por lo tanto, 500 moléculas son razonables como punto de partida para
una caja cercana a 5 nm por lado. En cambio, 500 moléculas en 10 × 10 ×
10 nm corresponden aproximadamente a:

\[
\rho \approx 0.108\ \mathrm{g\,cm^{-3}}
\]

Ese sistema está muy subdensificado y contiene grandes huecos. La
presión NPT inicial puede ser extrema y la caja tendría que contraerse
de manera drástica.

La densidad experimental utilizada debe corresponder a la temperatura
del protocolo. El valor inicial solo aproxima el volumen; la caja debe
equilibrarse con el modelo molecular elegido.

\subsubsection{5. Construir el líquido de
1-octanol}\label{construir-el-luxedquido-de-1-octanol}

Centre una molécula y asegúrese de que las coordenadas estén expresadas
en nanómetros:

\begin{Shaded}
\begin{Highlighting}[]
\ExtensionTok{gmx}\NormalTok{ editconf }\DataTypeTok{\textbackslash{}}
    \AttributeTok{{-}f}\NormalTok{ 00\_parametros/octanol.gro }\DataTypeTok{\textbackslash{}}
    \AttributeTok{{-}o}\NormalTok{ 00\_parametros/octanol\_centered.gro }\DataTypeTok{\textbackslash{}}
    \AttributeTok{{-}center}\NormalTok{ 0 0 0}
\end{Highlighting}
\end{Shaded}

Inserte las moléculas en una caja inicial de 5 × 5 × 5 nm:

\begin{Shaded}
\begin{Highlighting}[]
\FunctionTok{mkdir} \AttributeTok{{-}p}\NormalTok{ 01\_octanol\_liquido}

\ExtensionTok{gmx}\NormalTok{ insert{-}molecules }\DataTypeTok{\textbackslash{}}
    \AttributeTok{{-}ci}\NormalTok{ 00\_parametros/octanol\_centered.gro }\DataTypeTok{\textbackslash{}}
    \AttributeTok{{-}nmol}\NormalTok{ 500 }\DataTypeTok{\textbackslash{}}
    \AttributeTok{{-}box}\NormalTok{ 5 5 5 }\DataTypeTok{\textbackslash{}}
    \AttributeTok{{-}try}\NormalTok{ 500 }\DataTypeTok{\textbackslash{}}
    \AttributeTok{{-}seed}\NormalTok{ 2026 }\DataTypeTok{\textbackslash{}}
    \AttributeTok{{-}o}\NormalTok{ 01\_octanol\_liquido/octanol\_box.gro}
\end{Highlighting}
\end{Shaded}

\textbf{gmx insert-molecules} evita solapamientos mediante radios
atómicos, pero no garantiza que el número solicitado pueda insertarse.
Revise la línea final y use el número realmente añadido en
\texttt{{[}\ molecules\ {]}}.

Si no logra insertar todas las moléculas:

\begin{itemize}
\tightlist
\item
  aumente moderadamente la caja;
\item
  incremente \textbf{-try};
\item
  revise radios y nombres atómicos;
\item
  inserte una cantidad menor y comprima durante una equilibración
  controlada;
\item
  no reduzca agresivamente \textbf{-scale} sin inspeccionar los
  contactos creados.
\end{itemize}

No use una semilla aleatoria indefinida si desea reproducibilidad.

Topología inicial:

\begin{Shaded}
\begin{Highlighting}[]
\CommentTok{\#include "campo\_de\_fuerza.ff/forcefield.itp"}
\CommentTok{\#include "00\_parametros/octanol\_atomtypes.itp"}
\CommentTok{\#include "00\_parametros/octanol.itp"}

\KeywordTok{[ system ]}
\DataTypeTok{1{-}octanol líquido}

\KeywordTok{[ molecules ]}
\CommentTok{; molécula    cantidad}
\DataTypeTok{OCT           500}
\end{Highlighting}
\end{Shaded}

El nombre \textbf{OCT} debe ser idéntico al definido en
\texttt{{[}\ moleculetype\ {]}}.

\subsubsection{6. Minimizar el líquido de
1-octanol}\label{minimizar-el-luxedquido-de-1-octanol}

Archivo \textbf{em\_octanol.mdp}:

\begin{Shaded}
\begin{Highlighting}[]
\DataTypeTok{title           }\OtherTok{=}\StringTok{ Minimización del líquido de 1{-}octanol}
\DataTypeTok{integrator      }\OtherTok{=}\StringTok{ steep}
\DataTypeTok{nsteps          }\OtherTok{=}\StringTok{ }\DecValTok{50000}
\DataTypeTok{emtol           }\OtherTok{=}\StringTok{ }\FloatTok{1000.0}
\DataTypeTok{emstep          }\OtherTok{=}\StringTok{ }\FloatTok{0.01}

\DataTypeTok{cutoff{-}scheme   }\OtherTok{=}\StringTok{ Verlet}
\DataTypeTok{nstlist         }\OtherTok{=}\StringTok{ }\DecValTok{20}
\DataTypeTok{rlist           }\OtherTok{=}\StringTok{ }\FloatTok{1.2}
\DataTypeTok{coulombtype     }\OtherTok{=}\StringTok{ PME}
\DataTypeTok{rcoulomb        }\OtherTok{=}\StringTok{ }\FloatTok{1.2}
\DataTypeTok{vdwtype         }\OtherTok{=}\StringTok{ Cut{-}off}
\DataTypeTok{rvdw            }\OtherTok{=}\StringTok{ }\FloatTok{1.2}
\DataTypeTok{pbc             }\OtherTok{=}\StringTok{ xyz}
\end{Highlighting}
\end{Shaded}

Los cortes son ejemplos. Deben reemplazarse por los valores recomendados
para el campo de fuerza seleccionado.

\begin{Shaded}
\begin{Highlighting}[]
\FunctionTok{mkdir} \AttributeTok{{-}p}\NormalTok{ 03\_em}

\ExtensionTok{gmx}\NormalTok{ grompp }\DataTypeTok{\textbackslash{}}
    \AttributeTok{{-}f}\NormalTok{ em\_octanol.mdp }\DataTypeTok{\textbackslash{}}
    \AttributeTok{{-}c}\NormalTok{ 01\_octanol\_liquido/octanol\_box.gro }\DataTypeTok{\textbackslash{}}
    \AttributeTok{{-}p}\NormalTok{ topol.top }\DataTypeTok{\textbackslash{}}
    \AttributeTok{{-}o}\NormalTok{ 03\_em/octanol\_em.tpr }\DataTypeTok{\textbackslash{}}
    \AttributeTok{{-}pp}\NormalTok{ 03\_em/octanol\_processed.top}

\ExtensionTok{gmx}\NormalTok{ mdrun }\DataTypeTok{\textbackslash{}}
    \AttributeTok{{-}deffnm}\NormalTok{ 03\_em/octanol\_em }\DataTypeTok{\textbackslash{}}
    \AttributeTok{{-}v}
\end{Highlighting}
\end{Shaded}

Revise energía potencial, fuerza máxima, contactos anómalos y valores
NaN. La convergencia numérica de la minimización no demuestra que la
densidad o la estructura del líquido sean correctas.

\subsubsection{7. Equilibrar primero la fase de
1-octanol}\label{equilibrar-primero-la-fase-de-1-octanol}

Una trayectoria de 1 ps es insuficiente para equilibrar un líquido
construido por inserción aleatoria. Use primero NVT para estabilizar la
temperatura y luego NPT para ajustar densidad y volumen.

Ejemplo NVT:

\begin{Shaded}
\begin{Highlighting}[]
\DataTypeTok{title           }\OtherTok{=}\StringTok{ NVT del 1{-}octanol}
\DataTypeTok{integrator      }\OtherTok{=}\StringTok{ md}
\DataTypeTok{dt              }\OtherTok{=}\StringTok{ }\FloatTok{0.002}
\DataTypeTok{nsteps          }\OtherTok{=}\StringTok{ }\DecValTok{250000}
\DataTypeTok{continuation    }\OtherTok{=}\StringTok{ }\KeywordTok{no}
\DataTypeTok{gen{-}vel         }\OtherTok{=}\StringTok{ }\KeywordTok{yes}
\DataTypeTok{gen{-}temp        }\OtherTok{=}\StringTok{ }\DecValTok{300}
\DataTypeTok{gen{-}seed        }\OtherTok{=}\StringTok{ }\DecValTok{2026}

\DataTypeTok{constraints     }\OtherTok{=}\StringTok{ h{-}bonds}
\DataTypeTok{tcoupl          }\OtherTok{=}\StringTok{ V{-}rescale}
\DataTypeTok{tc{-}grps         }\OtherTok{=}\StringTok{ System}
\DataTypeTok{tau{-}t           }\OtherTok{=}\StringTok{ }\FloatTok{1.0}
\DataTypeTok{ref{-}t           }\OtherTok{=}\StringTok{ }\DecValTok{300}

\DataTypeTok{pcoupl          }\OtherTok{=}\StringTok{ }\KeywordTok{no}
\DataTypeTok{pbc             }\OtherTok{=}\StringTok{ xyz}
\end{Highlighting}
\end{Shaded}

Con 250000 pasos de 0.002 ps se simulan 500 ps.

\begin{Shaded}
\begin{Highlighting}[]
\FunctionTok{mkdir} \AttributeTok{{-}p}\NormalTok{ 04\_nvt}

\ExtensionTok{gmx}\NormalTok{ grompp }\DataTypeTok{\textbackslash{}}
    \AttributeTok{{-}f}\NormalTok{ nvt\_octanol.mdp }\DataTypeTok{\textbackslash{}}
    \AttributeTok{{-}c}\NormalTok{ 03\_em/octanol\_em.gro }\DataTypeTok{\textbackslash{}}
    \AttributeTok{{-}p}\NormalTok{ topol.top }\DataTypeTok{\textbackslash{}}
    \AttributeTok{{-}o}\NormalTok{ 04\_nvt/octanol\_nvt.tpr}

\ExtensionTok{gmx}\NormalTok{ mdrun }\DataTypeTok{\textbackslash{}}
    \AttributeTok{{-}deffnm}\NormalTok{ 04\_nvt/octanol\_nvt }\DataTypeTok{\textbackslash{}}
    \AttributeTok{{-}v}
\end{Highlighting}
\end{Shaded}

Ejemplo NPT isotrópico:

\begin{Shaded}
\begin{Highlighting}[]
\DataTypeTok{title            }\OtherTok{=}\StringTok{ NPT del 1{-}octanol}
\DataTypeTok{integrator       }\OtherTok{=}\StringTok{ md}
\DataTypeTok{dt               }\OtherTok{=}\StringTok{ }\FloatTok{0.002}
\DataTypeTok{nsteps           }\OtherTok{=}\StringTok{ }\DecValTok{2500000}
\DataTypeTok{continuation     }\OtherTok{=}\StringTok{ }\KeywordTok{yes}
\DataTypeTok{gen{-}vel          }\OtherTok{=}\StringTok{ }\KeywordTok{no}

\DataTypeTok{constraints      }\OtherTok{=}\StringTok{ h{-}bonds}
\DataTypeTok{tcoupl           }\OtherTok{=}\StringTok{ V{-}rescale}
\DataTypeTok{tc{-}grps          }\OtherTok{=}\StringTok{ System}
\DataTypeTok{tau{-}t            }\OtherTok{=}\StringTok{ }\FloatTok{1.0}
\DataTypeTok{ref{-}t            }\OtherTok{=}\StringTok{ }\DecValTok{300}

\DataTypeTok{pcoupl            }\OtherTok{=}\StringTok{ C{-}rescale}
\DataTypeTok{pcoupltype        }\OtherTok{=}\StringTok{ isotropic}
\DataTypeTok{tau{-}p             }\OtherTok{=}\StringTok{ }\FloatTok{5.0}
\DataTypeTok{ref{-}p             }\OtherTok{=}\StringTok{ }\FloatTok{1.0}
\DataTypeTok{compressibility   }\OtherTok{=}\StringTok{ 8.0e{-}5}

\DataTypeTok{pbc              }\OtherTok{=}\StringTok{ xyz}
\end{Highlighting}
\end{Shaded}

Aquí se simulan 5 ns. La compresibilidad es un valor inicial ilustrativo
y debe reemplazarse por un valor justificado para el líquido y las
condiciones elegidas.

\begin{Shaded}
\begin{Highlighting}[]
\FunctionTok{mkdir} \AttributeTok{{-}p}\NormalTok{ 05\_npt}

\ExtensionTok{gmx}\NormalTok{ grompp }\DataTypeTok{\textbackslash{}}
    \AttributeTok{{-}f}\NormalTok{ npt\_octanol.mdp }\DataTypeTok{\textbackslash{}}
    \AttributeTok{{-}c}\NormalTok{ 04\_nvt/octanol\_nvt.gro }\DataTypeTok{\textbackslash{}}
    \AttributeTok{{-}t}\NormalTok{ 04\_nvt/octanol\_nvt.cpt }\DataTypeTok{\textbackslash{}}
    \AttributeTok{{-}p}\NormalTok{ topol.top }\DataTypeTok{\textbackslash{}}
    \AttributeTok{{-}o}\NormalTok{ 05\_npt/octanol\_npt.tpr}

\ExtensionTok{gmx}\NormalTok{ mdrun }\DataTypeTok{\textbackslash{}}
    \AttributeTok{{-}deffnm}\NormalTok{ 05\_npt/octanol\_npt }\DataTypeTok{\textbackslash{}}
    \AttributeTok{{-}v}
\end{Highlighting}
\end{Shaded}

Controle densidad, volumen, presión y energía:

\begin{Shaded}
\begin{Highlighting}[]
\KeywordTok{(}\BuiltInTok{echo}\NormalTok{ Density}\KeywordTok{;} \BuiltInTok{echo}\NormalTok{ Volume}\KeywordTok{;} \BuiltInTok{echo}\NormalTok{ Pressure}\KeywordTok{;} \BuiltInTok{echo}\NormalTok{ Potential}\KeywordTok{;} \BuiltInTok{echo}\NormalTok{ 0}\KeywordTok{)} \KeywordTok{|} \DataTypeTok{\textbackslash{}}
\ExtensionTok{gmx}\NormalTok{ energy }\DataTypeTok{\textbackslash{}}
    \AttributeTok{{-}f}\NormalTok{ 05\_npt/octanol\_npt.edr }\DataTypeTok{\textbackslash{}}
    \AttributeTok{{-}o}\NormalTok{ 07\_analisis/octanol\_bulk\_properties.xvg}
\end{Highlighting}
\end{Shaded}

La presión instantánea tiene fluctuaciones grandes. Compare promedios
por bloques y compruebe que densidad y volumen hayan alcanzado una
región estacionaria.

\subsubsection{8. Preparar la lámina de
1-octanol}\label{preparar-la-luxe1mina-de-1-octanol}

Use la última configuración equilibrada, no un PDB extraído
arbitrariamente de una trayectoria. El formato GRO conserva la caja con
precisión suficiente y evita pérdidas innecesarias de información.

Primero consulte las dimensiones finales:

\begin{Shaded}
\begin{Highlighting}[]
\ExtensionTok{gmx}\NormalTok{ check }\AttributeTok{{-}f}\NormalTok{ 05\_npt/octanol\_npt.gro}
\end{Highlighting}
\end{Shaded}

Supóngase, solo como ejemplo, que la caja equilibrada mide
aproximadamente 5 × 5 × 5 nm. Amplíe únicamente z para crear espacio
para el agua:

\begin{Shaded}
\begin{Highlighting}[]
\FunctionTok{mkdir} \AttributeTok{{-}p}\NormalTok{ 02\_interfaz}

\ExtensionTok{gmx}\NormalTok{ editconf }\DataTypeTok{\textbackslash{}}
    \AttributeTok{{-}f}\NormalTok{ 05\_npt/octanol\_npt.gro }\DataTypeTok{\textbackslash{}}
    \AttributeTok{{-}o}\NormalTok{ 02\_interfaz/octanol\_slab\_box.gro }\DataTypeTok{\textbackslash{}}
    \AttributeTok{{-}box}\NormalTok{ 5 5 12 }\DataTypeTok{\textbackslash{}}
    \AttributeTok{{-}center}\NormalTok{ 2.5 2.5 6.0}
\end{Highlighting}
\end{Shaded}

Use los valores reales de \(L_x\) y \(L_y\) de la fase equilibrada.
Cambiarlos en este paso impone una deformación lateral. La región vacía
debe quedar distribuida a ambos lados de la lámina para generar dos
interfaces equivalentes.

Compruebe visualmente:

\begin{itemize}
\tightlist
\item
  que el octanol forme una única lámina continua;
\item
  que no quede dividido por una representación incorrecta de PBC;
\item
  que exista espacio suficiente para el agua;
\item
  que los grupos hidroxilo no hayan sido orientados artificialmente de
  forma uniforme.
\end{itemize}

\subsubsection{9. Incorporar el agua}\label{incorporar-el-agua}

Incluya la topología de agua compatible con el campo de fuerza antes de
la sección \texttt{{[}\ system\ {]}}:

\begin{Shaded}
\begin{Highlighting}[]
\CommentTok{\#include "campo\_de\_fuerza.ff/forcefield.itp"}
\CommentTok{\#include "00\_parametros/octanol\_atomtypes.itp"}
\CommentTok{\#include "00\_parametros/octanol.itp"}
\CommentTok{\#include "campo\_de\_fuerza.ff/modelo\_de\_agua.itp"}

\KeywordTok{[ system ]}
\DataTypeTok{Interfaz 1{-}octanol–agua}

\KeywordTok{[ molecules ]}
\CommentTok{; molécula    cantidad}
\DataTypeTok{OCT           500}
\end{Highlighting}
\end{Shaded}

Solvate usando explícitamente una caja de agua:

\begin{Shaded}
\begin{Highlighting}[]
\ExtensionTok{gmx}\NormalTok{ solvate }\DataTypeTok{\textbackslash{}}
    \AttributeTok{{-}cp}\NormalTok{ 02\_interfaz/octanol\_slab\_box.gro }\DataTypeTok{\textbackslash{}}
    \AttributeTok{{-}cs}\NormalTok{ spc216.gro }\DataTypeTok{\textbackslash{}}
    \AttributeTok{{-}o}\NormalTok{ 02\_interfaz/octanol\_water.gro }\DataTypeTok{\textbackslash{}}
    \AttributeTok{{-}p}\NormalTok{ topol.top}
\end{Highlighting}
\end{Shaded}

\textbf{spc216.gro} aporta una configuración geométrica de agua de tres
sitios; la interacción efectiva queda definida por la topología
incluida. Verifique que el modelo de agua sea el recomendado para el
campo de fuerza.

\textbf{gmx solvate} elimina moléculas que solapan con el octanol y
actualiza el número de agua en \texttt{{[}\ molecules\ {]}}. La salida
final debe quedar, por ejemplo:

\begin{Shaded}
\begin{Highlighting}[]
\KeywordTok{[ molecules ]}
\CommentTok{; molécula    cantidad}
\DataTypeTok{OCT           500}
\DataTypeTok{SOL           4442}
\end{Highlighting}
\end{Shaded}

El valor 4442 no es universal. Depende de las dimensiones finales, la
densidad de la lámina y los criterios geométricos de solvatación. Use el
número informado por su propia ejecución.

Revise el orden: las coordenadas contienen primero OCT y después SOL,
por lo que \texttt{{[}\ molecules\ {]}} debe seguir el mismo orden.

Compruebe el sistema:

\begin{Shaded}
\begin{Highlighting}[]
\ExtensionTok{gmx}\NormalTok{ check }\AttributeTok{{-}f}\NormalTok{ 02\_interfaz/octanol\_water.gro}
\end{Highlighting}
\end{Shaded}

\subsubsection{10. Minimización y equilibración de la
interfaz}\label{minimizaciuxf3n-y-equilibraciuxf3n-de-la-interfaz}

La interfaz recién construida contiene contactos y una distribución no
equilibrada de ambos líquidos. Ejecute nuevamente EM, NVT y una
equilibración apropiada de volumen o área.

\begin{Shaded}
\begin{Highlighting}[]
\ExtensionTok{gmx}\NormalTok{ grompp }\DataTypeTok{\textbackslash{}}
    \AttributeTok{{-}f}\NormalTok{ em\_interface.mdp }\DataTypeTok{\textbackslash{}}
    \AttributeTok{{-}c}\NormalTok{ 02\_interfaz/octanol\_water.gro }\DataTypeTok{\textbackslash{}}
    \AttributeTok{{-}p}\NormalTok{ topol.top }\DataTypeTok{\textbackslash{}}
    \AttributeTok{{-}o}\NormalTok{ 03\_em/interface\_em.tpr }\DataTypeTok{\textbackslash{}}
    \AttributeTok{{-}pp}\NormalTok{ 03\_em/interface\_processed.top}

\ExtensionTok{gmx}\NormalTok{ mdrun }\DataTypeTok{\textbackslash{}}
    \AttributeTok{{-}deffnm}\NormalTok{ 03\_em/interface\_em }\DataTypeTok{\textbackslash{}}
    \AttributeTok{{-}v}
\end{Highlighting}
\end{Shaded}

NVT:

\begin{Shaded}
\begin{Highlighting}[]
\ExtensionTok{gmx}\NormalTok{ grompp }\DataTypeTok{\textbackslash{}}
    \AttributeTok{{-}f}\NormalTok{ nvt\_interface.mdp }\DataTypeTok{\textbackslash{}}
    \AttributeTok{{-}c}\NormalTok{ 03\_em/interface\_em.gro }\DataTypeTok{\textbackslash{}}
    \AttributeTok{{-}p}\NormalTok{ topol.top }\DataTypeTok{\textbackslash{}}
    \AttributeTok{{-}o}\NormalTok{ 04\_nvt/interface\_nvt.tpr}

\ExtensionTok{gmx}\NormalTok{ mdrun }\DataTypeTok{\textbackslash{}}
    \AttributeTok{{-}deffnm}\NormalTok{ 04\_nvt/interface\_nvt }\DataTypeTok{\textbackslash{}}
    \AttributeTok{{-}v}
\end{Highlighting}
\end{Shaded}

Para una interfaz plana existen dos estrategias principales:

{\def\LTcaptype{none} % do not increment counter
\begin{longtable}[]{@{}
  >{\raggedright\arraybackslash}p{(\linewidth - 4\tabcolsep) * \real{0.2308}}
  >{\raggedright\arraybackslash}p{(\linewidth - 4\tabcolsep) * \real{0.3333}}
  >{\raggedright\arraybackslash}p{(\linewidth - 4\tabcolsep) * \real{0.4359}}@{}}
\toprule\noalign{}
\begin{minipage}[b]{\linewidth}\raggedright
Estrategia
\end{minipage} & \begin{minipage}[b]{\linewidth}\raggedright
Ventaja
\end{minipage} & \begin{minipage}[b]{\linewidth}\raggedright
Limitación
\end{minipage} \\
\midrule\noalign{}
\endhead
\bottomrule\noalign{}
\endlastfoot
NPT semiisotrópico & Permite ajustar por separado el plano xy y el eje
z. & Las fluctuaciones del área pueden modificar la interfaz. \\
NVT con caja previamente equilibrada & Mantiene fija el área
interfacial. & Requiere haber determinado antes dimensiones y densidades
adecuadas. \\
\end{longtable}
}

Ejemplo NPT semiisotrópico:

\begin{Shaded}
\begin{Highlighting}[]
\DataTypeTok{pcoupl            }\OtherTok{=}\StringTok{ C{-}rescale}
\DataTypeTok{pcoupltype        }\OtherTok{=}\StringTok{ semiisotropic}
\DataTypeTok{tau{-}p             }\OtherTok{=}\StringTok{ }\FloatTok{5.0}
\DataTypeTok{ref{-}p             }\OtherTok{=}\StringTok{ 1.0 1.0}
\DataTypeTok{compressibility   }\OtherTok{=}\StringTok{ 4.5e{-}5 4.5e{-}5}
\end{Highlighting}
\end{Shaded}

Los dos valores corresponden al plano xy y al eje z. No copie
automáticamente la compresibilidad del agua para todo el sistema: debe
justificarse según el protocolo y comprobarse que la caja no derive o
colapse.

El acoplamiento por tensión superficial también existe, pero no debe
utilizarse solo porque el sistema tenga una interfaz. Requiere una
tensión objetivo y una interpretación consistente del ensamble.

Para producción, continúe desde el checkpoint:

\begin{Shaded}
\begin{Highlighting}[]
\ExtensionTok{gmx}\NormalTok{ grompp }\DataTypeTok{\textbackslash{}}
    \AttributeTok{{-}f}\NormalTok{ md\_interface.mdp }\DataTypeTok{\textbackslash{}}
    \AttributeTok{{-}c}\NormalTok{ 05\_npt/interface\_npt.gro }\DataTypeTok{\textbackslash{}}
    \AttributeTok{{-}t}\NormalTok{ 05\_npt/interface\_npt.cpt }\DataTypeTok{\textbackslash{}}
    \AttributeTok{{-}p}\NormalTok{ topol.top }\DataTypeTok{\textbackslash{}}
    \AttributeTok{{-}o}\NormalTok{ 06\_md/interface\_md.tpr}

\ExtensionTok{gmx}\NormalTok{ mdrun }\DataTypeTok{\textbackslash{}}
    \AttributeTok{{-}deffnm}\NormalTok{ 06\_md/interface\_md }\DataTypeTok{\textbackslash{}}
    \AttributeTok{{-}v}
\end{Highlighting}
\end{Shaded}

\subsubsection{11. Controles mínimos de la
trayectoria}\label{controles-muxednimos-de-la-trayectoria}

Corrija únicamente la representación periódica necesaria para
visualizar. No centre el sistema de una manera que desplace la interfaz
entre fotogramas sin documentarlo.

Extraiga una configuración:

\begin{Shaded}
\begin{Highlighting}[]
\BuiltInTok{echo}\NormalTok{ System }\KeywordTok{|} \DataTypeTok{\textbackslash{}}
\ExtensionTok{gmx}\NormalTok{ trjconv }\DataTypeTok{\textbackslash{}}
    \AttributeTok{{-}s}\NormalTok{ 06\_md/interface\_md.tpr }\DataTypeTok{\textbackslash{}}
    \AttributeTok{{-}f}\NormalTok{ 06\_md/interface\_md.xtc }\DataTypeTok{\textbackslash{}}
    \AttributeTok{{-}o}\NormalTok{ 07\_analisis/interface\_last.gro }\DataTypeTok{\textbackslash{}}
    \AttributeTok{{-}dump}\NormalTok{ 100000}
\end{Highlighting}
\end{Shaded}

El tiempo de \textbf{-dump} se expresa en ps. En este ejemplo, 100000 ps
equivalen a 100 ns.

Perfil de densidad de masa a lo largo de z:

\begin{Shaded}
\begin{Highlighting}[]
\ExtensionTok{gmx}\NormalTok{ density }\DataTypeTok{\textbackslash{}}
    \AttributeTok{{-}s}\NormalTok{ 06\_md/interface\_md.tpr }\DataTypeTok{\textbackslash{}}
    \AttributeTok{{-}f}\NormalTok{ 06\_md/interface\_md.xtc }\DataTypeTok{\textbackslash{}}
    \AttributeTok{{-}n}\NormalTok{ index.ndx }\DataTypeTok{\textbackslash{}}
    \AttributeTok{{-}d}\NormalTok{ Z }\DataTypeTok{\textbackslash{}}
    \AttributeTok{{-}sl}\NormalTok{ 200 }\DataTypeTok{\textbackslash{}}
    \AttributeTok{{-}dens}\NormalTok{ mass }\DataTypeTok{\textbackslash{}}
    \AttributeTok{{-}o}\NormalTok{ 07\_analisis/density\_z.xvg}
\end{Highlighting}
\end{Shaded}

Cree grupos separados para OCT y SOL y analice ambos perfiles. Un
sistema bifásico equilibrado debe mostrar:

\begin{itemize}
\tightlist
\item
  una región rica en agua;
\item
  una región rica en octanol;
\item
  dos zonas interfaciales;
\item
  densidades aproximadamente constantes en el centro de cada fase, si el
  espesor es suficiente;
\item
  ausencia de deriva sistemática del espesor o de la posición de la
  lámina.
\end{itemize}

La anchura de la interfaz depende del binning. Repita el cálculo con
distintos valores de \textbf{-sl} para verificar que la conclusión no
sea un artefacto de discretización.

\subsubsection{12. Propiedades y análisis
adicionales}\label{propiedades-y-anuxe1lisis-adicionales}

\paragraph{Penetración mutua de los
solventes}\label{penetraciuxf3n-mutua-de-los-solventes}

Los perfiles de densidad de OCT y SOL permiten estimar la presencia de
agua en la fase rica en octanol y de octanol en la fase acuosa. Descarte
la etapa transitoria y compare bloques temporales independientes.

\paragraph{Orientación del
1-octanol}\label{orientaciuxf3n-del-1-octanol}

Puede analizarse el ángulo entre el vector C1--O del 1-octanol y el eje
z:

\begin{Shaded}
\begin{Highlighting}[]
\ExtensionTok{gmx}\NormalTok{ gangle }\DataTypeTok{\textbackslash{}}
    \AttributeTok{{-}s}\NormalTok{ 06\_md/interface\_md.tpr }\DataTypeTok{\textbackslash{}}
    \AttributeTok{{-}f}\NormalTok{ 06\_md/interface\_md.xtc }\DataTypeTok{\textbackslash{}}
    \AttributeTok{{-}n}\NormalTok{ index.ndx }\DataTypeTok{\textbackslash{}}
    \AttributeTok{{-}g1}\NormalTok{ vector }\DataTypeTok{\textbackslash{}}
    \AttributeTok{{-}group1} \StringTok{\textquotesingle{}vector connecting atomnr START END\textquotesingle{}} \DataTypeTok{\textbackslash{}}
    \AttributeTok{{-}g2}\NormalTok{ z }\DataTypeTok{\textbackslash{}}
    \AttributeTok{{-}oav}\NormalTok{ 07\_analisis/octanol\_orientation.xvg}
\end{Highlighting}
\end{Shaded}

La selección debe adaptarse a los índices reales. Para obtener una
distribución molecular completa puede ser necesario definir pares
equivalentes para todas las moléculas.

\paragraph{Tensión interfacial}\label{tensiuxf3n-interfacial}

Para una lámina con dos interfaces planas normales a z:

\[
\gamma =
\frac{L_z}{2}
\left[
P_{zz}-
\frac{P_{xx}+P_{yy}}{2}
\right]
\]

El factor 1/2 aparece porque la caja periódica contiene dos interfaces.
\(P_{xx}\), \(P_{yy}\) y \(P_{zz}\) son los componentes diagonales del
tensor de presión.

Extraiga los componentes con \textbf{gmx energy}:

\begin{Shaded}
\begin{Highlighting}[]
\KeywordTok{(}\BuiltInTok{echo}\NormalTok{ Pres{-}XX}\KeywordTok{;} \BuiltInTok{echo}\NormalTok{ Pres{-}YY}\KeywordTok{;} \BuiltInTok{echo}\NormalTok{ Pres{-}ZZ}\KeywordTok{;} \BuiltInTok{echo}\NormalTok{ Box{-}Z}\KeywordTok{;} \BuiltInTok{echo}\NormalTok{ 0}\KeywordTok{)} \KeywordTok{|} \DataTypeTok{\textbackslash{}}
\ExtensionTok{gmx}\NormalTok{ energy }\DataTypeTok{\textbackslash{}}
    \AttributeTok{{-}f}\NormalTok{ 06\_md/interface\_md.edr }\DataTypeTok{\textbackslash{}}
    \AttributeTok{{-}o}\NormalTok{ 07\_analisis/pressure\_tensor.xvg}
\end{Highlighting}
\end{Shaded}

La tensión interfacial converge lentamente porque el tensor de presión
es ruidoso. Deben usarse trayectorias suficientemente largas, promedios
por bloques y unidades consistentes. En GROMACS, presión se informa en
bar y longitud en nm; el resultado no queda automáticamente en mN·m⁻¹
sin conversión.

La equivalencia útil es:

\[
1\ \mathrm{bar\,nm}=0.1\ \mathrm{mN\,m^{-1}}
\]

\paragraph{Coeficiente de partición}\label{coeficiente-de-particiuxf3n}

Contar espontáneamente un soluto en cada fase puede servir si ocurren
muchas transiciones reversibles. Para moléculas con barreras altas, una
única trayectoria suele quedar atrapada en una fase y no permite estimar
un coeficiente de partición confiable. En esos casos se requieren
métodos de energía libre y réplicas, no solo una caja bifásica más
larga.

\subsubsection{13. Errores frecuentes}\label{errores-frecuentes}

\begin{itemize}
\tightlist
\item
  Usar 500 moléculas en una caja de 10 nm por lado sin calcular la
  densidad.
\item
  Parametrizar octanol con un campo de fuerza y combinarlo con otro.
\item
  Duplicar tipos atómicos con nombres iguales.
\item
  Confundir el solvente de coordenadas \textbf{spc216.gro} con el modelo
  definido en la topología.
\item
  Conservar el número de moléculas solicitado cuando
  \textbf{insert-molecules} insertó menos.
\item
  Escribir \texttt{{[}\ molecules\ {]}} en un orden distinto al archivo
  de coordenadas.
\item
  Extraer un PDB intermedio y perder precisión de caja o velocidades sin
  necesidad.
\item
  Equilibrar solo 1 ps.
\item
  Interpretar separación visual como equilibrio termodinámico.
\item
  Calcular promedios incluyendo la construcción y relajación inicial.
\item
  Aplicar un barostato isotrópico a una interfaz sin evaluar la
  deformación del área.
\item
  Interpretar ausencia de cruces del soluto como evidencia de partición
  estable.
\end{itemize}

\subsubsection{14. Criterios para iniciar la
producción}\label{criterios-para-iniciar-la-producciuxf3n}

Continúe únicamente cuando:

\begin{itemize}
\tightlist
\item
  la carga y los tipos atómicos sean correctos;
\item
  la topología expandida no contenga conflictos;
\item
  el número y orden de moléculas coincidan con las coordenadas;
\item
  no existan huecos ni solapamientos graves;
\item
  el octanol puro reproduzca razonablemente la densidad objetivo;
\item
  temperatura, energía y dimensiones de caja sean estacionarias;
\item
  se hayan formado dos regiones de densidad definidas;
\item
  la fase rica en cada solvente tenga espesor suficiente;
\item
  las advertencias de \textbf{gmx grompp} estén comprendidas y
  resueltas;
\item
  el protocolo de producción y los análisis se hayan definido de
  antemano.
\end{itemize}

\subsubsection{Fuentes}\label{fuentes-3}

\begin{enumerate}
\def\labelenumi{\arabic{enumi}.}
\tightlist
\item
  \href{https://manual.gromacs.org/current/user-guide/system-preparation.html}{Preparación
  de sistemas en GROMACS 2026.3}
\item
  \href{https://manual.gromacs.org/current/onlinehelp/gmx-insert-molecules.html}{Referencia
  de gmx insert-molecules}
\item
  \href{https://manual.gromacs.org/current/onlinehelp/gmx-solvate.html}{Referencia
  de gmx solvate}
\item
  \href{https://manual.gromacs.org/current/onlinehelp/gmx-editconf.html}{Referencia
  de gmx editconf}
\item
  \href{https://manual.gromacs.org/current/onlinehelp/gmx-trjconv.html}{Referencia
  de gmx trjconv}
\item
  \href{https://manual.gromacs.org/current/reference-manual/definitions.html}{Unidades
  de GROMACS}
\item
  \href{https://manual.gromacs.org/current/user-guide/mdp-options.html}{Opciones
  de archivos MDP}
\end{enumerate}

\subsection{Proteína de membrana: construcción de una bicapa e inserción
de una
acuaporina}\label{proteuxedna-de-membrana-construcciuxf3n-de-una-bicapa-e-inserciuxf3n-de-una-acuaporina}

\subsubsection{Introducción}\label{introducciuxf3n}

Las proteínas de membrana no deben simularse como proteínas solubles
rodeadas únicamente por agua. La bicapa aporta un entorno anisotrópico,
ejerce presión lateral, establece interfaces polares y apolares, modula
la conformación de la proteína y puede ocupar sitios de unión
específicos. La composición lipídica, la orientación, el ensamblado
oligomérico y el protocolo de equilibración forman parte del modelo
físico.

CHARMM-GUI Membrane Builder automatiza gran parte de esta preparación:
lectura y reparación de la estructura, orientación, construcción de
bicapas homogéneas o heterogéneas, agua en cavidades, solvatación,
incorporación de iones y generación de archivos para GROMACS. La
automatización reduce errores mecánicos, pero no decide qué unidad
biológica, estado de protonación, composición de membrana o ensamble
experimental son correctos.

Este tutorial usa una acuaporina como caso principal y contempla dos
resoluciones:

\begin{itemize}
\tightlist
\item
  \textbf{all-atom:} proteína, lípidos, agua e iones representados átomo
  por átomo;
\item
  \textbf{coarse-grained (CG):} varios átomos se agrupan en partículas
  efectivas.
\end{itemize}

Ambos modelos responden preguntas diferentes. No deben mezclarse
archivos de fuerza, topologías ni pasos de integración entre CHARMM36 y
Martini.

Flujo general:

\begin{Shaded}
\begin{Highlighting}[]
\NormalTok{estructura y pregunta biológica}
\NormalTok{    ↓}
\NormalTok{unidad biológica y protonación}
\NormalTok{    ↓}
\NormalTok{orientación respecto de la membrana}
\NormalTok{    ↓}
\NormalTok{composición y tamaño de la bicapa}
\NormalTok{    ↓}
\NormalTok{construcción en CHARMM{-}GUI}
\NormalTok{    ↓}
\NormalTok{inspección de topología y coordenadas}
\NormalTok{    ↓}
\NormalTok{EM y equilibración con restricciones decrecientes}
\NormalTok{    ↓}
\NormalTok{producción}
\NormalTok{    ↓}
\NormalTok{control de proteína, bicapa, agua, iones y poro}
\end{Highlighting}
\end{Shaded}

\subsubsection{1. Definir el objetivo antes de construir el
sistema}\label{definir-el-objetivo-antes-de-construir-el-sistema}

Registre antes de usar Membrane Builder:

\begin{itemize}
\tightlist
\item
  especie, tejido y membrana de origen;
\item
  isoforma y secuencia exacta;
\item
  ensamblado oligomérico;
\item
  estructura experimental o modelo utilizado;
\item
  residuos faltantes, mutaciones y modificaciones;
\item
  ligandos, cofactores, iones o lípidos estructurales;
\item
  orientación y topología transmembrana;
\item
  composición de cada monocapa;
\item
  pH, temperatura, presión y fuerza iónica;
\item
  variable principal que se analizará;
\item
  número de réplicas y duración prevista.
\end{itemize}

Una bicapa de POPC puro puede ser un modelo controlado útil, pero no
representa de forma universal una membrana plasmática, bacteriana,
mitocondrial o del cristalino. La composición lipídica puede modificar
estabilidad, inclinación, oligomerización y función de una proteína de
membrana.

\subsubsection{2. Particularidades de las
acuaporinas}\label{particularidades-de-las-acuaporinas}

Las acuaporinas suelen organizarse como homotetrámeros. Cada monómero
contiene su propio poro acuoso; el eje central del tetrámero no equivale
necesariamente a un quinto canal de agua funcional. La simulación de un
monómero aislado elimina contactos entre subunidades y expone
superficies que en el ensamblado biológico contactan con otros
monómeros.

Antes de construir el sistema:

\begin{itemize}
\tightlist
\item
  obtenga la unidad biológica, no solo la unidad asimétrica del PDB;
\item
  compruebe que estén presentes los cuatro monómeros cuando el objetivo
  sea el tetrámero;
\item
  revise los dos motivos NPA de cada monómero;
\item
  revise el filtro ar/R y los residuos que determinan selectividad;
\item
  identifique lípidos, metales u otras moléculas resueltas;
\item
  verifique que las hélices transmembrana estén completas;
\item
  decida si colas N- o C-terminales faltantes se modelarán o se dejarán
  truncadas.
\end{itemize}

CHARMM-GUI puede aplicar operaciones de simetría si la información está
disponible, pero el resultado debe compararse con el ensamblado
biológico informado por la fuente estructural.

\subsubsection{3. Preparación de la
estructura}\label{preparaciuxf3n-de-la-estructura}

Conserve el archivo original:

\begin{Shaded}
\begin{Highlighting}[]
\FunctionTok{mkdir} \AttributeTok{{-}p}\NormalTok{ 00\_entrada 01\_charmmgui 02\_em 03\_equilibracion 04\_md 05\_analisis}
\FunctionTok{cp}\NormalTok{ aquaporin\_input.pdb 00\_entrada/aquaporin\_original.pdb}
\end{Highlighting}
\end{Shaded}

Revise:

\begin{itemize}
\tightlist
\item
  cadenas y segmentos que se conservarán;
\item
  conformaciones alternativas;
\item
  residuos faltantes;
\item
  enlaces disulfuro;
\item
  terminales;
\item
  histidinas y otros grupos titulables;
\item
  aguas estructurales dentro del canal;
\item
  moléculas de cristalización que deben eliminarse;
\item
  ligandos o cofactores que necesitan parámetros.
\end{itemize}

No elimine automáticamente todas las aguas cristalográficas. Una cadena
de agua bien resuelta dentro de una acuaporina puede aportar una
configuración inicial razonable, aunque no debe conservarse si solapa
con el procedimiento de generación de agua del poro.

CHARMM-GUI puede modelar segmentos faltantes cortos. Las regiones largas
o poco determinadas requieren validación adicional; una geometría
generada por el servidor no constituye evidencia estructural.

\subsubsection{4. Orientación respecto de la
bicapa}\label{orientaciuxf3n-respecto-de-la-bicapa}

En Membrane Builder:

\begin{itemize}
\tightlist
\item
  el plano de la membrana es \textbf{xy};
\item
  la normal de la bicapa es el eje \textbf{z};
\item
  el centro hidrofóbico se sitúa aproximadamente en \(z=0\).
\end{itemize}

Una estructura obtenida de OPM o PPM puede usarse como orientación
inicial. Una estructura descargada directamente de RCSB no está
necesariamente orientada respecto de una membrana.

Las opciones basadas en ejes principales o en un vector son
aproximaciones geométricas. Pueden fallar cuando:

\begin{itemize}
\tightlist
\item
  existen dominios extramembrana grandes;
\item
  la proteína es irregular;
\item
  se cargó un oligómero incompleto;
\item
  el eje principal no coincide con el poro;
\item
  hay hélices anfipáticas o reentrantes;
\item
  la proteína es un barril beta asimétrico.
\end{itemize}

La orientación debe validarse visualmente y estructuralmente. Compruebe
que:

\begin{itemize}
\tightlist
\item
  la región hidrofóbica contacte con las colas;
\item
  los residuos cargados no queden enterrados sin justificación;
\item
  los dominios citosólicos y extracelulares estén del lado correcto;
\item
  el poro atraviese la bicapa en la dirección esperada;
\item
  no haya subunidades invertidas;
\item
  el centro del filtro selectivo quede cerca de la región prevista.
\end{itemize}

Para una acuaporina, la orientación del tetrámero debe evaluarse como
conjunto. Orientar una sola cadena y reconstruir después el tetrámero
puede introducir una geometría inconsistente.

\subsubsection{5. Agua del poro}\label{agua-del-poro}

Membrane Builder puede generar agua dentro de poros y cavidades. Este
paso es particularmente relevante para acuaporinas, canales, porinas y
transportadores.

Revise que:

\begin{itemize}
\tightlist
\item
  los cuatro poros monoméricos estén hidratados;
\item
  no haya agua aislada dentro de la región hidrofóbica fuera del canal;
\item
  no queden huecos extensos producidos por una mala construcción;
\item
  las aguas no solapen con ligandos o residuos;
\item
  la hidratación central del tetrámero no se confunda con permeación
  monomérica.
\end{itemize}

La ocupación inicial no debe interpretarse como un resultado. El patrón
de agua debe reequilibrarse durante la dinámica.

\subsubsection{6. Selección de la composición
lipídica}\label{selecciuxf3n-de-la-composiciuxf3n-lipuxeddica}

Elija la composición a partir del organismo y del compartimiento
celular. Si no existe información suficiente, declare explícitamente que
se utiliza una membrana modelo.

Ejemplos de decisiones justificables:

{\def\LTcaptype{none} % do not increment counter
\begin{longtable}[]{@{}
  >{\raggedright\arraybackslash}p{(\linewidth - 4\tabcolsep) * \real{0.2450}}
  >{\raggedright\arraybackslash}p{(\linewidth - 4\tabcolsep) * \real{0.3775}}
  >{\raggedright\arraybackslash}p{(\linewidth - 4\tabcolsep) * \real{0.3775}}@{}}
\toprule\noalign{}
\begin{minipage}[b]{\linewidth}\raggedright
Objetivo
\end{minipage} & \begin{minipage}[b]{\linewidth}\raggedright
Modelo posible
\end{minipage} & \begin{minipage}[b]{\linewidth}\raggedright
Limitación
\end{minipage} \\
\midrule\noalign{}
\endhead
\bottomrule\noalign{}
\endlastfoot
Estabilidad general de una acuaporina & POPC puro & No reproduce
asimetría ni diversidad celular. \\
Influencia del colesterol & POPC/colesterol & La proporción debe
justificarse experimentalmente. \\
Membrana bacteriana interna & Mezcla PE/PG/cardiolipina & Depende de
especie y condición de crecimiento. \\
Membrana externa Gram negativa & Monocapa externa con LPS y monocapa
interna fosfolipídica & Requiere parámetros, iones y equilibración más
exigentes. \\
Comparación entre isoformas & Misma bicapa controlada & Aumenta
comparabilidad, pero reduce realismo fisiológico. \\
\end{longtable}
}

En mezclas asimétricas, las dos monocapas pueden contener números y
áreas moleculares diferentes. CHARMM-GUI permite ajustar proporciones y
luego refinar el número de cada especie. No fuerce la misma cantidad de
lípidos en ambas monocapas si sus composiciones o áreas efectivas son
diferentes.

\subsubsection{7. Tamaño del sistema}\label{tamauxf1o-del-sistema}

La caja debe contener suficientes lípidos alrededor de la proteína para
reducir la interacción con sus imágenes periódicas. La extensión lateral
no debe calcularse solo como el rango máximo de coordenadas de la
proteína.

Como criterio inicial:

\[
L_x \gtrsim d_x + 2b
\]

\[
L_y \gtrsim d_y + 2b
\]

donde \(d_x\) y \(d_y\) son las dimensiones proyectadas de la proteína y
\(b\) es el espesor del anillo lipídico deseado. Para una proteína que
pueda inclinarse o cambiar de conformación, debe añadirse margen
adicional.

La cantidad aproximada de lípidos de una monocapa homogénea puede
estimarse como:

\[
N_{\mathrm{lip}} \approx
\frac{A_{xy}-A_{\mathrm{prot}}}
     {a_{\mathrm{lip}}}
\]

donde:

\begin{itemize}
\tightlist
\item
  \(A_{xy}=L_xL_y\);
\item
  \(A_{\mathrm{prot}}\) es el área proyectada ocupada por la proteína;
\item
  \(a_{\mathrm{lip}}\) es el área por lípido esperada en esas
  condiciones.
\end{itemize}

La ecuación es orientativa. En una mezcla lipídica, alrededor de una
proteína irregular o en una membrana asimétrica no existe una única área
por lípido que resuelva el empaquetamiento.

La altura de agua debe evitar contactos periódicos entre dominios
extramembrana. CHARMM-GUI expresa normalmente estas dimensiones en Å:

\[
10\ \text{Å}=1\ \text{nm}
\]

GROMACS utiliza nanómetros en coordenadas y parámetros espaciales.

\subsubsection{8. Temperatura y estado de
fase}\label{temperatura-y-estado-de-fase}

La temperatura no debe elegirse solo por el valor predeterminado del
servidor. Debe ser compatible con:

\begin{itemize}
\tightlist
\item
  la condición experimental;
\item
  el estado de fase esperado de los lípidos;
\item
  la parametrización;
\item
  la estabilidad de la proteína;
\item
  la pregunta biológica.
\end{itemize}

Conversión:

\[
T(\mathrm{K})=T(^{\circ}\mathrm{C})+273.15
\]

{\def\LTcaptype{none} % do not increment counter
\begin{longtable}[]{@{}rr@{}}
\toprule\noalign{}
°C & K \\
\midrule\noalign{}
\endhead
\bottomrule\noalign{}
\endlastfoot
20 & 293.15 \\
25 & 298.15 \\
30 & 303.15 \\
37 & 310.15 \\
\end{longtable}
}

Una bicapa destinada a representar una fase fluida debe simularse por
encima de la transición pertinente de la mezcla. En mezclas complejas,
la transición no puede inferirse únicamente a partir de un lípido
aislado.

\subsubsection{9. Iones y concentración}\label{iones-y-concentraciuxf3n}

CHARMM-GUI puede neutralizar el sistema y agregar NaCl, KCl u otras
sales. Distinga:

\begin{itemize}
\tightlist
\item
  \textbf{neutralización:} cantidad mínima de contraiones para carga
  neta cero;
\item
  \textbf{concentración salina nominal:} pares iónicos adicionales según
  el volumen accesible;
\item
  \textbf{fuerza iónica:} depende de carga y concentración de todas las
  especies.
\end{itemize}

\[
I=\frac{1}{2}\sum_i c_i z_i^2
\]

donde \(c_i\) es la concentración molar y \(z_i\) la carga del ion.

Una solución 0.15 M de NaCl se usa frecuentemente como aproximación
fisiológica, pero ``0.15 M'' e ``isotónica'' no son sinónimos generales.
MgCl₂ o CaCl₂ producen otra fuerza iónica y otra osmolaridad. Los iones
divalentes también pueden interactuar fuertemente con lípidos aniónicos.

No sustituya un ion estructural o catalítico por un contraion difusible.

\subsubsection{10. Construcción en CHARMM-GUI Membrane
Builder}\label{construcciuxf3n-en-charmm-gui-membrane-builder}

Secuencia conceptual de la interfaz:

\begin{enumerate}
\def\labelenumi{\arabic{enumi}.}
\tightlist
\item
  Seleccionar \textbf{Input Generator → Membrane Builder}.
\item
  Elegir un sistema proteína/membrana.
\item
  Cargar un PDB propio o recuperar una estructura.
\item
  Seleccionar modelo, cadenas, ligandos y componentes.
\item
  Aplicar terminales, protonación, enlaces disulfuro y reparaciones
  justificadas.
\item
  Orientar la proteína.
\item
  Generar y revisar agua del poro.
\item
  Definir forma y dimensiones de caja.
\item
  Elegir composición de ambas monocapas.
\item
  Seleccionar agua, iones, concentración y método de colocación.
\item
  Ensamblar y revisar el sistema.
\item
  Generar entradas para GROMACS con el campo de fuerza seleccionado.
\item
  Descargar el archivo TGZ y conservar el identificador del trabajo.
\end{enumerate}

La disponibilidad exacta de lípidos y opciones cambia con la versión del
servidor. Registre la fecha, el identificador del trabajo, el campo de
fuerza y las opciones elegidas.

\subsubsection{11. Inspección del archivo
descargado}\label{inspecciuxf3n-del-archivo-descargado}

Descomprima sin modificar el original:

\begin{Shaded}
\begin{Highlighting}[]
\FunctionTok{mkdir} \AttributeTok{{-}p}\NormalTok{ 01\_charmmgui}
\FunctionTok{tar} \AttributeTok{{-}xzf}\NormalTok{ charmm{-}gui.tgz }\AttributeTok{{-}C}\NormalTok{ 01\_charmmgui}
\BuiltInTok{cd}\NormalTok{ 01\_charmmgui}
\end{Highlighting}
\end{Shaded}

Localice la carpeta de GROMACS:

\begin{Shaded}
\begin{Highlighting}[]
\FunctionTok{find}\NormalTok{ . }\AttributeTok{{-}maxdepth}\NormalTok{ 3 }\AttributeTok{{-}type}\NormalTok{ f }\DataTypeTok{\textbackslash{}}
    \DataTypeTok{\textbackslash{}(} \AttributeTok{{-}name} \StringTok{\textquotesingle{}topol.top\textquotesingle{}} \AttributeTok{{-}o} \AttributeTok{{-}name} \StringTok{\textquotesingle{}index.ndx\textquotesingle{}} \AttributeTok{{-}o} \AttributeTok{{-}name} \StringTok{\textquotesingle{}*.mdp\textquotesingle{}} \AttributeTok{{-}o} \AttributeTok{{-}name} \StringTok{\textquotesingle{}README*\textquotesingle{}} \DataTypeTok{\textbackslash{})} \DataTypeTok{\textbackslash{}}
    \AttributeTok{{-}print}
\end{Highlighting}
\end{Shaded}

Los nombres pueden variar entre trabajos. Una salida típica contiene:

\begin{Shaded}
\begin{Highlighting}[]
\NormalTok{topol.top}
\NormalTok{index.ndx}
\NormalTok{step5\_input.gro o step5\_charmm2gmx.pdb}
\NormalTok{step6.0\_minimization.mdp}
\NormalTok{step6.1\_equilibration.mdp}
\NormalTok{...}
\NormalTok{step6.6\_equilibration.mdp}
\NormalTok{step7\_production.mdp}
\NormalTok{archivos .itp}
\NormalTok{directorio del campo de fuerza}
\NormalTok{README o script de ejecución}
\end{Highlighting}
\end{Shaded}

No regenere la topología con \textbf{pdb2gmx}. CHARMM-GUI ya produjo una
topología consistente con la proteína, los lípidos, el agua y los iones.
Ejecutar \textbf{pdb2gmx} sobre el sistema ensamblado elimina esa
coherencia.

Antes de correr:

\begin{Shaded}
\begin{Highlighting}[]
\ExtensionTok{gmx} \AttributeTok{{-}{-}version}
\FunctionTok{head} \AttributeTok{{-}n}\NormalTok{ 80 topol.top}
\FunctionTok{grep} \AttributeTok{{-}R} \StringTok{"\^{}define\textbackslash{}|\^{}integrator\textbackslash{}|\^{}dt\textbackslash{}|\^{}nsteps\textbackslash{}|\^{}tcoupl\textbackslash{}|\^{}pcoupl\textbackslash{}|\^{}pcoupltype"}\NormalTok{ ./}\PreprocessorTok{*}\NormalTok{.mdp}
\end{Highlighting}
\end{Shaded}

Compruebe:

\begin{itemize}
\tightlist
\item
  campo de fuerza esperado;
\item
  nombres y cantidades en \texttt{{[}\ molecules\ {]}};
\item
  existencia de todos los archivos incluidos;
\item
  grupos de \textbf{index.ndx};
\item
  temperatura;
\item
  paso de integración;
\item
  duración de cada etapa;
\item
  esquema de cortes;
\item
  termostato y barostato;
\item
  macros de restricciones;
\item
  frecuencia de salida;
\item
  tratamiento de dispersión y PME.
\end{itemize}

No reemplace parámetros del MDP por un archivo genérico. Los cortes y
modificadores de Lennard-Jones deben ser compatibles con CHARMM36 y con
la versión concreta usada por el servidor.

\subsubsection{12. CSH no es un requisito}\label{csh-no-es-un-requisito}

Los archivos generados pueden incluir un script C shell. GROMACS no
depende de CSH: el script solo encadena llamadas a \textbf{gmx grompp} y
\textbf{gmx mdrun}. Puede ejecutar cada etapa manualmente o usar
Bash/POSIX shell.

No ejecute el script sin leerlo. Los nombres de archivos y el número de
etapas pueden cambiar.

\subsubsection{13. Minimización manual}\label{minimizaciuxf3n-manual}

Adapte \textbf{step5\_input.gro} al nombre real del archivo descargado:

\begin{Shaded}
\begin{Highlighting}[]
\ExtensionTok{gmx}\NormalTok{ grompp }\DataTypeTok{\textbackslash{}}
    \AttributeTok{{-}f}\NormalTok{ step6.0\_minimization.mdp }\DataTypeTok{\textbackslash{}}
    \AttributeTok{{-}c}\NormalTok{ step5\_input.gro }\DataTypeTok{\textbackslash{}}
    \AttributeTok{{-}r}\NormalTok{ step5\_input.gro }\DataTypeTok{\textbackslash{}}
    \AttributeTok{{-}n}\NormalTok{ index.ndx }\DataTypeTok{\textbackslash{}}
    \AttributeTok{{-}p}\NormalTok{ topol.top }\DataTypeTok{\textbackslash{}}
    \AttributeTok{{-}o}\NormalTok{ step6.0\_minimization.tpr }\DataTypeTok{\textbackslash{}}
    \AttributeTok{{-}pp}\NormalTok{ step6.0\_processed.top}

\ExtensionTok{gmx}\NormalTok{ mdrun }\DataTypeTok{\textbackslash{}}
    \AttributeTok{{-}deffnm}\NormalTok{ step6.0\_minimization }\DataTypeTok{\textbackslash{}}
    \AttributeTok{{-}v}
\end{Highlighting}
\end{Shaded}

La opción \textbf{-r} proporciona las coordenadas de referencia para las
restricciones posicionales. Si el MDP o la topología usan esas
restricciones y \textbf{-r} se omite, \textbf{grompp} puede fallar o
aplicar una referencia incorrecta.

Revise:

\begin{Shaded}
\begin{Highlighting}[]
\ExtensionTok{gmx}\NormalTok{ check }\AttributeTok{{-}f}\NormalTok{ step6.0\_minimization.gro}

\KeywordTok{(}\BuiltInTok{echo}\NormalTok{ Potential}\KeywordTok{;} \BuiltInTok{echo}\NormalTok{ 0}\KeywordTok{)} \KeywordTok{|} \DataTypeTok{\textbackslash{}}
\ExtensionTok{gmx}\NormalTok{ energy }\DataTypeTok{\textbackslash{}}
    \AttributeTok{{-}f}\NormalTok{ step6.0\_minimization.edr }\DataTypeTok{\textbackslash{}}
    \AttributeTok{{-}o}\NormalTok{ ../05\_analisis/em\_potential.xvg}
\end{Highlighting}
\end{Shaded}

\subsubsection{14. Equilibración escalonada sin
CSH}\label{equilibraciuxf3n-escalonada-sin-csh}

CHARMM-GUI suele generar varias etapas que reducen gradualmente
restricciones sobre proteína, lípidos, agua y componentes internos. No
combine las seis etapas en una sola sin comparar sus MDP.

Script Bash, versión 4 o superior:

\begin{Shaded}
\begin{Highlighting}[]
\CommentTok{\#!/usr/bin/env bash}
\BuiltInTok{set} \AttributeTok{{-}euo}\NormalTok{ pipefail}

\VariableTok{GMX\_COMMAND}\OperatorTok{=}\StringTok{"}\VariableTok{$\{GMX\_COMMAND}\OperatorTok{:{-}}\NormalTok{gmx}\VariableTok{\}}\StringTok{"}
\VariableTok{REFERENCE}\OperatorTok{=}\StringTok{"step5\_input.gro"}
\VariableTok{PREVIOUS}\OperatorTok{=}\StringTok{"step6.0\_minimization"}

\ControlFlowTok{for}\NormalTok{ STEP }\KeywordTok{in}\NormalTok{ 1 2 3 4 5 6}
\ControlFlowTok{do}
    \VariableTok{NAME}\OperatorTok{=}\StringTok{"step6.}\VariableTok{$\{STEP\}}\StringTok{\_equilibration"}
    \VariableTok{MDP}\OperatorTok{=}\StringTok{"}\VariableTok{$\{NAME\}}\StringTok{.mdp"}

    \ControlFlowTok{if} \KeywordTok{[[} \OtherTok{!} \OtherTok{{-}f} \StringTok{"}\VariableTok{$\{MDP\}}\StringTok{"} \KeywordTok{]];} \ControlFlowTok{then}
        \BuiltInTok{printf} \StringTok{\textquotesingle{}Falta \%s\textbackslash{}n\textquotesingle{}} \StringTok{"}\VariableTok{$\{MDP\}}\StringTok{"} \OperatorTok{\textgreater{}\&}\DecValTok{2}
        \BuiltInTok{exit}\NormalTok{ 1}
    \ControlFlowTok{fi}

    \VariableTok{COMMAND}\OperatorTok{=}\VariableTok{(}
        \StringTok{"}\VariableTok{$\{GMX\_COMMAND\}}\StringTok{"}\NormalTok{ grompp}
\NormalTok{        {-}f }\StringTok{"}\VariableTok{$\{MDP\}}\StringTok{"}
\NormalTok{        {-}c }\StringTok{"}\VariableTok{$\{PREVIOUS\}}\StringTok{.gro"}
\NormalTok{        {-}r }\StringTok{"}\VariableTok{$\{REFERENCE\}}\StringTok{"}
\NormalTok{        {-}n index.ndx}
\NormalTok{        {-}p topol.top}
\NormalTok{        {-}o }\StringTok{"}\VariableTok{$\{NAME\}}\StringTok{.tpr"}
    \VariableTok{)}

    \ControlFlowTok{if} \KeywordTok{[[} \StringTok{"}\VariableTok{$\{STEP\}}\StringTok{"} \OtherTok{{-}gt}\NormalTok{ 1 }\KeywordTok{\&\&} \OtherTok{{-}f} \StringTok{"}\VariableTok{$\{PREVIOUS\}}\StringTok{.cpt"} \KeywordTok{]];} \ControlFlowTok{then}
        \VariableTok{COMMAND}\OperatorTok{+=}\VariableTok{(}\NormalTok{{-}t }\StringTok{"}\VariableTok{$\{PREVIOUS\}}\StringTok{.cpt"}\VariableTok{)}
    \ControlFlowTok{fi}

    \StringTok{"}\VariableTok{$\{COMMAND}\OperatorTok{[@]}\VariableTok{\}}\StringTok{"}

    \StringTok{"}\VariableTok{$\{GMX\_COMMAND\}}\StringTok{"}\NormalTok{ mdrun }\DataTypeTok{\textbackslash{}}
        \AttributeTok{{-}deffnm} \StringTok{"}\VariableTok{$\{NAME\}}\StringTok{"} \DataTypeTok{\textbackslash{}}
        \AttributeTok{{-}v}

    \VariableTok{PREVIOUS}\OperatorTok{=}\StringTok{"}\VariableTok{$\{NAME\}}\StringTok{"}
\ControlFlowTok{done}
\end{Highlighting}
\end{Shaded}

Guárdelo como \textbf{run\_equilibration.sh} y ejecute:

\begin{Shaded}
\begin{Highlighting}[]
\FunctionTok{chmod}\NormalTok{ +x run\_equilibration.sh}
\ExtensionTok{./run\_equilibration.sh}
\end{Highlighting}
\end{Shaded}

El script usa arreglos de Bash para preservar correctamente cada
argumento. No es compatible con un intérprete POSIX mínimo como
\textbf{dash}; debe ejecutarse con Bash.

Si el archivo inicial se llama \textbf{step5\_charmm2gmx.pdb}, cambie
\textbf{REFERENCE}. Si CHARMM-GUI genera más o menos etapas, modifique
la lista.

Para revisar sin ejecutar:

\begin{Shaded}
\begin{Highlighting}[]
\FunctionTok{bash} \AttributeTok{{-}n}\NormalTok{ run\_equilibration.sh}
\end{Highlighting}
\end{Shaded}

\subsubsection{15. Qué cambia entre las etapas de
equilibración}\label{quuxe9-cambia-entre-las-etapas-de-equilibraciuxf3n}

Inspeccione las diferencias:

\begin{Shaded}
\begin{Highlighting}[]
\FunctionTok{diff} \AttributeTok{{-}u}\NormalTok{ step6.1\_equilibration.mdp step6.2\_equilibration.mdp}
\FunctionTok{diff} \AttributeTok{{-}u}\NormalTok{ step6.5\_equilibration.mdp step6.6\_equilibration.mdp}
\end{Highlighting}
\end{Shaded}

Normalmente cambian una o más de estas variables:

\begin{itemize}
\tightlist
\item
  constantes de restricciones;
\item
  paso de integración;
\item
  temperatura;
\item
  acoplamiento de presión;
\item
  duración;
\item
  tratamiento de grupos;
\item
  frecuencias de salida.
\end{itemize}

Una etapa corta puede ser suficiente para resolver contactos, pero no
demuestra equilibrio de la bicapa. Extienda una etapa si persisten:

\begin{itemize}
\tightlist
\item
  huecos alrededor de la proteína;
\item
  deformación fuerte de la bicapa;
\item
  densidad no estacionaria;
\item
  área xy con deriva;
\item
  agua dentro del núcleo hidrofóbico fuera del poro;
\item
  errores LINCS;
\item
  proteína que se desplaza o inclina bruscamente;
\item
  contactos periódicos.
\end{itemize}

\subsubsection{16. Acoplamiento de presión en
membranas}\label{acoplamiento-de-presiuxf3n-en-membranas}

Para una bicapa en el plano xy se utiliza habitualmente acoplamiento
semiisotrópico:

\begin{Shaded}
\begin{Highlighting}[]
\DataTypeTok{pcoupltype       }\OtherTok{=}\StringTok{ semiisotropic}
\DataTypeTok{ref{-}p            }\OtherTok{=}\StringTok{ 1.0 1.0}
\DataTypeTok{compressibility  }\OtherTok{=}\StringTok{ 4.5e{-}5 4.5e{-}5}
\end{Highlighting}
\end{Shaded}

El primer valor corresponde conjuntamente a x/y y el segundo a z. El
acoplamiento isotrópico obliga a que todas las dimensiones respondan de
la misma manera y no suele ser apropiado para una bicapa plana.

\textbf{C-rescale} es adecuado para equilibración.
\textbf{Parrinello-Rahman} suele reservarse para producción una vez
estabilizado el sistema. Mantenga inicialmente los valores generados por
CHARMM-GUI y documente cualquier cambio.

No interprete fluctuaciones instantáneas de presión como fallo. Evalúe
deriva, densidad, área, espesor y promedios por bloques.

\subsubsection{17. Repartición de masa de hidrógeno
(HMR)}\label{reparticiuxf3n-de-masa-de-hidruxf3geno-hmr}

La \textbf{repartición de masa de hidrógeno} (\emph{hydrogen mass
repartitioning}, HMR) transfiere parte de la masa del átomo pesado a sus
hidrógenos enlazados sin cambiar la masa total de la molécula. Su
objetivo es disminuir las frecuencias más altas del sistema y permitir,
después de validarlo, un paso de integración mayor.

Para una vibración aproximadamente armónica:

\[
\omega=\sqrt{\frac{k}{\mu}},
\qquad
\mu=\frac{m_1m_2}{m_1+m_2}
\]

donde \(k\) es la constante de fuerza y \(\mu\) la masa reducida. Al
aumentar la masa del hidrógeno aumenta \(\mu\) y disminuye \(\omega\).
Para un factor de repartición \(f\):

\[
m'_{H}=f\,m_H,
\qquad
m'_X=m_X-\sum_H\left(m'_H-m_H\right)
\]

El segundo término se suma sobre los hidrógenos unidos al átomo pesado
\(X\). Con \(f=3\), un hidrógeno de aproximadamente
\(1.008\ \mathrm{u}\) pasa a aproximadamente \(3.024\ \mathrm{u}\). La
masa total se conserva. HMR no cambia las cargas, los parámetros de
Lennard-Jones ni la función de energía potencial; sí cambia la matriz de
masas y, por lo tanto, puede modificar propiedades dinámicas.

\paragraph{Qué permite y qué no permite
HMR}\label{quuxe9-permite-y-quuxe9-no-permite-hmr}

En sistemas atomísticos convencionales, GROMACS indica que combinar un
factor 3 con enlaces a hidrógeno restringidos permite
\textbf{habitualmente} usar
\(\Delta t=0.004\ \mathrm{ps}=4\ \mathrm{fs}\). No es una garantía de
estabilidad para cualquier topología. HMR:

\begin{itemize}
\tightlist
\item
  reduce aproximadamente a la mitad el número de pasos necesario para
  simular el mismo tiempo físico;
\item
  puede acercarse a una aceleración de dos veces, aunque el rendimiento
  real depende de GPU, CPU, PME, restricciones y comunicaciones;
\item
  no mejora el campo de fuerza ni genera más muestreo por nanosegundo
  simulado;
\item
  conserva, en principio, las distribuciones configuracionales de
  equilibrio;
\item
  altera escalas dinámicas: difusión, correlaciones temporales, cinética
  y tasas de transporte pueden resultar algo más lentas.
\end{itemize}

Por este último motivo, HMR es especialmente atractivo para propiedades
estructurales y termodinámicas de membranas, pero exige cautela si el
resultado principal es difusión lipídica, flujo de agua, permeabilidad
de una acuaporina, tiempos de residencia o cinética de apertura.

\paragraph{Método nativo en GROMACS
2026.3}\label{muxe9todo-nativo-en-gromacs-2026.3}

Desde GROMACS 2024, \textbf{gmx grompp} puede efectuar HMR a partir de
una topología ordinaria mediante una sola opción del MDP. Para GROMACS
2026.3, el bloque mínimo es:

\begin{Shaded}
\begin{Highlighting}[]
\DataTypeTok{integrator              }\OtherTok{=}\StringTok{ md}
\DataTypeTok{dt                      }\OtherTok{=}\StringTok{ }\FloatTok{0.004}
\DataTypeTok{constraints             }\OtherTok{=}\StringTok{ h{-}bonds}
\DataTypeTok{constraint{-}algorithm    }\OtherTok{=}\StringTok{ lincs}
\DataTypeTok{mass{-}repartition{-}factor }\OtherTok{=}\StringTok{ }\DecValTok{3}
\end{Highlighting}
\end{Shaded}

\textbf{No combine este método con \texttt{define\ =\ -DHEAVY\_H}, con
masas ya modificadas por CHARMM-GUI ni con una edición manual de los
ITP.} La opción nativa debe aplicarse al sistema completo. Reparticionar
solo lípidos o solo proteína deja otros hidrógenos con movimientos
rápidos y elimina la justificación física del paso de 4 fs.

\textbf{grompp} eleva al umbral correspondiente las masas de los átomos
ligeros y resta la diferencia al átomo enlazado. Emite una advertencia
si un átomo ligero no tiene vecino enlazado y un error si tiene más de
uno o si la resta haría que el vecino quedara por debajo de la masa
mínima. Esto obliga a revisar con especial cuidado moléculas no
estándar, modelos polarizables, sitios virtuales y topologías con
conectividad inusual.

\paragraph{Integración con los archivos de
CHARMM-GUI}\label{integraciuxf3n-con-los-archivos-de-charmm-gui}

Duplique el MDP de producción para conservar un control de 2 fs:

\begin{Shaded}
\begin{Highlighting}[]
\FunctionTok{cp}\NormalTok{ step7\_production.mdp step7\_production\_hmr.mdp}
\end{Highlighting}
\end{Shaded}

En \textbf{step7\_production\_hmr.mdp}, mantenga los parámetros de
CHARMM-GUI para electrostática, van der Waals, temperatura y presión, y
cambie únicamente lo necesario para la comparación:

\begin{Shaded}
\begin{Highlighting}[]
\DataTypeTok{dt                      }\OtherTok{=}\StringTok{ }\FloatTok{0.004}
\DataTypeTok{constraints             }\OtherTok{=}\StringTok{ h{-}bonds}
\DataTypeTok{mass{-}repartition{-}factor }\OtherTok{=}\StringTok{ }\DecValTok{3}
\end{Highlighting}
\end{Shaded}

No aproveche el cambio a HMR para acortar cortes, cambiar el barostato o
sustituir los parámetros no enlazados. El estudio de membranas que
evaluó HMR encontró resultados razonables con el protocolo validado,
pero no respaldó el corte de \(9\ \text{Å}=0.9\ \mathrm{nm}\) que
también ensayó. En CHARMM36 deben conservarse los cortes y el esquema de
conmutación recomendados para esa versión del campo de fuerza.

El tiempo simulado es:

\[
t_{\mathrm{total}}=n_{\mathrm{steps}}\,\Delta t
\]

Para 10 ns:

{\def\LTcaptype{none} % do not increment counter
\begin{longtable}[]{@{}lrr@{}}
\toprule\noalign{}
Paso & \texttt{nsteps} & Tiempo \\
\midrule\noalign{}
\endhead
\bottomrule\noalign{}
\endlastfoot
0.002 ps (2 fs) & 5 000 000 & 10 ns \\
0.004 ps (4 fs) & 2 500 000 & 10 ns \\
\end{longtable}
}

Los intervalos de salida también dependen de \(\Delta t\):

\[
\Delta t_{\mathrm{salida}}=\mathrm{nst}\times\Delta t
\]

Si se desea guardar cada 10 ps con 4 fs, use por ejemplo:

\begin{Shaded}
\begin{Highlighting}[]
\DataTypeTok{nstlog              }\OtherTok{=}\StringTok{ }\DecValTok{2500}
\DataTypeTok{nstenergy           }\OtherTok{=}\StringTok{ }\DecValTok{2500}
\DataTypeTok{nstxout{-}compressed  }\OtherTok{=}\StringTok{ }\DecValTok{2500}
\end{Highlighting}
\end{Shaded}

\paragraph{Generación y comprobación del
TPR}\label{generaciuxf3n-y-comprobaciuxf3n-del-tpr}

Ejemplo para producción después de una equilibración que ya empleó HMR:

\begin{Shaded}
\begin{Highlighting}[]
\ExtensionTok{gmx}\NormalTok{ grompp }\DataTypeTok{\textbackslash{}}
    \AttributeTok{{-}f}\NormalTok{ step7\_production\_hmr.mdp }\DataTypeTok{\textbackslash{}}
    \AttributeTok{{-}c}\NormalTok{ step6.6\_equilibration.gro }\DataTypeTok{\textbackslash{}}
    \AttributeTok{{-}t}\NormalTok{ step6.6\_equilibration.cpt }\DataTypeTok{\textbackslash{}}
    \AttributeTok{{-}r}\NormalTok{ step5\_input.gro }\DataTypeTok{\textbackslash{}}
    \AttributeTok{{-}n}\NormalTok{ index.ndx }\DataTypeTok{\textbackslash{}}
    \AttributeTok{{-}p}\NormalTok{ topol.top }\DataTypeTok{\textbackslash{}}
    \AttributeTok{{-}pp}\NormalTok{ processed\_hmr.top }\DataTypeTok{\textbackslash{}}
    \AttributeTok{{-}po}\NormalTok{ mdout\_hmr.mdp }\DataTypeTok{\textbackslash{}}
    \AttributeTok{{-}o}\NormalTok{ step7\_1\_hmr.tpr}
\end{Highlighting}
\end{Shaded}

Revise los parámetros efectivamente procesados:

\begin{Shaded}
\begin{Highlighting}[]
\FunctionTok{grep} \AttributeTok{{-}E} \StringTok{\textquotesingle{}\^{}(dt|nsteps|constraints|mass{-}repartition{-}factor)\textquotesingle{}}\NormalTok{ mdout\_hmr.mdp}
\ExtensionTok{gmx}\NormalTok{ dump }\AttributeTok{{-}s}\NormalTok{ step7\_1\_hmr.tpr }\OperatorTok{\textgreater{}}\NormalTok{ step7\_1\_hmr.dump}
\end{Highlighting}
\end{Shaded}

El archivo \textbf{processed\_hmr.top} permite auditar la topología
preprocesada, mientras que \textbf{gmx dump} permite inspeccionar el
TPR. No edite ninguno para ejecutar la simulación; corrija siempre el
MDP o la topología fuente y vuelva a ejecutar \textbf{grompp}.

Lo más limpio es usar HMR de manera consistente desde la primera etapa
de dinámica de la equilibración. Si el sistema fue equilibrado a 2 fs
con masas normales y se decide cambiar después, no trate la primera
ejecución HMR como una continuación bit a bit de la dinámica anterior.
Genere velocidades compatibles en una etapa NVT corta, vuelva a
equilibrar a NPT y recién después produzca. Mantenga separados los
nombres de archivos y no anexe un tramo HMR a un log o trayectoria de
masas normales.

\paragraph{Prueba escalonada de
estabilidad}\label{prueba-escalonada-de-estabilidad}

Antes de una producción larga:

\begin{enumerate}
\def\labelenumi{\arabic{enumi}.}
\tightlist
\item
  demuestre que el mismo sistema es estable con masas normales y 2 fs;
\item
  prepare HMR sin cambiar simultáneamente otros parámetros;
\item
  ejecute una prueba corta de 100 ps a 1 ns;
\item
  revise advertencias LINCS, NaN, temperatura, energía, volumen, área xy
  y espesor;
\item
  amplíe la prueba antes de lanzar réplicas largas.
\end{enumerate}

Una prueba de 1 ns a 4 fs contiene 250 000 pasos:

\begin{Shaded}
\begin{Highlighting}[]
\ExtensionTok{gmx}\NormalTok{ mdrun }\DataTypeTok{\textbackslash{}}
    \AttributeTok{{-}s}\NormalTok{ step7\_1\_hmr.tpr }\DataTypeTok{\textbackslash{}}
    \AttributeTok{{-}deffnm}\NormalTok{ hmr\_test }\DataTypeTok{\textbackslash{}}
    \AttributeTok{{-}nsteps}\NormalTok{ 250000 }\DataTypeTok{\textbackslash{}}
    \AttributeTok{{-}v}
\end{Highlighting}
\end{Shaded}

Busque problemas evidentes:

\begin{Shaded}
\begin{Highlighting}[]
\FunctionTok{grep} \AttributeTok{{-}Ei} \StringTok{\textquotesingle{}LINCS|constraint|nan|fatal|warning\textquotesingle{}}\NormalTok{ hmr\_test.log}
\ExtensionTok{gmx}\NormalTok{ energy }\AttributeTok{{-}f}\NormalTok{ hmr\_test.edr }\AttributeTok{{-}o}\NormalTok{ hmr\_test\_control.xvg}
\end{Highlighting}
\end{Shaded}

Seleccione al menos temperatura, energía potencial, presión, volumen y
dimensiones de caja. Si 4 fs falla, no oculte el problema aumentando
indiscriminadamente el orden de LINCS. Revise estructura, contactos,
topología, restricciones y equilibración. Un paso de 3 fs puede ser una
alternativa; si tampoco es estable o la topología no es compatible, use
2 fs.

\paragraph{Validación específica para una
membrana}\label{validaciuxf3n-especuxedfica-para-una-membrana}

Compare al menos una réplica de referencia a 2 fs y masas normales con
una réplica HMR a 4 fs, usando el mismo sistema y el mismo tiempo
físico. Evalúe por bloques:

\begin{itemize}
\tightlist
\item
  área por lípido o área xy;
\item
  espesor y perfiles de densidad;
\item
  parámetro de orden \(S_{CD}\);
\item
  compresibilidad areal;
\item
  RMSD, RMSF e inclinación de la proteína;
\item
  contactos lípido--proteína;
\item
  hidratación y radio del poro;
\item
  difusión lateral, indicando que HMR altera la dinámica;
\item
  estabilidad de temperatura, energía y presión;
\item
  rendimiento en ns/día sobre el mismo hardware.
\end{itemize}

Para una acuaporina, la recomendación práctica es usar HMR con confianza
solo después de demostrar que la estructura de la membrana y del canal
coincide con el control. Si el objetivo es una \textbf{permeabilidad
absoluta}, flujo de agua o una constante cinética, conserve producción
de 2 fs con masas normales como referencia principal o cuantifique
explícitamente el sesgo de HMR. Una ocupación o un perfil de densidad es
configuracional; una tasa de cruce depende del tiempo dinámico y no debe
asumirse invariante.

\paragraph{Decisión práctica}\label{decisiuxf3n-pruxe1ctica}

HMR es recomendable cuando se busca aumentar el muestreo de equilibrio
de estructura, contactos y termodinámica, siempre que el sistema pase
una validación apareada. No lo use de forma automática cuando:

\begin{itemize}
\tightlist
\item
  existen hidrógenos flexibles sin restringir;
\item
  la topología contiene conectividades o partículas no estándar que
  \textbf{grompp} rechaza;
\item
  se mezclan componentes con y sin HMR;
\item
  se pretende comparar directamente cinética con trayectorias de masas
  normales;
\item
  no hay un control estable a 2 fs;
\item
  se cambia al mismo tiempo HMR, cortes, termostato, barostato y grupos
  de acoplamiento.
\end{itemize}

La discusión histórica del foro incluye métodos dependientes del campo
de fuerza y casos de inestabilidad anteriores a la opción nativa. Son
útiles para entender los errores frecuentes, pero para GROMACS 2026.3
debe prevalecer el método documentado con
\textbf{mass-repartition-factor}.

\subsubsection{18. Producción segmentada sin
CSH}\label{producciuxf3n-segmentada-sin-csh}

Dividir una producción en segmentos no cambia por sí mismo la física.
Facilita checkpoints, colas HPC, copias de seguridad y reanudación. El
tiempo total es:

\[
t_{\mathrm{total}} =
N_{\mathrm{segmentos}}\,
n_{\mathrm{steps}}\,
\Delta t
\]

Si cada segmento contiene 5000000 pasos con
\(\Delta t=0.002\ \mathrm{ps}\):

\[
t_{\mathrm{segmento}}=10\ \mathrm{ns}
\]

Diez segmentos suman 100 ns.

Script Bash:

\begin{Shaded}
\begin{Highlighting}[]
\CommentTok{\#!/usr/bin/env bash}
\BuiltInTok{set} \AttributeTok{{-}euo}\NormalTok{ pipefail}

\VariableTok{GMX\_COMMAND}\OperatorTok{=}\StringTok{"}\VariableTok{$\{GMX\_COMMAND}\OperatorTok{:{-}}\NormalTok{gmx}\VariableTok{\}}\StringTok{"}
\VariableTok{FIRST\_SEGMENT}\OperatorTok{=}\StringTok{"}\VariableTok{$\{FIRST\_SEGMENT}\OperatorTok{:{-}}\NormalTok{1}\VariableTok{\}}\StringTok{"}
\VariableTok{LAST\_SEGMENT}\OperatorTok{=}\StringTok{"}\VariableTok{$\{LAST\_SEGMENT}\OperatorTok{:{-}}\NormalTok{10}\VariableTok{\}}\StringTok{"}
\VariableTok{REFERENCE}\OperatorTok{=}\StringTok{"step5\_input.gro"}

\ControlFlowTok{for}\NormalTok{ SEGMENT }\KeywordTok{in} \VariableTok{$(}\FunctionTok{seq} \StringTok{"}\VariableTok{$\{FIRST\_SEGMENT\}}\StringTok{"} \StringTok{"}\VariableTok{$\{LAST\_SEGMENT\}}\StringTok{"}\VariableTok{)}
\ControlFlowTok{do}
    \VariableTok{NAME}\OperatorTok{=}\StringTok{"step7\_}\VariableTok{$\{SEGMENT\}}\StringTok{"}

    \ControlFlowTok{if} \KeywordTok{[[} \StringTok{"}\VariableTok{$\{SEGMENT\}}\StringTok{"} \OtherTok{{-}eq}\NormalTok{ 1 }\KeywordTok{]];} \ControlFlowTok{then}
        \VariableTok{PREVIOUS}\OperatorTok{=}\StringTok{"step6.6\_equilibration"}
    \ControlFlowTok{else}
        \VariableTok{PREVIOUS}\OperatorTok{=}\StringTok{"step7\_}\VariableTok{$((SEGMENT} \OperatorTok{{-}} \DecValTok{1}\VariableTok{))}\StringTok{"}
    \ControlFlowTok{fi}

    \ControlFlowTok{if} \KeywordTok{[[} \OtherTok{!} \OtherTok{{-}f} \StringTok{"}\VariableTok{$\{NAME\}}\StringTok{.tpr"} \KeywordTok{]];} \ControlFlowTok{then}
        \StringTok{"}\VariableTok{$\{GMX\_COMMAND\}}\StringTok{"}\NormalTok{ grompp }\DataTypeTok{\textbackslash{}}
            \AttributeTok{{-}f}\NormalTok{ step7\_production.mdp }\DataTypeTok{\textbackslash{}}
            \AttributeTok{{-}c} \StringTok{"}\VariableTok{$\{PREVIOUS\}}\StringTok{.gro"} \DataTypeTok{\textbackslash{}}
            \AttributeTok{{-}t} \StringTok{"}\VariableTok{$\{PREVIOUS\}}\StringTok{.cpt"} \DataTypeTok{\textbackslash{}}
            \AttributeTok{{-}r} \StringTok{"}\VariableTok{$\{REFERENCE\}}\StringTok{"} \DataTypeTok{\textbackslash{}}
            \AttributeTok{{-}n}\NormalTok{ index.ndx }\DataTypeTok{\textbackslash{}}
            \AttributeTok{{-}p}\NormalTok{ topol.top }\DataTypeTok{\textbackslash{}}
            \AttributeTok{{-}o} \StringTok{"}\VariableTok{$\{NAME\}}\StringTok{.tpr"}
    \ControlFlowTok{fi}

    \ControlFlowTok{if} \KeywordTok{[[} \OtherTok{{-}f} \StringTok{"}\VariableTok{$\{NAME\}}\StringTok{.cpt"} \KeywordTok{]];} \ControlFlowTok{then}
        \StringTok{"}\VariableTok{$\{GMX\_COMMAND\}}\StringTok{"}\NormalTok{ mdrun }\DataTypeTok{\textbackslash{}}
            \AttributeTok{{-}deffnm} \StringTok{"}\VariableTok{$\{NAME\}}\StringTok{"} \DataTypeTok{\textbackslash{}}
            \AttributeTok{{-}cpi} \StringTok{"}\VariableTok{$\{NAME\}}\StringTok{.cpt"} \DataTypeTok{\textbackslash{}}
            \AttributeTok{{-}append} \DataTypeTok{\textbackslash{}}
            \AttributeTok{{-}v}
    \ControlFlowTok{else}
        \StringTok{"}\VariableTok{$\{GMX\_COMMAND\}}\StringTok{"}\NormalTok{ mdrun }\DataTypeTok{\textbackslash{}}
            \AttributeTok{{-}deffnm} \StringTok{"}\VariableTok{$\{NAME\}}\StringTok{"} \DataTypeTok{\textbackslash{}}
            \AttributeTok{{-}v}
    \ControlFlowTok{fi}
\ControlFlowTok{done}
\end{Highlighting}
\end{Shaded}

Guárdelo como \textbf{run\_production.sh}:

\begin{Shaded}
\begin{Highlighting}[]
\FunctionTok{chmod}\NormalTok{ +x run\_production.sh}
\FunctionTok{bash} \AttributeTok{{-}n}\NormalTok{ run\_production.sh}
\ExtensionTok{./run\_production.sh}
\end{Highlighting}
\end{Shaded}

Para ejecutar solo los segmentos 4 a 7:

\begin{Shaded}
\begin{Highlighting}[]
\VariableTok{FIRST\_SEGMENT}\OperatorTok{=}\NormalTok{4 }\VariableTok{LAST\_SEGMENT}\OperatorTok{=}\NormalTok{7 }\ExtensionTok{./run\_production.sh}
\end{Highlighting}
\end{Shaded}

Si la producción no usa restricciones, \textbf{-r} puede ser
innecesario, pero conservarlo no reemplaza la revisión del MDP y de las
macros activas.

En un clúster, el script debe ejecutarse dentro del sistema de colas. No
lance múltiples segmentos dependientes al mismo tiempo: cada uno
necesita el checkpoint final del anterior.

\subsubsection{19. Reanudar una ejecución
interrumpida}\label{reanudar-una-ejecuciuxf3n-interrumpida}

Si existe un checkpoint del mismo segmento:

\begin{Shaded}
\begin{Highlighting}[]
\ExtensionTok{gmx}\NormalTok{ mdrun }\DataTypeTok{\textbackslash{}}
    \AttributeTok{{-}deffnm}\NormalTok{ step7\_4 }\DataTypeTok{\textbackslash{}}
    \AttributeTok{{-}cpi}\NormalTok{ step7\_4.cpt }\DataTypeTok{\textbackslash{}}
    \AttributeTok{{-}append} \DataTypeTok{\textbackslash{}}
    \AttributeTok{{-}v}
\end{Highlighting}
\end{Shaded}

No vuelva a ejecutar \textbf{grompp} desde la última estructura si el
objetivo es continuar exactamente el mismo segmento. El checkpoint
conserva estado del integrador, velocidades y otra información necesaria
para una continuación reproducible.

\subsubsection{20. Preparación de la trayectoria para
análisis}\label{preparaciuxf3n-de-la-trayectoria-para-anuxe1lisis}

Conserve siempre la trayectoria original. Genere derivados para
visualización y análisis.

Primero haga moléculas enteras y centre la proteína:

\begin{Shaded}
\begin{Highlighting}[]
\KeywordTok{(}\BuiltInTok{echo}\NormalTok{ Protein}\KeywordTok{;} \BuiltInTok{echo}\NormalTok{ System}\KeywordTok{)} \KeywordTok{|} \DataTypeTok{\textbackslash{}}
\ExtensionTok{gmx}\NormalTok{ trjconv }\DataTypeTok{\textbackslash{}}
    \AttributeTok{{-}s}\NormalTok{ step7\_1.tpr }\DataTypeTok{\textbackslash{}}
    \AttributeTok{{-}f}\NormalTok{ md\_all.xtc }\DataTypeTok{\textbackslash{}}
    \AttributeTok{{-}o}\NormalTok{ 05\_analisis/md\_center.xtc }\DataTypeTok{\textbackslash{}}
    \AttributeTok{{-}pbc}\NormalTok{ mol }\DataTypeTok{\textbackslash{}}
    \AttributeTok{{-}center} \DataTypeTok{\textbackslash{}}
    \AttributeTok{{-}ur}\NormalTok{ compact}
\end{Highlighting}
\end{Shaded}

Luego ajuste la proteína si el análisis requiere eliminar traslación y
rotación:

\begin{Shaded}
\begin{Highlighting}[]
\KeywordTok{(}\BuiltInTok{echo}\NormalTok{ Backbone}\KeywordTok{;} \BuiltInTok{echo}\NormalTok{ System}\KeywordTok{)} \KeywordTok{|} \DataTypeTok{\textbackslash{}}
\ExtensionTok{gmx}\NormalTok{ trjconv }\DataTypeTok{\textbackslash{}}
    \AttributeTok{{-}s}\NormalTok{ step7\_1.tpr }\DataTypeTok{\textbackslash{}}
    \AttributeTok{{-}f}\NormalTok{ 05\_analisis/md\_center.xtc }\DataTypeTok{\textbackslash{}}
    \AttributeTok{{-}o}\NormalTok{ 05\_analisis/md\_fit.xtc }\DataTypeTok{\textbackslash{}}
    \AttributeTok{{-}fit}\NormalTok{ rot+trans}
\end{Highlighting}
\end{Shaded}

No use una trayectoria ajustada para difusión lateral de lípidos ni para
fluctuaciones de caja. El ajuste cambia las coordenadas colectivas que
esos análisis necesitan.

Si la producción está segmentada:

\begin{Shaded}
\begin{Highlighting}[]
\ExtensionTok{gmx}\NormalTok{ trjcat }\DataTypeTok{\textbackslash{}}
    \AttributeTok{{-}f}\NormalTok{ step7\_1.xtc step7\_2.xtc step7\_3.xtc step7\_4.xtc }\DataTypeTok{\textbackslash{}}
    \AttributeTok{{-}o}\NormalTok{ md\_all.xtc }\DataTypeTok{\textbackslash{}}
    \AttributeTok{{-}cat}
\end{Highlighting}
\end{Shaded}

Compruebe continuidad temporal y descarte marcos duplicados si los
segmentos se solapan.

\subsubsection{21. Controles estructurales de la
proteína}\label{controles-estructurales-de-la-proteuxedna}

RMSD del backbone:

\begin{Shaded}
\begin{Highlighting}[]
\KeywordTok{(}\BuiltInTok{echo}\NormalTok{ Backbone}\KeywordTok{;} \BuiltInTok{echo}\NormalTok{ Backbone}\KeywordTok{)} \KeywordTok{|} \DataTypeTok{\textbackslash{}}
\ExtensionTok{gmx}\NormalTok{ rms }\DataTypeTok{\textbackslash{}}
    \AttributeTok{{-}s}\NormalTok{ step7\_1.tpr }\DataTypeTok{\textbackslash{}}
    \AttributeTok{{-}f}\NormalTok{ 05\_analisis/md\_fit.xtc }\DataTypeTok{\textbackslash{}}
    \AttributeTok{{-}o}\NormalTok{ 05\_analisis/rmsd\_backbone.xvg }\DataTypeTok{\textbackslash{}}
    \AttributeTok{{-}tu}\NormalTok{ ns}
\end{Highlighting}
\end{Shaded}

RMSF por residuo:

\begin{Shaded}
\begin{Highlighting}[]
\BuiltInTok{echo}\NormalTok{ C{-}alpha }\KeywordTok{|} \DataTypeTok{\textbackslash{}}
\ExtensionTok{gmx}\NormalTok{ rmsf }\DataTypeTok{\textbackslash{}}
    \AttributeTok{{-}s}\NormalTok{ step7\_1.tpr }\DataTypeTok{\textbackslash{}}
    \AttributeTok{{-}f}\NormalTok{ 05\_analisis/md\_fit.xtc }\DataTypeTok{\textbackslash{}}
    \AttributeTok{{-}o}\NormalTok{ 05\_analisis/rmsf\_calpha.xvg }\DataTypeTok{\textbackslash{}}
    \AttributeTok{{-}res}
\end{Highlighting}
\end{Shaded}

Para una acuaporina tetramérica, calcule además:

\begin{itemize}
\tightlist
\item
  RMSD de cada monómero;
\item
  RMSD del tetrámero completo;
\item
  distancias entre centros de masa de subunidades;
\item
  contactos interfaciales entre monómeros;
\item
  inclinación del eje de cada poro respecto de z;
\item
  estabilidad de los motivos NPA;
\item
  geometría del filtro ar/R.
\end{itemize}

Un RMSD estable del tetrámero puede ocultar que un monómero se deforma;
cuatro RMSD monoméricos estables pueden ocultar una reorganización
cuaternaria.

\subsubsection{22. Área de la membrana}\label{uxe1rea-de-la-membrana}

El área instantánea proyectada es:

\[
A_{xy}(t)=L_x(t)L_y(t)
\]

Extraiga las dimensiones:

\begin{Shaded}
\begin{Highlighting}[]
\KeywordTok{(}\BuiltInTok{echo}\NormalTok{ Box{-}X}\KeywordTok{;} \BuiltInTok{echo}\NormalTok{ Box{-}Y}\KeywordTok{;} \BuiltInTok{echo}\NormalTok{ Box{-}Z}\KeywordTok{;} \BuiltInTok{echo}\NormalTok{ Volume}\KeywordTok{;} \BuiltInTok{echo}\NormalTok{ 0}\KeywordTok{)} \KeywordTok{|} \DataTypeTok{\textbackslash{}}
\ExtensionTok{gmx}\NormalTok{ energy }\DataTypeTok{\textbackslash{}}
    \AttributeTok{{-}f}\NormalTok{ step7\_1.edr }\DataTypeTok{\textbackslash{}}
    \AttributeTok{{-}o}\NormalTok{ 05\_analisis/box\_dimensions.xvg}
\end{Highlighting}
\end{Shaded}

Para una bicapa homogénea, simétrica y sin proteína:

\[
APL(t)=\frac{A_{xy}(t)}{N_{\mathrm{lip,leaflet}}}
\]

Con una proteína:

\[
APL_{\mathrm{aprox}}(t)=
\frac{A_{xy}(t)-A_{\mathrm{prot}}(t)}
     {N_{\mathrm{lip,leaflet}}}
\]

La corrección depende de cómo se defina el área proyectada de la
proteína. En bicapas mixtas no existe una única APL rigurosa para todas
las especies. Para análisis local conviene usar una partición de Voronoi
u otra herramienta específica.

No divida por el número total de lípidos de las dos monocapas: el
denominador es el número de una monocapa.

\subsubsection{23. Compresibilidad areal}\label{compresibilidad-areal}

En un ensamble apropiado puede estimarse:

\[
K_A=
\frac{k_\mathrm{B}T\langle A\rangle}
     {\langle A^2\rangle-\langle A\rangle^2}
\]

donde \(A=L_xL_y\). La estimación es sensible a:

\begin{itemize}
\tightlist
\item
  duración de la trayectoria;
\item
  barostato;
\item
  correlación temporal;
\item
  tamaño del sistema;
\item
  presencia de proteína;
\item
  mezcla y asimetría de lípidos.
\end{itemize}

Use promedios por bloques y no compare directamente valores obtenidos
con ensambles diferentes.

\subsubsection{24. Espesor y perfiles de
densidad}\label{espesor-y-perfiles-de-densidad}

Calcule perfiles a lo largo de z:

\begin{Shaded}
\begin{Highlighting}[]
\ExtensionTok{gmx}\NormalTok{ density }\DataTypeTok{\textbackslash{}}
    \AttributeTok{{-}s}\NormalTok{ step7\_1.tpr }\DataTypeTok{\textbackslash{}}
    \AttributeTok{{-}f}\NormalTok{ md\_all.xtc }\DataTypeTok{\textbackslash{}}
    \AttributeTok{{-}n}\NormalTok{ index.ndx }\DataTypeTok{\textbackslash{}}
    \AttributeTok{{-}d}\NormalTok{ Z }\DataTypeTok{\textbackslash{}}
    \AttributeTok{{-}sl}\NormalTok{ 200 }\DataTypeTok{\textbackslash{}}
    \AttributeTok{{-}dens}\NormalTok{ number }\DataTypeTok{\textbackslash{}}
    \AttributeTok{{-}o}\NormalTok{ 05\_analisis/density\_z.xvg}
\end{Highlighting}
\end{Shaded}

Cree grupos para:

\begin{itemize}
\tightlist
\item
  fósforos o grupos cabeza;
\item
  colas lipídicas;
\item
  agua;
\item
  proteína;
\item
  iones;
\item
  cada especie lipídica relevante.
\end{itemize}

Una definición simple del espesor entre cabezas es la distancia entre
los máximos de densidad de fósforo de ambas monocapas:

\[
d_{PP}=z_{P,\mathrm{sup}}-z_{P,\mathrm{inf}}
\]

Es una definición global. Una membrana deformada alrededor de la
proteína requiere mapas locales de espesor.

\subsubsection{25. Orden de las cadenas
lipídicas}\label{orden-de-las-cadenas-lipuxeddicas}

El parámetro de orden deuterio se expresa como:

\[
S_{CD}=
\frac{1}{2}
\left\langle 3\cos^2\theta-1 \right\rangle
\]

donde \(\theta\) es el ángulo entre un enlace de la cadena y la normal
de la membrana.

Consulte primero la sintaxis disponible:

\begin{Shaded}
\begin{Highlighting}[]
\ExtensionTok{gmx}\NormalTok{ order }\AttributeTok{{-}h}
\end{Highlighting}
\end{Shaded}

Ejemplo general:

\begin{Shaded}
\begin{Highlighting}[]
\ExtensionTok{gmx}\NormalTok{ order }\DataTypeTok{\textbackslash{}}
    \AttributeTok{{-}s}\NormalTok{ step7\_1.tpr }\DataTypeTok{\textbackslash{}}
    \AttributeTok{{-}f}\NormalTok{ md\_all.xtc }\DataTypeTok{\textbackslash{}}
    \AttributeTok{{-}n}\NormalTok{ lipid\_chains.ndx }\DataTypeTok{\textbackslash{}}
    \AttributeTok{{-}d}\NormalTok{ z }\DataTypeTok{\textbackslash{}}
    \AttributeTok{{-}od}\NormalTok{ 05\_analisis/order\_parameter.xvg}
\end{Highlighting}
\end{Shaded}

El archivo de índice debe definir los átomos de cadena en el orden
requerido. No mezcle carbonos de especies lipídicas diferentes en una
misma curva sin justificarlo.

\subsubsection{26. Difusión lateral de
lípidos}\label{difusiuxf3n-lateral-de-luxedpidos}

Para difusión bidimensional:

\[
MSD_{xy}(t)=
\left\langle
[x(t)-x(0)]^2+[y(t)-y(0)]^2
\right\rangle
\]

En el régimen difusivo:

\[
MSD_{xy}(t)=4D_{xy}t
\]

Comando:

\begin{Shaded}
\begin{Highlighting}[]
\ExtensionTok{gmx}\NormalTok{ msd }\DataTypeTok{\textbackslash{}}
    \AttributeTok{{-}s}\NormalTok{ step7\_1.tpr }\DataTypeTok{\textbackslash{}}
    \AttributeTok{{-}f}\NormalTok{ md\_all.xtc }\DataTypeTok{\textbackslash{}}
    \AttributeTok{{-}n}\NormalTok{ index.ndx }\DataTypeTok{\textbackslash{}}
    \AttributeTok{{-}sel} \StringTok{\textquotesingle{}res\_com of group "POPC"\textquotesingle{}} \DataTypeTok{\textbackslash{}}
    \AttributeTok{{-}lateral}\NormalTok{ z }\DataTypeTok{\textbackslash{}}
    \AttributeTok{{-}o}\NormalTok{ 05\_analisis/msd\_popc\_xy.xvg}
\end{Highlighting}
\end{Shaded}

Compruebe la sintaxis con \textbf{gmx msd -h}, porque cambió respecto de
versiones antiguas. El ajuste no debe incluir el régimen balístico
inicial ni una región donde el MSD todavía sea subdifusivo. La difusión
lipídica converge lentamente y depende del tamaño del sistema.

\subsubsection{27. Inclinación de la proteína y del
canal}\label{inclinaciuxf3n-de-la-proteuxedna-y-del-canal}

Defina un vector entre centros de masa de dos grupos situados en
extremos opuestos del poro y calcule su ángulo con z:

\begin{Shaded}
\begin{Highlighting}[]
\ExtensionTok{gmx}\NormalTok{ gangle }\DataTypeTok{\textbackslash{}}
    \AttributeTok{{-}s}\NormalTok{ step7\_1.tpr }\DataTypeTok{\textbackslash{}}
    \AttributeTok{{-}f}\NormalTok{ md\_all.xtc }\DataTypeTok{\textbackslash{}}
    \AttributeTok{{-}n}\NormalTok{ index.ndx }\DataTypeTok{\textbackslash{}}
    \AttributeTok{{-}g1}\NormalTok{ vector }\DataTypeTok{\textbackslash{}}
    \AttributeTok{{-}group1} \StringTok{\textquotesingle{}com of group "Pore\_lower" plus com of group "Pore\_upper"\textquotesingle{}} \DataTypeTok{\textbackslash{}}
    \AttributeTok{{-}g2}\NormalTok{ z }\DataTypeTok{\textbackslash{}}
    \AttributeTok{{-}oav}\NormalTok{ 05\_analisis/pore\_tilt.xvg}
\end{Highlighting}
\end{Shaded}

La selección debe adaptarse a residuos conservados de la acuaporina.
Para el tetrámero, genere una curva por monómero.

\subsubsection{28. Radio y continuidad del
poro}\label{radio-y-continuidad-del-poro}

GROMACS no proporciona por sí solo un perfil completo de radio de poro
equivalente a herramientas especializadas. Puede:

\begin{itemize}
\tightlist
\item
  medir distancias entre residuos del filtro;
\item
  calcular densidad de agua a lo largo del eje;
\item
  usar herramientas externas como HOLE o CHAP;
\item
  comparar perfiles por monómero y por bloque temporal.
\end{itemize}

Una disminución local del radio no implica cierre funcional si el agua
mantiene una cadena continua. Tampoco una cavidad geométricamente
abierta demuestra permeabilidad.

\subsubsection{29. Puentes de hidrógeno y orientación del
agua}\label{puentes-de-hidruxf3geno-y-orientaciuxf3n-del-agua}

Puentes de hidrógeno proteína--agua:

\begin{Shaded}
\begin{Highlighting}[]
\ExtensionTok{gmx}\NormalTok{ hbond }\DataTypeTok{\textbackslash{}}
    \AttributeTok{{-}s}\NormalTok{ step7\_1.tpr }\DataTypeTok{\textbackslash{}}
    \AttributeTok{{-}f}\NormalTok{ 05\_analisis/md\_fit.xtc }\DataTypeTok{\textbackslash{}}
    \AttributeTok{{-}n}\NormalTok{ index.ndx }\DataTypeTok{\textbackslash{}}
    \AttributeTok{{-}r} \StringTok{\textquotesingle{}group "Protein"\textquotesingle{}} \DataTypeTok{\textbackslash{}}
    \AttributeTok{{-}t} \StringTok{\textquotesingle{}group "Water"\textquotesingle{}} \DataTypeTok{\textbackslash{}}
    \AttributeTok{{-}num}\NormalTok{ 05\_analisis/protein\_water\_hbonds.xvg}
\end{Highlighting}
\end{Shaded}

Verifique la sintaxis exacta con:

\begin{Shaded}
\begin{Highlighting}[]
\ExtensionTok{gmx}\NormalTok{ hbond }\AttributeTok{{-}h}
\end{Highlighting}
\end{Shaded}

En acuaporinas importa no solo la ocupación, sino la orientación de las
moléculas de agua cerca de los motivos NPA. El cambio de orientación
contribuye al mecanismo de exclusión de protones. Este análisis requiere
vectores dipolares o enlaces O--H y una coordenada axial referida al
poro.

\subsubsection{30. Conteo de eventos de
permeación}\label{conteo-de-eventos-de-permeaciuxf3n}

Contar moléculas dentro de un cilindro en cada fotograma no equivale a
contar permeaciones. Un evento debe exigir una trayectoria completa
desde un reservorio hasta el opuesto.

Defina tres estados respecto del centro del canal:

\begin{Shaded}
\begin{Highlighting}[]
\NormalTok{A: z \textless{} {-}z0}
\NormalTok{P: {-}z0 ≤ z ≤ z0 y dentro del radio del poro}
\NormalTok{B: z \textgreater{} z0}
\end{Highlighting}
\end{Shaded}

Una permeación A→B requiere la secuencia A→P→B para la misma molécula,
sin reiniciar el conteo por fluctuaciones en el límite. Use coordenadas
relativas al centro de cada monómero y corrija PBC.

El flujo neto bajo un gradiente puede expresarse como:

\[
J=\frac{N_{A\rightarrow B}-N_{B\rightarrow A}}{t}
\]

Una estimación directa de permeabilidad osmótica necesita además el
gradiente de concentración:

\[
P_f=\frac{J}{\Delta c}
\]

Las unidades dependen de si \(J\) se expresa como moléculas por tiempo,
moles por tiempo o volumen por tiempo. En equilibrio, donde no hay flujo
neto sostenido, se utilizan formulaciones basadas en fluctuaciones
colectivas; no debe aplicarse la ecuación anterior con \(\Delta c=0\).

\subsubsection{31. Permeabilidad colectiva en
equilibrio}\label{permeabilidad-colectiva-en-equilibrio}

Para una coordenada colectiva \(n(t)\) que registra el desplazamiento
neto de agua a través del canal:

\[
\left\langle
[n(t)-n(0)]^2
\right\rangle
\xrightarrow[t\ \mathrm{grande}]{}
2D_n t
\]

La permeabilidad osmótica de canal único puede relacionarse con:

\[
p_f=v_wD_n
\]

donde \(v_w\) es el volumen molecular del agua y \(D_n\) es el
coeficiente de difusión de la coordenada colectiva. La definición exacta
de \(n(t)\), el tratamiento de PBC y el intervalo de ajuste deben
mantenerse constantes al comparar sistemas.

No calcule una permeabilidad confiable a partir de unos pocos cruces.
Use réplicas, intervalos de confianza y análisis por bloques.

\subsubsection{32. Contactos
lípido--proteína}\label{contactos-luxedpidoproteuxedna}

Los contactos persistentes pueden revelar sitios anulares o específicos.
Un contacto simple puede definirse por una distancia máxima:

\begin{Shaded}
\begin{Highlighting}[]
\ExtensionTok{gmx}\NormalTok{ select }\DataTypeTok{\textbackslash{}}
    \AttributeTok{{-}s}\NormalTok{ step7\_1.tpr }\DataTypeTok{\textbackslash{}}
    \AttributeTok{{-}f}\NormalTok{ md\_all.xtc }\DataTypeTok{\textbackslash{}}
    \AttributeTok{{-}n}\NormalTok{ index.ndx }\DataTypeTok{\textbackslash{}}
    \AttributeTok{{-}select} \StringTok{\textquotesingle{}resname POPC and within 0.45 of group "Protein"\textquotesingle{}} \DataTypeTok{\textbackslash{}}
    \AttributeTok{{-}os}\NormalTok{ 05\_analisis/lipid\_contacts\_size.xvg}
\end{Highlighting}
\end{Shaded}

Este resultado informa cuántos átomos o posiciones cumplen la selección,
según la salida elegida; no es automáticamente un número de lípidos
únicos. Para residencia lipídica deben seguirse identidades moleculares
y tolerar interrupciones cortas.

En una mezcla, analice enriquecimiento relativo:

\[
E_i=
\frac{x_i^{\mathrm{contacto}}}
     {x_i^{\mathrm{membrana}}}
\]

donde \(E_i>1\) sugiere enriquecimiento de la especie \(i\) alrededor de
la proteína. El resultado depende del corte, del área considerada y del
muestreo.

\subsubsection{33. Potencial electrostático y sistemas con
voltaje}\label{potencial-electrostuxe1tico-y-sistemas-con-voltaje}

El potencial promedio a lo largo de z puede calcularse con \textbf{gmx
potential} a partir de grupos de carga apropiados. Una sola bicapa
periódica no crea automáticamente dos reservorios independientes.

Para estudiar flujo iónico bajo potencial sostenido puede utilizarse el
protocolo de electrofisiología computacional de GROMACS, que emplea
típicamente dos bicapas, dos compartimientos y un desequilibrio de
carga:

\[
\Delta U=\frac{\Delta q}{C_{\mathrm{membrana}}}
\]

Este montaje no es necesario para permeación de agua en equilibrio y no
debe añadirse sin una pregunta electrofisiológica explícita.

\subsubsection{34. Réplicas, descarte y
convergencia}\label{ruxe9plicas-descarte-y-convergencia}

Una trayectoria larga no reemplaza réplicas independientes. Para
comparar acuaporinas, mutantes o composiciones:

\begin{itemize}
\tightlist
\item
  use varias semillas de velocidad;
\item
  mantenga idénticos los parámetros comunes;
\item
  descarte la etapa transitoria según observables;
\item
  compare bloques temporales;
\item
  informe incertidumbre entre réplicas;
\item
  evite seleccionar la trayectoria ``más estable''.
\end{itemize}

La estabilización del RMSD no demuestra que APL, espesor, difusión
lipídica o permeación hayan convergido.

\subsubsection{35. Simulación
coarse-grained}\label{simulaciuxf3n-coarse-grained}

Un modelo CG reduce grados de libertad y permite escalas espaciales o
temporales mayores. A cambio, pierde detalle atomístico, modifica la
cinética efectiva y puede requerir restricciones estructurales
adicionales.

Use un constructor y una versión de Martini compatibles. No convierta
simplemente las coordenadas all-atom y reutilice:

\begin{itemize}
\tightlist
\item
  \textbf{topol.top} de CHARMM36;
\item
  archivos ITP atomísticos;
\item
  modelo de agua atomístico;
\item
  cortes atomísticos;
\item
  paso de integración atomístico;
\item
  parámetros de restricciones atomísticas.
\end{itemize}

Registre:

\begin{itemize}
\tightlist
\item
  versión de Martini;
\item
  método de mapeo;
\item
  modelo de agua;
\item
  tratamiento electrostático;
\item
  red elástica;
\item
  radio y constante de la red;
\item
  lípidos CG;
\item
  paso de integración;
\item
  método de backmapping, si se utiliza.
\end{itemize}

\subsubsection{36. Redes elásticas en proteínas
CG}\label{redes-eluxe1sticas-en-proteuxednas-cg}

Una red elástica ayuda a conservar estructura terciaria o cuaternaria,
pero puede suprimir:

\begin{itemize}
\tightlist
\item
  apertura del canal;
\item
  inclinación helicoidal;
\item
  movimientos entre dominios;
\item
  respiración del poro;
\item
  reorganización oligomérica.
\end{itemize}

No use una red ``porque es el valor predeterminado'' si la variable de
interés depende del movimiento restringido. Realice análisis de
sensibilidad cambiando el corte o la constante, o compare con un modelo
sin red cuando sea estable.

La red elástica no reemplaza enlaces covalentes, disulfuros ni una
unidad biológica correcta.

\subsubsection{37. Paso de integración en
CG}\label{paso-de-integraciuxf3n-en-cg}

Los pasos CG suelen ser mayores que en all-atom, pero el máximo estable
depende del modelo, el agua, las restricciones y la geometría inicial.
Comience con el protocolo recomendado para la versión concreta y revise:

\begin{itemize}
\tightlist
\item
  energía;
\item
  errores de restricciones;
\item
  temperatura y presión;
\item
  estabilidad del poro;
\item
  comportamiento de los lípidos;
\item
  sensibilidad al paso de integración.
\end{itemize}

La aceleración aparente del tiempo en CG no debe interpretarse como una
correspondencia universal y exacta con tiempo experimental.

\subsubsection{38. Visualización CG}\label{visualizaciuxf3n-cg}

VMD y otros visores pueden no inferir correctamente enlaces entre
partículas CG porque usan reglas pensadas para geometrías atomísticas.
La ausencia visual de enlaces no significa que falten en la topología.

Las representaciones deben basarse en nombres reales de partículas y
residuos. Ejemplos habituales en Martini incluyen:

\begin{Shaded}
\begin{Highlighting}[]
\NormalTok{name BB}
\NormalTok{resname POPC}
\NormalTok{resname W}
\NormalTok{resname NA CL}
\end{Highlighting}
\end{Shaded}

Los nombres dependen de la versión y del constructor. Compruébelos en la
coordenada y en los ITP.

Para reducir memoria, genere una trayectoria de visualización:

\begin{Shaded}
\begin{Highlighting}[]
\ExtensionTok{gmx}\NormalTok{ trjconv }\DataTypeTok{\textbackslash{}}
    \AttributeTok{{-}s}\NormalTok{ cg\_md.tpr }\DataTypeTok{\textbackslash{}}
    \AttributeTok{{-}f}\NormalTok{ cg\_md.xtc }\DataTypeTok{\textbackslash{}}
    \AttributeTok{{-}o}\NormalTok{ cg\_visualization.xtc }\DataTypeTok{\textbackslash{}}
    \AttributeTok{{-}n}\NormalTok{ cg\_visualization.ndx }\DataTypeTok{\textbackslash{}}
    \AttributeTok{{-}dt}\NormalTok{ 1000}
\end{Highlighting}
\end{Shaded}

No use la trayectoria reducida para análisis que requieran resolución
temporal mayor.

La representación secundaria atomística no puede reconstruirse
exactamente a partir de unas pocas partículas por residuo. Una
caricatura CG es una interpretación gráfica apoyada en la asignación
secundaria, no una observación directa de todos los enlaces y ángulos
del esqueleto.

\subsubsection{39. Errores conceptuales
frecuentes}\label{errores-conceptuales-frecuentes}

\begin{itemize}
\tightlist
\item
  Simular una sola cadena de una acuaporina tetramérica sin
  justificarlo.
\item
  Interpretar el centro del tetrámero como el poro de cada monómero.
\item
  Orientar por eje principal sin revisar el cinturón hidrofóbico.
\item
  Usar POPC puro y describirlo como membrana fisiológica.
\item
  Considerar 0.15 M sinónimo universal de isotonicidad.
\item
  Regenerar con \textbf{pdb2gmx} una topología de CHARMM-GUI.
\item
  Sustituir los MDP generados por parámetros genéricos.
\item
  Usar presión isotrópica en una bicapa plana sin validación.
\item
  Retirar todas las restricciones en una sola etapa.
\item
  Ejecutar 1 ps y declarar la membrana equilibrada.
\item
  Ajustar rotación y traslación antes de calcular difusión lipídica.
\item
  Calcular APL dividiendo por todos los lípidos de ambas monocapas.
\item
  Contar ocupaciones del poro como eventos completos de permeación.
\item
  Interpretar una sola trayectoria como estimación convergida de
  permeabilidad.
\item
  Usar una red elástica CG que impide el movimiento que se quiere medir.
\item
  Suponer que CSH es necesario para ejecutar los archivos de CHARMM-GUI.
\end{itemize}

\subsubsection{40. Criterios para
producción}\label{criterios-para-producciuxf3n}

Inicie la producción cuando:

\begin{itemize}
\tightlist
\item
  el ensamblado oligomérico sea correcto;
\item
  la proteína esté orientada y centrada de forma razonable;
\item
  no existan residuos o componentes sin parametrizar;
\item
  las monocapas tengan la composición y asimetría previstas;
\item
  topología, coordenadas e índice sean coherentes;
\item
  los poros estén hidratados sin agua espuria en el núcleo;
\item
  no existan contactos graves ni errores LINCS;
\item
  temperatura, volumen, área xy y espesor sean estacionarios;
\item
  no haya deriva sistemática de la proteína;
\item
  las restricciones hayan alcanzado la etapa planificada;
\item
  el intervalo descartado y los análisis se hayan definido;
\item
  se hayan planificado réplicas independientes.
\end{itemize}

\subsubsection{Fuentes}\label{fuentes-4}

\begin{enumerate}
\def\labelenumi{\arabic{enumi}.}
\tightlist
\item
  \href{https://charmm-gui.org/?doc=tutorial&project=membrane&chapter=membrane_vdac_hetero}{Tutorial
  de CHARMM-GUI para complejos proteína--membrana heterogéneos}
\item
  \href{https://pmc.ncbi.nlm.nih.gov/articles/PMC8158057/}{Preparación
  de proteínas de membrana con CHARMM-GUI}
\item
  \href{https://www.charmm-gui.org/charmmdoc/usage.html}{Documentación
  de uso de CHARMM}
\item
  \href{https://manual.gromacs.org/current/user-guide/mdp-options.html}{Opciones
  MDP de GROMACS 2026.3}
\item
  \href{https://manual.gromacs.org/current/onlinehelp/gmx-msd.html}{Referencia
  de gmx msd}
\item
  \href{https://manual.gromacs.org/current/onlinehelp/gmx-order.html}{Referencia
  de gmx order}
\item
  \href{https://manual.gromacs.org/current/reference-manual/special/comp-electrophys.html}{Electrofisiología
  computacional en GROMACS}
\item
  \href{https://pmc.ncbi.nlm.nih.gov/articles/PMC7271963/}{Balusek et
  al., \emph{Accelerating Membrane Simulations with Hydrogen Mass
  Repartitioning}}
\item
  \href{https://gromacs.bioexcel.eu/t/hydrogen-mass-repartitioning/1171}{Discusión
  técnica de HMR en el foro de GROMACS}
\item
  \href{https://gitlab.com/gromacs/gromacs/-/work_items/5007}{Seguimiento
  de la documentación de HMR en GROMACS, work item 5007}
\item
  \href{https://manual.gromacs.org/2026.3/release-notes/2024/major/performance.html\#flexible-hydrogen-mass-repartitioning-using-grompp}{HMR
  flexible mediante grompp: notas de rendimiento de GROMACS 2026.3}
\item
  \href{https://manual.gromacs.org/2026.3/user-guide/mdp-options.html\#mdp-value-mass-repartition-factor}{Opciones
  MDP de GROMACS 2026.3: mass-repartition-factor}
\end{enumerate}

\subsection{Energía libre de perturbación: energía libre de
solvatación}\label{energuxeda-libre-de-perturbaciuxf3n-energuxeda-libre-de-solvataciuxf3n}

\subsubsection{Caso de estudio: formaldehído en
agua}\label{caso-de-estudio-formaldehuxeddo-en-agua}

La energía libre de solvatación describe el cambio de energía libre
asociado con transferir una especie química desde una fase de referencia
---habitualmente gas ideal--- hasta una solución a dilución infinita:

\[
\Delta G_{\mathrm{solv}} =
G_{\mathrm{soluto,solución}}-
G_{\mathrm{soluto,gas}}
\]

Una \(\Delta G_{\mathrm{solv}}<0\) indica que la transferencia al
solvente es favorable bajo la convención y los estados estándar
especificados. La magnitud depende de la identidad química, el estado de
protonación, el solvente, la temperatura, el campo de fuerza, los
estados estándar, el tratamiento electrostático, el tamaño de caja y el
muestreo.

Este tutorial muestra un cálculo alquímico de una molécula neutra. El
formaldehído sirve para ilustrar la mecánica de GROMACS, pero presenta
una limitación química crítica: en agua reacciona para formar
metanodiol. El resultado obtenido con una topología no reactiva
corresponde al formaldehído molecular \(\mathrm{H_2CO}\), no al
equilibrio químico completo de una solución de formaldehído.

\subsubsection{1. Fundamento del método}\label{fundamento-del-muxe9todo}

Los estados físicos A y B pueden tener distribuciones configuracionales
con poca superposición. Se construye entonces un camino alquímico
mediante \(\lambda\):

\[
H(\mathbf{x};\lambda)
\]

Una interpolación conceptual sencilla es:

\[
H(\lambda)=(1-\lambda)H_A+\lambda H_B
\]

Los estados intermedios no necesitan ser físicamente realizables. Solo
deben conectar los extremos mediante una ruta reversible con
superposición estadística suficiente.

Para un desacoplamiento en solución:

\begin{Shaded}
\begin{Highlighting}[]
\NormalTok{λ = 0                         λ = 1}
\NormalTok{soluto completamente          soluto sin interacciones}
\NormalTok{acoplado al agua              no enlazadas con el agua}
\end{Highlighting}
\end{Shaded}

Las interacciones internas del soluto se mantienen. La molécula
desacoplada sigue presente en las coordenadas, pero no ejerce
interacciones de Coulomb ni Lennard-Jones sobre el solvente.

\subsubsection{2. FEP, TI, BAR y MBAR}\label{fep-ti-bar-y-mbar}

La ecuación de Zwanzig es:

\[
\Delta G_{A\rightarrow B}
=
-k_\mathrm{B}T
\ln
\left\langle
\exp[-\beta(U_B-U_A)]
\right\rangle_A
\]

con \(\beta=1/(k_\mathrm{B}T)\), o \(\beta=1/(RT)\) para cantidades
molares. Es exacta en el límite de muestreo infinito, pero ineficiente
cuando A y B tienen poca superposición.

En integración termodinámica:

\[
\Delta G_{A\rightarrow B}
=
\int_0^1
\left\langle
\frac{\partial H}{\partial\lambda}
\right\rangle_\lambda
d\lambda
\]

TI requiere una grilla que resuelva la curvatura del integrando y una
integración numérica apropiada.

BAR combina información en ambas direcciones entre dos estados vecinos.
MBAR analiza todos los estados conjuntamente. Ninguno corrige una ruta
alquímico mal diseñada ni grados de libertad sin muestrear.

Este tutorial utiliza \textbf{gmx bar}, incluido en GROMACS.

\subsubsection{3. Convención de signo}\label{convenciuxf3n-de-signo}

Definamos:

\begin{itemize}
\tightlist
\item
  estado 0: formaldehído acoplado al agua;
\item
  estado 1: formaldehído desacoplado del agua.
\end{itemize}

La simulación calcula:

\[
\Delta G_{0\rightarrow1}
=
G_{\mathrm{desacoplado}}-
G_{\mathrm{acoplado}}
=
\Delta G_{\mathrm{desacoplamiento}}
\]

Por lo tanto:

\[
\boxed{
\Delta G_{\mathrm{solv}}
=
-\Delta G_{0\rightarrow1}
}
\]

Si se invierten \textbf{couple-lambda0} y \textbf{couple-lambda1},
también se invierte el signo. Antes de interpretar un número, escriba
qué representa cada extremo.

\subsubsection{4. Separar Coulomb y
Lennard-Jones}\label{separar-coulomb-y-lennard-jones}

Apagar simultáneamente cargas y Lennard-Jones puede servir en un
ejercicio corto, pero no es la ruta preferida para un resultado
cuantitativo.

Una secuencia habitual es:

\begin{enumerate}
\def\labelenumi{\arabic{enumi}.}
\tightlist
\item
  apagar electrostática mientras el volumen excluido permanece;
\item
  apagar Lennard-Jones mediante potenciales soft-core.
\end{enumerate}

Si desaparece primero el volumen excluido, el solvente puede ocupar la
posición del soluto mientras todavía existen cargas puntuales.

\[
\Delta G_{\mathrm{desacoplamiento}}
=
\Delta G_{\mathrm{Coulomb}}
+
\Delta G_{\mathrm{LJ}}
\]

Esta separación permite identificar qué componente necesita mayor
densidad de estados λ.

\subsubsection{5. Potenciales soft-core}\label{potenciales-soft-core}

La interpolación lineal de Lennard-Jones puede producir divergencias
cuando dos partículas se superponen. Los potenciales soft-core mantienen
finita la energía en estados intermedios.

\begin{Shaded}
\begin{Highlighting}[]
\DataTypeTok{sc{-}alpha       }\OtherTok{=}\StringTok{ }\FloatTok{0.5}
\DataTypeTok{sc{-}power       }\OtherTok{=}\StringTok{ }\DecValTok{1}
\DataTypeTok{sc{-}sigma       }\OtherTok{=}\StringTok{ }\FloatTok{0.3}
\DataTypeTok{sc{-}coul        }\OtherTok{=}\StringTok{ }\KeywordTok{no}
\end{Highlighting}
\end{Shaded}

Son valores iniciales, no constantes universales. \textbf{sc-power}
requiere un entero y \textbf{sc-sigma} se expresa en nm. Como Coulomb se
elimina antes que Lennard-Jones, no se activa soft-core electrostático.

\subsubsection{6. Selección de estados
λ}\label{selecciuxf3n-de-estados-ux3bb}

Agregar estados uniformemente no garantiza precisión. Conviene aumentar
su densidad donde disminuye la superposición, crece la varianza, entra
solvente en una cavidad o cambia rápidamente el potencial soft-core.

Grilla inicial de 19 estados:

\begin{Shaded}
\begin{Highlighting}[]
\DataTypeTok{coul{-}lambdas }\OtherTok{=}\StringTok{ 0.00 0.25 0.50 0.75 1.00 1.00 1.00 1.00 1.00 1.00 1.00 1.00 1.00 1.00 1.00 1.00 1.00 1.00 1.00}

\DataTypeTok{vdw{-}lambdas  }\OtherTok{=}\StringTok{ 0.00 0.00 0.00 0.00 0.00 0.05 0.10 0.15 0.20 0.30 0.40 0.50 0.60 0.70 0.80 0.85 0.90 0.95 1.00}
\end{Highlighting}
\end{Shaded}

Los índices 0--4 descargan el soluto y los índices 4--18 eliminan
Lennard-Jones. El estado 4 es común a ambas etapas. La grilla debe
reajustarse después de analizar solapamiento y errores por intervalo.

\subsubsection{7. Parametrización del
formaldehído}\label{parametrizaciuxf3n-del-formaldehuxeddo}

La topología debe ser compatible con el campo de fuerza y el modelo de
agua. No use automáticamente parámetros antiguos de CHARMM27 ni copie
tipos de SwissParam sin validación.

Compruebe:

\begin{itemize}
\tightlist
\item
  fórmula \(\mathrm{CH_2O}\);
\item
  geometría trigonal plana;
\item
  carga neta cero;
\item
  nombres y orden de átomos;
\item
  cargas parciales;
\item
  parámetros enlazados;
\item
  interacción del carbonilo;
\item
  compatibilidad con el campo de fuerza;
\item
  ausencia de tipos atómicos duplicados.
\end{itemize}

Organización sugerida:

\begin{Shaded}
\begin{Highlighting}[]
\NormalTok{00\_parametros/}
\NormalTok{    formaldehyde.gro}
\NormalTok{    formaldehyde.itp}
\NormalTok{    formaldehyde\_atomtypes.itp}
\NormalTok{01\_preparacion/}
\NormalTok{02\_equilibracion/}
\NormalTok{03\_fep/}
\NormalTok{04\_analisis/}
\NormalTok{topol.top}
\end{Highlighting}
\end{Shaded}

Topología general:

\begin{Shaded}
\begin{Highlighting}[]
\CommentTok{\#include "campo\_de\_fuerza.ff/forcefield.itp"}
\CommentTok{\#include "00\_parametros/formaldehyde\_atomtypes.itp"}
\CommentTok{\#include "00\_parametros/formaldehyde.itp"}
\CommentTok{\#include "campo\_de\_fuerza.ff/modelo\_de\_agua.itp"}

\KeywordTok{[ system ]}
\DataTypeTok{Formaldehído molecular en agua}

\KeywordTok{[ molecules ]}
\DataTypeTok{FOR    1}
\DataTypeTok{SOL    NUMERO\_DE\_AGUAS}
\end{Highlighting}
\end{Shaded}

\textbf{FOR} debe coincidir con \texttt{{[}\ moleculetype\ {]}}. Si
existen varias moléculas FOR, \textbf{couple-moltype = FOR} perturbará
todas. Para desacoplar una sola copia se necesita un tipo molecular
exclusivo.

\subsubsection{8. Construcción de la
caja}\label{construcciuxf3n-de-la-caja}

\begin{Shaded}
\begin{Highlighting}[]
\ExtensionTok{gmx}\NormalTok{ editconf }\DataTypeTok{\textbackslash{}}
    \AttributeTok{{-}f}\NormalTok{ 00\_parametros/formaldehyde.gro }\DataTypeTok{\textbackslash{}}
    \AttributeTok{{-}o}\NormalTok{ 01\_preparacion/formaldehyde\_box.gro }\DataTypeTok{\textbackslash{}}
    \AttributeTok{{-}c} \DataTypeTok{\textbackslash{}}
    \AttributeTok{{-}d}\NormalTok{ 1.2 }\DataTypeTok{\textbackslash{}}
    \AttributeTok{{-}bt}\NormalTok{ dodecahedron}
\end{Highlighting}
\end{Shaded}

La distancia de 1.2 nm es inicial. Debe reducir interacciones con
imágenes periódicas y ser coherente con los cortes.

\begin{Shaded}
\begin{Highlighting}[]
\ExtensionTok{gmx}\NormalTok{ solvate }\DataTypeTok{\textbackslash{}}
    \AttributeTok{{-}cp}\NormalTok{ 01\_preparacion/formaldehyde\_box.gro }\DataTypeTok{\textbackslash{}}
    \AttributeTok{{-}cs}\NormalTok{ spc216.gro }\DataTypeTok{\textbackslash{}}
    \AttributeTok{{-}o}\NormalTok{ 01\_preparacion/formaldehyde\_water.gro }\DataTypeTok{\textbackslash{}}
    \AttributeTok{{-}p}\NormalTok{ topol.top}
\end{Highlighting}
\end{Shaded}

\textbf{spc216.gro} aporta coordenadas de agua de tres sitios. El modelo
físico queda determinado por la topología incluida.

Para una sola molécula no hace falta insertarla en una caja vacía y
luego redefinir otra caja: \textbf{editconf} seguido de \textbf{solvate}
es suficiente.

\subsubsection{9. Minimización}\label{minimizaciuxf3n}

\begin{Shaded}
\begin{Highlighting}[]
\DataTypeTok{title           }\OtherTok{=}\StringTok{ EM de formaldehído en agua}
\DataTypeTok{integrator      }\OtherTok{=}\StringTok{ steep}
\DataTypeTok{nsteps          }\OtherTok{=}\StringTok{ }\DecValTok{50000}
\DataTypeTok{emtol           }\OtherTok{=}\StringTok{ }\FloatTok{1000.0}
\DataTypeTok{emstep          }\OtherTok{=}\StringTok{ }\FloatTok{0.01}

\DataTypeTok{cutoff{-}scheme   }\OtherTok{=}\StringTok{ Verlet}
\DataTypeTok{nstlist         }\OtherTok{=}\StringTok{ }\DecValTok{20}
\DataTypeTok{rlist           }\OtherTok{=}\StringTok{ }\FloatTok{1.2}
\DataTypeTok{coulombtype     }\OtherTok{=}\StringTok{ PME}
\DataTypeTok{rcoulomb        }\OtherTok{=}\StringTok{ }\FloatTok{1.2}
\DataTypeTok{vdwtype         }\OtherTok{=}\StringTok{ Cut{-}off}
\DataTypeTok{rvdw            }\OtherTok{=}\StringTok{ }\FloatTok{1.2}
\DataTypeTok{pbc             }\OtherTok{=}\StringTok{ xyz}
\end{Highlighting}
\end{Shaded}

Los cortes deben adaptarse al campo de fuerza.

\begin{Shaded}
\begin{Highlighting}[]
\FunctionTok{mkdir} \AttributeTok{{-}p}\NormalTok{ 02\_equilibracion}

\ExtensionTok{gmx}\NormalTok{ grompp }\DataTypeTok{\textbackslash{}}
    \AttributeTok{{-}f}\NormalTok{ em.mdp }\DataTypeTok{\textbackslash{}}
    \AttributeTok{{-}c}\NormalTok{ 01\_preparacion/formaldehyde\_water.gro }\DataTypeTok{\textbackslash{}}
    \AttributeTok{{-}p}\NormalTok{ topol.top }\DataTypeTok{\textbackslash{}}
    \AttributeTok{{-}o}\NormalTok{ 02\_equilibracion/em.tpr }\DataTypeTok{\textbackslash{}}
    \AttributeTok{{-}pp}\NormalTok{ 02\_equilibracion/processed.top}

\ExtensionTok{gmx}\NormalTok{ mdrun }\DataTypeTok{\textbackslash{}}
    \AttributeTok{{-}deffnm}\NormalTok{ 02\_equilibracion/em }\DataTypeTok{\textbackslash{}}
    \AttributeTok{{-}v}
\end{Highlighting}
\end{Shaded}

Terminar por alcanzar \textbf{nsteps} sin cumplir \textbf{emtol} no
significa convergencia. Una fuerza máxima de miles de kJ·mol⁻¹·nm⁻¹ debe
investigarse.

\subsubsection{10. Equilibración física
previa}\label{equilibraciuxf3n-fuxedsica-previa}

NVT:

\begin{Shaded}
\begin{Highlighting}[]
\DataTypeTok{title         }\OtherTok{=}\StringTok{ NVT previa a FEP}
\DataTypeTok{integrator    }\OtherTok{=}\StringTok{ md}
\DataTypeTok{dt            }\OtherTok{=}\StringTok{ }\FloatTok{0.002}
\DataTypeTok{nsteps        }\OtherTok{=}\StringTok{ }\DecValTok{250000}
\DataTypeTok{continuation  }\OtherTok{=}\StringTok{ }\KeywordTok{no}
\DataTypeTok{gen{-}vel       }\OtherTok{=}\StringTok{ }\KeywordTok{yes}
\DataTypeTok{gen{-}temp      }\OtherTok{=}\StringTok{ }\FloatTok{298.15}
\DataTypeTok{gen{-}seed      }\OtherTok{=}\StringTok{ }\DecValTok{2026}

\DataTypeTok{constraints   }\OtherTok{=}\StringTok{ h{-}bonds}
\DataTypeTok{tcoupl        }\OtherTok{=}\StringTok{ V{-}rescale}
\DataTypeTok{tc{-}grps       }\OtherTok{=}\StringTok{ System}
\DataTypeTok{tau{-}t         }\OtherTok{=}\StringTok{ }\FloatTok{1.0}
\DataTypeTok{ref{-}t         }\OtherTok{=}\StringTok{ }\FloatTok{298.15}
\DataTypeTok{pcoupl        }\OtherTok{=}\StringTok{ }\KeywordTok{no}
\DataTypeTok{pbc           }\OtherTok{=}\StringTok{ xyz}
\end{Highlighting}
\end{Shaded}

\begin{Shaded}
\begin{Highlighting}[]
\ExtensionTok{gmx}\NormalTok{ grompp }\DataTypeTok{\textbackslash{}}
    \AttributeTok{{-}f}\NormalTok{ nvt.mdp }\DataTypeTok{\textbackslash{}}
    \AttributeTok{{-}c}\NormalTok{ 02\_equilibracion/em.gro }\DataTypeTok{\textbackslash{}}
    \AttributeTok{{-}p}\NormalTok{ topol.top }\DataTypeTok{\textbackslash{}}
    \AttributeTok{{-}o}\NormalTok{ 02\_equilibracion/nvt.tpr}

\ExtensionTok{gmx}\NormalTok{ mdrun }\AttributeTok{{-}deffnm}\NormalTok{ 02\_equilibracion/nvt }\AttributeTok{{-}v}
\end{Highlighting}
\end{Shaded}

NPT:

\begin{Shaded}
\begin{Highlighting}[]
\DataTypeTok{title            }\OtherTok{=}\StringTok{ NPT previa a FEP}
\DataTypeTok{integrator       }\OtherTok{=}\StringTok{ md}
\DataTypeTok{dt               }\OtherTok{=}\StringTok{ }\FloatTok{0.002}
\DataTypeTok{nsteps           }\OtherTok{=}\StringTok{ }\DecValTok{1000000}
\DataTypeTok{continuation     }\OtherTok{=}\StringTok{ }\KeywordTok{yes}
\DataTypeTok{gen{-}vel          }\OtherTok{=}\StringTok{ }\KeywordTok{no}

\DataTypeTok{constraints      }\OtherTok{=}\StringTok{ h{-}bonds}
\DataTypeTok{tcoupl           }\OtherTok{=}\StringTok{ V{-}rescale}
\DataTypeTok{tc{-}grps          }\OtherTok{=}\StringTok{ System}
\DataTypeTok{tau{-}t            }\OtherTok{=}\StringTok{ }\FloatTok{1.0}
\DataTypeTok{ref{-}t            }\OtherTok{=}\StringTok{ }\FloatTok{298.15}

\DataTypeTok{pcoupl            }\OtherTok{=}\StringTok{ C{-}rescale}
\DataTypeTok{pcoupltype        }\OtherTok{=}\StringTok{ isotropic}
\DataTypeTok{tau{-}p             }\OtherTok{=}\StringTok{ }\FloatTok{5.0}
\DataTypeTok{ref{-}p             }\OtherTok{=}\StringTok{ }\FloatTok{1.0}
\DataTypeTok{compressibility   }\OtherTok{=}\StringTok{ 4.5e{-}5}
\DataTypeTok{pbc               }\OtherTok{=}\StringTok{ xyz}
\end{Highlighting}
\end{Shaded}

\begin{Shaded}
\begin{Highlighting}[]
\ExtensionTok{gmx}\NormalTok{ grompp }\DataTypeTok{\textbackslash{}}
    \AttributeTok{{-}f}\NormalTok{ npt.mdp }\DataTypeTok{\textbackslash{}}
    \AttributeTok{{-}c}\NormalTok{ 02\_equilibracion/nvt.gro }\DataTypeTok{\textbackslash{}}
    \AttributeTok{{-}t}\NormalTok{ 02\_equilibracion/nvt.cpt }\DataTypeTok{\textbackslash{}}
    \AttributeTok{{-}p}\NormalTok{ topol.top }\DataTypeTok{\textbackslash{}}
    \AttributeTok{{-}o}\NormalTok{ 02\_equilibracion/npt.tpr}

\ExtensionTok{gmx}\NormalTok{ mdrun }\AttributeTok{{-}deffnm}\NormalTok{ 02\_equilibracion/npt }\AttributeTok{{-}v}
\end{Highlighting}
\end{Shaded}

Cuarenta picosegundos no constituyen una duración universal de
equilibración. Evalúe densidad, volumen, energía y relajación del
solvente.

\subsubsection{11. Bloque MDP para
desacoplamiento}\label{bloque-mdp-para-desacoplamiento}

Mantenga los parámetros no enlazados del protocolo validado y añada:

\begin{Shaded}
\begin{Highlighting}[]
\DataTypeTok{free{-}energy            }\OtherTok{=}\StringTok{ }\KeywordTok{yes}
\DataTypeTok{couple{-}moltype         }\OtherTok{=}\StringTok{ FOR}

\CommentTok{; λ=0 acoplado; λ=1 desacoplado}
\DataTypeTok{couple{-}lambda0         }\OtherTok{=}\StringTok{ vdw{-}q}
\DataTypeTok{couple{-}lambda1         }\OtherTok{=}\StringTok{ none}
\DataTypeTok{couple{-}intramol        }\OtherTok{=}\StringTok{ }\KeywordTok{no}

\DataTypeTok{coul{-}lambdas }\OtherTok{=}\StringTok{ 0.00 0.25 0.50 0.75 1.00 1.00 1.00 1.00 1.00 1.00 1.00 1.00 1.00 1.00 1.00 1.00 1.00 1.00 1.00}
\DataTypeTok{vdw{-}lambdas  }\OtherTok{=}\StringTok{ 0.00 0.00 0.00 0.00 0.00 0.05 0.10 0.15 0.20 0.30 0.40 0.50 0.60 0.70 0.80 0.85 0.90 0.95 1.00}

\DataTypeTok{init{-}lambda{-}state      }\OtherTok{=}\StringTok{ LAMBDA\_STATE}
\DataTypeTok{calc{-}lambda{-}neighbors  }\OtherTok{=}\StringTok{ }\DecValTok{1}

\DataTypeTok{sc{-}alpha               }\OtherTok{=}\StringTok{ }\FloatTok{0.5}
\DataTypeTok{sc{-}power               }\OtherTok{=}\StringTok{ }\DecValTok{1}
\DataTypeTok{sc{-}sigma               }\OtherTok{=}\StringTok{ }\FloatTok{0.3}
\DataTypeTok{sc{-}coul                }\OtherTok{=}\StringTok{ }\KeywordTok{no}

\DataTypeTok{nstdhdl                }\OtherTok{=}\StringTok{ }\DecValTok{100}
\DataTypeTok{dhdl{-}derivatives       }\OtherTok{=}\StringTok{ }\KeywordTok{yes}
\DataTypeTok{separate{-}dhdl{-}file     }\OtherTok{=}\StringTok{ }\KeywordTok{yes}
\end{Highlighting}
\end{Shaded}

\textbf{calc-lambda-neighbors = 1} calcula diferencias con estados
vecinos para BAR. \textbf{nstdhdl} debe ser múltiplo de
\textbf{nstcalcenergy}.

\textbf{init-lambda-state} es el índice entero de las listas, no el
valor de λ. No use simultáneamente \textbf{fep-lambdas} y
\textbf{coul-lambdas/vdw-lambdas} para describir la misma ruta sin
controlar cómo se forma el vector λ.

\subsubsection{12. Equilibrar cada estado}\label{equilibrar-cada-estado}

Prepare:

\begin{itemize}
\tightlist
\item
  \textbf{fep\_equilibration.template.mdp} para relajar cada
  Hamiltoniano;
\item
  \textbf{fep\_production.template.mdp} para acumular datos.
\end{itemize}

Ambos contienen \textbf{LAMBDA\_STATE}.

Equilibración por ventana:

\begin{Shaded}
\begin{Highlighting}[]
\DataTypeTok{nsteps       }\OtherTok{=}\StringTok{ }\DecValTok{250000}
\DataTypeTok{nstdhdl      }\OtherTok{=}\StringTok{ }\DecValTok{0}
\DataTypeTok{gen{-}vel      }\OtherTok{=}\StringTok{ }\KeywordTok{yes}
\DataTypeTok{gen{-}temp     }\OtherTok{=}\StringTok{ }\FloatTok{298.15}
\DataTypeTok{gen{-}seed     }\OtherTok{=}\StringTok{ }\DecValTok{{-}1}
\end{Highlighting}
\end{Shaded}

Producción por ventana:

\begin{Shaded}
\begin{Highlighting}[]
\DataTypeTok{nsteps       }\OtherTok{=}\StringTok{ }\DecValTok{2500000}
\DataTypeTok{nstdhdl      }\OtherTok{=}\StringTok{ }\DecValTok{100}
\DataTypeTok{gen{-}vel      }\OtherTok{=}\StringTok{ }\KeywordTok{no}
\DataTypeTok{continuation }\OtherTok{=}\StringTok{ }\KeywordTok{yes}
\end{Highlighting}
\end{Shaded}

Con \(dt=0.002\ \mathrm{ps}\), corresponden a 0.5 ns y 5 ns. Son puntos
de partida, no garantías de convergencia.

\subsubsection{13. Crear y ejecutar las
ventanas}\label{crear-y-ejecutar-las-ventanas}

Script para Bash 4 o superior:

\begin{Shaded}
\begin{Highlighting}[]
\CommentTok{\#!/usr/bin/env bash}
\BuiltInTok{set} \AttributeTok{{-}euo}\NormalTok{ pipefail}

\VariableTok{GMX\_COMMAND}\OperatorTok{=}\StringTok{"}\VariableTok{$\{GMX\_COMMAND}\OperatorTok{:{-}}\NormalTok{gmx}\VariableTok{\}}\StringTok{"}
\VariableTok{START\_STRUCTURE}\OperatorTok{=}\StringTok{"02\_equilibracion/npt.gro"}
\VariableTok{TOPOLOGY}\OperatorTok{=}\StringTok{"topol.top"}

\ControlFlowTok{for}\NormalTok{ STATE }\KeywordTok{in} \VariableTok{$(}\FunctionTok{seq}\NormalTok{ 0 18}\VariableTok{)}
\ControlFlowTok{do}
    \VariableTok{DIRECTORY}\OperatorTok{=}\VariableTok{$(}\BuiltInTok{printf} \StringTok{"03\_fep/lambda\_\%02d"} \StringTok{"}\VariableTok{$\{STATE\}}\StringTok{"}\VariableTok{)}
    \FunctionTok{mkdir} \AttributeTok{{-}p} \StringTok{"}\VariableTok{$\{DIRECTORY\}}\StringTok{"}

    \FunctionTok{sed} \StringTok{"s/LAMBDA\_STATE/}\VariableTok{$\{STATE\}}\StringTok{/g"} \DataTypeTok{\textbackslash{}}
\NormalTok{        fep\_equilibration.template.mdp }\DataTypeTok{\textbackslash{}}
        \OperatorTok{\textgreater{}} \StringTok{"}\VariableTok{$\{DIRECTORY\}}\StringTok{/equilibration.mdp"}

    \FunctionTok{sed} \StringTok{"s/LAMBDA\_STATE/}\VariableTok{$\{STATE\}}\StringTok{/g"} \DataTypeTok{\textbackslash{}}
\NormalTok{        fep\_production.template.mdp }\DataTypeTok{\textbackslash{}}
        \OperatorTok{\textgreater{}} \StringTok{"}\VariableTok{$\{DIRECTORY\}}\StringTok{/production.mdp"}

    \StringTok{"}\VariableTok{$\{GMX\_COMMAND\}}\StringTok{"}\NormalTok{ grompp }\DataTypeTok{\textbackslash{}}
        \AttributeTok{{-}f} \StringTok{"}\VariableTok{$\{DIRECTORY\}}\StringTok{/equilibration.mdp"} \DataTypeTok{\textbackslash{}}
        \AttributeTok{{-}c} \StringTok{"}\VariableTok{$\{START\_STRUCTURE\}}\StringTok{"} \DataTypeTok{\textbackslash{}}
        \AttributeTok{{-}p} \StringTok{"}\VariableTok{$\{TOPOLOGY\}}\StringTok{"} \DataTypeTok{\textbackslash{}}
        \AttributeTok{{-}o} \StringTok{"}\VariableTok{$\{DIRECTORY\}}\StringTok{/equilibration.tpr"}

    \StringTok{"}\VariableTok{$\{GMX\_COMMAND\}}\StringTok{"}\NormalTok{ mdrun }\DataTypeTok{\textbackslash{}}
        \AttributeTok{{-}deffnm} \StringTok{"}\VariableTok{$\{DIRECTORY\}}\StringTok{/equilibration"} \DataTypeTok{\textbackslash{}}
        \AttributeTok{{-}v}

    \StringTok{"}\VariableTok{$\{GMX\_COMMAND\}}\StringTok{"}\NormalTok{ grompp }\DataTypeTok{\textbackslash{}}
        \AttributeTok{{-}f} \StringTok{"}\VariableTok{$\{DIRECTORY\}}\StringTok{/production.mdp"} \DataTypeTok{\textbackslash{}}
        \AttributeTok{{-}c} \StringTok{"}\VariableTok{$\{DIRECTORY\}}\StringTok{/equilibration.gro"} \DataTypeTok{\textbackslash{}}
        \AttributeTok{{-}t} \StringTok{"}\VariableTok{$\{DIRECTORY\}}\StringTok{/equilibration.cpt"} \DataTypeTok{\textbackslash{}}
        \AttributeTok{{-}p} \StringTok{"}\VariableTok{$\{TOPOLOGY\}}\StringTok{"} \DataTypeTok{\textbackslash{}}
        \AttributeTok{{-}o} \StringTok{"}\VariableTok{$\{DIRECTORY\}}\StringTok{/production.tpr"}

    \StringTok{"}\VariableTok{$\{GMX\_COMMAND\}}\StringTok{"}\NormalTok{ mdrun }\DataTypeTok{\textbackslash{}}
        \AttributeTok{{-}deffnm} \StringTok{"}\VariableTok{$\{DIRECTORY\}}\StringTok{/production"} \DataTypeTok{\textbackslash{}}
        \AttributeTok{{-}dhdl} \StringTok{"}\VariableTok{$\{DIRECTORY\}}\StringTok{/dhdl.xvg"} \DataTypeTok{\textbackslash{}}
        \AttributeTok{{-}v}
\ControlFlowTok{done}
\end{Highlighting}
\end{Shaded}

\begin{Shaded}
\begin{Highlighting}[]
\FunctionTok{chmod}\NormalTok{ +x run\_fep.sh}
\FunctionTok{bash} \AttributeTok{{-}n}\NormalTok{ run\_fep.sh}
\ExtensionTok{./run\_fep.sh}
\end{Highlighting}
\end{Shaded}

Las ventanas son independientes después de generar sus TPR y pueden
ejecutarse como un array de trabajos. No use el mismo directorio o
nombre de salida para ventanas diferentes.

\subsubsection{14. Ejecución en un array
HPC}\label{ejecuciuxf3n-en-un-array-hpc}

Ejemplo mínimo basado en \textbf{SLURM\_ARRAY\_TASK\_ID}:

\begin{Shaded}
\begin{Highlighting}[]
\VariableTok{STATE}\OperatorTok{=}\StringTok{"}\VariableTok{$\{SLURM\_ARRAY\_TASK\_ID\}}\StringTok{"}
\VariableTok{DIRECTORY}\OperatorTok{=}\VariableTok{$(}\BuiltInTok{printf} \StringTok{"03\_fep/lambda\_\%02d"} \StringTok{"}\VariableTok{$\{STATE\}}\StringTok{"}\VariableTok{)}

\ExtensionTok{gmx}\NormalTok{ mdrun }\DataTypeTok{\textbackslash{}}
    \AttributeTok{{-}deffnm} \StringTok{"}\VariableTok{$\{DIRECTORY\}}\StringTok{/production"} \DataTypeTok{\textbackslash{}}
    \AttributeTok{{-}dhdl} \StringTok{"}\VariableTok{$\{DIRECTORY\}}\StringTok{/dhdl.xvg"} \DataTypeTok{\textbackslash{}}
    \AttributeTok{{-}v}
\end{Highlighting}
\end{Shaded}

Los TPR deben haberse generado previamente. Los recursos, GPU, partición
y tiempo dependen del clúster.

Registre:

\begin{Shaded}
\begin{Highlighting}[]
\NormalTok{estado | coul{-}lambda | vdw{-}lambda | semilla | duración | estado del trabajo}
\end{Highlighting}
\end{Shaded}

\subsubsection{15. Continuar una ventana}\label{continuar-una-ventana}

\begin{Shaded}
\begin{Highlighting}[]
\ExtensionTok{gmx}\NormalTok{ mdrun }\DataTypeTok{\textbackslash{}}
    \AttributeTok{{-}deffnm}\NormalTok{ 03\_fep/lambda\_07/production }\DataTypeTok{\textbackslash{}}
    \AttributeTok{{-}cpi}\NormalTok{ 03\_fep/lambda\_07/production.cpt }\DataTypeTok{\textbackslash{}}
    \AttributeTok{{-}append} \DataTypeTok{\textbackslash{}}
    \AttributeTok{{-}dhdl}\NormalTok{ 03\_fep/lambda\_07/dhdl.xvg }\DataTypeTok{\textbackslash{}}
    \AttributeTok{{-}v}
\end{Highlighting}
\end{Shaded}

No regenere el TPR si solo continúa el mismo Hamiltoniano. Si cambia
MDP, duración o λ, documente el nuevo segmento.

\subsubsection{16. Análisis con BAR}\label{anuxe1lisis-con-bar}

Los directorios llevan ceros iniciales para que el glob preserve el
orden:

\begin{Shaded}
\begin{Highlighting}[]
\FunctionTok{mkdir} \AttributeTok{{-}p}\NormalTok{ 04\_analisis}

\ExtensionTok{gmx}\NormalTok{ bar }\DataTypeTok{\textbackslash{}}
    \AttributeTok{{-}b}\NormalTok{ 500 }\DataTypeTok{\textbackslash{}}
    \AttributeTok{{-}f}\NormalTok{ 03\_fep/lambda\_}\PreprocessorTok{*}\NormalTok{/dhdl.xvg }\DataTypeTok{\textbackslash{}}
    \AttributeTok{{-}o}\NormalTok{ 04\_analisis/bar\_deltaG.xvg }\DataTypeTok{\textbackslash{}}
    \AttributeTok{{-}oi}\NormalTok{ 04\_analisis/bar\_integral.xvg }\DataTypeTok{\textbackslash{}}
    \AttributeTok{{-}oh}\NormalTok{ 04\_analisis/bar\_histograms.xvg }\DataTypeTok{\textbackslash{}}
    \AttributeTok{{-}prec}\NormalTok{ 3}
\end{Highlighting}
\end{Shaded}

\textbf{-b 500} descarta 500 ps de cada archivo. Determine el descarte
mediante la relajación real de las ventanas, no por copia del ejemplo.

Con la convención usada:

\[
\Delta G_{\mathrm{solv}}
=
-\Delta G_{\mathrm{BAR}}
\]

\subsubsection{17. Diagnóstico de BAR}\label{diagnuxf3stico-de-bar}

Revise por cada par:

\begin{itemize}
\tightlist
\item
  \(\Delta G\) parcial;
\item
  incertidumbre;
\item
  entropías relativas \textbf{s\_A} y \textbf{s\_B};
\item
  desviación esperada por muestra;
\item
  histogramas;
\item
  consistencia entre ventanas.
\end{itemize}

Las entropías relativas son indicadores de distancia entre
distribuciones; valores grandes sugieren peor superposición.

El error de BAR no incluye automáticamente error del campo de fuerza,
modelo de agua, especie química, estado estándar, tamaño finito o
muestreo conformacional ausente. Precisión estadística no implica
exactitud química.

\subsubsection{18. Convergencia temporal y
réplicas}\label{convergencia-temporal-y-ruxe9plicas}

Repita BAR con distintos descartes:

\begin{Shaded}
\begin{Highlighting}[]
\ExtensionTok{gmx}\NormalTok{ bar }\AttributeTok{{-}b}\NormalTok{ 250  }\AttributeTok{{-}f}\NormalTok{ 03\_fep/lambda\_}\PreprocessorTok{*}\NormalTok{/dhdl.xvg}
\ExtensionTok{gmx}\NormalTok{ bar }\AttributeTok{{-}b}\NormalTok{ 500  }\AttributeTok{{-}f}\NormalTok{ 03\_fep/lambda\_}\PreprocessorTok{*}\NormalTok{/dhdl.xvg}
\ExtensionTok{gmx}\NormalTok{ bar }\AttributeTok{{-}b}\NormalTok{ 1000 }\AttributeTok{{-}f}\NormalTok{ 03\_fep/lambda\_}\PreprocessorTok{*}\NormalTok{/dhdl.xvg}
\end{Highlighting}
\end{Shaded}

Una estimación defendible debe mostrar:

\begin{itemize}
\tightlist
\item
  ausencia de deriva;
\item
  superposición bilateral;
\item
  contribuciones sin saltos dominantes;
\item
  compatibilidad entre mitades temporales;
\item
  compatibilidad entre réplicas;
\item
  menor incertidumbre al aumentar datos efectivos.
\end{itemize}

Las filas consecutivas están correlacionadas. El número de filas no es
el número de muestras independientes.

Use semillas distintas y vuelva a equilibrar las ventanas para cada
réplica. Trayectorias con el mismo checkpoint y velocidades no son
independientes.

\subsubsection{19. Refinar la grilla}\label{refinar-la-grilla}

Agregue estados entre \(\lambda_i\) y \(\lambda_{i+1}\) si los
histogramas apenas se superponen, el error parcial domina,
\textbf{s\_A/s\_B} aumentan abruptamente o el solvente entra y sale de
forma discontinua.

No agregue ventanas donde el solapamiento ya es alto si el cuello de
botella está en otra región. Redistribuya el cálculo.

\subsubsection{20. Unidades de energía}\label{unidades-de-energuxeda}

GROMACS informa energías molares en kJ·mol⁻¹.

\[
1\ \mathrm{kcal\,mol^{-1}}
=
4.184\ \mathrm{kJ\,mol^{-1}}
\]

\[
1\ \mathrm{eV\ por\ molécula}
=
96.485332\ \mathrm{kJ\,mol^{-1}}
\]

\[
1\ E_h
=
2625.50\ \mathrm{kJ\,mol^{-1}}
\]

Por tanto:

\[
-0.151256\ \mathrm{eV}
\times
96.485332
=
-14.59399\ \mathrm{kJ\,mol^{-1}}
\]

El valor anterior, −14.5934028 kJ·mol⁻¹, difiere ligeramente por el
factor de conversión empleado. Esa diferencia es despreciable frente a
la incertidumbre física, pero debe usarse una constante consistente.

\subsubsection{21. Energía térmica}\label{energuxeda-tuxe9rmica}

A 298.15 K:

\[
RT=
(8.314462618\times10^{-3})
(298.15)
=
2.47896\ \mathrm{kJ\,mol^{-1}}
\]

Una unidad \(k_\mathrm{B}T\) por partícula equivale numéricamente a una
unidad \(RT\) por mol.

\[
-14.594\ \mathrm{kJ\,mol^{-1}}
\approx
-5.89\ k_\mathrm{B}T
\]

No multiplique por el número de Avogadro una energía ya expresada por
mol.

\subsubsection{22. Tiempo y λ}\label{tiempo-y-ux3bb}

\[
1000\ \mathrm{ps}=1\ \mathrm{ns}
\]

Con:

\begin{Shaded}
\begin{Highlighting}[]
\DataTypeTok{dt     }\OtherTok{=}\StringTok{ }\FloatTok{0.002}
\DataTypeTok{nsteps }\OtherTok{=}\StringTok{ }\DecValTok{2500000}
\end{Highlighting}
\end{Shaded}

\[
t=0.002\ \mathrm{ps}\times2500000
=5000\ \mathrm{ps}
=5\ \mathrm{ns}
\]

λ es adimensional; \textbf{init-lambda-state} es un índice entero.

\subsubsection{23. Estados estándar}\label{estados-estuxe1ndar}

Una referencia experimental puede usar gas ideal a 1 atm o 1 bar,
solución a 1 mol·L⁻¹, dilución infinita o una convención de ley de
Henry.

La conversión gas--solución contiene:

\[
\Delta G^\circ_{\mathrm{corr}}
=
RT\ln
\left(
\frac{C^\circ RT}{p^\circ}
\right)
\]

A 298.15 K, la magnitud entre 1 atm y 1 mol·L⁻¹ es aproximadamente 7.9
kJ·mol⁻¹. El signo depende de la dirección y de la definición original.

No aplique esta corrección sin identificar las convenciones experimental
y computacional.

\subsubsection{24. Tamaño finito y
carga}\label{tamauxf1o-finito-y-carga}

Para solutos con carga neta, PME y una caja periódica con fondo
neutralizante introducen correcciones dependientes de la carga, la caja,
la constante dieléctrica y la convención electrostática.

El formaldehído aquí es neutro y no requiere la corrección principal de
carga neta. Persisten posibles efectos de tamaño y dipolo.

No compare transformaciones que cambian la carga neta sin un protocolo
específico.

\subsubsection{25. Limitación química del
formaldehído}\label{limitaciuxf3n-quuxedmica-del-formaldehuxeddo}

En agua:

\[
\mathrm{H_2CO + H_2O
\rightleftharpoons
CH_2(OH)_2}
\]

El producto es metanodiol o metilenglicol. Según la concentración
también pueden existir oligómeros.

Un campo de fuerza clásico de topología fija no rompe el enlace C=O, no
forma enlaces C--O, no transfiere protones y no convierte formaldehído
en metanodiol.

Este cálculo estima:

\[
\mathrm{H_2CO(g)}
\rightarrow
\mathrm{H_2CO(aq)}
\]

para formaldehído mantenido artificialmente como tal. No estima
directamente:

\[
\mathrm{H_2CO(g)+H_2O(l)}
\rightarrow
\mathrm{CH_2(OH)_2(aq)}
\]

ni la solubilidad total de formaldehído.

\subsubsection{26. Ciclo químico
correcto}\label{ciclo-quuxedmico-correcto}

El proceso efectivo puede descomponerse:

\[
\Delta G^\circ_{\mathrm{global}}
=
\Delta G^\circ_{\mathrm{solv}}(\mathrm{H_2CO})
+
\Delta G^\circ_{\mathrm{hidratación,aq}}
\]

Alternativas:

\begin{itemize}
\tightlist
\item
  calcular solo solvatación física de \(\mathrm{H_2CO}\) y declararlo;
\item
  parametrizar metanodiol como especie diferente;
\item
  calcular la hidratación covalente mediante QM/MM o estructura
  electrónica;
\item
  combinar solvatación y reacción en un ciclo termodinámico;
\item
  usar etanol o metanol como ejemplo pedagógico no reactivo.
\end{itemize}

El tutorial oficial de GROMACS usa etanol, una elección más clara para
enseñar el procedimiento general. El formaldehído resulta útil para
mostrar que la especie química debe definirse antes del cálculo.

\subsubsection{27. Comparación con el valor atribuido a
NIST}\label{comparaciuxf3n-con-el-valor-atribuido-a-nist}

El texto anterior cita:

\[
-0.151256\ \mathrm{eV}
=
-14.59399\ \mathrm{kJ\,mol^{-1}}
\]

Antes de compararlo con BAR deben coincidir:

\begin{itemize}
\tightlist
\item
  especie;
\item
  sentido solvatación/desolvatación;
\item
  temperatura;
\item
  estado estándar;
\item
  solvente;
\item
  definición de dilución;
\item
  inclusión de hidratación covalente;
\item
  convención de la base.
\end{itemize}

La página general del proyecto NIST no reemplaza el registro específico
con sus metadatos. Un acuerdo numérico accidental puede ocultar signos
opuestos o procesos químicos distintos.

\subsubsection{28. Interpretación de
resultados}\label{interpretaciuxf3n-de-resultados}

Un resultado como:

\[
-22.94\pm1.04\ \mathrm{kJ\,mol^{-1}}
\]

solo es interpretable si se informa:

\begin{itemize}
\tightlist
\item
  extremos de λ;
\item
  signo aplicado a BAR;
\item
  tiempo descartado;
\item
  duración por ventana;
\item
  grilla;
\item
  réplicas;
\item
  método de incertidumbre;
\item
  campo de fuerza;
\item
  agua;
\item
  temperatura;
\item
  correcciones;
\item
  especie química.
\end{itemize}

El ± suele representar incertidumbre estadística condicionada al
muestreo. No incluye errores de parametrización o química ausente.

\subsubsection{29. Controles adicionales}\label{controles-adicionales}

\begin{Shaded}
\begin{Highlighting}[]
\ExtensionTok{gmx}\NormalTok{ trjconv }\DataTypeTok{\textbackslash{}}
    \AttributeTok{{-}s}\NormalTok{ 03\_fep/lambda\_00/production.tpr }\DataTypeTok{\textbackslash{}}
    \AttributeTok{{-}f}\NormalTok{ 03\_fep/lambda\_00/production.xtc }\DataTypeTok{\textbackslash{}}
    \AttributeTok{{-}o}\NormalTok{ 04\_analisis/lambda\_00\_snapshot.gro }\DataTypeTok{\textbackslash{}}
    \AttributeTok{{-}dump}\NormalTok{ 5000}

\ExtensionTok{gmx}\NormalTok{ trjconv }\DataTypeTok{\textbackslash{}}
    \AttributeTok{{-}s}\NormalTok{ 03\_fep/lambda\_18/production.tpr }\DataTypeTok{\textbackslash{}}
    \AttributeTok{{-}f}\NormalTok{ 03\_fep/lambda\_18/production.xtc }\DataTypeTok{\textbackslash{}}
    \AttributeTok{{-}o}\NormalTok{ 04\_analisis/lambda\_18\_snapshot.gro }\DataTypeTok{\textbackslash{}}
    \AttributeTok{{-}dump}\NormalTok{ 5000}
\end{Highlighting}
\end{Shaded}

Controle temperatura, volumen, errores LINCS, integridad del soluto,
respuesta del solvente, ausencia de singularidades, producción de
archivos de diferencias de energía y estado λ informado en el log.

En el extremo desacoplado el solvente puede ocupar el espacio del
soluto. Es un comportamiento esperado con soft-core.

\subsubsection{30. Errores frecuentes}\label{errores-frecuentes-1}

\begin{itemize}
\tightlist
\item
  No definir los estados extremos.
\item
  Confundir solvatación con desolvatación.
\item
  Usar \textbf{vdwq} en vez de la sintaxis actual \textbf{vdw-q}.
\item
  Eliminar Lennard-Jones antes que las cargas.
\item
  Usar cinco λ uniformes y asumir convergencia.
\item
  Desacoplar Coulomb y Lennard-Jones simultáneamente sin validación.
\item
  No equilibrar cada ventana.
\item
  Copiar arbitrariamente el descarte \textbf{-b}.
\item
  Reutilizar un directorio para varias ventanas.
\item
  Usar archivos sin orden definido.
\item
  Tratar muestras correlacionadas como independientes.
\item
  Informar solamente el error de BAR.
\item
  Confundir eV por molécula con kJ·mol⁻¹.
\item
  Multiplicar otra vez por \(N_A\).
\item
  Ignorar el estado estándar.
\item
  Perturbar varias copias del mismo \textbf{moleculetype}.
\item
  Cambiar carga neta sin correcciones.
\item
  Comparar formaldehído molecular con formaldehído total acuoso.
\item
  Aceptar una minimización que no alcanzó su criterio.
\item
  Usar \textbf{-maxwarn} para producir el TPR.
\end{itemize}

\subsubsection{31. Reproducibilidad
mínima}\label{reproducibilidad-muxednima}

Informe:

\begin{itemize}
\tightlist
\item
  versión exacta de GROMACS;
\item
  origen y versión de parámetros;
\item
  modelo de agua;
\item
  topología molecular;
\item
  caja y número de aguas;
\item
  temperatura y presión;
\item
  integrador y paso;
\item
  lista completa de λ;
\item
  parámetros soft-core;
\item
  equilibración y producción;
\item
  frecuencia de diferencias de energía;
\item
  método de análisis y descarte;
\item
  réplicas y semillas;
\item
  correcciones;
\item
  signo de \(\Delta G\);
\item
  especie química calculada.
\end{itemize}

\subsubsection{32. Criterios de
aceptación}\label{criterios-de-aceptaciuxf3n}

El cálculo es técnicamente defendible cuando:

\begin{itemize}
\tightlist
\item
  los extremos están definidos;
\item
  la topología es químicamente correcta;
\item
  todas las ventanas terminaron;
\item
  existe superposición entre vecinos;
\item
  las contribuciones parciales son estables;
\item
  el resultado converge temporalmente;
\item
  las réplicas son compatibles;
\item
  unidades y signo son inequívocos;
\item
  se aplicaron las correcciones necesarias;
\item
  el dato experimental usa una convención comparable;
\item
  la especie simulada coincide con la experimental.
\end{itemize}

Para formaldehído acuoso en equilibrio, el último criterio no se cumple
si solo se modela \(\mathrm{H_2CO}\) mediante una topología clásica
fija.

\subsubsection{Fuentes}\label{fuentes-5}

\begin{enumerate}
\def\labelenumi{\arabic{enumi}.}
\tightlist
\item
  \href{https://tutorials.gromacs.org/docs/free-energy-of-solvation.html}{Tutorial
  oficial de energía libre de solvatación de GROMACS}
\item
  \href{https://manual.gromacs.org/current/user-guide/mdp-options.html}{Opciones
  MDP de energía libre en GROMACS 2026.3}
\item
  \href{https://manual.gromacs.org/current/onlinehelp/gmx-bar.html}{Referencia
  de gmx bar}
\item
  \href{https://pmc.ncbi.nlm.nih.gov/articles/PMC8479811/}{Perturbation
  Free-Energy Toolkit}
\item
  \href{https://link.springer.com/article/10.1007/s10822-015-9840-9}{Guidelines
  for the analysis of free energy calculations}
\item
  \href{https://link.springer.com/article/10.1007/s10822-015-9873-0}{Python
  tool for relative free-energy calculations in GROMACS}
\item
  \href{http://www.mdtutorials.com/gmx/2018/free_energy/index.html}{Tutorial
  histórico de FEP con GROMACS}
\item
  \href{https://pmc.ncbi.nlm.nih.gov/articles/PMC12163693/}{Revisión
  adicional proporcionada}
\item
  \href{https://www.nist.gov/programs-projects/solvation-free-energies}{Proyecto
  NIST sobre energías libres de solvatación}
\end{enumerate}

\subsection{\texorpdfstring{Energía de unión mediante
\texttt{gmx\_MMPBSA}}{Energía de unión mediante gmx\_MMPBSA}}\label{energuxeda-de-uniuxf3n-mediante-gmx_mmpbsa}

\subsubsection{1. Qué calcula este
método}\label{quuxe9-calcula-este-muxe9todo}

\textbf{gmx\_MMPBSA} es la implementación actual para realizar cálculos
de energía libre de unión de estados finales usando trayectorias y
topologías de GROMACS. Está basado en \textbf{MMPBSA.py} de AmberTools y
reemplaza el flujo antiguo basado en \textbf{g\_mmpbsa}, un archivo
\textbf{pbsa.mdp}, APBS ejecutado externamente y el script
\textbf{pbsa.py}.

MM/PBSA y MM/GBSA no son FEP. No contienen estados λ, no hacen
desaparecer el ligando y no necesitan umbrella sampling para impedir que
el ligando se aleje durante un desacoplamiento. Analizan configuraciones
ya muestreadas de los estados unido y libre mediante un modelo de
solvente implícito.

El esquema básico es:

\[
\Delta G_{\mathrm{bind}}
=
G_{\mathrm{complejo}}
-
G_{\mathrm{receptor}}
-
G_{\mathrm{ligando}}
\]

Para cada especie:

\[
G=
E_{\mathrm{MM}}
+
G_{\mathrm{solv}}
-
TS
\]

La energía de mecánica molecular puede escribirse:

\[
E_{\mathrm{MM}}
=
E_{\mathrm{bonded}}
+
E_{\mathrm{vdW}}
+
E_{\mathrm{elec}}
\]

y la contribución de solvatación:

\[
G_{\mathrm{solv}}
=
G_{\mathrm{polar}}
+
G_{\mathrm{nonpolar}}
\]

Por tanto:

\[
\Delta G_{\mathrm{bind}}
=
\Delta E_{\mathrm{MM}}
+
\Delta G_{\mathrm{polar}}
+
\Delta G_{\mathrm{nonpolar}}
-
T\Delta S
\]

Si no se calcula entropía, el resultado contiene solo:

\[
\Delta G_{\mathrm{bind}}^{*}
=
\Delta E_{\mathrm{MM}}
+
\Delta G_{\mathrm{solv}}
\]

El asterisco recuerda que no es una energía libre absoluta completa.
Muchos informes la llaman ``binding free energy'', pero físicamente es
una estimación end-state sin el término entrópico explícito.

\subsubsection{2. PB frente a GB}\label{pb-frente-a-gb}

MM/PBSA calcula la contribución polar resolviendo numéricamente la
ecuación de Poisson--Boltzmann. MM/GBSA utiliza una aproximación
Generalized Born.

{\def\LTcaptype{none} % do not increment counter
\begin{longtable}[]{@{}
  >{\raggedright\arraybackslash}p{(\linewidth - 4\tabcolsep) * \real{0.0486}}
  >{\raggedright\arraybackslash}p{(\linewidth - 4\tabcolsep) * \real{0.4167}}
  >{\raggedright\arraybackslash}p{(\linewidth - 4\tabcolsep) * \real{0.5347}}@{}}
\toprule\noalign{}
\begin{minipage}[b]{\linewidth}\raggedright
Método
\end{minipage} & \begin{minipage}[b]{\linewidth}\raggedright
Ventaja
\end{minipage} & \begin{minipage}[b]{\linewidth}\raggedright
Limitación
\end{minipage} \\
\midrule\noalign{}
\endhead
\bottomrule\noalign{}
\endlastfoot
MM/GBSA & Más rápido; útil para explorar protocolos y series. & Mayor
dependencia del modelo GB y del conjunto de radios. \\
MM/PBSA & Tratamiento continuo electrostático más explícito. & Más
costoso y sensible a malla, radios, dieléctricos y convergencia
numérica. \\
3D-RISM & Describe estructura promedio del solvente con mayor detalle. &
Coste y preparación superiores; escalamiento paralelo limitado. \\
\end{longtable}
}

No existe una regla universal por la cual PB sea siempre más exacto que
GB. La comparación debe validarse para la familia química y el objetivo.

\subsubsection{3. Qué puede y qué no puede
inferirse}\label{quuxe9-puede-y-quuxe9-no-puede-inferirse}

Usos razonables:

\begin{itemize}
\tightlist
\item
  comparar poses del mismo ligando;
\item
  ordenar ligandos estrechamente relacionados;
\item
  comparar mutantes con un protocolo idéntico;
\item
  detectar residuos que contribuyen de forma persistente;
\item
  estudiar sensibilidad a parámetros del continuo;
\item
  complementar estabilidad estructural y contactos.
\end{itemize}

No debe interpretarse automáticamente como:

\begin{itemize}
\tightlist
\item
  afinidad experimental absoluta;
\item
  constante de disociación exacta;
\item
  mecanismo de unión;
\item
  tasa de asociación o disociación;
\item
  sustituto de FEP para transformaciones pequeñas;
\item
  prueba causal de que un residuo ``produce'' la unión;
\item
  evidencia de convergencia de la dinámica.
\end{itemize}

La aproximación omite o simplifica agua explícita, reorganización del
solvente, cambios de protonación, polarización, estados alternativos y,
con frecuencia, entropía conformacional.

\subsubsection{4. Protocolo de trayectoria
única}\label{protocolo-de-trayectoria-uxfanica}

En el protocolo de trayectoria única, \textbf{single trajectory, ST},
receptor y ligando se extraen de cada instantánea del complejo:

\begin{Shaded}
\begin{Highlighting}[]
\NormalTok{trayectoria del complejo}
\NormalTok{        ├── complejo}
\NormalTok{        ├── receptor extraído}
\NormalTok{        └── ligando extraído}
\end{Highlighting}
\end{Shaded}

Esto mantiene correspondencia conformacional entre los tres términos y
favorece la cancelación:

\[
\Delta E_{\mathrm{bonded}}\approx0
\]

Ventajas:

\begin{itemize}
\tightlist
\item
  menor ruido;
\item
  menor costo;
\item
  correspondencia exacta entre marcos;
\item
  buena reproducibilidad operativa.
\end{itemize}

Supuesto fuerte:

\begin{itemize}
\tightlist
\item
  receptor y ligando libres adoptan las mismas conformaciones que en el
  complejo.
\end{itemize}

ST no cuantifica adecuadamente una reorganización grande inducida por
unión.

\subsubsection{5. Protocolo de trayectorias
múltiples}\label{protocolo-de-trayectorias-muxfaltiples}

En el protocolo \textbf{multiple trajectory, MT}, complejo, receptor y
ligando provienen de simulaciones separadas:

\[
\Delta G_{\mathrm{bind}}
=
\langle G_C\rangle_C
-
\langle G_R\rangle_R
-
\langle G_L\rangle_L
\]

Puede incorporar reorganización conformacional, pero aumenta mucho la
varianza porque se pierde cancelación entre configuraciones
correlacionadas.

Use MT solo cuando:

\begin{itemize}
\tightlist
\item
  existan trayectorias libres suficientemente muestreadas;
\item
  la reorganización sea parte de la pregunta;
\item
  se disponga de varias réplicas;
\item
  se evalúe la convergencia de cada estado por separado.
\end{itemize}

Un resultado MT más ``completo'' puede ser menos preciso que ST si el
muestreo es insuficiente.

\subsubsection{6. Requisitos actuales}\label{requisitos-actuales}

La documentación de desarrollo consultada en septiembre de 2026
recomienda rangos probados:

\begin{itemize}
\tightlist
\item
  Python 3.11 o 3.12;
\item
  AmberTools desde 24.8 y menor que 27;
\item
  GROMACS desde 2022 y menor que 2027;
\item
  ParmEd 4.2.2 o posterior dentro de la serie 4;
\item
  MPI opcional para paralelización.
\end{itemize}

GROMACS 2026 está dentro del rango probado. Esto no garantiza que
cualquier topología se convierta correctamente.

Archivos necesarios para ST:

\begin{Shaded}
\begin{Highlighting}[]
\NormalTok{md.tpr              estructura y masas del complejo}
\NormalTok{md\_fit.xtc          trayectoria corregida}
\NormalTok{index.ndx           grupos receptor y ligando}
\NormalTok{topol.top           topología completa de GROMACS}
\NormalTok{reference.pdb       referencia con cadenas y numeración, recomendada}
\NormalTok{mmpbsa.in           opciones del cálculo}
\end{Highlighting}
\end{Shaded}

La topología del ligando debe estar incluida en \textbf{topol.top}. La
ruta actual no reconstruye ligandos desde un PDB sin parámetros.

\subsubsection{7. Instalación
reproducible}\label{instalaciuxf3n-reproducible-1}

Conda o Mamba, Python 3.12:

\begin{Shaded}
\begin{Highlighting}[]
\ExtensionTok{conda}\NormalTok{ create }\AttributeTok{{-}n}\NormalTok{ gmxMMPBSA python=3.12 }\AttributeTok{{-}y}
\ExtensionTok{conda}\NormalTok{ activate gmxMMPBSA}

\ExtensionTok{conda}\NormalTok{ install }\AttributeTok{{-}c}\NormalTok{ conda{-}forge }\DataTypeTok{\textbackslash{}}
    \StringTok{"mpi4py\textgreater{}=4.0.1,\textless{}5"} \DataTypeTok{\textbackslash{}}
    \StringTok{"ambertools\textgreater{}=24.8,\textless{}27"} \DataTypeTok{\textbackslash{}}
    \AttributeTok{{-}y}

\ExtensionTok{python} \AttributeTok{{-}m}\NormalTok{ pip install gmx\_MMPBSA}
\ExtensionTok{python} \AttributeTok{{-}m}\NormalTok{ pip check}
\end{Highlighting}
\end{Shaded}

Si GROMACS 2026 ya está instalado en el sistema y disponible en
\textbf{PATH}, no instale otra copia sin necesidad. Verifique:

\begin{Shaded}
\begin{Highlighting}[]
\ExtensionTok{gmx} \AttributeTok{{-}{-}version}
\ExtensionTok{gmx\_MMPBSA} \AttributeTok{{-}{-}version}
\ExtensionTok{gmx\_MMPBSA} \AttributeTok{{-}h}
\ExtensionTok{python} \AttributeTok{{-}m}\NormalTok{ pip check}
\end{Highlighting}
\end{Shaded}

Para una instalación completamente aislada:

\begin{Shaded}
\begin{Highlighting}[]
\ExtensionTok{conda}\NormalTok{ install }\AttributeTok{{-}c}\NormalTok{ conda{-}forge }\StringTok{"gromacs\textgreater{}=2022,\textless{}2027"}\NormalTok{ pocl }\AttributeTok{{-}y}
\end{Highlighting}
\end{Shaded}

La interfaz gráfica requiere PyQt6:

\begin{Shaded}
\begin{Highlighting}[]
\ExtensionTok{conda}\NormalTok{ install }\AttributeTok{{-}c}\NormalTok{ conda{-}forge pyqt6 }\AttributeTok{{-}y}
\end{Highlighting}
\end{Shaded}

No es necesaria en un nodo HPC sin entorno gráfico.

\subsubsection{8. Probar la instalación}\label{probar-la-instalaciuxf3n}

\begin{Shaded}
\begin{Highlighting}[]
\ExtensionTok{gmx\_MMPBSA\_test} \AttributeTok{{-}h}
\end{Highlighting}
\end{Shaded}

Ejecute primero una prueba pequeña proporcionada por el paquete. Esto
permite separar problemas de instalación de problemas propios de la
topología.

Compruebe además:

\begin{Shaded}
\begin{Highlighting}[]
\BuiltInTok{command} \AttributeTok{{-}v}\NormalTok{ gmx}
\BuiltInTok{command} \AttributeTok{{-}v}\NormalTok{ gmx\_MMPBSA}
\BuiltInTok{command} \AttributeTok{{-}v}\NormalTok{ ante{-}MMPBSA.py}
\BuiltInTok{echo} \StringTok{"}\VariableTok{$\{AMBERHOME}\OperatorTok{:{-}}\NormalTok{AMBERHOME no definido}\VariableTok{\}}\StringTok{"}
\end{Highlighting}
\end{Shaded}

La instalación debe encontrar AmberTools y GROMACS dentro del entorno
activo o mediante las opciones de configuración correspondientes.

\subsubsection{9. Preparación de los
grupos}\label{preparaciuxf3n-de-los-grupos}

Cree un índice que contenga grupos separados y sin solapamiento:

\begin{itemize}
\tightlist
\item
  receptor;
\item
  ligando;
\item
  complejo receptor--ligando.
\end{itemize}

Forma interactiva:

\begin{Shaded}
\begin{Highlighting}[]
\ExtensionTok{gmx}\NormalTok{ make\_ndx }\DataTypeTok{\textbackslash{}}
    \AttributeTok{{-}f}\NormalTok{ md.tpr }\DataTypeTok{\textbackslash{}}
    \AttributeTok{{-}o}\NormalTok{ index.ndx}
\end{Highlighting}
\end{Shaded}

Ejemplo dentro de \textbf{make\_ndx}, si el ligando se llama LIG:

\begin{Shaded}
\begin{Highlighting}[]
\NormalTok{r LIG}
\NormalTok{Protein | r LIG}
\NormalTok{name NOMBRE\_GRUPO\_RECEPTOR Receptor}
\NormalTok{name NOMBRE\_GRUPO\_LIGANDO Ligando}
\NormalTok{name NOMBRE\_GRUPO\_COMPLEJO Complejo}
\NormalTok{q}
\end{Highlighting}
\end{Shaded}

Los números reales dependen del índice creado. Verifique:

\begin{Shaded}
\begin{Highlighting}[]
\ExtensionTok{gmx}\NormalTok{ make\_ndx }\DataTypeTok{\textbackslash{}}
    \AttributeTok{{-}f}\NormalTok{ md.tpr }\DataTypeTok{\textbackslash{}}
    \AttributeTok{{-}n}\NormalTok{ index.ndx }\DataTypeTok{\textbackslash{}}
    \AttributeTok{{-}o}\NormalTok{ index\_check.ndx}
\end{Highlighting}
\end{Shaded}

No use números copiados de otro sistema. La opción actual \textbf{-cg}
acepta números de grupo basados en cero o nombres.

\subsubsection{10. Corregir la
trayectoria}\label{corregir-la-trayectoria}

La documentación requiere una trayectoria sin problemas de PBC y
ajustada. El ligando debe permanecer junto al receptor.

Primero centre el complejo:

\begin{Shaded}
\begin{Highlighting}[]
\KeywordTok{(}\BuiltInTok{echo}\NormalTok{ Complejo}\KeywordTok{;} \BuiltInTok{echo}\NormalTok{ System}\KeywordTok{)} \KeywordTok{|} \DataTypeTok{\textbackslash{}}
\ExtensionTok{gmx}\NormalTok{ trjconv }\DataTypeTok{\textbackslash{}}
    \AttributeTok{{-}s}\NormalTok{ md.tpr }\DataTypeTok{\textbackslash{}}
    \AttributeTok{{-}f}\NormalTok{ md.xtc }\DataTypeTok{\textbackslash{}}
    \AttributeTok{{-}n}\NormalTok{ index.ndx }\DataTypeTok{\textbackslash{}}
    \AttributeTok{{-}o}\NormalTok{ md\_center.xtc }\DataTypeTok{\textbackslash{}}
    \AttributeTok{{-}pbc}\NormalTok{ mol }\DataTypeTok{\textbackslash{}}
    \AttributeTok{{-}center} \DataTypeTok{\textbackslash{}}
    \AttributeTok{{-}ur}\NormalTok{ compact}
\end{Highlighting}
\end{Shaded}

Luego elimine rotación y traslación:

\begin{Shaded}
\begin{Highlighting}[]
\KeywordTok{(}\BuiltInTok{echo}\NormalTok{ Backbone}\KeywordTok{;} \BuiltInTok{echo}\NormalTok{ System}\KeywordTok{)} \KeywordTok{|} \DataTypeTok{\textbackslash{}}
\ExtensionTok{gmx}\NormalTok{ trjconv }\DataTypeTok{\textbackslash{}}
    \AttributeTok{{-}s}\NormalTok{ md.tpr }\DataTypeTok{\textbackslash{}}
    \AttributeTok{{-}f}\NormalTok{ md\_center.xtc }\DataTypeTok{\textbackslash{}}
    \AttributeTok{{-}n}\NormalTok{ index.ndx }\DataTypeTok{\textbackslash{}}
    \AttributeTok{{-}o}\NormalTok{ md\_fit.xtc }\DataTypeTok{\textbackslash{}}
    \AttributeTok{{-}fit}\NormalTok{ rot+trans}
\end{Highlighting}
\end{Shaded}

Compruebe visualmente que:

\begin{itemize}
\tightlist
\item
  la proteína esté entera;
\item
  el ligando no salte por PBC;
\item
  el orden atómico no cambie;
\item
  no se hayan descartado componentes necesarios;
\item
  los tiempos sean correctos.
\end{itemize}

Para MM/PBSA se elimina el solvente durante la preparación interna. No
genere una trayectoria manualmente desolvatada si todavía no verificó
que la selección mantiene todos los componentes del complejo.

\subsubsection{11. Estructura de
referencia}\label{estructura-de-referencia}

Genere un PDB de referencia con el complejo completo:

\begin{Shaded}
\begin{Highlighting}[]
\BuiltInTok{echo}\NormalTok{ Complejo }\KeywordTok{|} \DataTypeTok{\textbackslash{}}
\ExtensionTok{gmx}\NormalTok{ trjconv }\DataTypeTok{\textbackslash{}}
    \AttributeTok{{-}s}\NormalTok{ md.tpr }\DataTypeTok{\textbackslash{}}
    \AttributeTok{{-}f}\NormalTok{ md\_fit.xtc }\DataTypeTok{\textbackslash{}}
    \AttributeTok{{-}n}\NormalTok{ index.ndx }\DataTypeTok{\textbackslash{}}
    \AttributeTok{{-}o}\NormalTok{ reference.pdb }\DataTypeTok{\textbackslash{}}
    \AttributeTok{{-}dump}\NormalTok{ 0}
\end{Highlighting}
\end{Shaded}

Revise cadenas, numeración, nombres de residuos y átomos. La opción
\textbf{-cr reference.pdb} es recomendada porque evita asignaciones
automáticas incorrectas, especialmente en complejos con varias cadenas.

\subsubsection{12. Generar un archivo de
entrada}\label{generar-un-archivo-de-entrada}

La versión actual puede generar plantillas:

\begin{Shaded}
\begin{Highlighting}[]
\ExtensionTok{gmx\_MMPBSA} \AttributeTok{{-}{-}create\_input}\NormalTok{ gb}
\ExtensionTok{gmx\_MMPBSA} \AttributeTok{{-}{-}create\_input}\NormalTok{ pb}
\ExtensionTok{gmx\_MMPBSA} \AttributeTok{{-}{-}create\_input}\NormalTok{ gb decomp}
\ExtensionTok{gmx\_MMPBSA} \AttributeTok{{-}{-}create\_input}\NormalTok{ gb nmode}
\end{Highlighting}
\end{Shaded}

Revise siempre la plantilla generada por la versión instalada:

\begin{Shaded}
\begin{Highlighting}[]
\ExtensionTok{gmx\_MMPBSA} \AttributeTok{{-}{-}input{-}file{-}help}
\end{Highlighting}
\end{Shaded}

No reutilice el archivo \textbf{pbsa.mdp} de \textbf{g\_mmpbsa}. La
sintaxis actual usa bloques namelist como \textbf{\&general},
\textbf{\&gb}, \textbf{\&pb} y \textbf{\&decomp}.

\subsubsection{13. Entrada mínima
MM/GBSA}\label{entrada-muxednima-mmgbsa}

Archivo \textbf{mmpbsa\_gb.in}:

\begin{Shaded}
\begin{Highlighting}[]
\NormalTok{\&general}
\NormalTok{  sys\_name="Complejo\_proteina\_ligando",}
\NormalTok{  startframe=1,}
\NormalTok{  endframe=5000,}
\NormalTok{  interval=10,}
\NormalTok{  PBRadii=4,}
\NormalTok{  temperature=300.0,}
\NormalTok{/}

\NormalTok{\&gb}
\NormalTok{  igb=8,}
\NormalTok{  saltcon=0.150,}
\NormalTok{/}
\end{Highlighting}
\end{Shaded}

Con \textbf{igb=8}, la documentación recomienda \textbf{PBRadii=4},
correspondiente a mbondi3.

\textbf{startframe}, \textbf{endframe} e \textbf{interval} se refieren a
índices de marcos procesados, no necesariamente a ps. El número
analizado es aproximadamente:

\[
N_{\mathrm{frames}}
=
\left\lfloor
\frac{f_{\mathrm{final}}-f_{\mathrm{inicial}}}
     {\mathrm{interval}}
\right\rfloor+1
\]

No elija 500 marcos consecutivos y los trate como 500 observaciones
independientes.

\subsubsection{14. Ejecutar MM/GBSA}\label{ejecutar-mmgbsa}

\begin{Shaded}
\begin{Highlighting}[]
\ExtensionTok{gmx\_MMPBSA} \AttributeTok{{-}O} \DataTypeTok{\textbackslash{}}
    \AttributeTok{{-}i}\NormalTok{ mmpbsa\_gb.in }\DataTypeTok{\textbackslash{}}
    \AttributeTok{{-}cs}\NormalTok{ md.tpr }\DataTypeTok{\textbackslash{}}
    \AttributeTok{{-}ct}\NormalTok{ md\_fit.xtc }\DataTypeTok{\textbackslash{}}
    \AttributeTok{{-}ci}\NormalTok{ index.ndx }\DataTypeTok{\textbackslash{}}
    \AttributeTok{{-}cg}\NormalTok{ Receptor Ligando }\DataTypeTok{\textbackslash{}}
    \AttributeTok{{-}cp}\NormalTok{ topol.top }\DataTypeTok{\textbackslash{}}
    \AttributeTok{{-}cr}\NormalTok{ reference.pdb }\DataTypeTok{\textbackslash{}}
    \AttributeTok{{-}o}\NormalTok{ FINAL\_RESULTS\_MMGBSA.dat }\DataTypeTok{\textbackslash{}}
    \AttributeTok{{-}eo}\NormalTok{ FINAL\_RESULTS\_MMGBSA.csv }\DataTypeTok{\textbackslash{}}
    \AttributeTok{{-}nogui}
\end{Highlighting}
\end{Shaded}

Significado de las opciones principales:

{\def\LTcaptype{none} % do not increment counter
\begin{longtable}[]{@{}
  >{\raggedright\arraybackslash}p{(\linewidth - 2\tabcolsep) * \real{0.1316}}
  >{\raggedright\arraybackslash}p{(\linewidth - 2\tabcolsep) * \real{0.8684}}@{}}
\toprule\noalign{}
\begin{minipage}[b]{\linewidth}\raggedright
Opción
\end{minipage} & \begin{minipage}[b]{\linewidth}\raggedright
Función
\end{minipage} \\
\midrule\noalign{}
\endhead
\bottomrule\noalign{}
\endlastfoot
\textbf{-O} & Permite sobrescribir resultados existentes. Úsela
deliberadamente. \\
\textbf{-i} & Archivo de parámetros. \\
\textbf{-cs} & TPR o PDB del complejo; se recomienda el TPR de
producción. \\
\textbf{-ct} & Una o más trayectorias del complejo. \\
\textbf{-ci} & Índice del complejo. \\
\textbf{-cg} & Grupos receptor y ligando. \\
\textbf{-cp} & Topología completa de GROMACS, obligatoria. \\
\textbf{-cr} & PDB de referencia recomendado. \\
\textbf{-o} & Resumen estadístico. \\
\textbf{-eo} & Energías por marco en CSV. \\
\textbf{-nogui} & No abre el analizador gráfico al terminar. \\
\end{longtable}
}

\subsubsection{15. Entrada mínima
MM/PBSA}\label{entrada-muxednima-mmpbsa}

Genere primero la plantilla de su versión:

\begin{Shaded}
\begin{Highlighting}[]
\ExtensionTok{gmx\_MMPBSA} \AttributeTok{{-}{-}create\_input}\NormalTok{ pb}
\end{Highlighting}
\end{Shaded}

Ejemplo inicial \textbf{mmpbsa\_pb.in}:

\begin{Shaded}
\begin{Highlighting}[]
\NormalTok{\&general}
\NormalTok{  sys\_name="Complejo\_proteina\_ligando\_PB",}
\NormalTok{  startframe=1,}
\NormalTok{  endframe=1000,}
\NormalTok{  interval=10,}
\NormalTok{  temperature=300.0,}
\NormalTok{/}

\NormalTok{\&pb}
\NormalTok{  istrng=0.150,}
\NormalTok{  fillratio=4.0,}
\NormalTok{/}
\end{Highlighting}
\end{Shaded}

Ejecución:

\begin{Shaded}
\begin{Highlighting}[]
\ExtensionTok{gmx\_MMPBSA} \AttributeTok{{-}O} \DataTypeTok{\textbackslash{}}
    \AttributeTok{{-}i}\NormalTok{ mmpbsa\_pb.in }\DataTypeTok{\textbackslash{}}
    \AttributeTok{{-}cs}\NormalTok{ md.tpr }\DataTypeTok{\textbackslash{}}
    \AttributeTok{{-}ct}\NormalTok{ md\_fit.xtc }\DataTypeTok{\textbackslash{}}
    \AttributeTok{{-}ci}\NormalTok{ index.ndx }\DataTypeTok{\textbackslash{}}
    \AttributeTok{{-}cg}\NormalTok{ Receptor Ligando }\DataTypeTok{\textbackslash{}}
    \AttributeTok{{-}cp}\NormalTok{ topol.top }\DataTypeTok{\textbackslash{}}
    \AttributeTok{{-}cr}\NormalTok{ reference.pdb }\DataTypeTok{\textbackslash{}}
    \AttributeTok{{-}o}\NormalTok{ FINAL\_RESULTS\_MMPBSA.dat }\DataTypeTok{\textbackslash{}}
    \AttributeTok{{-}eo}\NormalTok{ FINAL\_RESULTS\_MMPBSA.csv }\DataTypeTok{\textbackslash{}}
    \AttributeTok{{-}nogui}
\end{Highlighting}
\end{Shaded}

PB suele ser más costoso. Pruebe primero pocos marcos y revise
convergencia numérica antes de lanzar el conjunto completo.

\subsubsection{16. Dieléctricos y fuerza
iónica}\label{dieluxe9ctricos-y-fuerza-iuxf3nica}

En PB aparecen:

\begin{itemize}
\tightlist
\item
  \textbf{indi}: dieléctrico interno;
\item
  \textbf{exdi}: dieléctrico externo;
\item
  \textbf{istrng}: fuerza iónica molar.
\end{itemize}

En GB:

\begin{itemize}
\tightlist
\item
  \textbf{intdiel}: dieléctrico interno;
\item
  \textbf{extdiel}: dieléctrico externo;
\item
  \textbf{saltcon}: concentración salina molar.
\end{itemize}

No ajuste el dieléctrico interno solo para acercar el resultado a un
dato experimental. Puede hacerse un análisis de sensibilidad, por
ejemplo con 1, 2 y 4, manteniendo el resto constante.

La fuerza iónica se define:

\[
I=\frac{1}{2}\sum_i c_i z_i^2
\]

Para NaCl ideal 1:1, 0.15 M produce \(I=0.15\ \mathrm{M}\). Para una sal
divalente la concentración molar y la fuerza iónica no son iguales.

\subsubsection{17. Radios atómicos}\label{radios-atuxf3micos}

Los radios influyen en la frontera dieléctrica y en el término polar. No
son una opción cosmética.

Correspondencias documentadas:

{\def\LTcaptype{none} % do not increment counter
\begin{longtable}[]{@{}rll@{}}
\toprule\noalign{}
PBRadii & Conjunto & Uso típico \\
\midrule\noalign{}
\endhead
\bottomrule\noalign{}
\endlastfoot
1 & bondi & Asociado habitualmente con igb=7. \\
2 & mbondi & Asociado habitualmente con igb=1. \\
3 & mbondi2 & Asociado con igb=2 o 5. \\
4 & mbondi3 & Asociado con igb=8. \\
7 & charmm\_radii & Solo PB y topologías preparadas con CHARMM. \\
\end{longtable}
}

El campo de fuerza de la dinámica y el conjunto de radios del solvente
implícito son conceptos distintos. \textbf{PBRadii} no cambia los
parámetros enlazados ni Lennard-Jones de la topología original.

Para topologías CHARMM, consulte las limitaciones de conversión y valore
\textbf{charmm\_radii} en PB. La documentación advierte que no toda
topología producida por CHARMM-GUI queda validada automáticamente.

\subsubsection{18. Compatibilidad de
topologías}\label{compatibilidad-de-topologuxedas}

\textbf{gmx\_MMPBSA} convierte la topología GROMACS usando ParmEd.
Compruebe cuidadosamente sistemas con:

\begin{itemize}
\tightlist
\item
  parámetros CHARMM y CMAP;
\item
  átomos virtuales;
\item
  dummy atoms;
\item
  puntos de carga extra;
\item
  metales coordinados;
\item
  enlaces covalentes proteína--ligando;
\item
  topologías híbridas;
\item
  lípidos;
\item
  residuos modificados;
\item
  restricciones no convencionales.
\end{itemize}

La versión actual documenta rutas probadas para topologías
representativas Amber, OPLS y CHARMM, pero esto no implica
compatibilidad universal.

Siempre inspeccione el log y las topologías intermedias creadas.

\subsubsection{19. Sistemas de membrana}\label{sistemas-de-membrana}

No elimine la membrana conceptualmente y aplique un modelo acuoso
homogéneo sin evaluar el efecto. La versión actual incluye opciones PB
para membranas:

\begin{Shaded}
\begin{Highlighting}[]
\NormalTok{\&pb}
\NormalTok{  memopt=1,}
\NormalTok{  emem=7.0,}
\NormalTok{  indi=4.0,}
\NormalTok{  mctrdz=automatic,}
\NormalTok{  mthick=automatic,}
\NormalTok{  membrane\_atoms="P",}
\NormalTok{  poretype=1,}
\NormalTok{  radiopt=0,}
\NormalTok{  istrng=0.150,}
\NormalTok{/}
\end{Highlighting}
\end{Shaded}

Es un esquema orientativo. Debe adaptarse a la orientación de la bicapa,
nombres de átomos, espesor, presencia de poro y modelo experimental. Una
membrana implícita mal centrada puede producir un resultado peor que un
modelo acuoso simplificado claramente declarado.

\subsubsection{20. Entropía}\label{entropuxeda}

Sin entropía:

\[
\Delta G_{\mathrm{estimada}}
=
\Delta E_{\mathrm{MM}}
+
\Delta G_{\mathrm{solv}}
\]

Con entropía:

\[
\Delta G_{\mathrm{bind}}
=
\Delta H_{\mathrm{aprox}}
-
T\Delta S
\]

Métodos disponibles:

\begin{itemize}
\tightlist
\item
  Interaction Entropy;
\item
  C2 Entropy;
\item
  modos normales, NMODE.
\end{itemize}

NMODE es costoso y sensible a minimización. El consumo total de memoria
crece aproximadamente como:

\[
RAM_{\mathrm{total}}
=
RAM_{\mathrm{por\ marco}}
\times
N_{\mathrm{procesos}}
\]

No paralelice NMODE hasta agotar memoria.

\subsubsection{21. Interaction Entropy}\label{interaction-entropy}

La aproximación IE usa fluctuaciones de la energía de interacción:

\[
-T\Delta S_{\mathrm{IE}}
=
k_\mathrm{B}T
\ln
\left\langle
\exp
\left[
\beta\,
\Delta E_{\mathrm{int}}
\right]
\right\rangle
\]

con:

\[
\Delta E_{\mathrm{int}}
=
E_{\mathrm{int}}
-
\left\langle E_{\mathrm{int}}\right\rangle
\]

Entrada:

\begin{Shaded}
\begin{Highlighting}[]
\NormalTok{\&general}
\NormalTok{  sys\_name="Complejo\_IE",}
\NormalTok{  startframe=1,}
\NormalTok{  endframe=5000,}
\NormalTok{  interval=1,}
\NormalTok{  temperature=300.0,}
\NormalTok{  PBRadii=4,}
\NormalTok{  interaction\_entropy=1,}
\NormalTok{  ie\_segment=25,}
\NormalTok{/}

\NormalTok{\&gb}
\NormalTok{  igb=8,}
\NormalTok{  saltcon=0.150,}
\NormalTok{/}
\end{Highlighting}
\end{Shaded}

Debe informarse \(\sigma_{IE}\), la desviación de la energía de
interacción. La documentación desaconseja IE cuando:

\[
\sigma_{IE}\gtrsim3.6\ \mathrm{kcal\,mol^{-1}}
\]

porque el promedio exponencial se vuelve difícil de converger. Esto
equivale aproximadamente a:

\[
3.6\ \mathrm{kcal\,mol^{-1}}
=
15.1\ \mathrm{kJ\,mol^{-1}}
\]

\textbf{ie\_segment} es un diagnóstico de la cola de la curva
acumulativa; no reemplaza el resultado calculado sobre todo el conjunto
seleccionado.

\subsubsection{22. Descomposición por
residuos}\label{descomposiciuxf3n-por-residuos}

Genere una plantilla:

\begin{Shaded}
\begin{Highlighting}[]
\ExtensionTok{gmx\_MMPBSA} \AttributeTok{{-}{-}create\_input}\NormalTok{ gb decomp}
\end{Highlighting}
\end{Shaded}

Ejemplo:

\begin{Shaded}
\begin{Highlighting}[]
\NormalTok{\&general}
\NormalTok{  sys\_name="Descomposicion",}
\NormalTok{  startframe=1,}
\NormalTok{  endframe=5000,}
\NormalTok{  interval=10,}
\NormalTok{  PBRadii=4,}
\NormalTok{/}

\NormalTok{\&gb}
\NormalTok{  igb=8,}
\NormalTok{  saltcon=0.150,}
\NormalTok{/}

\NormalTok{\&decomp}
\NormalTok{  idecomp=2,}
\NormalTok{  dec\_verbose=1,}
\NormalTok{  print\_res="within 4",}
\NormalTok{/}
\end{Highlighting}
\end{Shaded}

Ejecución:

\begin{Shaded}
\begin{Highlighting}[]
\ExtensionTok{gmx\_MMPBSA} \AttributeTok{{-}O} \DataTypeTok{\textbackslash{}}
    \AttributeTok{{-}i}\NormalTok{ mmpbsa\_decomp.in }\DataTypeTok{\textbackslash{}}
    \AttributeTok{{-}cs}\NormalTok{ md.tpr }\DataTypeTok{\textbackslash{}}
    \AttributeTok{{-}ct}\NormalTok{ md\_fit.xtc }\DataTypeTok{\textbackslash{}}
    \AttributeTok{{-}ci}\NormalTok{ index.ndx }\DataTypeTok{\textbackslash{}}
    \AttributeTok{{-}cg}\NormalTok{ Receptor Ligando }\DataTypeTok{\textbackslash{}}
    \AttributeTok{{-}cp}\NormalTok{ topol.top }\DataTypeTok{\textbackslash{}}
    \AttributeTok{{-}cr}\NormalTok{ reference.pdb }\DataTypeTok{\textbackslash{}}
    \AttributeTok{{-}o}\NormalTok{ FINAL\_RESULTS\_DECOMP.dat }\DataTypeTok{\textbackslash{}}
    \AttributeTok{{-}do}\NormalTok{ FINAL\_DECOMP\_MMPBSA.dat }\DataTypeTok{\textbackslash{}}
    \AttributeTok{{-}eo}\NormalTok{ FINAL\_RESULTS\_DECOMP.csv }\DataTypeTok{\textbackslash{}}
    \AttributeTok{{-}deo}\NormalTok{ FINAL\_DECOMP\_MMPBSA.csv }\DataTypeTok{\textbackslash{}}
    \AttributeTok{{-}nogui}
\end{Highlighting}
\end{Shaded}

\textbf{print\_res=``within 4''} selecciona residuos próximos dentro de
4 Å en la sintaxis del programa. El conjunto impreso debe contener al
menos un residuo del receptor y uno del ligando.

La descomposición:

\begin{itemize}
\tightlist
\item
  depende del esquema elegido;
\item
  no es estrictamente única;
\item
  reparte términos colectivos;
\item
  no equivale a una mutación experimental;
\item
  no demuestra causalidad;
\item
  puede variar con radios y dieléctricos.
\end{itemize}

Úsela para priorizar residuos y formular hipótesis.

\subsubsection{23. Alanine scanning}\label{alanine-scanning}

El alanine scanning computacional estima el cambio al mutar un residuo:

\[
\Delta\Delta G_{\mathrm{bind}}
=
\Delta G_{\mathrm{bind}}^{\mathrm{mutante}}
-
\Delta G_{\mathrm{bind}}^{\mathrm{WT}}
\]

Interpretación:

\begin{itemize}
\tightlist
\item
  \(\Delta\Delta G>0\): la mutación debilita la unión en esta
  convención;
\item
  \(\Delta\Delta G<0\): la mutación la favorece.
\end{itemize}

Es una mutación end-state sobre estructuras existentes. No reemplaza una
dinámica completa del mutante si este reorganiza la proteína.

\subsubsection{24. Paralelización}\label{paralelizaciuxf3n}

Ejecución serial:

\begin{Shaded}
\begin{Highlighting}[]
\ExtensionTok{gmx\_MMPBSA} \AttributeTok{{-}O} \DataTypeTok{\textbackslash{}}
    \AttributeTok{{-}i}\NormalTok{ mmpbsa\_gb.in }\DataTypeTok{\textbackslash{}}
    \AttributeTok{{-}cs}\NormalTok{ md.tpr }\DataTypeTok{\textbackslash{}}
    \AttributeTok{{-}ct}\NormalTok{ md\_fit.xtc }\DataTypeTok{\textbackslash{}}
    \AttributeTok{{-}ci}\NormalTok{ index.ndx }\DataTypeTok{\textbackslash{}}
    \AttributeTok{{-}cg}\NormalTok{ Receptor Ligando }\DataTypeTok{\textbackslash{}}
    \AttributeTok{{-}cp}\NormalTok{ topol.top }\DataTypeTok{\textbackslash{}}
    \AttributeTok{{-}cr}\NormalTok{ reference.pdb }\DataTypeTok{\textbackslash{}}
    \AttributeTok{{-}nogui}
\end{Highlighting}
\end{Shaded}

Con MPI:

\begin{Shaded}
\begin{Highlighting}[]
\ExtensionTok{mpirun} \AttributeTok{{-}np}\NormalTok{ 4 gmx\_MMPBSA }\AttributeTok{{-}O} \DataTypeTok{\textbackslash{}}
    \AttributeTok{{-}i}\NormalTok{ mmpbsa\_gb.in }\DataTypeTok{\textbackslash{}}
    \AttributeTok{{-}cs}\NormalTok{ md.tpr }\DataTypeTok{\textbackslash{}}
    \AttributeTok{{-}ct}\NormalTok{ md\_fit.xtc }\DataTypeTok{\textbackslash{}}
    \AttributeTok{{-}ci}\NormalTok{ index.ndx }\DataTypeTok{\textbackslash{}}
    \AttributeTok{{-}cg}\NormalTok{ Receptor Ligando }\DataTypeTok{\textbackslash{}}
    \AttributeTok{{-}cp}\NormalTok{ topol.top }\DataTypeTok{\textbackslash{}}
    \AttributeTok{{-}cr}\NormalTok{ reference.pdb }\DataTypeTok{\textbackslash{}}
    \AttributeTok{{-}nogui}
\end{Highlighting}
\end{Shaded}

No use \textbf{gmx\_mpi} dentro de esta ejecución MPI. La documentación
indica usar \textbf{gmx}, porque las herramientas auxiliares de GROMACS
no se benefician del MPI de \textbf{mdrun} y pueden entrar en conflicto
con \textbf{mpirun}.

El escalamiento deja de mejorar cuando el número de procesos se aproxima
al número de marcos o domina la carga de topologías.

\subsubsection{25. Análisis gráfico}\label{anuxe1lisis-gruxe1fico}

\begin{Shaded}
\begin{Highlighting}[]
\ExtensionTok{gmx\_MMPBSA\_ana} \DataTypeTok{\textbackslash{}}
    \AttributeTok{{-}f}\NormalTok{ FINAL\_RESULTS\_MMPBSA.dat}
\end{Highlighting}
\end{Shaded}

El analizador permite revisar términos, evolución temporal,
descomposición y exportar gráficos. En HPC use los archivos DAT y CSV, y
abra el analizador en una estación con entorno gráfico.

Conserve:

\begin{Shaded}
\begin{Highlighting}[]
\NormalTok{FINAL\_RESULTS\_MMPBSA.dat}
\NormalTok{FINAL\_RESULTS\_MMPBSA.csv}
\NormalTok{FINAL\_DECOMP\_MMPBSA.dat}
\NormalTok{FINAL\_DECOMP\_MMPBSA.csv}
\NormalTok{gmx\_MMPBSA.log}
\NormalTok{mmpbsa.in}
\end{Highlighting}
\end{Shaded}

No conserve únicamente una captura del valor final.

\subsubsection{26. Unidades}\label{unidades}

Las salidas heredadas de AmberTools suelen expresarse en kcal·mol⁻¹.
Verifique siempre el encabezado del archivo.

\[
1\ \mathrm{kcal\,mol^{-1}}
=
4.184\ \mathrm{kJ\,mol^{-1}}
\]

Ejemplo:

\[
-28.6\ \mathrm{kcal\,mol^{-1}}
=
-119.7\ \mathrm{kJ\,mol^{-1}}
\]

No mezcle resultados antiguos de \textbf{g\_mmpbsa}, que normalmente se
informaban en kJ·mol⁻¹, con resultados de \textbf{gmx\_MMPBSA} sin
convertir unidades.

En parámetros PB y GB aparecen además:

\begin{itemize}
\tightlist
\item
  concentración y fuerza iónica en mol·L⁻¹;
\item
  radios y distancias en Å;
\item
  tensión superficial en kcal·mol⁻¹·Å⁻².
\end{itemize}

\subsubsection{27. Relación con afinidad
experimental}\label{relaciuxf3n-con-afinidad-experimental}

Bajo estados estándar compatibles:

\[
\Delta G^\circ_{\mathrm{bind}}
=
RT\ln K_d
=
-RT\ln K_a
\]

A 298.15 K:

\[
RT=2.47896\ \mathrm{kJ\,mol^{-1}}
=0.59248\ \mathrm{kcal\,mol^{-1}}
\]

Entonces:

\[
K_d=
\exp
\left(
\frac{\Delta G^\circ_{\mathrm{bind}}}{RT}
\right)
\]

Esta conversión solo es válida para una energía libre estándar completa.
No convierta directamente un MM/GBSA sin entropía ni correcciones en un
\(K_d\) ``predicho''.

Para rankings, compare correlación y error frente a datos experimentales
usando la misma serie química y protocolo.

\subsubsection{28. Muestreo de marcos}\label{muestreo-de-marcos}

Más marcos muy correlacionados aportan menos información que marcos
separados por encima del tiempo de autocorrelación.

Evalúe:

\begin{itemize}
\tightlist
\item
  distancia temporal entre marcos;
\item
  estabilidad del ligando;
\item
  cambios de pose;
\item
  RMSD del sitio;
\item
  bloques tempranos y tardíos;
\item
  réplicas independientes.
\end{itemize}

Ejemplo conceptual:

\begin{Shaded}
\begin{Highlighting}[]
\NormalTok{0–20 ns}
\NormalTok{20–40 ns}
\NormalTok{40–60 ns}
\NormalTok{60–80 ns}
\NormalTok{80–100 ns}
\end{Highlighting}
\end{Shaded}

Calcule MM/PBSA por bloque. Una media estable con bloques discrepantes
no implica convergencia.

La incertidumbre entre réplicas suele ser más informativa que el error
interno calculado sobre marcos correlacionados.

\subsubsection{29. Sensibilidad del
protocolo}\label{sensibilidad-del-protocolo}

Repita una fracción representativa cambiando una variable por vez:

\begin{itemize}
\tightlist
\item
  PB frente a GB;
\item
  conjunto de radios;
\item
  dieléctrico interno;
\item
  fuerza iónica;
\item
  inclusión de agua explícita estructural;
\item
  selección temporal;
\item
  intervalo entre marcos;
\item
  estado de protonación.
\end{itemize}

Si el ranking cambia completamente ante modificaciones pequeñas y
razonables, el resultado no es robusto.

No elija retrospectivamente los parámetros que maximizan correlación sin
validación externa: eso equivale a ajustar el método al conjunto.

\subsubsection{30. Agua estructural}\label{agua-estructural}

El protocolo estándar elimina agua explícita antes de evaluar el
solvente implícito. Algunas aguas del sitio pueden mediar la unión y
formar parte del receptor efectivo.

La versión actual incluye ejemplos de ST MM/GBSA con aguas explícitas
del receptor. Si se conservan:

\begin{itemize}
\tightlist
\item
  seleccione aguas con criterios reproducibles;
\item
  mantenga identidades o sitios coherentes;
\item
  compruebe ocupación;
\item
  evite incluir aguas transitorias arbitrarias;
\item
  documente si pertenecen al receptor o al solvente.
\end{itemize}

Una selección fija basada solo en una instantánea puede sesgar el
resultado.

\subsubsection{31. Interpretación de
componentes}\label{interpretaciuxf3n-de-componentes}

Una salida típica contiene:

\begin{itemize}
\tightlist
\item
  \textbf{VDWAALS}: Lennard-Jones;
\item
  \textbf{EEL}: electrostática molecular;
\item
  \textbf{EGB} o \textbf{EPB}: solvatación polar;
\item
  \textbf{ESURF/ENPOLAR/EDISPER}: términos no polares según el modelo;
\item
  \textbf{GGAS}: contribución molecular;
\item
  \textbf{GSOLV}: contribución de solvatación;
\item
  \textbf{TOTAL}: suma reportada.
\end{itemize}

No interprete cada componente como una magnitud experimental separable.
Electrostatica molecular y solvatación polar suelen compensarse
fuertemente.

En algunos modos PB, la documentación advierte que la partición entre
\textbf{GGAS} y \textbf{GSOLV} no conserva la interpretación habitual
aunque \textbf{TOTAL} siga siendo válido. Lea el log antes de
interpretar columnas.

\subsubsection{32. Comparación entre
ligandos}\label{comparaciuxf3n-entre-ligandos}

Para comparar una serie:

\begin{itemize}
\tightlist
\item
  mismo receptor y estado de protonación;
\item
  mismo campo de fuerza;
\item
  mismo modelo PB/GB;
\item
  mismas opciones y radios;
\item
  temperatura equivalente;
\item
  ventanas temporales comparables;
\item
  número de réplicas similar;
\item
  mismo tratamiento de aguas;
\item
  misma definición del receptor.
\end{itemize}

Reporte:

\[
\Delta\Delta G_i
=
\Delta G_i-
\Delta G_{\mathrm{referencia}}
\]

Las diferencias relativas suelen ser más útiles que valores absolutos,
pero no eliminan sesgos específicos de cada grupo funcional.

\subsubsection{33. Errores frecuentes}\label{errores-frecuentes-2}

\begin{itemize}
\tightlist
\item
  Confundir g\_mmpbsa con gmx\_MMPBSA.
\item
  Descargar ejecutables antiguos y copiar \textbf{pbsa.py}.
\item
  Usar \textbf{pbsa.mdp} con la interfaz nueva.
\item
  Describir MM/PBSA como FEP con λ.
\item
  Afirmar que umbrella sampling es obligatorio.
\item
  Omitir \textbf{-cp topol.top} en la versión actual.
\item
  Usar un PDB sin parámetros para el ligando.
\item
  Seleccionar grupos solapados o equivocados.
\item
  Copiar números de grupo de otro sistema.
\item
  Analizar una trayectoria rota por PBC.
\item
  Permitir que el ligando quede lejos del receptor.
\item
  Mezclar topología y TPR con distinto orden atómico.
\item
  Suponer compatibilidad automática de cualquier topología CHARMM.
\item
  Tratar concentración de sal como fuerza iónica para sales
  multivalentes.
\item
  Elegir dieléctricos para reproducir un resultado esperado.
\item
  Interpretar TOTAL sin conocer sus unidades.
\item
  Convertir un resultado sin entropía directamente en \(K_d\).
\item
  Tratar marcos correlacionados como réplicas.
\item
  Usar descomposición como prueba causal.
\item
  Paralelizar NMODE hasta agotar memoria.
\item
  Usar \textbf{gmx\_mpi} bajo \textbf{mpirun} con gmx\_MMPBSA.
\item
  Informar muchos decimales sin incertidumbre entre réplicas.
\end{itemize}

\subsubsection{34. Controles mínimos}\label{controles-muxednimos}

Antes de aceptar el resultado:

\begin{itemize}
\tightlist
\item
  instalación verificada;
\item
  topología convertida sin errores;
\item
  ligando parametrizado;
\item
  grupos comprobados;
\item
  referencia con cadenas correctas;
\item
  trayectoria entera y ajustada;
\item
  marcos posteriores a la equilibración;
\item
  resultado estable por bloques;
\item
  réplicas compatibles;
\item
  sensibilidad a PB/GB o dieléctricos evaluada;
\item
  entropía incluida o ausencia declarada;
\item
  unidades confirmadas;
\item
  componentes interpretados con cautela.
\end{itemize}

\subsubsection{35. Información mínima para
reproducibilidad}\label{informaciuxf3n-muxednima-para-reproducibilidad}

Informe:

\begin{itemize}
\tightlist
\item
  versión de gmx\_MMPBSA;
\item
  versiones de GROMACS, AmberTools y ParmEd;
\item
  sistema operativo o entorno;
\item
  campo de fuerza y modelo de agua;
\item
  parametrización del ligando;
\item
  archivos y grupos usados;
\item
  protocolo ST o MT;
\item
  cantidad y separación temporal de marcos;
\item
  bloques namelist completos;
\item
  radios, dieléctricos y fuerza iónica;
\item
  método entrópico;
\item
  número de procesos;
\item
  réplicas;
\item
  unidades;
\item
  criterio de incertidumbre;
\item
  tratamiento de aguas y membranas.
\end{itemize}

\subsubsection{Fuentes}\label{fuentes-6}

\begin{enumerate}
\def\labelenumi{\arabic{enumi}.}
\tightlist
\item
  \href{https://valdes-tresanco-ms.github.io/gmx_MMPBSA/dev/getting-started/}{Inicio
  y requisitos de gmx\_MMPBSA}
\item
  \href{https://valdes-tresanco-ms.github.io/gmx_MMPBSA/dev/installation/}{Instalación
  actual}
\item
  \href{https://valdes-tresanco-ms.github.io/gmx_MMPBSA/dev/howworks/}{Funcionamiento
  y preparación de topologías}
\item
  \href{https://valdes-tresanco-ms.github.io/gmx_MMPBSA/dev/gmx_MMPBSA_command-line/}{Interfaz
  de línea de comandos}
\item
  \href{https://valdes-tresanco-ms.github.io/gmx_MMPBSA/dev/examples/}{Ejemplos
  oficiales}
\item
  \href{https://valdes-tresanco-ms.github.io/gmx_MMPBSA/dev/examples/Protein_ligand/ST/}{Ejemplo
  proteína--ligando ST}
\item
  \href{https://valdes-tresanco-ms.github.io/gmx_MMPBSA/dev/examples/Decomposition_analysis/}{Ejemplo
  de descomposición}
\item
  \href{https://valdes-tresanco-ms.github.io/gmx_MMPBSA/dev/examples/Entropy_calculations/Interaction_Entropy/}{Ejemplo
  de Interaction Entropy}
\item
  \href{https://pubs.acs.org/doi/10.1021/acs.jctc.1c00645}{Artículo
  original de gmx\_MMPBSA}
\end{enumerate}

\section{Interacción energética lineal (Linear Interaction Energy,
LIE)}\label{interacciuxf3n-energuxe9tica-lineal-linear-interaction-energy-lie}

La \textbf{energía de interacción lineal} (LIE) es un método de estado
final para estimar afinidades de unión. Requiere muestrear
explícitamente dos estados mediante dinámica molecular:

\begin{enumerate}
\def\labelenumi{\arabic{enumi}.}
\tightlist
\item
  el ligando unido al receptor y rodeado por solvente e iones;
\item
  el mismo ligando libre en solvente, con el mismo estado de protonación
  y un protocolo compatible.
\end{enumerate}

LIE es más económico que una transformación alquímica completa porque no
introduce estados intermedios de \(\lambda\). A cambio, es un modelo
semiempírico: no debe interpretarse como una energía libre rigurosa ni
utilizarse con coeficientes tomados arbitrariamente de otros sistemas.
Su utilidad principal es comparar ligandos relacionados para un mismo
receptor y dentro del dominio químico e interaccional empleado para
calibrar el modelo.

\subsection{Fundamento teórico}\label{fundamento-teuxf3rico}

En la forma más habitual,

\[
\Delta G_{\mathrm{bind}}^{\mathrm{LIE}}
=
\alpha\left(
\left\langle V_{\mathrm{vdW}}^{L-E}\right\rangle_{\mathrm{bound}}
-
\left\langle V_{\mathrm{vdW}}^{L-E}\right\rangle_{\mathrm{free}}
\right)
+
\beta\left(
\left\langle V_{\mathrm{elec}}^{L-E}\right\rangle_{\mathrm{bound}}
-
\left\langle V_{\mathrm{elec}}^{L-E}\right\rangle_{\mathrm{free}}
\right)
+
\gamma .
\]

Aquí:

\begin{itemize}
\tightlist
\item
  \(L\) es el ligando y \(E\) es su entorno;
\item
  en el estado unido, el entorno incluye receptor, agua, iones y
  cualquier cofactor que no forme parte del ligando;
\item
  en el estado libre, el entorno incluye agua e iones;
\item
  los corchetes \(\langle\cdots\rangle\) representan promedios de
  conjunto, aproximados mediante promedios temporales sobre trayectorias
  equilibradas;
\item
  \(\alpha\) y \(\beta\) son coeficientes adimensionales para las
  contribuciones de Lennard-Jones y electrostática;
\item
  \(\gamma\) es un intercepto opcional, expresado en kJ mol\(^{-1}\).
\end{itemize}

La contribución electrostática procede de una aproximación de respuesta
lineal. El valor histórico \(\beta=0{,}5\) corresponde al caso ideal de
respuesta lineal, pero no es universal. La documentación de GROMACS
2026.3 conserva como valores predeterminados \(\alpha=0{,}181\) y
\(\beta=0{,}5\); deben considerarse valores iniciales o de referencia,
no una validación del modelo para cualquier ligando.

LIE no calcula de manera explícita todos los términos entrópicos,
reorganizaciones internas, cambios conformacionales o contribuciones de
estado estándar. Se presupone que una parte de esos efectos queda
representada de forma efectiva por los coeficientes y el intercepto. Por
eso un resultado aislado obtenido con los valores predeterminados debe
describirse como \textbf{estimación LIE no calibrada}.

\subsection{Unidades y comparación con datos
experimentales}\label{unidades-y-comparaciuxf3n-con-datos-experimentales}

GROMACS informa las energías de interacción en \textbf{kJ mol\(^{-1}\)}.
Como \(\alpha\) y \(\beta\) son adimensionales,
\(\Delta G_{\mathrm{bind}}^{\mathrm{LIE}}\) y \(\gamma\) conservan esas
unidades.

\[
1\ \mathrm{kcal\ mol^{-1}}=4{,}184\ \mathrm{kJ\ mol^{-1}}.
\]

Para transformar una constante experimental de disociación en una
energía libre estándar,

\[
\Delta G^\circ_{\mathrm{bind}}
=
RT\ln\left(\frac{K_d}{C^\circ}\right)
=
-RT\ln\left(K_a C^\circ\right),
\]

donde \(C^\circ=1\ \mathrm{mol\ L^{-1}}\). La razón dentro del logaritmo
debe ser adimensional. Deben registrarse la temperatura, la fuerza
iónica, el pH y el estado de protonación del ligando; comparar
directamente valores obtenidos bajo condiciones experimentales distintas
puede introducir un error mayor que el que se intenta modelar.

\subsection{Diseño de las
simulaciones}\label{diseuxf1o-de-las-simulaciones}

Los estados unido y libre deben usar, en lo posible:

\begin{itemize}
\tightlist
\item
  el mismo campo de fuerzas para el ligando;
\item
  el mismo modelo de agua, concentración salina y temperatura;
\item
  idéntico tratamiento de interacciones no enlazantes;
\item
  el mismo estado de protonación, tautomería y carga;
\item
  longitudes de producción y criterios de descarte comparables;
\item
  varias réplicas independientes cuando el coste lo permita.
\end{itemize}

No debe extraerse el ligando de la caja del complejo y analizarse sin
una simulación libre propia: la relajación y la reorganización del
solvente son parte esencial del estado libre. Para sistemas flexibles o
con poses alternativas conviene iniciar réplicas desde conformaciones
distintas. Si el ligando abandona el sitio, cambia de pose de forma
irreversible o el receptor experimenta una transición grande, esa
trayectoria no debe promediarse ciegamente con las demás.

\subsection{Preparación de grupos de
energía}\label{preparaciuxf3n-de-grupos-de-energuxeda}

La orden \textbf{gmx lie} necesita términos de interacción no enlazante
entre el ligando y el resto del sistema en el archivo EDR. La separación
más clara es usar dos grupos no superpuestos que cubran todo el sistema:
\textbf{LIG} y \textbf{Environment}.

Para crear un índice estático a partir del TPR del complejo:

\begin{Shaded}
\begin{Highlighting}[]
\ExtensionTok{gmx}\NormalTok{ select }\AttributeTok{{-}s}\NormalTok{ md\_bound.tpr }\AttributeTok{{-}on}\NormalTok{ lie\_bound.ndx }\DataTypeTok{\textbackslash{}}
  \AttributeTok{{-}select} \StringTok{\textquotesingle{}"LIG" resname LIG; "Environment" not resname LIG\textquotesingle{}}
\end{Highlighting}
\end{Shaded}

Para el ligando libre:

\begin{Shaded}
\begin{Highlighting}[]
\ExtensionTok{gmx}\NormalTok{ select }\AttributeTok{{-}s}\NormalTok{ md\_free.tpr }\AttributeTok{{-}on}\NormalTok{ lie\_free.ndx }\DataTypeTok{\textbackslash{}}
  \AttributeTok{{-}select} \StringTok{\textquotesingle{}"LIG" resname LIG; "Environment" not resname LIG\textquotesingle{}}
\end{Highlighting}
\end{Shaded}

Se debe reemplazar \textbf{LIG} por el nombre real del residuo o por una
selección inequívoca. Si existen varias moléculas con ese mismo nombre,
hay que decidir si constituyen un único ligando termodinámico o si deben
analizarse por separado. Verifique siempre los grupos:

\begin{Shaded}
\begin{Highlighting}[]
\ExtensionTok{gmx}\NormalTok{ check }\AttributeTok{{-}f}\NormalTok{ md\_bound.xtc}
\ExtensionTok{gmx}\NormalTok{ make\_ndx }\AttributeTok{{-}f}\NormalTok{ md\_bound.tpr }\AttributeTok{{-}n}\NormalTok{ lie\_bound.ndx}
\end{Highlighting}
\end{Shaded}

Dentro de una copia del MDP de producción, agregue:

\begin{Shaded}
\begin{Highlighting}[]
\CommentTok{; Grupos no superpuestos que cubren el sistema}
\DataTypeTok{energygrps }\OtherTok{=}\StringTok{ LIG Environment}
\end{Highlighting}
\end{Shaded}

El resto del MDP usado para recalcular energías debe conservar el campo
de fuerzas, los cortes, las reglas de dispersión, el modificador de
potencial y el tratamiento electrostático de la simulación de
producción. Cambiar esos parámetros durante el análisis define otro
descriptor y rompe la comparabilidad con un modelo calibrado
previamente.

\subsubsection{Limitación importante de
PME}\label{limitaciuxf3n-importante-de-pme}

Los términos que utiliza \textbf{gmx lie} son, entre otros,
\textbf{Coul-SR} y \textbf{LJ-SR}. La parte recíproca de PME no se
descompone en pares de grupos como una energía ligando--entorno
independiente. Por tanto, una LIE basada en esos términos es dependiente
del protocolo y no contiene una asignación completa de la electrostática
de largo alcance al ligando.

Esto no significa que deba eliminarse PME de una simulación moderna.
Significa que todos los ligandos, los dos estados y el conjunto de
calibración deben tratarse de manera idéntica, y que la limitación debe
documentarse. Recalcular con reacción de campo u otro modelo
electrostático sólo es defendible si todo el modelo LIE fue calibrado y
validado con ese mismo protocolo.

\subsection{Obtención de los archivos EDR para
LIE}\label{obtenciuxf3n-de-los-archivos-edr-para-lie}

Incluir grupos de energía durante una producción puede limitar la
aceleración por GPU. Una alternativa práctica es recalcular las energías
sobre las trayectorias ya producidas. Genere TPR de análisis separados
para el complejo y el ligando libre:

\begin{Shaded}
\begin{Highlighting}[]
\ExtensionTok{gmx}\NormalTok{ grompp }\AttributeTok{{-}f}\NormalTok{ lie\_bound.mdp }\AttributeTok{{-}c}\NormalTok{ md\_bound.gro }\AttributeTok{{-}t}\NormalTok{ md\_bound.cpt }\DataTypeTok{\textbackslash{}}
  \AttributeTok{{-}p}\NormalTok{ topol\_bound.top }\AttributeTok{{-}n}\NormalTok{ lie\_bound.ndx }\AttributeTok{{-}o}\NormalTok{ lie\_bound.tpr}

\ExtensionTok{gmx}\NormalTok{ grompp }\AttributeTok{{-}f}\NormalTok{ lie\_free.mdp }\AttributeTok{{-}c}\NormalTok{ md\_free.gro }\AttributeTok{{-}t}\NormalTok{ md\_free.cpt }\DataTypeTok{\textbackslash{}}
  \AttributeTok{{-}p}\NormalTok{ topol\_free.top }\AttributeTok{{-}n}\NormalTok{ lie\_free.ndx }\AttributeTok{{-}o}\NormalTok{ lie\_free.tpr}
\end{Highlighting}
\end{Shaded}

Luego efectúe el recálculo en CPU:

\begin{Shaded}
\begin{Highlighting}[]
\ExtensionTok{gmx}\NormalTok{ mdrun }\AttributeTok{{-}s}\NormalTok{ lie\_bound.tpr }\AttributeTok{{-}rerun}\NormalTok{ md\_bound.xtc }\DataTypeTok{\textbackslash{}}
  \AttributeTok{{-}deffnm}\NormalTok{ lie\_bound }\AttributeTok{{-}nb}\NormalTok{ cpu }\AttributeTok{{-}pme}\NormalTok{ cpu}

\ExtensionTok{gmx}\NormalTok{ mdrun }\AttributeTok{{-}s}\NormalTok{ lie\_free.tpr }\AttributeTok{{-}rerun}\NormalTok{ md\_free.xtc }\DataTypeTok{\textbackslash{}}
  \AttributeTok{{-}deffnm}\NormalTok{ lie\_free }\AttributeTok{{-}nb}\NormalTok{ cpu }\AttributeTok{{-}pme}\NormalTok{ cpu}
\end{Highlighting}
\end{Shaded}

\textbf{mdrun -rerun} evalúa la energía de cada marco de la trayectoria
suministrada; no genera un nuevo muestreo. Los archivos TPR deben
conservar la misma topología, orden de átomos y parámetros no enlazantes
usados para producir cada trayectoria. Los valores cinéticos y de
temperatura obtenidos en un recálculo de este tipo no son relevantes
para LIE.

Antes de continuar, compruebe que existen los términos esperados:

\begin{Shaded}
\begin{Highlighting}[]
\ExtensionTok{gmx}\NormalTok{ energy }\AttributeTok{{-}f}\NormalTok{ lie\_free.edr }\AttributeTok{{-}o}\NormalTok{ /dev/null}
\ExtensionTok{gmx}\NormalTok{ energy }\AttributeTok{{-}f}\NormalTok{ lie\_bound.edr }\AttributeTok{{-}o}\NormalTok{ /dev/null}
\end{Highlighting}
\end{Shaded}

Deben aparecer nombres equivalentes a:

\begin{Shaded}
\begin{Highlighting}[]
\NormalTok{Coul{-}SR:LIG{-}Environment}
\NormalTok{LJ{-}SR:LIG{-}Environment}
\end{Highlighting}
\end{Shaded}

Si no aparecen, el TPR no contenía una definición válida de
\textbf{energygrps} o el cálculo se ejecutó con una ruta que no produjo
la descomposición solicitada.

\subsection{Cálculo con gmx lie en GROMACS
2026}\label{cuxe1lculo-con-gmx-lie-en-gromacs-2026}

Primero obtenga los promedios del ligando libre en solvente. El inicio
del intervalo productivo se expresa en picosegundos:

\begin{Shaded}
\begin{Highlighting}[]
\BuiltInTok{printf} \StringTok{"Coul{-}SR:LIG{-}Environment\textbackslash{}nLJ{-}SR:LIG{-}Environment\textbackslash{}n0\textbackslash{}n"} \KeywordTok{|} \DataTypeTok{\textbackslash{}}
  \ExtensionTok{gmx}\NormalTok{ energy }\AttributeTok{{-}f}\NormalTok{ lie\_free.edr }\AttributeTok{{-}o}\NormalTok{ free\_interactions.xvg }\AttributeTok{{-}b}\NormalTok{ 20000}
\end{Highlighting}
\end{Shaded}

Tome de la salida los promedios \textbf{Average} y asígnelos, sin
cambiar signos ni unidades:

\begin{Shaded}
\begin{Highlighting}[]
\VariableTok{FREE\_COUL}\OperatorTok{=}\NormalTok{{-}104.226}
\VariableTok{FREE\_LJ}\OperatorTok{=}\NormalTok{{-}183.449}
\end{Highlighting}
\end{Shaded}

Los números anteriores son únicamente ejemplos. No deben copiarse para
otro ligando.

Calcule después la estimación sobre el estado unido:

\begin{Shaded}
\begin{Highlighting}[]
\ExtensionTok{gmx}\NormalTok{ lie }\AttributeTok{{-}f}\NormalTok{ lie\_bound.edr }\AttributeTok{{-}o}\NormalTok{ lie.xvg }\AttributeTok{{-}b}\NormalTok{ 20000 }\DataTypeTok{\textbackslash{}}
  \AttributeTok{{-}Elj} \StringTok{"}\VariableTok{$FREE\_LJ}\StringTok{"} \AttributeTok{{-}Eqq} \StringTok{"}\VariableTok{$FREE\_COUL}\StringTok{"} \DataTypeTok{\textbackslash{}}
  \AttributeTok{{-}Clj}\NormalTok{ 0.181 }\AttributeTok{{-}Cqq}\NormalTok{ 0.5 }\AttributeTok{{-}ligand}\NormalTok{ LIG}
\end{Highlighting}
\end{Shaded}

Significado de las opciones principales:

{\def\LTcaptype{none} % do not increment counter
\begin{longtable}[]{@{}
  >{\raggedright\arraybackslash}p{(\linewidth - 4\tabcolsep) * \real{0.1538}}
  >{\raggedright\arraybackslash}p{(\linewidth - 4\tabcolsep) * \real{0.5495}}
  >{\raggedright\arraybackslash}p{(\linewidth - 4\tabcolsep) * \real{0.2967}}@{}}
\toprule\noalign{}
\begin{minipage}[b]{\linewidth}\raggedright
Opción
\end{minipage} & \begin{minipage}[b]{\linewidth}\raggedright
Significado
\end{minipage} & \begin{minipage}[b]{\linewidth}\raggedright
Unidad
\end{minipage} \\
\midrule\noalign{}
\endhead
\bottomrule\noalign{}
\endlastfoot
\textbf{-Elj} & Promedio LJ ligando--solvente del estado libre & kJ
mol\(^{-1}\) \\
\textbf{-Eqq} & Promedio Coulomb ligando--solvente del estado libre & kJ
mol\(^{-1}\) \\
\textbf{-Clj} & Coeficiente \(\alpha\) & adimensional \\
\textbf{-Cqq} & Coeficiente \(\beta\) & adimensional \\
\textbf{-ligand} & Nombre del grupo de energía del ligando & --- \\
\textbf{-b}, \textbf{-e} & Inicio y final del intervalo analizado & ps,
de forma predeterminada \\
\textbf{-dt} & Separación temporal entre marcos utilizados & ps \\
\end{longtable}
}

El archivo \textbf{lie.xvg} contiene la evolución de la estimación. No
se debe elegir el comienzo del promedio sólo porque la curva ``parece
estable'' a partir de un punto conveniente. El descarte debe
establecerse antes de comparar ligandos o justificarse mediante
diagnósticos de equilibrio y convergencia aplicados de manera uniforme.

También es recomendable extraer los cuatro promedios por separado y
verificar manualmente la ecuación. Esto detecta con facilidad grupos
invertidos, signos erróneos o valores libres copiados de otro sistema.

\subsection{Calibración de los
coeficientes}\label{calibraciuxf3n-de-los-coeficientes}

Para predicción cuantitativa, construya una tabla con un conjunto de
ligandos de afinidad experimental conocida:

{\def\LTcaptype{none} % do not increment counter
\begin{longtable}[]{@{}
  >{\raggedright\arraybackslash}p{(\linewidth - 6\tabcolsep) * \real{0.1146}}
  >{\raggedleft\arraybackslash}p{(\linewidth - 6\tabcolsep) * \real{0.2708}}
  >{\raggedleft\arraybackslash}p{(\linewidth - 6\tabcolsep) * \real{0.2812}}
  >{\raggedleft\arraybackslash}p{(\linewidth - 6\tabcolsep) * \real{0.3333}}@{}}
\toprule\noalign{}
\begin{minipage}[b]{\linewidth}\raggedright
Ligando
\end{minipage} & \begin{minipage}[b]{\linewidth}\raggedleft
\(\Delta V_{\mathrm{vdW}}\)
\end{minipage} & \begin{minipage}[b]{\linewidth}\raggedleft
\(\Delta V_{\mathrm{elec}}\)
\end{minipage} & \begin{minipage}[b]{\linewidth}\raggedleft
\(\Delta G^\circ_{\mathrm{exp}}\)
\end{minipage} \\
\midrule\noalign{}
\endhead
\bottomrule\noalign{}
\endlastfoot
compuesto 1 & unido − libre & unido − libre & kJ mol\(^{-1}\) \\
compuesto 2 & unido − libre & unido − libre & kJ mol\(^{-1}\) \\
\end{longtable}
}

Ajuste entonces

\[
\Delta G^\circ_{\mathrm{exp}}
=
\alpha\Delta V_{\mathrm{vdW}}
+
\beta\Delta V_{\mathrm{elec}}
+
\gamma .
\]

No es correcto promediar valores de \(\alpha\) o \(\beta\) publicados
para proteínas, campos de fuerzas o familias químicas diferentes. Con
pocos compuestos, ajustar simultáneamente tres parámetros produce
sobreajuste; en ese caso conviene fijar uno de los coeficientes con una
justificación previa o ampliar el conjunto. Deben informarse, como
mínimo, validación cruzada o conjunto externo, MAE, RMSE, correlación y
dominio de aplicabilidad.

Una predicción nueva es más confiable cuando el ligando se parece al
conjunto de calibración no sólo en estructura, sino también en sus
patrones de interacción con el receptor. Un compuesto con carga, pose o
química diferente puede estar fuera del dominio aunque su esqueleto
molecular parezca similar.

\subsection{Convergencia e
incertidumbre}\label{convergencia-e-incertidumbre}

Use series temporales y promedios por bloques para cada uno de los
cuatro términos energéticos. Para réplicas independientes, informe la
media entre réplicas y su incertidumbre; una trayectoria larga no
sustituye necesariamente varias inicializaciones cuando existen poses o
estados conformacionales separados.

Una aproximación útil para la propagación de la incertidumbre es

\[
\sigma^2_{\Delta G}
\approx
\alpha^2\sigma^2_{\Delta V_{\mathrm{vdW}}}
+
\beta^2\sigma^2_{\Delta V_{\mathrm{elec}}}
+
2\alpha\beta\operatorname{Cov}
\left(\Delta V_{\mathrm{vdW}},\Delta V_{\mathrm{elec}}\right).
\]

Esta expresión debe emplear errores de los promedios que tengan en
cuenta la autocorrelación, no la desviación estándar de los marcos como
si fueran observaciones independientes. Si \(\alpha\), \(\beta\) y
\(\gamma\) fueron ajustados, su incertidumbre también contribuye al
error predictivo.

Como diagnóstico mínimo:

\begin{itemize}
\tightlist
\item
  grafique promedios acumulativos y por bloques;
\item
  compare mitades de la región productiva;
\item
  examine por separado los estados libre y unido;
\item
  compruebe la estabilidad de la pose, contactos, hidratación y estado
  conformacional;
\item
  repita el cálculo con varias réplicas y documente cualquier exclusión.
\end{itemize}

\subsection{Múltiples poses y métodos
híbridos}\label{muxfaltiples-poses-y-muxe9todos-huxedbridos}

Cuando varias poses o conformaciones del receptor sean plausibles, no
deben concatenarse y promediarse sin criterio. Los esquemas LIE
iterativos pueden combinar estados mediante pesos de tipo Boltzmann,
pero esos pesos dependen de los parámetros del modelo y exigen
calibración iterativa.

También se han propuesto métodos híbridos que combinan LIE con energías
de solvatación alquímicas. Pueden mejorar la representación del estado
libre o permitir estrategias de anclaje, pero no convierten
automáticamente un cálculo LIE convencional en una energía libre
rigurosa. Cada término adicional, coeficiente y protocolo debe
calibrarse como parte del mismo modelo y validarse fuera del conjunto de
entrenamiento.

\subsection{Errores frecuentes}\label{errores-frecuentes-3}

\begin{itemize}
\tightlist
\item
  Usar una sola simulación del complejo y omitir el ligando libre.
\item
  Copiar \textbf{-Elj} y \textbf{-Eqq} de otro ligando.
\item
  Promediar coeficientes tomados de artículos incompatibles.
\item
  Mezclar estados de protonación o parámetros no enlazantes entre los
  dos estados.
\item
  Confundir \textbf{Coul-SR} con la energía electrostática total bajo
  PME.
\item
  Usar grupos superpuestos o dejar átomos fuera de \textbf{LIG +
  Environment}.
\item
  Interpretar cada marco de una trayectoria como una muestra
  independiente.
\item
  Seleccionar retrospectivamente el intervalo que produce el resultado
  esperado.
\item
  Reportar demasiadas cifras significativas sin incertidumbre.
\item
  Presentar una estimación con coeficientes predeterminados como
  afinidad absoluta validada.
\end{itemize}

\subsection{Qué debe informarse}\label{quuxe9-debe-informarse}

Un resultado LIE reproducible debe incluir: versión de GROMACS, campo de
fuerzas, modelo de agua, estados de protonación, parámetros
electrostáticos y de dispersión, composición de los grupos, duración
descartada y analizada, número de réplicas, promedios libre y unido,
coeficientes e intercepto, procedimiento de calibración, métricas de
validación, incertidumbre y dominio de aplicabilidad.

\subsubsection{Fuentes}\label{fuentes-7}

\begin{enumerate}
\def\labelenumi{\arabic{enumi}.}
\tightlist
\item
  \href{https://manual.gromacs.org/current/onlinehelp/gmx-lie.html}{gmx
  lie --- documentación de GROMACS 2026.3}
\item
  \href{https://manual.gromacs.org/current/reference-manual/functions/free-energy-interactions.html}{Interacciones
  de energía libre --- manual de referencia de GROMACS}
\item
  \href{https://www.frontiersin.org/journals/molecular-biosciences/articles/10.3389/fmolb.2020.00114/full}{Recent
  Developments in Linear Interaction Energy Based Binding Free Energy
  Calculations}
\item
  \href{https://pubs.acs.org/jctcce/article/16/2/1300/606208/Combined-Linear-Interaction-Energy-and-Alchemical}{Combined
  Linear Interaction Energy and Alchemical Solvation Free-Energy
  Approach}
\item
  \href{https://pubs.acs.org/jcisd8/article/59/9/4018/987937/A-Comparative-Linear-Interaction-Energy-and-MM}{A
  Comparative Linear Interaction Energy and MM/PBSA Study}
\end{enumerate}

\end{document}
